\documentclass[11pt,oneside,openany]{book}
\usepackage[T1]{fontenc}
\usepackage[utf8]{inputenc}
\usepackage{lmodern,microtype,amsmath,amssymb,amsthm,mathtools,mathrsfs}
\usepackage[paperwidth=195mm,paperheight=238mm,textwidth=158mm,top=22mm,bottom=22mm,headheight=15pt]{geometry}
\usepackage{booktabs,longtable,tabularx,array,enumitem,fancyhdr,xcolor}
\usepackage[numbers,sort&compress]{natbib}
\usepackage[colorlinks=true,linkcolor=blue!45!black,citecolor=blue!45!black,urlcolor=blue!45!black,breaklinks=true]{hyperref}
\usepackage{bookmark}
\newtheorem{theorem}{Theorem}[section]
\newtheorem{lemma}[theorem]{Lemma}
\newtheorem{proposition}[theorem]{Proposition}
\newtheorem{corollary}[theorem]{Corollary}
\theoremstyle{definition}
\newtheorem{definition}[theorem]{Definition}
\newtheorem{example}[theorem]{Example}
\newtheorem{openproblem}[theorem]{Open problem}
\newtheorem{conjecture}[theorem]{Conjecture}
\theoremstyle{remark}
\newtheorem{remark}[theorem]{Remark}
\providecommand{\doi}[1]{\href{https://doi.org/#1}{doi:#1}}
\providecommand{\natexlab}[1]{#1}
\setlist{itemsep=2pt,topsep=4pt}
\setlength{\emergencystretch}{3em}
\setcounter{tocdepth}{1}
\allowdisplaybreaks[2]
\pagestyle{fancy}\fancyhf{}
\fancyhead[L]{\small\itshape DRAGA\enspace /\enspace Multiletter Hankel Operators and Executable Memory}
\fancyhead[R]{\small\thepage}
\renewcommand{\headrulewidth}{0.3pt}
\hypersetup{pdftitle={DRAGA: Dimension Reduction and Guaranteed Approximation -- Multiletter Hankel Operators and Executable Memory},pdfauthor={Vashista Nobaub},pdfsubject={Multiletter Hankel Operators and Executable Memory},bookmarksdepth=2,pdfdisplaydoctitle=true}

\newcommand{\SR}{\mathbb R}\newcommand{\SC}{\mathbb C}
\newcommand{\SK}{\mathbb K}\newcommand{\SN}{\mathbb N}
\newcommand{\Seps}{\varepsilon}\newcommand{\Swords}{\Sigma^*}
\newcommand{\ScH}{\mathcal H}\newcommand{\ScA}{\mathcal A}
\newcommand{\Snorm}[1]{\left\lVert#1\right\rVert}
\newcommand{\Sabs}[1]{\left|#1\right|}
\newcommand{\Sip}[2]{\left\langle#1,#2\right\rangle}
\DeclareMathOperator{\Srank}{rank}\DeclareMathOperator{\Str}{tr}
\DeclareMathOperator{\Sdiag}{diag}\DeclareMathOperator{\SSpan}{span}
\DeclareMathOperator{\Sran}{ran}\DeclareMathOperator{\Ssuff}{suff}
\DeclareMathOperator{\Salg}{alg}\DeclareMathOperator{\SLL}{LL}
\DeclareMathOperator{\SEnd}{End}\DeclareMathOperator{\Sdist}{dist}
\newcommand{\SdeltaF}[1]{\delta_{#1}^{\mathbb F}}
\newcommand{\UHH}{\ell^2(\Sigma^*)}
\newcommand{\Unorm}[1]{\left\lVert#1\right\rVert}
\newcommand{\Urank}{\operatorname{rank}}
\newcommand{\UId}{I}
\newcommand{\UFF}{\mathbb F}
\DeclareMathOperator{\Udiag}{diag}
\DeclareMathOperator{\Utr}{tr}
\DeclareMathOperator{\Ucodim}{codim}
\newcommand{\UTV}{\operatorname{TV}}
\newcommand{\Bnorm}[1]{\left\lVert#1\right\rVert}
\newcommand{\BC}{\mathbb C}
\newcommand{\BR}{\mathbb R}
\newcommand{\BHH}{\ell^2(\Sigma^*)}
\newcommand{\BHtwo}{H^2}
\newcommand{\Beps}{\varepsilon}
\newcommand{\Bran}{\operatorname{ran}}
\newcommand{\Brank}{\operatorname{rank}}
\newcommand{\Btr}{\operatorname{tr}}
\newcommand{\Bdiag}{\operatorname{diag}}
\newcommand{\Fnorm}[1]{\left\lVert#1\right\rVert}
\newcommand{\FC}{\mathbb C}
\newcommand{\FR}{\mathbb R}
\newcommand{\FHH}{\ell^2(\Sigma^*)}
\newcommand{\Feps}{\varepsilon}
\newcommand{\Fran}{\operatorname{ran}}
\newcommand{\Frank}{\operatorname{rank}}
\newcommand{\Ftr}{\operatorname{tr}}
\newcommand{\Fdiag}{\operatorname{diag}}
\newcommand{\Rnorm}[1]{\left\lVert#1\right\rVert}
\newcommand{\RC}{\mathbb C}\newcommand{\RR}{\mathbb R}
\newcommand{\Rrank}{\operatorname{rank}}
\newcommand{\RHerm}{\operatorname{Herm}}
\newcommand{\Rtr}{\operatorname{Tr}}
\newcommand{\Rdiag}{\operatorname{diag}}
\newcommand{\RHH}{\ell^2(\Sigma^*)}
\newcommand{\MHH}{\ell^2(\Sigma^*)}
\newcommand{\Mnorm}[1]{\left\lVert#1\right\rVert}
\newcommand{\Mrank}{\operatorname{rank}}
\newcommand{\MId}{I}
\newcommand{\MFF}{\mathbb F}
\DeclareMathOperator{\Mdiag}{diag}
\DeclareMathOperator{\Mtr}{tr}
\DeclareMathOperator{\Mcodim}{codim}
\newcommand{\MTV}{\operatorname{TV}}
\DeclareMathOperator{\Mdist}{dist}
\DeclareMathOperator{\Mtrim}{trim}
\DeclareMathOperator{\Mrange}{range}
\newcommand{\MN}{\mathbb N}
\newcommand{\MR}{\mathbb R}
\newcommand{\Msaf}{\operatorname{stop}}

\linespread{1.31094}\begin{document}
\frontmatter
\begin{titlepage}
\vspace*{18mm}{\sffamily\large DATAR CONSULTING}\par
\vspace{17mm}\noindent{\color{blue!45!black}\rule{\linewidth}{1.4pt}}\par\vspace{6mm}
{\fontsize{46}{52}\selectfont\bfseries DRAGA}\par\vspace{6mm}
{\LARGE Dimension Reduction and\\[2pt] Guaranteed Approximation}\par\vspace{8mm}
{\Large\itshape Multiletter Hankel Operators\\ and Executable Memory}\par
\vfill{\Large Vashista Nobaub}\par\vspace{3mm}
{\sffamily Monograph edition\enspace /\enspace November 2026}\par
\end{titlepage}
\chapter*{Preface}\addcontentsline{toc}{chapter}{Preface}
Can a small sequential model preserve the responses of a larger one after any continuation?
This question has three answers in this volume, depending on what the reader charges:
operator error, error between probability laws, or bits of persistent state.
A change of criterion can reverse a conclusion about compressibility.

This edition is based on the manuscript dated 7 November 2026 and the 189-page submission
proof stamped 9 November. The three mathematical parts retain the order and numbering of
the dated manuscript. The volume adds this reading guide and a guide to the companion paper.
The companion combines two of the source constructions on one normalized trace law.
References to theorem numbers in the articles should be checked against this edition's
chapter, section, and item numbers.
\section*{Synopsis}
Sequential reasoning systems need compact states that preserve useful predictions and support certified updates. We study this requirement through multiletter word-Hankel operators, probability laws, and finite-state memory. For every fixed rank budget, an explicit binary family has a vanishing singular-value tail while its structured approximation error stays positive. We determine its boundary error by scalar rational approximation, solve the finite-parameter rank-one problem, and derive first-order corrections. Joint Stein certificates yield exact decision and synthesis formulations, with polynomial bit complexity for rational stable sources and real competitors when alphabet size and target rank are fixed. For proper finite-word laws, we prove sharp spectral-to-Kullback--Leibler rates under length-tail assumptions and a discontinuity at every fixed positive-state budget under a mean bound. Occupation flows certify total-variation error from a supplied probability dictionary. For executable memory, we characterize the residual state behind a supplied suffix window, establish exact bit and Fibonacci state optima for specified systems, and construct replay through irreversible mode changes. These results give explicit conditions for reducing sequential representations, reasoning under uncertainty, and maintaining predictive memory with charged storage and execution costs. Agent memory and recurrent-state compression for operator-controlled deployments motivate the framework.
\section*{A route through the volume}
\begin{description}
\item[Approximation theory.] Read Chapter 1, Sections 2, 8 and 19 for the binary
obstruction, its elementary proof, and the exact boundary problem. Sections 28 and
32 describe the finite-parameter optimizer and the first boundary correction.
\item[Probabilistic modeling.] Start with Chapter 2, Sections 2 and 7.
The one-state information projection and fixed-state lower bound concern laws
on completed words. Section 8 introduces positive approximants when a dictionary
and its closure errors are given.
\item[Executable memory.] Chapter 3 begins with the external suffix interface.
Sections 7 and 13 give, respectively, exact residual dimensions and a finite Fibonacci
state count. Sections 14--19 treat persistent information, arithmetic and coding costs.
\end{description}
\section*{Conventions at a glance}
\begin{tabularx}{\linewidth}{@{}lX@{}}\toprule
Symbol or unit & Meaning\\\midrule
$H_f[u,v]=f(uv)$ & Counting-measure word-Hankel operator. A split is made after the prefix $u$.\\
$\sigma_j$ & Singular values in decreasing order, numbered from one.\\
$\delta_k$ & Least operator-norm error among bounded word-Hankels of rank at most $k$.\\
$D_{\mathrm{KL}}(p\Vert q)$ & Completed-word log-loss in nats, with the source in the first argument.\\
Persistent bits & Base-two logarithmic cost of a state identifier, rounded up for fixed-width storage.\\\bottomrule
\end{tabularx}
\tableofcontents
\mainmatter
\chapter{Multiletter Hankel Approximation}\label{sec:hankel}

\section{Introduction and the approximation question}
An intelligent system that receives observations over time must decide
what to retain from its history. The retained state should support future
predictions, remain consistent after new events, and fit the resources
available for its execution. These requirements arise in knowledge
representation and reasoning under uncertainty whenever an inference
procedure replaces an event history by a compact sequential model.
The relevant mathematical question is which distinctions between
histories must survive so that the specified future responses remain
accurate after further updates.

Predictive representations of state express this idea through predictions
of future observations~\citep{littman2001}. Weighted automata provide a
finite linear model of word-dependent responses, with a classical
equivalence between finite Hankel rank and minimal realization
dimension~\citep{carlyle1971,fliess1974}. Spectral methods connect these
representations to data, including learning from a single trajectory
under finite-state policies~\citep{balle2017}. Once a model is available,
reducing its state raises further questions about which output errors are
controlled, whether the reduced state defines a valid probability law,
and how many bits its recurrent implementation requires.

Recent work on agent memory makes the operational questions concrete.
Martin~\citep{martin2026memory} describes persistent external records
organized as transcripts, notes, linked representations, and context
loaded into later sessions. Offline consolidation can check notes against
the source history and revise their organization. This notion of memory
is distinct from updating model parameters or retaining activations
within an inference run. FrontierMemory~\citep{letta2026frontiermemoryv1}
evaluates whether memory derived from prior trajectories improves later
tasks, distinguishing task conventions, user preferences, and knowledge
of the working environment. Its first release studies offline memory
formation. These systems motivate a mathematical question: which
history distinctions must be retained to preserve specified future
answers through subsequent updates? Empirical task utility and a
uniform predictive certificate are different evaluation objects; their
relationship requires an explicit task model.

The distinction matters even in a simple sequence of revisions.
An event that grants a permission followed by one that revokes it has a
different effect from the reverse order. A representation intended to
answer later permission queries must preserve the distinction relevant
to those queries. The multiletter formalism keeps each event transition
and the order of its composition. For a scalar word response $f$, its
ordinary word-Hankel operator has entries $H_f[u,v]=f(uv)$, so a row
records the continuations associated with a past word $u$. The choice of
responses, transition model, and error metric determines which histories
may be merged.

This paper develops three levels of certified reduction. The first
concerns linear realizations and structured operator error. The second
requires a proper probabilistic model and controls discrepancies between
laws of completed words. The third requires an autonomous executable
state with an explicit update, readout, and storage contract.
These levels are connected by additional assumptions that are stated
where they are used. In particular, conditioning on a rare history and
implementing a real coordinate with finitely many bits each introduce
requirements beyond a small linear approximation error.

The contributions are organized into three sections.
{\clubpenalty=10000
\begin{enumerate}[label=(\roman*)]
\item \begin{samepage}Section~\ref{sec:hankel} gives an explicit binary family whose
singular-value tail vanishes while its best fixed-rank word-Hankel error
stays positive. It identifies the exact boundary limit, solves the
finite-parameter rank-one problem, and constructs matching first-order
responses. Joint Stein certificates yield finite algebraic descriptions.
For rational mean-square-stable sources and real competitors, decision
is polynomial in input bit length when alphabet size and target rank
are fixed.
\end{samepage}
\item \begin{samepage}Section~\ref{sec:probability} establishes sharp one-state
Kullback--Leibler approximation rates under exponential length tails
and higher moments. A mean bound alone permits a discontinuity, with a
separate construction for every fixed positive-state budget. Occupation
flows turn a supplied probability dictionary into a proper decoded law
with an explicit total-variation certificate. The acquisition results
specify the positivity, conditioning, sampling, and closure premises
under which that certificate can be obtained.
\end{samepage}
\item \begin{samepage}Section~\ref{sec:memory} characterizes residual state behind a
supplied suffix window and proves exact persistent-bit optima for
explicit systems. Its Fibonacci state count applies to the specified
initial leaf union. Mode checkpoints support replay through irreversible
transitions. Pressure bounds and finite-precision execution results
include the separation, observability, and arithmetic hypotheses needed
for their claims.
\end{samepage}
\end{enumerate}}

For artificial intelligence, these results give conditions under which a
compressed sequential representation supports a stated inference or
memory contract. Operator certificates concern the modeled continuation
responses. Probability certificates concern complete-word laws, with
additional conditions for rare-prefix predictions. Executable-memory
certificates account for persistent state and distinguish it from program
size, workspace, and update time. The dictionary and interface are part of
the mathematical input. The results are theoretical, and the final
recurrent-state application states a separate requirement for measured
gains in complete encoded length.

\paragraph{Research programme and deployment context.}
This work forms part of Datar Consulting's research on memory and
compression for large language models (LLMs), with sovereign model
deployments as an application objective. The programme studies
persistent agent memory, numerical recurrent state---including Gated
DeltaNet (GDN)~\citep{yang2025gated}---and exact residual encodings.
Here a sovereign deployment means that the operator controls the
infrastructure, retained data, and execution policy within a chosen
deployment boundary. The research objective is to reduce storage and
execution costs while preserving specified predictive behavior within
that boundary. An agent's external records and a decoder's recurrent
matrix state have different interfaces and correctness criteria.
The results below supply contracts for those interfaces under explicit
hypotheses; practical compression gains and complete-system behavior
remain implementation and evaluation questions.

For the operator problem, a weighted finite automaton represents a word
series by $f(w)=c^*A_wb$, where
$A_{a_1\cdots a_t}=A_{a_1}\cdots A_{a_t}$.
We use counting-measure $\ell^2(\Sigma^*)$ and define
\[
 \delta_k^{\mathbb F}(H_f)=
 \inf_{\substack{H_g\text{ bounded word-Hankel over }\mathbb F\\
                  \operatorname{rank}H_g\le k}}
                  \|H_f-H_g\|,\qquad \mathbb F\in\{\mathbb R,\mathbb C\}.
\]
The unconstrained lower bound is $\sigma_{k+1}(H_f)$ for compact $H_f$.
Classical one-letter AAK theory, including its automata formulation
\citep{lacroce2024}, gives exact optimal reduction under its hypotheses.
Singular value automata give canonical realizations and truncation bounds
\citep{sva2019}. A multiletter noncommutative approach was proposed in
\citep{lacroce2022}. Here the target is the ordinary counting-measure
word-Hankel structure. Subsection~\ref{I:U:sec:free-aak-obstruction}
checks the row-contraction premise needed for that cited approach.

The binary obstruction shows that a vanishing spectral tail alone cannot
control the structured optimum uniformly at fixed rank. Its stability
margin closes at the boundary, so uniform stability assumptions matter
when comparing it with other approximation guarantees. The exact
boundary constant has a scalar rational-approximation characterization.
Its identification with a reciprocal Hardy coefficient norm uses
published low-degree extremals through budget four and the stated
external outer-extremal input at general budget. The proofs identify
these inputs and keep the real and complex competitor classes explicit.

\section{A binary obstruction with an exact limit}
For distinct letters $a,b$, $k\ge1$ and $r\in(0,1)$, put $q=1-r^2$ and
\begin{equation}\label{I:eq:binary-source}
 f_{r,k}(a^j b^k)=qr^j\quad(j\ge0),\qquad
 f_{r,k}(w)=0\text{ otherwise}.
\end{equation}
The same notation is used throughout the proofs. Let
\[
 E_m^{\mathbb F}=\inf_{\substack{h\in H^2\text{ over }\mathbb F\\
                 \operatorname{rank}(h_{i+j})_{i,j\ge0}\le m}}
                 \|z^m-h\|_{H^2},\qquad E_0^{\mathbb F}=1,
\]
where circle measure is normalized. Write $E_m=E_m^{\mathbb C}$.

\begin{theorem}[Structural separation and boundary reduction]
\label{I:I:main:separation}
The source in \eqref{I:eq:binary-source} has rank $k+1$, norm one, and nonzero
singular values $1,\sqrt q,\ldots,\sqrt q$, with $k$ copies of $\sqrt q$.
For either field and every $0<r<1$,
\[
 \delta_k^{\mathbb F}(H_{f_{r,k}})\ge E_k>0.
\]
For each fixed budget and each field separately,
\[
 \lim_{r\uparrow1}\delta_k^{\mathbb F}(H_{f_{r,k}})=E_{k-1}^{\mathbb F}>0.
\]
Consequently no finite factor depending only on $(k,|\Sigma|)$ bounds
$\delta_k^{\mathbb F}/\sigma_{k+1}$ on all bounded word-Hankel sources.
\end{theorem}

The singular decomposition is explicit. Let
$z_j=\sqrt q\sum_{n\ge0}r^ne_{a^nb^j}$ for $0\le j\le k$.
These vectors are orthonormal and direct enumeration of the cuts of a word gives
\[
 H_{f_{r,k}}=z_0z_k^*+\sqrt q\sum_{j=0}^{k-1}z_{k-j}e_{b^j}^*.
\]
Both input and output systems in this decomposition are orthonormal.
The singular-value statement follows, and every usual finite-rank
spectral-tail norm tends to zero with $q$.

For the uniform floor, the row space of a rank-$k$ word-Hankel is
invariant under each backward letter shift. Its overlap with $z_k$ is
bounded by a scalar backward-shift invariant subspace problem. The
constant previously denoted by $\kappa_k$ has the single interpretation
\[
 \kappa_k=\Gamma_k^2=1-E_k^2,
 \qquad
 \Gamma_k=\max_{D^*D+l^*l\preceq I_k}\|lD^k\|.
\]
Thus the uniform floor at budget $k$ is $E_k$, while the boundary
limit is $E_{k-1}$. Full proofs are in
Theorems~\ref{I:U:thm:separation}, \ref{I:M:rev13:thm:uniform} and
\ref{I:B:thm:main}.

\subsection{Why one dimension is lost at the boundary}
Let $A,B$ be the backward shifts for $a,b$ restricted to a competitor's
row space, and let $lx=x[\varepsilon]$. Then
$A^*A+B^*B+l^*l\preceq I$. The exact identity
\[
 (I-rA^*)(I-rA)-q(I-A^*A)=(A-rI)^*(A-rI)
\]
implies $\|\sqrt q\,B(I-rA)^{-1}\|\le1$. By compactness, a sequence
of passive tuples has a limit. If $I-A_0$ is invertible, the relevant
singular-vector overlap tends to zero. Otherwise the fixed subspace of
the contraction $A_0$ is reducing, and passivity forces both $B_0$ and
$l_0$ to vanish on it. At least one dimension is therefore unavailable
to the remaining scalar extremal problem. This gives the lower bound
$E_{k-1}^{\mathbb F}$. A block realization lifting a scalar minimizer
provides the matching upper bound. The argument is over all bounded
rank-$k$ competitors and does not impose a common stability margin.

\section{The Hardy coefficient bridge and its inputs}
For $0<p<\infty$, define
$C(m,p)=\sup_{\|F\|_{H^p}=1}|[z^m]F|$.
The scalar rational class includes polynomial parts; its rank convention
is the rank of the coefficient Hankel $(h_{i+j})$, not a rational-degree
convention that discards those parts.

\begin{theorem}[Coefficient-norm identification]
\label{I:I:main:hardy}
For every $m\ge0$, $E_m^{\mathbb F}\ge1/C(m,2/3)$.
Equality over both fields for $0\le m\le3$ follows from the published
low-degree extremals in~\citet{brevig2020}. For general $m$, equality
over both fields uses the outer-extremal theorem of
\citet{outer2026}, together with the polynomial structure of coefficient
extremals in~\citet{brevig2020}.
\end{theorem}

The lower inequality has a short proof independent of that external
outer-extremal input. A degree-$m$ model space is $K_\Theta$, with
$\Theta=D_m^*/D$. Let $b=P_{\le m}\Theta$. Taylor coefficient extraction
and H\"older's inequality give
\[
 \|D\|_2^2=[z^m](bD^2)
 \le C(m,2/3)\|bD^2\|_{2/3}
 \le C(m,2/3)\|b\|_2\|D\|_2^2.
\]
The squared projection error for $z^m$ is
$\sum_{j=0}^m|\theta_j|^2=\|b\|_2^2$, proving the bound. Smaller
model spaces may be enlarged, so the same inequality covers all ranks
at most $m$.

For the upper bound, an outer coefficient extremal has the form
$P^3$, normalized by $\|P\|_2=1$ and $[z^m]P^3=C$.
Stationarity on the polynomial unit sphere yields
$P^2=CP_m^*+O(z^{m+1})$. Therefore
$\Theta=P_m^*/P=P/C+O(z^{m+1})$, and the projection error is $1/C$.
Nonnegative real coefficient extremals give a real attaining lift.
Theorem~\ref{I:B:thm:hardy-bridge} includes the full construction and the
precise external input; no numerical optimization is used as a proof.

\begin{center}
\begin{tabular}{cll}\toprule
Budget $k$ & Boundary error & Source of scalar input\\\midrule
1 & $1$ & Elementary\\
2 & $\sqrt3/2$ & Published degree-one extremal\\
3 & $3/4$ & Published degree-two extremal\\
4 & $\sqrt{(1103-33\sqrt{33})/2048}$ & Published degree-three extremal\\
General $k$ & $1/C(k-1,2/3)$ & Outer-extremal theorem~\cite{outer2026}\\\bottomrule
\end{tabular}
\end{center}
The classical model-space formulation is also used in rational
approximation~\citep{rational2026}. A numerical coincidence with a
constant in a different rational-degree problem is not an identification
of the approximation classes.

\section{Finite-parameter optimality and first-order response}
Write $D_k^{\mathbb F}(q)=\delta_k^{\mathbb F}(H_{f_{r,k}})^2$.
At rank one, let $u(q)$ be the unique root in $(q,\sqrt q)$ of
\[
 q^2(1-u^2)(u-q)^2=(1-q)^3u^2(q-u^2).
\]
The exact optimum is
\[
 D_1^{\mathbb R}(q)=D_1^{\mathbb C}(q)
 =\frac{q^2+(1-2q)u(q)^2}{q(1-u(q)^2)}.
\]
Theorem~\ref{I:F:thm:F1-exact} constructs and classifies every optimizer.
The phase reduction is an explicit two-sided unitary conjugation; the
real and complex equality here is proved, rather than inferred from
real-valued input data.

At every fixed budget $k\ge2$, the general response theorem gives
\[
 D_k^{\mathbb F}(q)=(E_{k-1}^{\mathbb F})^2-b_k^{\mathbb F}q+o(q).
\]
The coefficient is the attained finite variational maximum in
\eqref{I:R:eq:main-response-coefficient}; no universal closed formula or
uniform efficient optimizer is asserted. The diagonal real/complex
comparison uses the scalar-extremal classification identified in
Section~\ref{I:R:subsec:real-diagonal-orbits}. Explicit cases are
\begin{align*}
 D_2^{\mathbb F}(q)&=\frac34-\frac14q+O(q^2),\\
 D_3^{\mathbb F}(q)&=\frac9{16}-\frac17q+O(q^2),\\
 D_4^{\mathbb F}(q)&=\frac{1103-33\sqrt{33}}{2048}
       -\frac{855+103\sqrt{33}}{15872}q+o(q).
\end{align*}
The theorem also permits independent state budget and terminal delay,
and extensions to arbitrary terminal languages.

The proof separates two tasks. A passive-envelope calculation bounds
all structured competitors from below. A lossless construction then
produces an actual word-Hankel operator attaining that response to first
order. Passivity alone is insufficient for this second task. Joint
stability of the constructed realization follows from a finite
observation Gramian, and the improved $O(q^2)$ upper remainder follows
from an exact sign-conjugation symmetry of the analytic trial. These
steps are proved in Theorem~\ref{I:R:thm:v4-matching}. The global lower
remainder is only $o(q)$ without the additional structure established
in the explicit low-order cases.

\section{Finite certificates and computational scope}
Given a stable finite realization, a candidate's error realization has
joint reachability and observability Gramians. Their Stein equations and
a positive-semidefinite norm constraint give an exact finite certificate
for $\delta_k(H_f)\le t$. Attainment is essential for the non-strict
threshold. Factoring positive-semidefinite matrices yields a degree-three
existential real formula and algebraic optimal witnesses
(Theorem~\ref{I:U:thm:exact} and Corollary~\ref{I:U:cor:algebraic}).
This is an application of established realization, matrix-positivity and
real-algebraic tools, not a hardness result.

For rational mean-square-stable sources and real competitors, when
$d=|\Sigma|$ and $k$ are fixed, source-dependent Gramian variables
can be eliminated before applying real-algebraic decision algorithms.
There remain
$p=dk^2+2k+k(k+1)/2$ candidate variables; the number and degrees of the
polynomials grow polynomially with the source dimension. The resulting
bit-time bound has a polynomial exponent $O(p)$ for fixed parameters.
This is an XP result in $(d,k)$, not a fixed-parameter tractability claim
and not joint polynomiality in all dimensions
(Theorem~\ref{I:U:thm:fixed}). The constructive global net in
Theorem~\ref{I:F:thm:g1-net} gives a separate certified approximation scheme
with explicit but generally large costs.

\section{Limits and organization of the proofs}
The target space and norm are part of every claim. Ordinary counting-measure
backward letter shifts have row norm $\sqrt d$, so the row-contraction
hypothesis required by the cited noncommutative AAK route is not verified
for those shifts. This observation does not identify that model with all
Fock-space or bilateral-shift formulations.

Three further questions remain open here: useful complexity dependence
on varying rank; finite-parameter optimizers and field comparison beyond
the solved examples; and a uniform multiletter balanced-truncation tail
bound under explicit geometric hypotheses. The Stein-witness truncation
bounds proved below do not establish the classical
$2\sum_{j>k}\sigma_j$ estimate.

The remaining subsections contain complete proofs and constructions, including the
operator preliminaries, obstruction family, exact boundary reduction,
Hardy bridge, finite-parameter and response analysis, structured examples,
and numerical certificates. Repeated foundational results are cited
within the article rather than reproved under different names.
\label{I:I:main-end}

\section*{Detailed proofs and supporting results}

\section{Notation and finite stable sources}\label{I:U:sec:preliminaries}
Let $\Sigma$ be a finite alphabet of size $d\ge1$, let $\Sigma^*$ include
the empty word $\varepsilon$, and let $\UFF$ be either $\mathbb R$ or
$\mathbb C$. All operators act on the corresponding Hilbert space
$\UHH$ unless another space is specified. Its orthonormal basis is
$(e_w)_{w\in\Sigma^*}$. A word-Hankel operator has entries
\[
H_f[u,v]=f(uv),\qquad u,v\in\Sigma^*.
\]
The real and complex optimization problems are treated separately; no
equality between their optimal values is assumed. The notation $X^*$
means transpose in the real case and conjugate transpose in the complex
case. Positive semidefinite order is denoted by $\preceq$.

A finite realization of dimension $n$ consists of matrices
$A_a\in\UFF^{n\times n}$ and vectors $b_f,c_f\in\UFF^n$ such that
\begin{equation}\label{I:U:eq:realization}
f(w)=c_f^*A_wb_f,
\qquad A_{a_1\cdots a_j}=A_{a_1}\cdots A_{a_j},
\qquad A_\varepsilon=\UId_n.
\end{equation}
It is \emph{mean-square stable} when the positive linear map
\[
\Phi_A(X)=\sum_{a\in\Sigma}A_aXA_a^*
\]
has spectral radius less than one. This is stability of a finite linear
map, not the generally different condition on the joint spectral radius
of arbitrary products of the $A_a$.

Under this assumption the observation and reachability maps, with rows
$c_f^*A_u$ and columns $A_vb_f$, respectively, are bounded and satisfy
\begin{align}
P_f&=\sum_{w\in\Sigma^*}A_wb_fb_f^*A_w^*,
&P_f&=b_fb_f^*+\sum_aA_aP_fA_a^*,\label{I:U:eq:sourceP}\\
Q_f&=\sum_{w\in\Sigma^*}A_w^*c_fc_f^*A_w,
&Q_f&=c_fc_f^*+\sum_aA_a^*Q_fA_a.\label{I:U:eq:sourceQ}
\end{align}
These are convergent series of finite matrices; the Stein equations have
unique solutions. Writing the maps as $O_f,C_f$ gives
$H_f=O_fC_f$, $P_f=C_fC_f^*$, and $Q_f=O_f^*O_f$.
The nonzero squared singular values of $H_f$ are the nonzero eigenvalues
of $P_f^{1/2}Q_fP_f^{1/2}$. In particular,
\begin{equation}\label{I:U:eq:trace}
\Unorm{H_f}^2\le\Unorm{H_f}_{\mathrm{HS}}^2
=\Utr(P_fQ_f).
\end{equation}
These identities do not require invertible Gramians.

Throughout the algorithmic results, the source is supplied by a finite
stable realization, not by a table of finitely many coefficients or an
arbitrary infinite-dimensional oracle. Algebraic input entries are given
by effective exact representations. Bit-complexity statements concern
rational input, with numerators and denominators encoded in binary and
total description length denoted by $L$. For complex rational data, both
real and imaginary parts have this encoding. The threshold is rational
and included in $L$. A requested rank $k\ge n$ is trivial since the
source itself is feasible; $k=0$ means that the only competitor is zero.

For later use define the left and right creation operators by
\begin{equation}\label{I:U:eq:shifts}
L_ae_w=e_{aw},\qquad R_ae_w=e_{wa}.
\end{equation}
They obey
$\sum_aL_aL_a^*=\sum_aR_aR_a^*=\UId-e_\varepsilon e_\varepsilon^*$
and $R_a^*R_a=L_a^*L_a=\UId$. The word-Hankel identities are equivalently
\begin{equation}\label{I:U:eq:intertwining}
R_a^*H_f=H_fL_a\qquad(a\in\Sigma).
\end{equation}
Taking adjoints yields the dual identity
\begin{equation}\label{I:U:eq:adjoint-intertwining}
L_a^*H_f^*=H_f^*R_a\qquad(a\in\Sigma).
\end{equation}
In particular, the prefix residual $v\mapsto p(uv)$ of a real probability
law is the column $H_p^*e_u$, rather than $H_pe_u$.
For a complex series, $(H_f^*e_u)[v]=\overline{f(uv)}$ instead.


The structured approximation problem, with scalar field fixed, is
\begin{equation}\label{I:U:eq:problem}
 \delta_k(H_f)=\inf\{\|H_f-H_g\|:H_g\text{ bounded word-Hankel},\ \Urank H_g\le k\}.
\end{equation}
The empty word is denoted by $\varepsilon$; a positive error tolerance is
written $\epsilon$ or explicitly named when both occur.

\section{A binary obstruction at every rank budget}\label{I:U:sec:obstruction}
Use the source in \eqref{I:eq:binary-source}, with $q=1-r^2$.
The source changes with $r$, but the alphabet and target rank stay fixed.

\begin{theorem}[No uniform scalar AAK factor]\label{I:U:thm:separation}
The source has rank $k+1$ and singular values $1,\sqrt q,\ldots,\sqrt q$,
where $\sqrt q$ occurs $k$ times. Let $C_k=\binom{2k}{k}$. Over either
scalar field,
\[
 \delta_k(H_{f_{r,k}})\ge
 \max\left\{\sqrt q,\sqrt{\frac{1+q^2(C_k-1)}{C_k}}\right\}.
\]
In particular, $\sigma_{k+1}\to0$ but
$\liminf_{r\uparrow1}\delta_k\ge C_k^{-1/2}>0$. No constant depending
only on $k$ and the alphabet size can bound $\delta_k$ by that singular value.
\end{theorem}
\begin{proof}
Let $z_j=\sqrt q\sum_{n\ge0}r^n e_{a^n b^j}$ for $0\le j\le k$.
These are orthonormal. Directly sorting the possible cuts of $a^n b^k$ gives
\[
 H_{f_{r,k}}=z_0z_k^*+\sqrt q\sum_{j=0}^{k-1}z_{k-j}e_{b^j}^*.
\]
The right vectors $z_k,e_\varepsilon,e_b,\ldots,e_{b^{k-1}}$ are
orthonormal, as are the corresponding left vectors. This proves the rank and
spectrum, and the usual unstructured rank lower bound gives $\sqrt q$.

Consider any bounded rank-$s$ word-Hankel $G$, where $s\le k$. Its finite
range is invariant under the right backward shifts. The canonical realization
on that range (proved in Lemma~\ref{I:U:lem:compact}) has letter matrices of norm
at most one. The characteristic polynomial of its $b$ matrix, multiplied
by $z^{k-s}$, is a monic annihilating polynomial
$p(z)=z^k+\sum_{j=0}^{k-1}c_jz^j$. Its roots have modulus at most one.
The elementary symmetric coefficient bounds and Vandermonde's identity yield
\[
 S:=\sum_{j=0}^{k-1}|c_j|^2\le\sum_{j=0}^{k-1}\binom{k}{j}^2=C_k-1.
\]
For each $n$, the columns indexed by $a^n b^j$ obey the annihilating
relation. Boundedness permits summation of the convergent vectors $z_j$,
and therefore $Gx=0$ for $x=z_k+\sum_{j<k}c_jz_j$.
The displayed source decomposition gives
\[
 \|x\|^2=1+S,\qquad
 H_{f_{r,k}}x=z_0+q\sum_{j<k}c_jz_{k-j},\qquad
 \|H_{f_{r,k}}x\|^2=1+q^2S.
\]
Since $(1+q^2S)/(1+S)$ decreases with $S\ge0$, its value is at least
$(1+q^2(C_k-1))/C_k$. Taking the infimum over $G$ proves the claim.
The zero competitor is included by the polynomial $z^k$.
\end{proof}


\begin{remark}[Scope of the obstruction]
For rational $r$, the target has a nonnegative rational realization: a waiting
state loops on $a$ with weight $r$, followed by a chain of $k$ transitions on
$b$, with output scale $q$. It is minimal by the spectrum above and is
mean-square stable. Its stability margin tends to zero as $r\uparrow1$.
Thus a positive theorem depending on a fixed stability margin or additional
geometry is not excluded. The target is not a normalized word probability
law. The theorem is neither a computational hardness statement nor a
persistent-bit lower bound.
\end{remark}
\subsection{The row-contraction hypothesis in noncommutative AAK theory}\label{I:U:sec:free-aak-obstruction}

The preceding separation does not contradict Popescu's multivariable
AAK theorem \cite{popescu2003}.  The required hypothesis is stronger than
the Hankel intertwining identity alone.  The following calculation makes
the distinction explicit.

\begin{proposition}[Target-shift obstruction]\label{I:U:prop:rowobstruction}
On counting-measure word space, let
$L_ae_w=e_{aw}$ and $R_ae_w=e_{wa}$.  Every bounded word-Hankel satisfies
\[
 H_fL_a=R_a^*H_f.
\]
For an alphabet of size $d$, however,
\begin{equation}\label{I:U:eq:row-shift-norm}
 \Unorm{[R_1^*\ \cdots\ R_d^*]}=\sqrt d.
\end{equation}
Thus the tuple $(R_a^*)_{a\in\Sigma}$ is not a row contraction when
$d\ge2$.
\end{proposition}
\begin{proof}
The intertwining identity follows by evaluating both sides at row $u$
and column $v$: each entry is $f(uav)$.  Each $R_a$ is an isometry, so
\[
 [R_1^*\ \cdots\ R_d^*]
 [R_1^*\ \cdots\ R_d^*]^*
 =\sum_aR_a^*R_a=d\UId.
\]
This proves \eqref{I:U:eq:row-shift-norm}.  Equivalently, taking
$h_a=R_ax$ with $x\ne0$ gives
$\Unorm{\sum_aR_a^*h_a}^2=d^2\Unorm{x}^2$, whereas
$\sum_a\Unorm{h_a}^2=d\Unorm{x}^2$.
\end{proof}


Popescu's hypotheses require a row-expansive source tuple $X$ and a
row-contractive target tuple $T$, namely
\[
 \Unorm{\sum_aX_ay_a}^2\ge\sum_a\Unorm{y_a}^2,
 \qquad
 \Unorm{\sum_aT_ah_a}^2\le\sum_a\Unorm{h_a}^2.
\]
The resulting theorem is stated in terms of restrictions to
$X$-invariant subspaces of bounded codimension and, equivalently,
multipliers on those subspaces
\cite[Theorems~5.3--5.4]{popescu2003}.
The left creators satisfy the first inequality, but
Proposition~\ref{I:U:prop:rowobstruction} rules out the second for the
ordinary backward target shifts.  Unitary word reversal preserves
the norm in \eqref{I:U:eq:row-shift-norm}; a dilation cannot fix this by
compression either, since compressions of row contractions remain
row contractions.

The comparison with the noncommutative framework of
\cite[Theorems~7--10]{lacroce2022} must keep its operator conventions
separate from the counting-measure model used here. For the latter model,
the hypotheses of Popescu's theorem are not verified by the ordinary
backward target shifts: their row norm is $\sqrt d$, and their ranges
are not pairwise orthogonal. Proposition~\ref{I:U:prop:rowobstruction}
and Theorem~\ref{I:U:thm:separation} are the precise conclusions needed
here. They do not address a different Fock-space or bilateral-shift
formulation without a separate identification of its norm and rank classes.

\section{Attainment and finite-dimensional realizations}\label{I:U:sec:realization}
We first separate existence from computation. Existence uses only the operator topology; the realization results subsequently turn finite-rank competitors into finite certificates.

\begin{theorem}[Attainment for every bounded source]\label{I:U:thm:attainment}\label{I:S:thm:attainment}
Let $H_f$ be any bounded operator on word space (in particular a bounded word-Hankel operator), and let $k\ge0$. There exists a bounded word-Hankel operator $G_\star$ with $\Urank G_\star\le k$ such that
\[
 \Unorm{H_f-G_\star}=\delta_k(H_f).
\]
The assertion holds separately over $\mathbb R$ and $\mathbb C$, without any finite-rank or stability assumption on the source.
\end{theorem}
\begin{proof}
Choose a minimizing sequence $(G_j)$ of admissible operators with
$\Unorm{H_f-G_j}\le\delta_k(H_f)+1$. Since the zero operator is admissible,
$\delta_k(H_f)\le\Unorm{H_f}$ and hence
$\Unorm{G_j}\le 2\Unorm{H_f}+1=:M$.
The closed ball of radius $M$ in $B(\UHH)$ is compact in the weak operator topology. In the present separable space, a subsequence can be obtained directly by diagonal extraction of all matrix coefficients: the uniform norm bound extends the limiting coefficients to a bounded operator $G$, and then finite-support approximation gives weak operator convergence on arbitrary pairs of vectors.

The word-Hankel identities are equalities between matrix coefficients and therefore pass to this limit. Every $(k+1)\times(k+1)$ minor also passes to the limit and remains zero, so $\Urank G\le k$. Indeed, rank greater than $k$ would give $k+1$ independent columns and then a nonzero coordinate minor of that order. Finally, for any unit vectors $x,y$,
\[
 |\langle(H_f-G)x,y\rangle|
 =\lim_j|\langle(H_f-G_j)x,y\rangle|
 \le\liminf_j\Unorm{H_f-G_j}.
\]
Taking the supremum over $x,y$ proves the result. For real operators, perform the same argument in the real Hilbert space.
\end{proof}


The following normalization packages bounded realizations into a compact parameter set. It uses the familiar reachability and observability factorization of a rational Hankel operator; see \cite{sva2019} for the weighted-automata setting. Its formulation with non-strict inequalities also admits singular and redundant realizations.

For $M>0$ and $k\ge1$, let $\mathcal K_{k,M}$ consist of tuples
$\theta=((T_a)_{a\in\Sigma},b,c,P)$, with $T_a\in\mathbb F^{k\times k}$,
$b,c\in\mathbb F^k$ and $P=P^*$, satisfying
\begin{align}
 \sum_a T_a^*T_a+cc^*&\preceq\UId_k,\label{I:U:eq:normal-observation}\\
 \sum_a T_aPT_a^*+bb^*&\preceq P,\label{I:U:eq:normal-reachability}\\
 0\preceq P&\preceq M^2\UId_k.\label{I:U:eq:normal-bound}
\end{align}
Here $\mathbb F$ is the chosen scalar field. Define $g_\theta(w)=c^*T_wb$.

\begin{lemma}[Compact and complete parametrization]\label{I:U:lem:compact}
The set $\mathcal K_{k,M}$ is compact, and
\[
 \{H_{g_\theta}:\theta\in\mathcal K_{k,M}\}
 =\{G:G\text{ is bounded word-Hankel},\ \Urank G\le k,\ \Unorm G\le M\}.
\]
\end{lemma}
\begin{proof}
For the forward inclusion, introduce the positive maps
\[
 \Psi(X)=\sum_aT_a^*XT_a,
 \qquad \Phi_T(X)=\sum_aT_aXT_a^*.
\]
Iterating \eqref{I:U:eq:normal-observation} and
\eqref{I:U:eq:normal-reachability} gives the explicit partial-sum bounds
\begin{align*}
 \sum_{j=0}^{N}\Psi^j(cc^*)
 &\preceq \UId_k-\Psi^{N+1}(\UId_k)\preceq\UId_k,\\
 \sum_{j=0}^{N}\Phi_T^j(bb^*)
 &\preceq P-\Phi_T^{N+1}(P)\preceq P.
\end{align*}
The summands are respectively
$\sum_{|u|=j}T_u^*cc^*T_u$ and
$\sum_{|v|=j}T_vbb^*T_v^*$. Thus the matrix with rows $c^*T_u$ defines a bounded observation operator $O:\mathbb F^k\to\UHH$, and the matrix with columns $T_vb$ defines a bounded reachability operator $C:\UHH\to\mathbb F^k$, with
\[
 O^*O\preceq\UId_k,\qquad CC^*\preceq P.
\]
For $C$, one may first define $C^*x=(b^*T_v^*x)_v$ and use the second bound, avoiding any assumption of absolute convergence of an arbitrary input sum. Since
$(OC)[u,v]=c^*T_uT_vb=g_\theta(uv)$, the product is the desired bounded Hankel, and its rank is at most $k$ and its norm at most $M$.

Conversely, let $G$ be a nonzero bounded word-Hankel of rank $s\le k$. Its range $S$ is finite-dimensional and closed. Choose an isometry
$O:\mathbb F^s\to\UHH$ onto $S$. Define the left and right creation operators by
\[
 L_ae_v=e_{av},\qquad R_ae_u=e_{ua}.
\]
The Hankel identity is $R_a^*G=GL_a$, so $S$ is invariant under $R_a^*$. Set
\[
 T_a=O^*R_a^*O,\quad c=O^*e_\varepsilon,
 \quad C=O^*G,\quad b=Ce_\varepsilon.
\]
Then $OT_a=R_a^*O$ and $T_aC=CL_a$. To check the word order explicitly, if $o_u$ denotes row $u$ of $O$, then $o_{ua}=o_uT_a$ and $o_\varepsilon=c^*$. Thus $o_u=c^*T_u$. Similarly, column $av$ of $C$ equals $T_a$ times column $v$, so column $v$ is $T_vb$. Hence $G=OC$ has the prescribed realization.

The ranges of the $R_a$ partition the nonempty words. Therefore
\[
 \sum_aT_a^*T_a+cc^*
 =O^*\left(\sum_aR_aR_a^*+e_\varepsilon e_\varepsilon^*\right)O
 =\UId_s.
\]
Taking $P=CC^*$ and partitioning the columns of $C$ by their first letter gives
\[
 P=bb^*+\sum_aT_aPT_a^*,\qquad
 P\preceq\Unorm{G}^2\UId_s\preceq M^2\UId_s.
\]
Pad all state variables by zeros to reach dimension $k$. The zero Hankel is represented by the zero tuple.

Finally, the constraints imply
$\Unorm{T_a}\le1$, $\Unorm c\le1$, $\Unorm b\le M$, and
$\Unorm P\le M^2$. They are closed conditions in a finite-dimensional real vector space, including after separating real and imaginary parts. Hence the parameter set is compact.
\end{proof}


Although the compact set permits boundary modes, every output it represents also has a strictly stable realization of the same allowed dimension.

\begin{lemma}[Stable realization of a bounded finite-rank Hankel]\label{I:U:lem:stable}
Every bounded word-Hankel of rank at most $k$ has a realization of dimension $k$ satisfying $\rho(\Phi_T)<1$. A nonzero operator of rank $s$ has such a realization of dimension $s$.
\end{lemma}
\begin{proof}
Use the canonical $s$-dimensional realization in the converse proof of Lemma~\ref{I:U:lem:compact}. Its observation operator is an isometry, so
\[
 \sum_{w\in\Sigma^*}T_w^*cc^*T_w=\UId_s.
\]
The identity $\UId_s=cc^*+\Psi(\UId_s)$ then yields
\[
 \Psi^N(\UId_s)=\UId_s-\sum_{|w|<N}T_w^*cc^*T_w\longrightarrow0
\]
in matrix norm. For a Hermitian matrix $X$, positivity and
$-\Unorm X\UId_s\preceq X\preceq\Unorm X\UId_s$ imply
$\Unorm{\Psi^N(X)}\le\Unorm X\,\Unorm{\Psi^N(\UId_s)}$.
Splitting a general complex matrix into Hermitian real and imaginary parts shows that the powers of $\Psi$ tend to zero as linear maps on the entire matrix space. In finite dimension this is equivalent to $\rho(\Psi)<1$.
For real realizations, apply the same argument to the complexification of the real matrix map.
The map $\Phi_T$ is the adjoint of $\Psi$ for the Hilbert--Schmidt inner product, and hence has the same spectral radius. Zero padding preserves strict stability; the zero operator has a zero-transition realization.
\end{proof}


\begin{remark}
Theorem~\ref{I:U:thm:attainment} is a compactness argument, not an algorithm. Lemma~\ref{I:U:lem:stable} concerns the existence of a stable representation, not the stability of every redundant representation of the same series. In particular, the closed constraints defining $\mathcal K_{k,M}$ intentionally allow unreachable or unobservable boundary modes.
\end{remark}

\section{An exact finite certificate}\label{I:U:sec:exact}
Assume now that the source is supplied by the finite mean-square-stable realization
$f(w)=c_f^*A_wb_f$ of dimension $n$. Let the candidate have $k$ states and put $m=n+k$. Its joint error realization is
\begin{equation}\label{I:U:eq:error-realization}
 B_a=\begin{pmatrix}A_a&0\\0&T_a\end{pmatrix},\qquad
 \beta=\begin{pmatrix}b_f\\b\end{pmatrix},\qquad
 \gamma=\begin{pmatrix}c_f\\-c\end{pmatrix}.
\end{equation}
Thus $\gamma^*B_w\beta=f(w)-g(w)$. No damping is applied to either block.

The next theorem is a Hankel-norm certificate built from reachability and observability Stein inequalities. The two ingredients have standard precedents in weighted-automata Gramians \cite{sva2019} and realization theory over free semigroups \cite{ball2005}. We prove the equivalence directly, including singular Gramians, rather than identifying it with a transfer-function bounded-real lemma.

\begin{theorem}[Exact joint Stein certificate]\label{I:U:thm:exact}
For $t\ge0$, the inequality $\delta_k(H_f)\le t$ holds if and only if there exist candidate parameters $(T_a,b,c)$, Hermitian matrices $P,Q\in\mathbb F^{m\times m}$ and $W\in\mathbb F^{m\times m}$ such that
\begin{align}
 P&\succeq0, & Q&\succeq0,\label{I:U:eq:exact-positive}\\
 P-\beta\beta^*-\sum_aB_aPB_a^*&\succeq0,\label{I:U:eq:exact-reach}\\
 Q-\gamma\gamma^*-\sum_aB_a^*QB_a&\succeq0,\label{I:U:eq:exact-observe}\\
 P&=WW^*,\label{I:U:eq:exact-factor}\\
 t^2\UId_m-W^*QW&\succeq0.\label{I:U:eq:exact-norm}
\end{align}
Every feasible tuple defines a bounded word-Hankel $H_g$ of rank at most $k$ with $\Unorm{H_f-H_g}\le t$. At $t=\delta_k(H_f)$ a feasible tuple exists.
\end{theorem}
\begin{proof}
Suppose first that a tuple is feasible. Iterating
\eqref{I:U:eq:exact-reach}--\eqref{I:U:eq:exact-observe} gives
\[
 \sum_{|w|\le N}B_w\beta\beta^*B_w^*\preceq P,
 \qquad
 \sum_{|w|\le N}B_w^*\gamma\gamma^*B_w\preceq Q
\]
for every $N$. Exactly as in Lemma~\ref{I:U:lem:compact}, these inequalities define bounded reachability and observation maps $C_E$ and $O_E$, with exact Gramians
\[
 P_E=C_EC_E^*\preceq P,\qquad Q_E=O_E^*O_E\preceq Q.
\]
Restriction to the candidate state coordinates also gives bounded candidate maps $C_g,O_g$: the lower diagonal blocks of the partial sums are precisely their partial Gramians. Thus $H_g=O_gC_g$ is bounded and has rank at most $k$, and
$E:=H_f-H_g=O_EC_E$.

For bounded operators factoring through an $m$-dimensional space,
\begin{equation}\label{I:U:eq:gram-norm}
 \Unorm E^2=\lambda_{\max}(Q_E^{1/2}P_EQ_E^{1/2}).
\end{equation}
Indeed, $E^*E=C_E^*Q_EC_E$ has the same nonzero eigenvalues as $Q_E^{1/2}C_EC_E^*Q_E^{1/2}$; the equality also holds when both operators are zero. Congruence preserves the positive-semidefinite order, so
\begin{align*}
 \Unorm E^2
 &\le\lambda_{\max}(Q_E^{1/2}P Q_E^{1/2})\\
 &=\lambda_{\max}(W^*Q_EW)
 \le\lambda_{\max}(W^*QW)\le t^2.
\end{align*}
The middle equality follows from $P=WW^*$ and equality of the nonzero spectra of $ZZ^*$ and $Z^*Z$, with $Z=Q_E^{1/2}W$. No inverse or positive definiteness is used.

Conversely, Theorem~\ref{I:U:thm:attainment} supplies an optimal bounded rank-$\le k$ Hankel. Choose the canonical realization of Lemma~\ref{I:U:lem:compact}, padded with zeros, and form \eqref{I:U:eq:error-realization}. The source maps are bounded by its mean-square stability; the candidate maps are bounded by the lemma. Consequently
\[
 C_E=\begin{pmatrix}C_f\\ C_g\end{pmatrix},\qquad
 O_E=\begin{pmatrix}O_f&-O_g\end{pmatrix}
\]
are bounded. Their finite-dimensional Gramians $P_E,Q_E$ satisfy the exact Stein equations obtained by splitting words at the first and last letter, respectively. Set $P=P_E$, $Q=Q_E$ and choose $W=P_E^{1/2}$. Equation~\eqref{I:U:eq:gram-norm} then implies
$W^*QW\preceq\delta_k(H_f)^2\UId_m\preceq t^2\UId_m$.
This proves necessity and feasibility at the optimum.
\end{proof}


\begin{corollary}[Algebraic optimization and an existential formulation]\label{I:U:cor:algebraic}\label{I:S:thm:etr}
For a finite mean-square-stable source with algebraic coefficients, $\delta_k(H_f)$ is a nonnegative real algebraic number, and an optimal candidate exists with algebraic coefficients. Both can be computed by real quantifier elimination and algebraic sample-point extraction.
For rational source data and a rational threshold, the decision problem $\delta_k(H_f)\le t$ has a polynomial-size existential formula over the reals, of degree at most three, using
\[
 O\bigl(dk^2+(n+k)^2\bigr)
\]
real scalar variables. In particular, the rational-input decision problem belongs to $\exists\mathbb R$.
\end{corollary}
Here $\exists\mathbb R$ denotes the class of decision problems admitting
a polynomial-time many-one reduction to existential feasibility of
polynomial equalities and inequalities over the reals; see
\citet{schaefer2024}. Membership is an upper bound, not a hardness claim.
\begin{proof}
Replace every matrix inequality $Y\succeq0$ in Theorem~\ref{I:U:thm:exact} by an equation $Y=UU^*$ with a fresh square matrix $U$ of the appropriate size; Hermitian variables can be represented by their independent real coordinates. The condition $P\succeq0$ is already implied by $P=WW^*$. For complex data, separate real and imaginary parts. All resulting scalar polynomial equations and inequalities have degree at most three. The cubic terms are $B_aPB_a^*$, $B_a^*QB_a$, and $W^*QW$; all PSD factorizations themselves are quadratic. The number of matrices of order $m$ is constant, yielding the stated variable count. Expanding the sums and products produces polynomially many terms in $d,n,k$ and in the binary size of rational input coefficients.

By Theorem~\ref{I:U:thm:exact}, the projection of this feasible set onto $t$ is exactly the closed ray $[\delta_k(H_f),\infty)$. Quantifier elimination over the real algebraic numbers gives a univariate semialgebraic description of this ray, whose finite endpoint is algebraic. Applying sample-point extraction at that endpoint produces an algebraic feasible tuple and hence an optimal algebraic realization. These are standard effective operations over real closed fields \cite{basu2006,renegar1992}.

For rational input the threshold formula is an existential real formula of polynomial size, which proves membership in $\exists\mathbb R$. The stipulated source stability can itself be checked in polynomial bit time: solve $Y-\sum_aA_aYA_a^*=\UId_n$ over the rationals and test $Y\succ0$. A singular linear system is rejected; a positive-definite solution implies stability by conjugating the positive map with $Y^{1/2}$, and stability gives a unique such solution. Thus invalid source encodings can be rejected before the reduction. If $k$ is supplied in binary, the case $k\ge n$ is trivial for $t\ge0$, because the source itself is admissible. Otherwise $k<n$, so the construction is polynomial in the input size. Negative thresholds are rejected before constructing the formula. The case $k=0$ has the sole candidate $G=0$ and is handled by the same Gramian norm test without candidate variables.
\end{proof}


\begin{remark}\label{I:U:rem:exact-scope}
The existential formulation is an exact decision and synthesis description, not a claim of polynomial-time solvability when all parameters vary. It also does not assert equality between the real and complex optima: the chosen scalar field is fixed throughout each instance. Algebraic output describes an exact optimizer; rounding it to a prescribed finite precision requires a separate error analysis. No finite-horizon or numerical optimization oracle is assumed in the certificate.
\end{remark}

\section{Polynomial-time decision for fixed alphabet and target rank}
\label{I:U:sec:complexity}

The cubic certificate of Theorem~\ref{I:U:thm:exact} uses a number of
variables that grows with the source dimension.  For a fixed alphabet and
target rank, one can eliminate the Gramian variables before applying real
algebraic decision algorithms.  The resulting formulation has a fixed number
of variables and polynomially growing degrees.  Throughout this section the
source, threshold, and requested accuracy are rational, and competitors are
real.  The complexity model counts bit operations on binary encodings of
rational numbers.

\begin{theorem}[Fixed-parameter slices]\label{I:U:thm:fixed}
Fix integers $d\ge1$ and $k\ge0$.  Given a rational, mean-square-stable
$n$-state realization of $f$ and a rational threshold $t$, deciding whether
$\delta_k(H_f)\le t$ takes deterministic time polynomial in the binary
input length.  The exponent may depend on $d$ and $k$.
\end{theorem}

The theorem concerns polynomial time on each slice with $d,k$ fixed.
The construction below gives an XP algorithm with respect to $(d,k)$;
it does not give a fixed-parameter tractable algorithm or a running time
polynomial jointly in $d,k,n$. With $p=dk^2+2k+k(k+1)/2$ candidate variables, the semialgebraic decision step has exponent $O(p)=O(dk^2)$ in the polynomial-size input parameters; constants and prefactors may depend on $(d,k)$.

\begin{lemma}[A real-rooted polynomial test]\label{I:U:lem:coefftest}
Suppose a real $m\times m$ matrix $R$ has only real eigenvalues.
All of its eigenvalues are nonnegative if and only if all nonleading
coefficients of $\det(z\UId_m+R)$ are nonnegative.
\end{lemma}
\begin{proof}
Write the eigenvalues, with algebraic multiplicity, as
$\mu_1,\ldots,\mu_m$.  If each $\mu_i\ge0$, the coefficients of
$\prod_i(z+\mu_i)$ are nonnegative.  Conversely, if some $\mu_i<0$,
this polynomial has the positive root $-\mu_i$.  A monic polynomial
whose remaining coefficients are nonnegative is strictly positive at
every positive real number.  This is a contradiction.  Diagonalizability
of $R$ is not needed.
\end{proof}


\begin{proof}[Proof of Theorem~\ref{I:U:thm:fixed}]
Negative thresholds are immediately rejected.  We give the proof for
$k\ge1$; for $k=0$ the sole competitor is zero and the same calculation
below has no candidate variables.

\paragraph{A fixed number of candidate variables.}
By Theorem~\ref{I:U:thm:attainment} and Lemma~\ref{I:U:lem:stable}, the inequality
$\delta_k(H_f)\le t$ is equivalent to the existence of a stable
$k$-state realization $g(w)=c^\top T_wb$ with error at most $t$.
A smaller stable realization can be padded with zero states.
Introduce the entries of $T_1,\ldots,T_d,b,c$, and a symmetric
$k\times k$ matrix $X$ as variables.  Impose
\begin{equation}\label{I:U:eq:fixed-stability}
 X\succ0,\qquad X-\sum_{a\in\Sigma}T_aXT_a^\top=\UId_k.
\end{equation}
These conditions are equivalent to mean-square stability of the tuple
$T$.  For the forward implication, positivity and
$\UId_k\succeq \Unorm{X}^{-1}X$ give a strict order contraction in the
order-unit norm determined by $X$; the spectral radius is therefore less
than one.  For the reverse implication, take
$X=\sum_{j\ge0}\Phi_T^j(\UId_k)$.
Positive definiteness in~\eqref{I:U:eq:fixed-stability} is expressed by
positivity of the $k$ leading principal minors of $X$.
The number of real variables is
\begin{equation}\label{I:U:eq:fixed-varcount}
 p=dk^2+2k+\frac{k(k+1)}2,
\end{equation}
which is independent of $n$.

\paragraph{Eliminating the joint Gramians.}
Let $m=n+k$, and form
\[
 B_a=\begin{pmatrix}A_a&0\\0&T_a\end{pmatrix},\qquad
 \beta=\begin{pmatrix}b_f\\b\end{pmatrix},\qquad
 \gamma=\begin{pmatrix}c_f\\-c\end{pmatrix}.
\]
The joint tuple $B$ is stable: a positive definite Stein solution for
the source and the matrix $X$ in~\eqref{I:U:eq:fixed-stability}, placed on
the diagonal, give a positive definite Stein solution for $B$.
Hence its exact reachability and observability Gramians are uniquely
defined by
\[
 P=\beta\beta^\top+\sum_aB_aPB_a^\top,\qquad
 Q=\gamma\gamma^\top+\sum_aB_a^\top QB_a.
\]
These are the usual finite-dimensional Gramian equations for weighted
automata; see also~\cite{sva2019}.
Using column vectorization, put
\begin{align}
 \mathcal L&=\UId_{m^2}-\sum_a B_a\otimes B_a,
 &\Delta&=\det\mathcal L,\label{I:U:eq:fixed-L}\\
 \widehat P
 &=\operatorname{unvec}\!\left(
       \operatorname{adj}(\mathcal L)\operatorname{vec}(\beta\beta^\top)
     \right),&&\nonumber\\
 \widehat Q
 &=\operatorname{unvec}\!\left(
       \operatorname{adj}(\mathcal L^\top)\operatorname{vec}(\gamma\gamma^\top)
     \right).&&
 \label{I:U:eq:fixed-adj}
\end{align}
Stability guarantees $\Delta\ne0$ throughout the feasible set, and
\begin{equation}\label{I:U:eq:fixed-rational-grams}
 P=\widehat P/\Delta,\qquad Q=\widehat Q/\Delta.
\end{equation}
Thus the Gramians are rational functions of the fixed collection of
variables; they need not be introduced as additional variables.

\paragraph{A polynomial norm constraint without square roots.}
The squared norm of the error Hankel is
$\lambda_{\max}(P^{1/2}QP^{1/2})$.
The matrices $PQ$ and $P^{1/2}QP^{1/2}$ have the same characteristic
polynomial: apply
$\det(z\UId-UV)=\det(z\UId-VU)$ with
$U=P^{1/2}$ and $V=P^{1/2}Q$.
In particular, all eigenvalues of $PQ$ are real and nonnegative,
including when $P$ is singular.
Define the polynomial matrix
\begin{equation}\label{I:U:eq:fixed-R}
 R=t^2\Delta^2\UId_m-\widehat P\widehat Q
   =\Delta^2(t^2\UId_m-PQ).
\end{equation}
Its eigenvalues are real.  Since $\Delta^2>0$, the error norm is at
most $t$ if and only if these eigenvalues are nonnegative.  By
Lemma~\ref{I:U:lem:coefftest}, this is equivalent to
\begin{equation}\label{I:U:eq:fixed-coeff}
 \det(z\UId_m+R)=z^m+a_1z^{m-1}+\cdots+a_m,
 \qquad a_1\ge0,\ldots,a_m\ge0.
\end{equation}
Each $a_i$ is a polynomial in the variables counted
in~\eqref{I:U:eq:fixed-varcount}.  Squaring the common denominator avoids
any assumption about its sign.  No matrix square root, eigenvalue
variable, or list of exponentially many principal minors is used.

\paragraph{Degrees, coefficient sizes, and construction time.}
Each entry of the $m^2\times m^2$ matrix $\mathcal L$ has degree at most two.
Consequently
\[
 \deg\Delta\le2m^2,\qquad
 \deg\widehat P,\deg\widehat Q\le2m^2,\qquad
 \deg a_i\le4m^3,
\]
where $t$ is the fixed rational input threshold.
Together,~\eqref{I:U:eq:fixed-stability} and~\eqref{I:U:eq:fixed-coeff}
give $O(k^2+m)$ polynomial equalities and inequalities of degree
$O(m^3)$ in $p$ variables.

Their coefficients have bit lengths polynomial in $m$ and the binary
input length.  To see this, clear the input denominators and use the
determinant expansion: each monomial coefficient is a sum of products
of polynomially many input coefficients, and the logarithm of the
number of summands is polynomial in $m$.  The same bound applies to
adjugate entries, their products, and the final characteristic
polynomial.  This argument bounds bit length; it does not call for
enumerating the determinant expansion.

For completeness, the polynomials themselves can be constructed in
polynomial bit time when $p$ is fixed.  Their degree bounds imply only
$m^{O(p)}$ possible monomials.  Evaluate determinants, adjugates, and
characteristic polynomials over the rationals on an integer
interpolation grid of polynomial side length, and perform exact
multivariate interpolation.  The grid contains polynomially many
points for fixed $p$, and every evaluation has polynomial bit
complexity.  Adjugates are evaluated directly by cofactors or a
division-free matrix algorithm, so singular interpolation points cause
no difficulty.  Their intermediate and output bit lengths are
polynomial by the same determinant bounds.  The interpolation step
also has polynomial bit complexity.

\paragraph{Decision.}
Real algebraic feasibility with a fixed number of variables is
decidable in time polynomial in the number of polynomials, their
degrees, and their coefficient bit lengths
\cite{renegar1992,basu2006}.  Apply this result to the formula above.
The source dimension enters through polynomial degree and coefficient
size, while the number of variables is the fixed number $p$.
This proves the claimed bit-complexity bound.
\end{proof}


\begin{corollary}[Additive synthesis on a fixed slice]\label{I:U:cor:fixed-synthesis}
For fixed $d,k$, a rational stable source and a positive rational
$\varepsilon$ admit an algorithm that returns real algebraic
coefficients for a stable realization $g$ of dimension at most $k$
such that
\[
 \delta_k(H_f)\le\Unorm{H_f-H_g}
 \le\delta_k(H_f)+\varepsilon.
\]
Its bit complexity is polynomial in the binary encodings of the
source and $\varepsilon$, with an exponent depending on $d,k$.
\end{corollary}
\begin{proof}
Compute the rational source Gramians $P_f,Q_f$ by their Stein
equations.  The rational number
$F=1+\operatorname{tr}(P_fQ_f)$ has polynomial encoding length and
satisfies
$\Unorm{H_f}\le\Unorm{H_f}_{\mathrm{HS}}
=\sqrt{\operatorname{tr}(P_fQ_f)}\le F$.
The zero competitor is feasible at threshold $F$.
Bisect $[0,F]$ using Theorem~\ref{I:U:thm:fixed} until the interval width
is at most $\varepsilon$, retaining a feasible upper endpoint $u$.
The number of decisions is
$O(1+\log_+(F/\varepsilon))$, and all intermediate thresholds have
polynomial encoding length.
Standard real algebraic sample-point extraction
\cite{basu2006} applied to the feasible formula at $u$ returns
algebraic $T_a,b,c,X$ in polynomial time for fixed $p$.
They define a stable candidate with error at most
$u\le\delta_k(H_f)+\varepsilon$.
\end{proof}


\begin{remark}[Output model and scope]
The corollary returns exact real algebraic data, for example through
isolating intervals and algebraic sample-point encodings.  It does not
assert that independently rounding the coefficients to floating point
preserves stability or the error certificate.  Nor does it make a
practical running-time claim for the resulting algebraic algorithms.
For Gaussian-rational sources and complex competitors, the same
fixed-slice argument applies after writing real and imaginary parts
as separate variables and expressing the two Stein equations in real
coordinates on Hermitian matrices.  Their denominators can be
combined and squared in the norm test.  The number of candidate
variables still depends only on $d,k$, and all degree and bit-size
bounds remain polynomial in the source input length.
\end{remark}

\section{Sharper bounds for the binary obstruction family}\label{I:M:rev13:sec:sep}

Theorem~\ref{I:U:thm:separation} gives, for the source
$f_{r,k}(w)=q\,r^j$ on $w=a^jb^k$ ($q=1-r^2$),
\begin{equation}\label{I:M:rev13:eq:old-bound}
\delta_k(H_{f_{r,k}})\ge\max\Bigl\{\sqrt q,\ \sqrt{\tfrac{1+q^2(C_k-1)}{C_k}}\Bigr\},\qquad
C_k=\tbinom{2k}k,
\end{equation}
so that $\liminf_{r\uparrow1}\delta_k\ge C_k^{-1/2}$. We prove a bound that is uniform in
$r$ and much larger, and for $k=1$ we identify the limit. Throughout,
$z_j=\sqrt q\sum_{n\ge0}r^ne_{a^nb^j}$ ($0\le j\le k$) are the orthonormal vectors of the proof of Theorem~\ref{I:U:thm:separation} and
\[
H_f=z_0z_k^{*}+\sqrt q\sum_{j=0}^{k-1}z_{k-j}e_{b^j}^{*},\qquad
H_f^{*}z_0=z_k ,
\]
the last identity because $z_{k-j}\perp z_0$ for $j<k$.

\subsection{A model-space constant}
For a monic polynomial $\chi$ of degree $k$ with all roots in the open unit disk, let
$V_\chi\subset\ell^2(\MN)$ be the $k$-dimensional space of $\chi$-recurrent square-summable
sequences ($\sum_i c_iu_{j+i}=0$ for all $j\ge0$, $\chi=\sum_ic_iz^i$). Define
\begin{equation}\label{I:M:rev13:eq:kappa}
\kappa_k=\sup_\chi\ \sup_{0\ne u\in V_\chi}\frac{|u_k|^2}{\sum_{j\ge0}|u_j|^2}
=\sup_\chi\Mnorm{P_{V_\chi}e_k}^2 .
\end{equation}
\begin{lemma}[Blaschke form of $\kappa_k$]\label{I:M:rev13:lem:kappa}
Under $\ell^2(\MN)\cong H^2$, $V_\chi$ is the model space $K_\Theta=H^2\ominus\Theta H^2$ of
the degree-$k$ Blaschke product $\Theta$ with zeros at the conjugates of the roots of
$\chi$, and
\[
\Mnorm{P_{K_\Theta}z^k}^2=1-\sum_{m=0}^{k}|\theta_m|^2,\qquad
\Theta(z)=\sum_{m\ge0}\theta_mz^m .
\]
Hence $\kappa_k=1-\inf_\Theta\sum_{m\le k}|\theta_m|^2$ over degree-$k$ Blaschke products.
Moreover $\kappa_k\le1-1/C_k$, $\kappa_1=1/4$, and $\kappa_2=7/16$.
\end{lemma}
\begin{proof}
A $\chi$-recurrent sequence is $u_j=\sum_i p_i(j)\lambda_i^j$ over the roots $\lambda_i$
with polynomial coefficients; under $u\mapsto\sum_ju_jz^j$ the sequence $(\lambda^j)_j$
becomes $1/(1-\lambda z)=k_{\bar\lambda}$, the reproducing kernel at $\bar\lambda$, and
derivatives for multiple roots; these kernels span $K_\Theta$. Since
$P_{K_\Theta}f=f-\Theta P_+(\bar\Theta f)$ and multiplication by the inner function $\Theta$
is isometric, $\Mnorm{P_{K_\Theta}z^k}^2=1-\Mnorm{P_+(\bar\Theta z^k)}^2$, and
$\langle\bar\Theta z^k,z^j\rangle=\langle z^k,\Theta z^j\rangle=\overline{\theta_{k-j}}$
for $0\le j\le k$, zero for $j>k$. For the crude bound, the coefficients of $\chi$ satisfy
$|c_i|\le\binom ki$, so $|u_k|^2\le(\sum_{i<k}|c_i|^2)\sum_{j<k}|u_j|^2\le(C_k-1)
(\Mnorm u^2-|u_k|^2)$. For $k=1$, $\Theta=B_\lambda$ gives $\theta_0=-\lambda$,
$\theta_1=1-|\lambda|^2$ and $1-|\theta_0|^2-|\theta_1|^2=|\lambda|^2(1-|\lambda|^2)\le1/4$,
attained at $|\lambda|^2=1/2$. For $k=2$ apply two steps of the Schur algorithm to $\Theta$.
Writing $\alpha=\Theta(0)$, there are $\beta\in\mathbb D$ and $|\gamma|=1$
such that
\[
 \Theta(z)=\frac{\alpha+zF(z)}{1+\bar\alpha zF(z)},\qquad
 F(z)=\frac{\beta+z\gamma}{1+\bar\beta z\gamma}.
\]
Indeed $(\Theta-\alpha)/(1-\bar\alpha\Theta)$ is a degree-two inner
rational function vanishing at zero; dividing by $z$ leaves degree one.
Repeating this step leaves a unimodular constant. This is the elementary
finite Schur algorithm; see \citet{schur_nevanlinna}.
The first coefficients are
\[
 \theta_0=\alpha,\quad \theta_1=(1-|\alpha|^2)\beta,\quad
 \theta_2=(1-|\alpha|^2)\bigl((1-|\beta|^2)\gamma-\bar\alpha\beta^2\bigr).
\]
Set $a=|\alpha|$ and $t=|\beta|^2$. By the reverse triangle inequality,
\begin{align*}
 \sum_{j=0}^2|\theta_j|^2
 &\ge a^2+(1-a^2)^2\bigl[t+(1-(1+a)t)^2\bigr]\\
 &=\frac9{16}+(a-\tfrac12)^2(a+\tfrac34)\\
 &\quad +(1-a^2)^2(1+a)^2
 \left(t-\frac{1+2a}{2(1+a)^2}\right)^2
 \ \ge\ \frac9{16}.
\end{align*}
Equality holds for $\alpha=1/2$, $\beta=2/3$, $\gamma=1$, giving
\[
 \Theta(z)=\frac{z^2+z+1/2}{1+z+z^2/2},\qquad
 (\theta_0,\theta_1,\theta_2)=(1/2,1/2,1/4).
\]
Its zeros $(-1\pm i)/2$ lie in $\mathbb D$, so this is an admissible
degree-two Blaschke product and $\kappa_2=7/16$. No global numerical
minimisation is needed.
\end{proof}


\subsection{A uniform lower bound}
\begin{theorem}[Uniform obstruction constant]\label{I:M:rev13:thm:uniform}
For every $k\ge1$, every $0<r<1$ and either scalar field,
\[
\delta_k(H_{f_{r,k}})\ \ge\ \sqrt{1-\kappa_k}=E_k
\quad\bigl(\ge C_k^{-1/2}\bigr).
\]
Here $E_k$ is the complex scalar distance of
Lemma~\ref{I:B:lem:one-letter-projection}; that lemma gives
$\kappa_k=\Gamma_k^2=1-E_k^2$ without the outer-extremal input.
In particular $\inf_{0<r<1}\delta_1\ge\sqrt3/2$ and $\inf_{0<r<1}\delta_2\ge3/4$, while
$\sigma_{k+1}(H_{f_{r,k}})=\sqrt q$ ranges over $(0,1)$.
\end{theorem}
\begin{proof}
Let $G$ be a bounded word-Hankel operator of rank at most $k$ and $A=H_f-G$. Then
$\Mnorm A=\Mnorm{A^*}\ge\Mnorm{A^*z_0}=\Mnorm{z_k-G^*z_0}\ge\Mdist(z_k,\Mrange G^*)$.
By Lemma~\ref{I:U:lem:compact}, $G=OC$ with $O$ an
isometry of $\MFF^s$ ($s=\Mrank G\le k$), $C=O^*G$, $Ce_v=T_vb$ and $\Mnorm{T_a}\le1$; hence
$\overline{\Mrange G^*}=\overline{\Mrange C^*}=\{h_\eta:\eta\in\MFF^s\}$ with $h_\eta(v)=\eta^*T_vb$. Complex conjugation preserves norms and fixes the real vector $z_k$,
so it does not affect the desired projection estimate. Fix $\eta$ with
$h_\eta\ne0$. For each $n\ge0$ the sequence $u^{(n)}_j=\eta^*T_a^nT_b^jb$, $j\ge0$,
satisfies the recurrence of the characteristic polynomial of $T_b$ (Cayley--Hamilton),
whose roots lie in the closed unit disk, and is square-summable because
$\sum_j|u^{(n)}_j|^2\le\Mnorm{h_\eta}^2<\infty$. A square-summable sequence recurrent for a
polynomial with unimodular roots has no component on those roots, so it is recurrent for a
degree-$k$ polynomial with roots in the open disk (replace the unimodular roots by $0$).
Padding with zero states if $s<k$ changes nothing. Therefore
$|u^{(n)}_k|^2\le\kappa_k\sum_j|u^{(n)}_j|^2$ for every $n$, and
\[
|\langle z_k,h_\eta\rangle|^2\le\sum_n|h_\eta(a^nb^k)|^2
=\sum_n|u^{(n)}_k|^2\le\kappa_k\sum_n\sum_j|h_\eta(a^nb^j)|^2\le\kappa_k\Mnorm{h_\eta}^2,
\]
using Cauchy--Schwarz on the support of the unit vector $z_k$ and the fact that the words
$a^nb^j$ are distinct. Thus $\Mnorm{P_{\Mrange G^*}z_k}^2\le\kappa_k$ and
$\Mdist(z_k,\Mrange G^*)^2\ge1-\kappa_k$. The crude bound of Lemma~\ref{I:M:rev13:lem:kappa} gives the
parenthetical comparison.
\end{proof}


\begin{remark}
The proof of Theorem~\ref{I:U:thm:separation} bounds the Rayleigh quotient of $H_f$ on one kernel vector of $G$;
Theorem~\ref{I:M:rev13:thm:uniform} instead bounds how much of the top right singular vector $z_k$
can lie in the $k$-dimensional backward-shift-invariant row space of $G$, using only the
words $a^nb^j$. The resulting constant does not depend on $r$, so the obstruction holds at
every $r$. The diverging ratio to $\sigma_{k+1}=\sqrt q$ occurs as $r\uparrow1$, when the source stability margin closes; a margin-dependent positive bound is not excluded.\end{remark}

\subsection{The rank-one case: the obstruction is asymptotically total}
\begin{lemma}[Rank-one word-Hankel operators]\label{I:M:rev13:lem:rank-one}\label{I:S:lem:rankone}
Over $\MFF\in\{\MR,\mathbb C\}$, every nonzero bounded word-Hankel operator of rank one is
$G=\gamma\,mm^{\top}$ (transpose, also over $\mathbb C$), where
$m_w=\prod_{i}t_{w_i}$ for scalars $t_a\in\MFF$ with $\sum_a|t_a|^2<1$ and
$\gamma\in\MFF\setminus\{0\}$; conversely each such $G$ is a bounded rank-one word-Hankel.
\end{lemma}
\begin{proof}
Write $G=xy^*$ and $g(uv)=G[u,v]=x_u\bar y_v$. Then $g(u)=x_u\bar y_\varepsilon$,
$g(v)=x_\varepsilon\bar y_v$ and $g(\varepsilon)=x_\varepsilon\bar y_\varepsilon$. If
$g(\varepsilon)=0$ then $x_\varepsilon=0$ or $y_\varepsilon=0$, forcing $g\equiv0$.
Otherwise $g(uv)=g(u)g(v)/g(\varepsilon)$, so $w\mapsto g(w)/g(\varepsilon)$ is
multiplicative: $g(w)=\gamma\prod_it_{w_i}$ with $\gamma=g(\varepsilon)$,
$t_a=g(a)/g(\varepsilon)$, and $G[u,v]=\gamma m_um_v$. Finally
$\Mnorm G=|\gamma|\Mnorm m^2$ and $\Mnorm m^2=\sum_n(\sum_a|t_a|^2)^n$.
\end{proof}


\begin{theorem}[Sharpened rank-one obstruction]\label{I:M:rev13:thm:rank-one-sharp}
For $k=1$, $0<r<1$, $q=1-r^2$, over either scalar field,
\begin{equation}\label{I:M:rev13:eq:rank-one-sharp}
\delta_1(H_{f_{r,1}})\ \ge\ \frac{\sqrt{1+2\sqrt q}}{1+\sqrt q}\ =\ 1-\frac q{2}+O(q^{3/2}).
\end{equation}
Consequently $\lim_{r\uparrow1}\delta_1(H_{f_{r,1}})=1=\Mnorm{H_{f_{r,1}}}$ while
$\sigma_2(H_{f_{r,1}})=\sqrt q\to0$: as the stability margin closes, no rank-one
word-Hankel operator approximates the source better than the zero operator, up to
$O(q)$. The right-hand side of \eqref{I:M:rev13:eq:rank-one-sharp} is at least $\sqrt3/2$ for all
$q\in(0,1]$, so it dominates Theorem~\ref{I:M:rev13:thm:uniform} at $k=1$, and it exceeds
\eqref{I:M:rev13:eq:old-bound} for $q<0.76$ approximately.
\end{theorem}
\begin{proof}
Let $G=\gamma mm^{\top}$ be any bounded rank-one competitor (Lemma~\ref{I:M:rev13:lem:rank-one};
$G=0$ gives error $\Mnorm{H_f}=1$) and $A=H_f-G$, with
$H_f=z_0z_1^{\top}+\sqrt q\,z_1e_\varepsilon^{\top}$. Put
\[
p=m^{\top}z_0=\sqrt q\sum_nr^nt_a^n=\frac{\sqrt q}{1-rt_a},\qquad
s=m^{\top}z_1=\sqrt q\sum_nr^nt_a^nt_b=t_bp .
\]
Then $Az_1=z_0-\gamma s\,m$ and, since $H_f^{\top}z_0=z_1$ and $G^{\top}=G$,
$A^{\top}z_0=z_1-\gamma p\,m$. With $c=\gamma sp$,
\[
\Mnorm{Az_1}^2=1-2\operatorname{Re}c+|c|^2\frac{\Mnorm m^2}{|p|^2},\qquad
\Mnorm{A^{\top}z_0}^2=1-2\operatorname{Re}c+|c|^2\frac{\Mnorm m^2}{|s|^2}.
\]
Because $|s|=|t_b||p|\le|p|$, the second expression dominates, so
$\Mnorm A^2\ge1-2|c|+|c|^2\Mnorm m^2/|s|^2\ge1-|s|^2/\Mnorm m^2$ after minimising over
$|c|\ge0$ (if $t_b=0$ then $s=0$ and $\Mnorm{Az_1}=\Mnorm{z_0}=1$). Now
\[
\frac{|s|^2}{\Mnorm m^2}=\frac{q\,|t_b|^2(1-|t_a|^2-|t_b|^2)}{|1-rt_a|^2}
\le\frac{q\,\beta(1-\alpha^2-\beta)}{(1-r\alpha)^2}\le\frac q4\,\frac{(1-\alpha^2)^2}{(1-r\alpha)^2},
\]
with $\alpha=|t_a|$, $\beta=|t_b|^2\le1-\alpha^2$. The function
$\varphi(\alpha)=(1-\alpha^2)/(1-r\alpha)$ on $[0,1)$ has derivative proportional to
$r\alpha^2-2\alpha+r$ and attains its maximum $2/(1+\sqrt q)$ at
$\alpha=(1-\sqrt q)/r$ (where $1-\alpha^2=2\sqrt q/(1+\sqrt q)$ and $1-r\alpha=\sqrt q$).
Hence $|s|^2/\Mnorm m^2\le q/(1+\sqrt q)^2$ and
$\delta_1^2\ge1-q/(1+\sqrt q)^2=(1+2\sqrt q)/(1+\sqrt q)^2$. The expansion and the limit
are immediate, and $\delta_1\le\Mnorm{H_f}=1$ because $G=0$ is admissible.
\end{proof}




\begin{remark}The rank-two boundary limit is proved in Theorem~\ref{I:B:thm:main}, and its first correction is proved in Theorem~\ref{I:R:thm:main-general}.\end{remark}

\section{Exact behavioral realization and the marked algebra}\label{I:S:sec:exact}
Fix a finite nonempty alphabet $\Sigma$, $d=|\Sigma|$, and the free monoid $\Swords$ with empty word $\Seps$. A linear weighted finite automaton (WFA) over a field $\SK$ is
\[
 f(w)=\alpha^\top A_w\beta,\qquad A_{a_1\cdots a_m}=A_{a_1}\cdots A_{a_m},\quad A_\Seps=I_n.
\]
We use row states $s(w)=\alpha^\top A_w$. Here \emph{row-span minimality} means that the $s(w)$ span $\SK^{1\times n}$; \emph{column-span minimality} means that $\{A_w\beta\}$ spans $\SK^n$. A realization is minimal when both hold. In complex analytic arguments $\alpha^\top$ remains the algebraic output row; Hermitian adjoints in Gramians are written separately.

\begin{theorem}[Hankel realization and exact finite witnesses]\label{I:S:thm:realization}
The algebraic Hankel matrix $H_f[u,v]=f(uv)$ has finite rank $n$ if and only if $f$ has a finite linear realization. Its rank is the minimum linear state dimension. Two realizations of dimensions $n_1,n_2$ agree on all words if they agree on words of length at most $n_1+n_2-1$.
\end{theorem}
\begin{proof}
A realization gives $H_f=OC$, with $O[u,:]=\alpha^\top A_u$ and $C[:,v]=A_v\beta$, so its rank is at most $n$. Conversely take the span of residual rows $r_u(v)=f(uv)$. The shift $r_u\mapsto r_{ua}$ is well-defined: if $\sum c_u r_u=0$, evaluation at $av$ gives $\sum c_u f(uav)=0$ for all $v$. The initial residual is $r_\Seps$ and the output is evaluation at $\Seps$. A basis supplies a realization of dimension $\Srank H_f$.

For a realization of dimension $N$, let $V_j=\SSpan\{\alpha^\top A_w:|w|\le j\}$. Once $V_j=V_{j+1}$, it is invariant and contains all reachable rows. If the initial row is nonzero, $\dim V_0=1$, so stabilization occurs by $j=N-1$; the zero case is immediate. Vanishing output on words of length at most $N-1$ therefore implies vanishing on all words. Apply this to the direct-sum difference realization of dimension $n_1+n_2$.
\end{proof}


An explicit minimization first restricts to the generated row space and then quotients by the rows annihilating every future column. Exact basis closure terminates in finite dimension. The same argument applies to finitely many output coordinates: the future tests include both a word and an output coordinate.

\begin{proposition}[Discrete and affine memory have different costs]\label{I:S:prop:types}
The minimum number of deterministic states realizing an observable $f$ is the number of residual classes for $u\sim v$ iff $f(ut)=f(vt)$ for every $t$. It can be infinite while $\Srank H_f$ is finite. If a vector observable has residual affine-hull dimension $r_{\rm aff}$, the minimum dimension of an affine-state realization with affine readout is $r_{\rm aff}$, and
\[
 \Srank H_f-1\le r_{\rm aff}\le\Srank H_f.
\]
\end{proposition}
\begin{proof}
Equal deterministic states must give equal future outputs; the residual classes themselves form an attaining deterministic automaton because the relation is a right congruence. For $f(a^m)=m$ in characteristic zero, residuals $t\mapsto m+|t|$ are all distinct but their linear span has dimension two. For the affine statement, every future observable is affine in an affine state, so the residual affine dimension is a lower bound. Shift maps preserve affine relations by the same evaluation argument as in Theorem~\ref{I:S:thm:realization}; their restriction to the residual affine hull attains the bound. An affine hull of dimension $r$ has linear span of dimension $r$ or $r+1$, proving the inequalities.
\end{proof}


Define the syntactic ideal and the generated algebra by
\[
 I_f=\{p\in\SK\langle\Sigma\rangle:f(upv)=0\ \forall u,v\},\quad
 \ScA_f=\SK\langle\Sigma\rangle/I_f,\quad B=\Salg(I,(A_a)_a).
\]
The letters, unit, and tested observable are retained as markings. For a finite-dimensional algebra, let $J(B)$ denote its Jacobson radical and $\SLL(B)=\min\{j\ge1:J(B)^j=0\}$, with $\SLL(0)=0$.

\begin{theorem}[Universal marked quotient and Loewy dominance]\label{I:S:thm:marked}
Every finite realization induces a letter-marked surjection $B\twoheadrightarrow\ScA_f$. It is an isomorphism for a minimal realization. For any finite-dimensional algebra epimorphism $\pi:B\twoheadrightarrow A$,
\[
 \pi(J(B)^j)=J(A)^j\quad(j\ge0).
\]
Thus dimensions of the algebra and of every graded radical layer, and Loewy length, cannot increase under the quotient.
\end{theorem}
\begin{proof}
The representation kernel is contained in $I_f$: a zero matrix polynomial remains zero between every prefix row and future column. This yields the quotient. Under both span conditions, $\alpha^\top A_u p(A)A_v\beta=0$ for all $u,v$ implies $p(A)=0$, giving kernel equality. The image of $J(B)$ is a nilpotent ideal, hence contained in $J(A)$. Also $A/\pi(J(B))$ is a quotient of the semisimple algebra $B/J(B)$ and is semisimple. Its radical vanishes, giving the reverse inclusion. Taking products proves the power identities; passage to successive quotients gives the graded surjections.
\end{proof}


\begin{remark}[The charged algebra matters]
This theorem concerns a quotient of the algebra actually generated by the transitions. It is false with an arbitrary containing algebra in its place: $\SK[t]/(t^2)$ embeds in $M_2(\SK)$, which is semisimple. Loewy length is also not neural depth. A single recurrent layer with an $(m+1)$-state nilpotent shift realizes a length-$m$ impulse. For a unary matrix over $\SC$, the Chinese remainder decomposition of $\SC[t]/(\mu_A)$ shows that the Loewy length of its generated algebra is the largest Jordan-block size, not the number of network layers.
\end{remark}

\section{Bounded Hankels, strict stability, and exact norms}\label{I:S:sec:analytic}
Now $\ScH=\ell^2(\Swords)$ has its unweighted orthonormal word basis. The analytic theorems are stated over $\SR$; the Hermitian versions hold over $\SC$. Define
\[
 \Phi_A(X)=\sum_a A_a^\top X A_a,\qquad
 \Psi_A(X)=\sum_a A_a X A_a^\top
\]
on real symmetric matrices. They are adjoint for the trace pairing. The observability and controllability feature maps, in their operator-factorization convention, are
\[
 (Ox)(u)=\alpha^\top A_u x,\qquad Cz=\sum_v A_v\beta z(v).
\]

\begin{theorem}[Minimal boundedness is strict all-word stability]\label{I:S:thm:stability}
For a minimal finite real realization the following are equivalent:
\[
 H_f\text{ bounded};\quad f\in\ell^2(\Swords);\quad
 \rho(\Phi_A)<1;\quad
 \exists X\succ0:\ X-\Phi_A(X)\succ0.
\]
The Gramians are finite and positive definite and satisfy
\begin{align}
 P&=\sum_w A_w^\top\alpha\alpha^\top A_w
   =\alpha\alpha^\top+\Phi_A(P),\label{I:S:eq:P}\\
 Q&=\sum_w A_w\beta\beta^\top A_w^\top
   =\beta\beta^\top+\Psi_A(Q).\label{I:S:eq:Q}
\end{align}
Without minimality, strict stability remains sufficient but is not necessary for boundedness of the realized Hankel.
\end{theorem}
\begin{proof}
Boundedness implies $f=H_f e_\Seps\in\ell^2$. If $f\in\ell^2$, all shifted sequences $u\mapsto f(uv)$ and $v\mapsto f(uv)$, with the other word fixed, are square summable, because concatenation with a fixed word is injective. Choose a basis among $A_v\beta$; its images under $O$ are such sequences. Hence $O$ is bounded. Choose a basis among $A_u^\top\alpha$ to prove boundedness of $C^*$. Their Gramians are precisely \eqref{I:S:eq:P}--\eqref{I:S:eq:Q}; the span conditions make them positive definite.

Unique decomposition by the last letter gives the first Stein equation. Its tail identity is
\[
 \Phi_A^m(P)=\sum_{|v|=m}A_v^\top P A_v
           =\sum_{|w|\ge m}A_w^\top\alpha\alpha^\top A_w\longrightarrow0.
\]
In the order-unit norm $\Snorm{Y}_P=\Snorm{P^{-1/2}YP^{-1/2}}$, positivity implies
$\Snorm{\Phi_A^m}_{P\to P}\le\Snorm{\Phi_A^m(P)}_P$. Some power has norm less than one, so $\rho(\Phi_A)<1$.

Conversely, a positive map with spectral radius below one has $X=\sum_{j\ge0}\Phi_A^j(I)\succ0$ and $X-\Phi_A(X)=I$. If a strict Lyapunov inequality holds, choose $\mu>0$ with $\Phi_A(X)\preceq(1-\mu)X$. The same order-unit argument gives spectral radius below one. Since $\Psi_A$ is adjoint to $\Phi_A$, both Neumann sums in \eqref{I:S:eq:P}--\eqref{I:S:eq:Q} converge. They define bounded $O,C$, and $H_f=OC$ first on finitely supported vectors and then by continuity. This converse does not use minimality.
\end{proof}


\begin{corollary}[Complete contractive cover]
Every bounded finite-rank real word Hankel has a minimal realization similar to one satisfying $\sum_a A_a^\top A_a\prec I$. Any smaller rank can be embedded in a prescribed larger stable dimension by zero-padding.
\end{corollary}
\begin{proof}
Apply the realization theorem and Theorem~\ref{I:S:thm:stability}; conjugate $A_a$ to $X^{1/2}A_aX^{-1/2}$. Zero-padding adds a strictly positive identity slack on the unused coordinates.
\end{proof}


\begin{theorem}[Finite norm formula, including singular Gramians]\label{I:S:thm:norm}
For any stable, possibly nonminimal realization, $H_f=OC$ and
\[
 \Snorm{H_f}^2=\lambda_{\max}(P^{1/2}QP^{1/2})
              =\lambda_{\max}(L^\top QL),\qquad P=LL^\top.
\]
Its nonzero squared singular values are the nonzero eigenvalues of these finite matrices. For $t\ge0$,
\begin{equation}\label{I:S:eq:factorfree}
 \Snorm{H_f}\le t\quad\Longleftrightarrow\quad t^2P-PQP\succeq0.
\end{equation}
\end{theorem}
\begin{proof}
$H_f^*H_f=C^*PC$; the nonzero spectra of $T^*T$ and $TT^*$ agree for $T=P^{1/2}C$, proving the formula. The factor $L$ gives the same nonzero spectrum, as does any square factor of $P$. Put $S=P^{1/2}QP^{1/2}$. It vanishes on $\ker P$ and preserves $\Sran P$. Congruence by the invertible restriction of $P^{1/2}$ on that range identifies $S\preceq t^2I$ with $t^2P-PQP\succeq0$ there. Both latter terms vanish on the kernel, and its additional zero eigenvalues are harmless because $t\ge0$.
\end{proof}


For two models use $D_a=\Sdiag(A_a,B_a)$, $c=(\alpha,-\gamma)$, $b=(\beta,\nu)$ as columns. Their difference is $c^\top D_wb$. Its Gramians certify the error on \emph{all} words; a finite word section alone normally supplies only a lower bound.

\begin{proposition}[Elementary norm and tail bounds]\label{I:S:prop:normbounds}
For any series $h$,
\[
 \Snorm{H_h}_{\mathrm{HS}}^2=\sum_w(|w|+1)|h(w)|^2
\]
whenever either side is finite. If $h\in\ell^1$, then $\Snorm{H_h}\le\Snorm h_1$. If $h=\sum_m h_m$ with $h_m$ homogeneous and $\sum_m\Snorm{h_m}_2<\infty$, then $\Snorm{H_h}\le\sum_m\Snorm{h_m}_2$.
\end{proposition}
\begin{proof}
Every word of length $m$ has exactly $m+1$ prefix--suffix splits, giving the Hilbert--Schmidt identity by nonnegative summation. Each absolute row and column sum is at most $\Snorm h_1$; weighted Cauchy--Schwarz proves the Schur bound. A homogeneous symbol has norm $\Snorm{h_m}_2$ by the orthogonal block argument in Lemma~\ref{I:S:lem:homogeneous}; summing operators proves the last assertion and identifies all entries of the limit.
\end{proof}


\section{Existence, exact algebraicity, and certificate logic}\label{I:S:sec:global}
For a real target, $\mathbb F\in\{\SR,\SC\}$, and $k\ge0$, define
\[
 \mathscr H_k^{\mathbb F}=\{H_g:H_g\text{ bounded},\ g:\Swords\to\mathbb F,\ \Srank H_g\le k\},
 \qquad \SdeltaF{k}(H)=\inf_{G\in\mathscr H_k^{\mathbb F}}\Snorm{H-G}.
\]
Ranks are linear WFA dimensions. Competitors are not required to be commuting, nonnegative, or probabilistically normalized. For compact $H$, the unstructured lower bound is $\sigma_{k+1}(H)\le\SdeltaF{k}(H)$. Real targets satisfy $\delta_k^\SC\le\delta_k^\SR$.

Attainment over either scalar field is Theorem~\ref{I:U:thm:attainment};
its proof applies to any bounded target operator. The algebraic endpoint
and degree-three existential formulation are Corollary~\ref{I:U:cor:algebraic}.

The same conclusion holds for a bounded real-algebraic WFA after minimal
stable realization. Strict thresholds can be encoded by requiring the
factor of the norm slack to be invertible. No practical running-time
bound follows from general quantifier elimination.

\begin{remark}[Strict versus non-strict certificates]
For singular $P$, strict positive definiteness of $t^2P-PQP$ is impossible even if the error is strictly below $t$. The $t^2I-L^\top QL$ encoding handles strict thresholds correctly. Alternatively, proving error at most a smaller rational threshold suffices for a strict public upper endpoint. A feasible candidate is an upper certificate, not a certificate of its optimality; a solver timeout is not an infeasibility proof.
\end{remark}

\begin{proposition}[Perturbation transfer]\label{I:S:prop:transfer}
For one fixed competitor class, $|\SdeltaF{k}(H)-\SdeltaF{k}(\widetilde H)|\le\Snorm{H-\widetilde H}$. If $L\le\SdeltaF{k}(H_f)\le\Snorm{H_f-H_g}\le U$ and $\Snorm{H_T-H_f}\le\eta$, then
\[
 \max(0,L-\eta)\le\SdeltaF{k}(H_T)\le\Snorm{H_T-H_g}\le U+\eta.
\]
For compact operators, the spectral gap $\SdeltaF{k}-\sigma_{k+1}$ decreases by at most $2\eta$ under a perturbation of size $\eta$.
\end{proposition}
\begin{proof}
Apply the triangle inequality against each competitor and take infima in both directions. Individual singular values are 1-Lipschitz, proving the last assertion.
\end{proof}


\section{Rank-one approximation and the free-word obstruction}\label{I:S:sec:rankone}
Let $L_a e_w=e_{aw}$ and $R_a e_w=e_{wa}$. Then $H_f L_a=R_a^*H_f$. Consequently the range of a finite-rank Hankel is invariant under every $R_a^*$, and its adjoint range is invariant under every $L_a^*$. Restricting the usual subspace-width minimization to these invariant subspaces gives lower bounds on $\SdeltaF{k}$. This restriction remembers the letter action which a free matrix SVD discards.

By Lemma~\ref{I:M:rev13:lem:rank-one}, write $G=c k_zk_z^\top$,
where $k_z(w)=\prod_i z_{w_i}$, $\sum_a|z_a|^2<1$, and
$\|k_z\|^2=(1-\|z\|_2^2)^{-1}$. The coefficient transpose is
intentional also for complex parameters.

\begin{lemma}[One-corner completion]\label{I:S:lem:corner}
For a bounded operator $T$ and unit output/input vectors $u,v$,
\[
 \min_c\Snorm{T-cuv^*}
   =\max\{\Snorm{T(I-vv^*)},\Snorm{(I-uu^*)T}\}.
\]
The minimizing scalar can be real for real data.
\end{lemma}
\begin{proof}
Use the decompositions along $u,v$. Only the upper-left scalar block is free; the bottom row and right column give the lower bound $\gamma$. If $\gamma=0$, cancel the only nonzero block. Otherwise divide by $\gamma$ and denote the three fixed blocks by $b,c,d$, so their row and column norms are at most one. The inequalities give contractions $B,C$ with
$b=B(I-d^*d)^{1/2}$ and $c=(I-dd^*)^{1/2}C$. These follow by defining the contraction on the range of the defect square root and extending continuously. The block matrix with free corner $-Bd^*C$ factors as
\[
 \begin{pmatrix}B&0\\0&I\end{pmatrix}
 \begin{pmatrix}-d^*&(I-d^*d)^{1/2}\\(I-dd^*)^{1/2}&d\end{pmatrix}
 \begin{pmatrix}C&0\\0&I\end{pmatrix}.
\]
The middle matrix is unitary by the defect-operator identities; the others are contractions. Rescaling proves attainment. This is the classical Parrott completion argument~\cite{parrott}.
\end{proof}


\begin{theorem}[An exact real rank-one variational formula]\label{I:S:thm:rankone}
For a bounded real Hankel and $P_b$ the orthogonal projection onto $\SR k_b$,
\[
 \delta_1^\SR(H_f)=\min_{b\in\SR^d,\ \Snorm b<1}
 \max\{\Snorm{H_f(I-P_b)},\Snorm{(I-P_b)H_f}\}.
\]
The minimum over the open ball is attained. If $H_f$ has finite rank with singular values $\sigma_j$ and right/left singular vectors $\xi_j,\eta_j$, each squared term is a finite rank-one downdate:
\[
 \lambda_{\max}\bigl(D^{1/2}(I-c^Rc^{R\top})D^{1/2}\bigr),\qquad
 \lambda_{\max}\bigl(D^{1/2}(I-c^Lc^{L\top})D^{1/2}\bigr),
\]
where $D=\Sdiag(\sigma_j^2)$, $c^R_j=\Sip{\xi_j}{k_b/\Snorm{k_b}}$ and $c^L_j=\Sip{\eta_j}{k_b/\Snorm{k_b}}$.
\end{theorem}
\begin{proof}
Fix $b$ and apply Lemma~\ref{I:S:lem:corner} with $u=v=k_b/\Snorm{k_b}$. The zero candidate cannot improve on this minimum. The objective is continuous on the ball and bounded above by $\Snorm{H_f}$. As $\Snorm b\to1$, the normalized kernels converge weakly to zero, uniformly on every fixed finite set of coordinates: each such coordinate is bounded in absolute value by $\sqrt{1-\Snorm b^2}$. Choose a finitely supported unit vector nearly attaining $\Snorm{H_f}$ to see that $\Snorm{H_f(I-P_b)}\to\Snorm{H_f}$, uniformly in boundary direction. Apply the same argument to $H_f^*$ for the other term. If the infimum equals $\Snorm{H_f}$, every point minimizes; otherwise it lies in a compact sub-ball and is attained. Finally apply the identity between the nonzero spectra of $TT^*$ and $T^*T$ to $(I-P_b)[\xi_j]D^{1/2}$, and similarly on the left.
\end{proof}


For a real minimal stable realization the kernel evaluations are rational:
\[
 Ck_b=(I-\sum_a b_aA_a)^{-1}\beta,\qquad
 O^*k_b=(I-\sum_a b_aA_a^\top)^{-1}\alpha.
\]
The inverses exist on the ball: in the contractive gauge,
$\Snorm{\sum b_aA_a}\le\Snorm b<1$. The transpose sum has the same norm. The power series converge absolutely by Cauchy--Schwarz against the finite Gramians.

\begin{corollary}[Rank-two secular equality criterion]\label{I:S:cor:secular}
Suppose $\Srank H_f=2$ and $\sigma_1>\sigma_2>0$. A one-sided term in Theorem~\ref{I:S:thm:rankone} equals $\sigma_2$ exactly when its corresponding normalized coefficients satisfy
\[
 c_2=0,\qquad c_1^2\ge1-\sigma_2^2/\sigma_1^2.
\]
Both the left and the right requirements must hold at the same $b$ for spectral equality. If $\sigma_1=\sigma_2=s$, every one-sided term equals $s$ and this zero test is inapplicable.
\end{corollary}
\begin{proof}
Subtract $\sigma_2^2I$ from the $2\times2$ downdate. Its determinant is
$-\sigma_2^2c_2^2(\sigma_1^2-\sigma_2^2)$. Negative semidefiniteness therefore forces $c_2=0$, after which its remaining diagonal inequality is the stated condition. For a repeated value, $s^2(I-cc^\top)$ has an eigenvalue $s^2$ and is bounded above by $s^2I$.
\end{proof}


\subsection{A commuting isospectral separation}
\begin{theorem}[Spectrum and abstract algebra do not determine the optimum]\label{I:S:thm:commutative}
Let $f=\mathbf1_{\Seps}+\mathbf1_b+\mathbf1_{aa}$ over $\{a,b\}$. It has nonnegative commuting minimal transitions $A_a=J$, $A_b=J^2$, where $J$ is the $3\times3$ upper shift. Its nonzero singular values are $(2,1,1)$, and
\[
 \sqrt{727/544}\le\delta_1^\SC(H_f)\le\delta_1^\SR(H_f)<29/25.
\]
There is a unary-supported nonnegative target with exactly the same singular values and abstract syntactic algebra $\mathbb F[t]/(t^3)$ whose optimal rank-one error is exactly one.
\end{theorem}
\begin{proof}
Take $\alpha=e_1$, $\beta=(1,0,1)^\top$. The only nonzero rows and columns are indexed by $(\Seps,a,b,aa)$, and the entire nonzero operator is
\[
 H=\begin{pmatrix}1&0&1&1\\0&1&0&0\\1&0&0&0\\1&0&0&0\end{pmatrix}.
\]
Its eigenvectors $(2,0,1,1)/\sqrt6$, $(0,1,0,0)$ and $(1,0,-1,-1)/\sqrt3$ have eigenvalues $2,1,-1$. Thus the rank is three and the realization is minimal.

For a complex rank-one candidate with parameters $(p,q)$, put $x=|p|^2$, $t=|q|$, $m=1-x-t^2>0$, and $z=q+p^2$. Projecting $H$ away from its possible output range gives the finite matrix
$M=H^2-mww^*$, $w=(1+z,p,1,1)^\top$. On $\Sran H$ one eigenvalue is one; the other two, $0\le\mu_-\le1\le\mu_+\le4$, satisfy
\[
 Q_0(\lambda)=(\lambda-4)(\lambda-1)+4A(\lambda-1)+B(\lambda-4),
\]
where $A=m|2+z|^2/6$ and $B=m(x+|1-z|^2/3)$. Direct evaluation gives
\[
 Q_0(4/3)=\tfrac23\{m(4\Re z-|z|^2-4x)-4/3\}.
\]
Since $t+x<1+t-t^2\le5/4$, the reverse triangle inequality yields
$4\Re z-|z|^2-4x\le4t-(t+x)^2$. If the right side is negative the next bound is immediate; otherwise both factors decrease with $x$. Therefore
$m(4\Re z-|z|^2-4x)\le(1-t^2)(4t-t^2)$.
For $y=t-1/2\in[-1/2,1/2]$, the exact identity
\[
 4-3(1-t^2)(4t-t^2)
 =\tfrac{51}{4}(y-1/17)^2+\tfrac5{272}
   +3y^2(y+1/2)(5/2-y)
\]
makes $Q_0(4/3)\le-5/1224$. Hence
$\mu_+\ge4/3+(5/1224)/(4/3)=727/544$, proving the lower bound. Subsection~\ref{I:S:app:commupper} gives a rational candidate and its strict upper certificate.

For the comparison put $r=2^{1/3}$ and let $h(\Seps)=2-r$, $h(a)=r^2-1$, $h(aa)=r$, with all other coefficients zero. Its unary block
\[
 K=\begin{pmatrix}2-r&r^2-1&r\\r^2-1&r&0\\r&0&0\end{pmatrix}
\]
has characteristic polynomial $(\lambda-2)(\lambda-1)(\lambda+1)$. Let $b_0=r^{-1}$ and $h_*(a^n)=(2-r)b_0^n$, zero on words containing $b$. The vector $(1,b_0,b_0^2)$ is a top eigenvector, $b_0^6=1/4$, and $\Snorm{H_{h_*}}=2$. Its normalized infinite kernel has squared overlap $3/4$ with the normalized top eigenvector of $K$. The difference of the two rank-one terms of size two has norm $2\sqrt{1-3/4}=1$; the orthogonal eigenvectors of eigenvalues $1,-1$ are unchanged. Thus $\Snorm{H_h-H_{h_*}}=1$, which is optimal by the singular-value bound. Both targets have generated algebra $\mathbb F[J]\simeq\mathbb F[t]/(t^3)$, but the letter $b$ is marked by $J^2$ in one and zero in the other.
\end{proof}


The identification of $b$ with two unary steps is not an isometry: $b$ and $aa$ are distinct words with different splits. More generally, if a symmetric series depends only on its count vector $\mathbf m$, define $c(\mathbf m)=|\mathbf m|!/\prod_i m_i!$. Normalized sums over words with a fixed count give an isometry under which its nonzero block is
\[
 \widehat H_f[\mathbf m,\mathbf n]
 =f(\mathbf m+\mathbf n)\sqrt{c(\mathbf m)c(\mathbf n)}.
\]
This follows by summing the identical entries over each fiber. Dropping the multinomial factors changes the norm. In a minimal rational realization, symmetry is equivalent to commuting transitions, since $f(uabv)=f(ubav)$ for all contexts forces $A_aA_b=A_bA_a$. At rank one all competitors are symmetric; at larger ranks, restricting competitors to symmetric series can increase the optimum.

\section{Coefficient and trace bounds; a sharp homogeneous constant}\label{I:S:sec:homogeneous}
For $p,q\ge0$, write $C_{p,q}(f)=(f(uv))_{|u|=p,|v|=q}$, and let $E_k(C)^2=\sum_{j>k}\sigma_j(C)^2$.

\begin{theorem}[Coefficient lower bound across lengths]\label{I:S:thm:coefficient}
For every bounded Hankel and $N\ge0$,
\[
 \delta_k^\SC(H_f)^2\ge
 \sum_{m=0}^N\max_{0\le p\le m}E_k(C_{p,m-p}(f))^2.
\]
If $H_f$ is compact, this combines with the independent lower bound $\sigma_{k+1}(H_f)$.
\end{theorem}
\begin{proof}
Every $C_{p,q}(g)$ of a rank-$k$ competitor has rank at most $k$. Frobenius Eckart--Young gives an error of at least $E_k(C_{p,q}(f))^2$. For fixed $m=p+q$, that squared error is exactly $\sum_{|w|=m}|f(w)-g(w)|^2$, independent of the split. Maximize over the splits, and sum over different lengths. The result is bounded above by $\Snorm{(H_f-H_g)e_\Seps}^2\le\Snorm{H_f-H_g}^2$. Take infima.
\end{proof}

One must maximize, not sum, over splits of the same length; otherwise the same coefficients are counted repeatedly without justification.

\begin{lemma}[Homogeneous norm and spectrum]\label{I:S:lem:homogeneous}
If $f\ne0$ is supported at length $m$, then $\Snorm{H_f}=a:=\Snorm f_2$. If $m\ge2$, then $\sigma_1=\sigma_2=a$ and $\sigma_3\ge a/\sqrt d$.
\end{lemma}
\begin{proof}
The word-length decomposition makes each nonzero block map degree $q$ to degree $m-q$, with block $C_{m-q,q}(f)$. Their input and output spaces are mutually orthogonal, so their singular-value lists combine. Each Frobenius norm is $a$. The two endpoint blocks have singular value $a$, and every other block has norm at most $a$. The block $C_{1,m-1}$ has at most $d$ rows and Frobenius norm $a$, hence has a singular value at least $a/\sqrt d$.
\end{proof}


\begin{theorem}[A trace certificate at one length]\label{I:S:thm:trace}
Put $A=\Snorm{H_f}$ and $a_m=(\sum_{|w|=m}|f(w)|^2)^{1/2}$. Then
\[
 \delta_k^\SC(H_f)\ge
 \max\left\{0,\frac{(m+1)a_m-kA}{m+1+k}\right\}.
\]
For a homogeneous degree-$m$ target this becomes
$\delta_k^\SC(H_f)\ge A\max\{0,(m+1-k)/(m+1+k)\}$.
\end{theorem}
\begin{proof}
If $a_m=0$ the claim is immediate. Let $T=H_{f\mathbf1_{|w|=m}}$, so $\Snorm T=a_m$. For a candidate $G=H_g$ of rank at most $k$ and error $t$, trace duality and $\Snorm G_1\le k\Snorm G\le k(A+t)$ give
\[
 (m+1)\left|\sum_{|w|=m}\overline{f(w)}g(w)\right|
 =|\Sip T G_{\rm HS}|\le a_m k(A+t).
\]
The empty column, paired with $f\mathbf1_{|w|=m}/a_m$, gives $a_m-t$ as a lower bound on the real part of that normalized coefficient inner product. Combining and rearranging proves the result. The factor $m+1$ comes from the exact number of splits and is justified here by trace duality.
\end{proof}


\begin{theorem}[Unbounded alphabet-dependent gap and sharp fixed-alphabet constant]\label{I:S:thm:sharp}
Let $f_{d,m}=\sum_{i=1}^d\mathbf1_{a_i^m}$ with $m\ge2$. It has a nonnegative commuting nilpotent realization of rank $d(m-1)+2$ and spectrum
\[
 (\sqrt d,\sqrt d,\underbrace{1,\ldots,1}_{d(m-1)}).
\]
For $2\le k\le d$, at $m=2$,
\[
 \max\{1,\sqrt{d-k}\}\le\delta_k^\SC(H_{f_{d,2}})
 \le\delta_k^\SR(H_{f_{d,2}})\le\sqrt{d-k+2}.
\]
Thus no multiplicative constant independent of $d$ bounds $\delta_k/\sigma_{k+1}$ for a fixed $k\ge2$. For each fixed $d\ge1$ and $\mathbb F\in\{\SR,\SC\}$,
\begin{equation}\label{I:S:eq:sharp}
 \sup_{\substack{f\ne0\ \text{real and homogeneous}\\ \deg f\ge2}}
 \frac{\SdeltaF{2}(H_f)}{\sigma_3(H_f)}=\sqrt d.
\end{equation}
The same supremum holds within the nonnegative commuting nilpotent target class.
\end{theorem}
\begin{proof}
Give each letter a chain of $m$ edges sharing only its source and sink with the other chains. The $d(m-1)+2$ states realize the indicated series; mixed products are zero and the letter matrices commute. The endpoint Hankel blocks have singular values $\sqrt d$, while each of the $m-1$ interior blocks is an identity on its $d$ pure words. This proves the rank and spectrum. For $m=2$, $C_{1,1}=I_d$ gives $\sqrt{d-k}$ by Theorem~\ref{I:S:thm:coefficient}. Keep $k-2$ pure-square terms to obtain a rank-$k$ candidate; the remaining $d-k+2$ terms have norm $\sqrt{d-k+2}$. At $k=2$ use zero. Letting $d\to\infty$ proves the unbounded ratio.

For every homogeneous target, zero and Lemma~\ref{I:S:lem:homogeneous} give $\delta_2\le\sqrt d\,\sigma_3$. For the family above, Theorem~\ref{I:S:thm:trace} and the spectral bound give
\[
 \max\{1,\sqrt d(m-1)/(m+3)\}\le\delta_2\le\sqrt d,
 \qquad\sigma_3=1.
\]
Let $m\to\infty$ with $d$ fixed. This proves \eqref{I:S:eq:sharp}, including its restricted-class assertion.
\end{proof}


For $d=100$, tolerance two, the spectral state lower bound is only two, while Theorem~\ref{I:S:thm:coefficient} excludes all ranks at most 95. Keeping 96 terms gives a rank-98 candidate with error two. Thus the true tolerance budget lies between 96 and 98. Equation~\eqref{I:S:eq:sharp} is a supremum, not an equality for every homogeneous target. It does not establish a finite factor for all binary mixed-length targets.

\subsection{A strict binary rank-two lower bound without a solver}
Let $f_0=\mathbf1_{aa}+\mathbf1_{bb}$, $h=(e_{aa}+e_{bb})/\sqrt2$, $T=\SSpan\{e_\Seps,h\}$ and $M=\SSpan\{e_a,e_b\}$. Its exact support is $S=T\oplus M$, and
$\Snorm{H_{f_0}x}^2=\Snorm x^2+\Snorm{P_Tx}^2$ on $S$.

\begin{proposition}[Uniform binary geometric obstruction]\label{I:S:prop:binary}
Let $\tau$ be the unique root in $(1,\sqrt2)$ of
\[
 \Phi(t)=\sqrt2-t-\sqrt2(\sqrt2+t)(t^2-1).
\]
Then $\delta_2^\SC(H_{f_0})\ge\tau>2627/2500=1.0508$.
\end{proposition}
\begin{proof}
For a rank-at-most-two candidate $G$, write $t=\Snorm{H_{f_0}-G}$. Only $1\le t<\sqrt2$ needs consideration. There are two-dimensional spaces $K_R\subset\ker G\cap S$ and $K_L\subset\ker G^*\cap S$. On either, $\Snorm{P_Tx}\le q\Snorm x$, $q^2=t^2-1<1$. If the columns of $\binom XY$ are an orthonormal basis of such a space, $X^*X\preceq q^2I$ and $Y^*Y\succeq(1-q^2)I$. The square matrix $Y$ has $YY^*\succeq(1-q^2)I$, proving $\Snorm{(I-P_K)P_M}\le q$. Therefore
\[
 \Snorm{P_MGP_M}\le q^2\Snorm G\le(t^2-1)(\sqrt2+t).
\]
By the Hankel equalities, the four coefficients of $F=P_MGP_M$ are exactly the degree-two coefficients of $Ge_\Seps$. Hence
\[
 \sqrt2-t\le\Snorm F_F\le\sqrt2\Snorm F
 \le\sqrt2(t^2-1)(\sqrt2+t).
\]
Thus $\Phi(t)\le0$. On $[1,\sqrt2]$, $\Phi$ is strictly decreasing and changes sign. Exact rational squaring with checked positive signs gives
$2627/2500<\tau<105083/100000$; the endpoint inequalities follow by direct rational substitution. The bound is uniform and includes complex candidates.
\end{proof}


\section{Constructive equality and controlled-gap families}\label{I:S:sec:constructive}
\subsection{The length-only equality locus}
\begin{proposition}[Radial unary reduction]\label{I:S:prop:radial}
If $f(w)=\varphi(|w|)$ and $\psi(n)=d^{n/2}\varphi(n)$ has bounded finite-rank unary Hankel, then
$H_f\simeq H_\psi\oplus0$ and $\SdeltaF{k}(H_f)=\sigma_{k+1}(H_f)$ whenever scalar AAK applies over the chosen field.
\end{proposition}
\begin{proof}
The isometry $Ve_n=d^{-n/2}\sum_{|w|=n}e_w$ has $V^*H_fV=H_\psi$ by direct summation. $H_f$ vanishes on the complement of the length-constant subspace and has range in it. A unary rank-$k$ AAK optimizer $h$ lifts to $g(w)=d^{-|w|/2}h(|w|)$ with every letter transition equal to the unary transition divided by $\sqrt d$. The lift preserves error and rank; the singular-value lower bound supplies the reverse inequality. Only the unary optimizer is imported from classical AAK~\cite{lacroce2024}.
\end{proof}


\subsection{An exact mixed-length optimum with seven target states}
Take real parameters $s,\eta,\rho$, with $|\rho|<1$, and let
\[
 g_s(w)=s\{(1/2)^{\#a(w)}(1/4)^{\#b(w)}
             +(1/4)^{\#a(w)}(1/2)^{\#b(w)}\},\qquad s\ne0,
\]
and define $h_\rho(a^3)=-\rho$, $h_\rho(a^j)=(1-\rho^2)\rho^{j-4}$ for $j\ge4$, zero on other words, with $|\rho|<1$. At $\rho=0$ it is $\mathbf1_{a^4}$.

\begin{theorem}[A globally optimal explicit two-state model]\label{I:S:thm:mixedexact}
For $s\eta\ne0$, $f=g_s+\eta h_\rho$ has rank seven, commuting nonproportional transitions, and
\[
 \delta_2^\SC(H_f)=\delta_2^\SR(H_f)
  =\sigma_3(H_f)=\Snorm{H_f-H_{g_s}}=|\eta|.
\]
If $|\eta|<2|s|/33$, one state cannot meet tolerance $|\eta|$.
\end{theorem}
\begin{proof}
The two-state realization of $g_s$ has diagonal transitions
$B_a=\Sdiag(1/2,1/4)$, $B_b=\Sdiag(1/4,1/2)$, initial row $(s,s)$ and final column $(1,1)^\top$. The residual has $A_b=0$ and
\[
 T_\rho=\begin{pmatrix}0&1&0&0&0\\0&0&1&0&0\\0&0&0&1&0\\0&0&0&0&1\\0&0&0&0&\rho\end{pmatrix},
 \quad c=e_1,\quad v=(0,0,0,-\rho,1-\rho^2)^\top.
\]
Its exact Stein solutions are
$P=\Sdiag(1,1,1,1,(1-\rho^2)^{-1})$ and
$Q=\Sdiag(1,1,1,1,1-\rho^2)$, so $PQ=I_5$. Thus $H_{h_\rho}$ has norm one and is an isometry on a five-dimensional input subspace. For any rank-two operator $G$, the difference $H_{g_s}-G$ has rank at most four and has a unit kernel vector in that subspace. The error on that vector is $|\eta|$. The candidate $H_{g_s}$ attains it, even among unstructured rank-two operators, proving the singular-value equality as well.

The direct-sum realization has seven states. The principal section indexed by $\Seps$, $a$, $a^2$, $a^3$, $a^4$, $b$, $b^2$ has determinant
$s^2\eta^5(1-\rho^2)/1024\ne0$: eliminate the invertible $(b,b^2)$ block of the rank-two part, whose determinant is $s^2/1024$, leaving the residual five-block of determinant $1-\rho^2$. This establishes minimality. The two geometric kernels have Gram matrix $\begin{psmallmatrix}16/11&4/3\\4/3&16/11\end{psmallmatrix}$, so the smaller singular value of $H_{g_s}$ is $4|s|/33$. Weyl's bound gives $\sigma_2(H_f)\ge4|s|/33-|\eta|>|\eta|$ under the final assumption.
\end{proof}


More generally, if $R$ has norm one and is an isometry on a space of dimension $2k+1$, then every $G_0$ of rank at most $k$ is a globally optimal rank-$k$ approximation to $G_0+\eta R$, with error $|\eta|$. The preceding kernel argument proves this general plateau principle. The family is deliberately solvable; it does not solve arbitrary mixed-length approximation.

\subsection{Recognizable low-rank background with one homogeneous residual}
For $r\ge1$, $m\ge2r$, let $\mathcal C_{r,m}$ consist of bounded rational real series $f$ such that some bounded $g_0$ of rank at most $r$ satisfies $f(w)=g_0(w)$ whenever $|w|\ne m$.

\begin{theorem}[Reconstruction with a certified additive gap]\label{I:S:thm:constructgap}
The background $g_0$ is unique and reconstructible from coefficients of lengths at most $2r-1$. Given the full exact WFA, membership is decidable by finite linear algebra. With $h=f-g_0$, $U=\Snorm h_2$, one obtains
\[
 L:=U\frac{m+1-2r}{m+1+2r}\le\delta_r^\SC(H_f)
 \le\delta_r^\SR(H_f)\le\Snorm{H_f-H_{g_0}}=U,
 \qquad U-L=\frac{4rU}{m+1+2r}.
\]
\end{theorem}
\begin{proof}
Two rank-at-most-$r$ backgrounds agree on every coefficient up to length $2r-1$, since $m\ge2r$. Their difference is identically zero by Theorem~\ref{I:S:thm:realization}. For reconstruction set $V=\Sigma^{\le r-1}$ and $K=(f(pq))_{p,q\in V}$. Under the promise, $K$ has rank $s=\Srank H_{g_0}\le r$. Select an invertible $s$-square submatrix $M=H_f[P,Q]$ and form $M_a=(f(paq))_{P,Q}$. The realization
\[
 \widehat B_a=M^{-1}M_a,\qquad
 \widehat\alpha^\top=(f(q))_{q\in Q},\qquad
 \widehat\beta=M^{-1}(f(p))_{p\in P}
\]
is similar to the minimal background realization: factor $M=O_PC_Q$ and cancel the invertible factors. The zero-rank case gives zero. Reject if $\Srank K>r$.

Without the promise, form $D=F-\widehat G$. For each length less than $m$, generate its exact-length row space and test the output. Generate the row space at length $m+1$, close it under all letters, and test the output on that closure. These finite tests are equivalent to vanishing outside length $m$. If they pass, the candidate is bounded as the difference of a bounded Hankel and a polynomial Hankel. Its residual norm is $U$ by Lemma~\ref{I:S:lem:homogeneous}. For any competitor $q$ of rank at most $r$, the Hankel $H_q-H_{g_0}$ has rank at most $2r$. Apply the homogeneous trace bound at budget $2r$ to $H_h$ and take infima.
\end{proof}


The residual norm can be computed without enumerating its degree-$m$ words: for $h(w)=c^\top D_wb$, start $Q_0=bb^\top$, iterate $Q_{j+1}=\sum_aD_aQ_jD_a^\top$, and take $U^2=c^\top Q_mc$. A perturbation of the target by Hankel norm at most $e$ transfers the bracket to $[\max(0,L-e),U+e]$. For a fixed $m$, the displayed gap is not arbitrarily small merely by rerunning the construction.

\begin{proposition}[Directional trace improvement]\label{I:S:prop:directional}
If $h=f-g_0$ is homogeneous of degree $m$, $U=\Snorm h_2>0$, $A_+\ge\Snorm{H_f}$, and $c_m=|\sum_{|w|=m}\overline{h(w)}f(w)|/U$, then
\[
 \delta_r^\SC(H_f)\ge
 \max\{0,((m+1)c_m-rA_+)/(m+1+r)\}.
\]
\end{proposition}
\begin{proof}
For an error $t$, Cauchy--Schwarz on the empty column gives $U(c_m-t)\le|\Sip h q|$. The trace identity and rank bound give $(m+1)|\Sip h q|\le Ur(A_++t)$. Rearrange.
\end{proof}


For $g_0=g_s$, $s=33/736$, and $h=\mathbf1_{a^{397}}+\mathbf1_{b^{397}}$, one has $\Snorm{H_{g_0}}=1/8$, $U=\sqrt2$, and $\sigma_3(H_f)\le9/8$. Positivity gives $c_m\ge U$, so
\[
 \frac98\ge\sigma_3(H_f),\qquad
 \frac{99}{100}\sqrt2-\frac1{1600}\le\delta_2^\SC(H_f)
 \le\delta_2^\SR(H_f)\le\sqrt2.
\]
The left error endpoint exceeds $1.399445$, and the certified width is below $0.014768$. This is a different target from the fixed small target of the next section. Increasing $m$ changes the target and cannot be described as closing the interval for one fixed target.

\section{The question and the exact answer}
Let $\Sigma$ be any finite alphabet containing distinct letters $a,b$. All word spaces use counting measure. For a scalar series $f:\Sigma^*\to\BC$, put
\[
H_f[u,v]=f(uv),\qquad
\delta_k^{\BC}(H_f)=\inf\{\Bnorm{H_f-H_g}:H_g\text{ bounded word-Hankel},\ \Brank H_g\le k\}.
\]
The norm is the operator norm on $\BHH$. For real $f$, let $\delta_k^{\BR}$ denote the same infimum restricted to real coefficients. Without a superscript, $\delta_k$ means $\delta_k^{\BC}$. In this part use the source \eqref{I:eq:binary-source}, with fixed $k\ge1$, $0<r<1$, and $q=1-r^2$.
The asymptotic is always at fixed $k$, rather than uniformly over growing budgets.

For $0<p<\infty$, the classical Hardy quasi-norm is
\[
\Bnorm{F}_{H^p}=
\sup_{0<\rho<1}\left(\int_0^{2\pi}|F(\rho e^{it})|^p\frac{dt}{2\pi}\right)^{1/p},
\qquad C(m,p)=\sup_{\Bnorm{F}_{H^p}=1}|[z^m]F|.
\]
The circle measure is normalized. In particular $C(0,p)=1$. This is the convention of Brevig--Saksman~\cite{brevig2020}; it is fixed independently of the limit we wish to prove.

\begin{theorem}[Exact boundary reduction and its Hardy identification]
\label{I:B:thm:main}
For $\mathbb F\in\{\mathbb R,\mathbb C\}$ define
\[
 E_m^{\mathbb F}=\inf\{\|z^m-h\|_{H^2}:h\in H^2\text{ over }\mathbb F,
                    \ \operatorname{rank}(h_{i+j})_{i,j\ge0}\le m\}.
\]
For every fixed integer $k\ge1$, the self-contained reduction gives
\[
 \boxed{\displaystyle\lim_{r\uparrow1}\delta_k^{\mathbb F}(H_{f_{r,k}})
       =E_{k-1}^{\mathbb F}\ge\frac1{C(k-1,2/3)}>0.}
\]
For $k\le4$, the published low-degree extremals of
Brevig--Saksman~\cite{brevig2020} give
$E_{k-1}^{\mathbb R}=E_{k-1}^{\mathbb C}=1/C(k-1,2/3)$.
For arbitrary $k$, the same equality uses the outer-extremal theorem
of Brevig--S\o vik~\cite{outer2026}, together with the classical
polynomial structure of coefficient extremals. This is an external
theorem input; the reduction to $E_{k-1}^{\mathbb F}$ does not use it.
Both infima use ordinary word-Hankel rank. No equality of the real and
complex optima at finite $r$ is asserted.
\end{theorem}

The proof has two independent reductions. First, with $m=k-1$,
\[
\lim_{r\uparrow1}\delta_k^{\BC}(H_{f_{r,k}})
=E_m:=\inf_{\Brank H_h\le m}\Bnorm{z^m-h}_{H^2},
\]
where $H_h[i,j]=h_{i+j}$ is a one-letter coefficient Hankel and $h\in H^2$. Second,
\[
E_m=\frac1{C(m,2/3)}.
\]
The first reduction follows from finite-dimensional shift geometry and an explicit construction. The second uses classical finite model spaces and the outer polynomial extremal supplied by~\cite{outer2026}. The proof below distinguishes those imported ingredients from the identities established here.

The first four limits are obtained from the primary low-order constants in~\cite{brevig2020}:
\begin{center}
\begin{tabular}{cl}
\toprule Budget $k$ & Exact limiting structured error\\\midrule
1 & $1$\\
2 & $\sqrt3/2$\\
3 & $3/4$\\
4 & $\sqrt{1153/[2(1103+33\sqrt{33})]}$\\\bottomrule
\end{tabular}
\end{center}
Remark 2 of~\cite{rational2026} concerns degree-one rational approximation of $z^N$. Its value at $N=2$ also equals $\sqrt3/2$, but this is a different approximation problem from $E_1$, where the target is $z$. The contribution here is the precise word-Hankel boundary reduction and coefficient-norm identification.

\section{Exact geometry of the target}
Let $s=\sqrt q$. Define vectors in separate input and output copies of $\BHH$ by
\[
u_j=s\sum_{n\ge0}r^n e_{a^n b^j},\quad 0\le j\le k,
\qquad v_{r,k}=s\sum_{n\ge0}r^n e_{a^n b^k}.
\]
All these normalized vectors with distinct $j$ have disjoint supports. Splitting the possible column words gives
\begin{equation}\label{I:B:eq:target-svd}
H_{f_{r,k}}=u_0v_{r,k}^*+s\sum_{j=0}^{k-1}u_{k-j}e_{b^j}^*.
\end{equation}
The input vectors $v_{r,k},e_\Beps,e_b,\ldots,e_{b^{k-1}}$ are orthonormal, as are the corresponding output vectors. Thus
\[
\Bnorm{H_{f_{r,k}}}=1,\quad \Brank H_{f_{r,k}}=k+1,\quad
\sigma_1=1,\quad\sigma_2=\cdots=\sigma_{k+1}=s,
\quad u_0^*H_{f_{r,k}}=v_{r,k}^*.
\]
Equation~\eqref{I:B:eq:target-svd} proves boundedness as well as the spectral assertions. It also exhibits the difference between the unstructured tail, which tends to zero, and the positive structured boundary limit in Theorem~\ref{I:B:thm:main}.

\section{A lower bound for all complex competitors}
\label{I:B:sec:complex-lower}

We reduce the boundary obstruction to a finite-dimensional one-letter
extremal problem. All coefficients and comparison realizations in this
section may be complex. The alphabet may contain additional letters besides
$a,b$; their nonnegative energy terms can be discarded in the estimates.

\subsection{A one-letter distance and its passive realization formula}

Identify $\BHtwo$ with $\ell^2(\mathbb N_0)$ by Taylor coefficients, and let
$S^*$ denote its backward shift. For an integer $m\ge0$, define
\[
 \mathcal R_m=\left\{h\in\BHtwo:
 \dim\operatorname{span}\{(S^*)^j h:j\ge0\}\le m\right\},
 \qquad
 E_m=\inf_{h\in\mathcal R_m}\Bnorm{z^m-h}_{H^2}.
\]
Thus $\mathcal R_0=\{0\}$ and $E_0=1$. A finite backward-shift cyclic
span is precisely the scalar stable rational-series condition with its
dimension equal to the one-letter Hankel rank. To see boundedness directly,
restrict $S^*$ to this finite invariant span. It is a contraction with no
unit-circle eigenvector, since the corresponding geometric sequence would
not be square summable. Its spectral radius is therefore below one, and
the generated column sequence defines a bounded finite-rank Hankel operator.

For $m\ge1$, also define
\begin{equation}
 \Gamma_m=
 \max_{\substack{D\in\BC^{m\times m},\ l\in\BC^{1\times m}\\
                 D^*D+l^*l\preceq I_m}}
              \Bnorm{lD^m},
 \qquad \Gamma_0=0.
 \label{I:B:eq:passive-gamma}
\end{equation}
The feasible set is compact and the objective continuous, so the maximum
exists. Smaller state spaces are included by padding with zero coordinates.

\begin{lemma}[One-letter projection identity]
\label{I:B:lem:one-letter-projection}
For every $m\ge0$,
\begin{equation}
 E_m^2=1-\Gamma_m^2.
 \label{I:B:eq:E-gamma}
\end{equation}
The infimum defining $E_m$ is attained. Moreover,
$\kappa_m=\Gamma_m^2=1-E_m^2$ for $m\ge1$, with $\kappa_m$ from Lemma~\ref{I:M:rev13:lem:kappa}.
The projection and attainment argument also holds over $\mathbb R$, with field-specific $E_m^{\mathbb R}$ and $\Gamma_m^{\mathbb R}$.
\end{lemma}

\begin{proof}
The statement is immediate for $m=0$. Let $m\ge1$, and let $K$ range over
the finite-dimensional $S^*$-invariant subspaces of $\BHtwo$ with dimension
at most $m$. Every $h\in\mathcal R_m$ lies in such a space, and every
vector in such a space belongs to $\mathcal R_m$. Consequently
\begin{equation}
 E_m^2=1-\sup_K\Bnorm{P_Kz^m}_{H^2}^2.
 \label{I:B:eq:E-model-projection}
\end{equation}
On a particular $K$, set $D=S^*|_K$ and $lx=x(0)$. The identity
$D^*D+l^*l=I_K$ holds, and
\[
 \Bnorm{P_Kz^m}_{H^2}
   =\sup_{x\in K,\,\Bnorm{x}_{H^2}\le1}|x_m|
   =\Bnorm{lD^m}.
\]
Thus the supremum in \eqref{I:B:eq:E-model-projection} is at most $\Gamma_m$.

Conversely, any passive tuple in \eqref{I:B:eq:passive-gamma} has a contraction
observation map
\[
 O:\BC^m\longrightarrow\BHtwo,\qquad
 (Ox)(z)=\sum_{j\ge0}lD^jx\,z^j.
\]
Indeed, telescoping $D^*D+l^*l\preceq I_m$ bounds every finite output
energy by $\Bnorm{x}^2$. Its range $K=\operatorname{ran}O$ is closed,
finite dimensional and backward-shift invariant, since $S^*O=OD$.
For every unit vector $x$,
\[
 |lD^mx|=|(Ox)_m|
 \le\Bnorm{P_Kz^m}_{H^2}\Bnorm{Ox}_{H^2}
 \le\Bnorm{P_Kz^m}_{H^2}.
\]
This proves the reverse inequality and hence \eqref{I:B:eq:E-gamma}. The spaces used here are exactly the finite backward-shift invariant model spaces in Lemma~\ref{I:M:rev13:lem:kappa}; spaces of smaller dimension may be enlarged, proving $\kappa_m=\Gamma_m^2$. Over the real field the same observation-map argument uses real invariant spaces and real tuples.
Choose a maximizing tuple in \eqref{I:B:eq:passive-gamma}; its associated
$K$ has projection norm exactly $\Gamma_m$. Therefore $P_Kz^m$ attains
$E_m$. Any unobservable unitary part of $D$ disappears in $O$, so no strict
stability assumption on the feasible tuple is needed.
\end{proof}


\begin{lemma}[An attained minimizer uses its full rank budget]
\label{I:B:lem:exact-rank-minimizer}
If $m\ge1$, every minimizer of $E_m$ has backward-shift cyclic dimension
exactly $m$. In particular $E_m>0$.
\end{lemma}

\begin{proof}
Let $h$ be a minimizer and $K_h$ its cyclic span. Projection onto $K_h$
is admissible, so minimality forces $h=P_{K_h}z^m$. The residual
$R=z^m-h$ is nonzero, because $z^m$ itself has cyclic dimension $m+1$.
If $\dim K_h<m$, choose $\lambda\in\mathbb D$ with $R(\lambda)\ne0$.
The reproducing kernel $k_\lambda(z)=(1-\overline\lambda z)^{-1}$ is
a backward-shift eigenvector. The invariant space
$K_h+\operatorname{span}\{k_\lambda\}$ has dimension at most $m$,
and projection onto it strictly improves the approximation, since
$R\perp K_h$ but $\langle R,k_\lambda\rangle\ne0$. This is a
contradiction. Attainment and $h\ne z^m$ also give $E_m>0$.
\end{proof}


\subsection{The exact resolvent defect inequality}

Let $G$ be any bounded word-Hankel operator with rank at most $k$, and
let $\mathcal S_G=\operatorname{ran}G^*$ be its input row space.
For the left creation operators $L_c e_w=e_{cw}$, the Hankel relation
implies that $\mathcal S_G$ is invariant under $L_c^*$. Set
\[
 A=L_a^*|_{\mathcal S_G},\qquad
 B=L_b^*|_{\mathcal S_G},\qquad
 \ell x=x[\Beps].
\]
The partition of nonempty words according to first letter gives
\begin{equation}
 A^*A+B^*B+\ell^*\ell\preceq I_{\mathcal S_G}.
 \label{I:B:eq:boundary-passivity}
\end{equation}
The inequality becomes equality when the alphabet is exactly $\{a,b\}$.

Write $q=1-r^2$. For every contraction $A$,
\begin{equation}
 (I-rA^*)(I-rA)-q(I-A^*A)
       =(A-rI)^*(A-rI)\succeq0.
 \label{I:B:eq:resolvent-defect-exact}
\end{equation}
Since $B^*B\preceq I-A^*A$, congruence by $(I-rA)^{-1}$ yields
\begin{equation}
 \Bnorm{\sqrt q\,B(I-rA)^{-1}}\le1.
 \label{I:B:eq:resolvent-defect-contraction}
\end{equation}
This is a joint passivity consequence; it does not follow merely from
separate spectral-radius bounds on the letter matrices.

\subsection{Compactness at the boundary}

\begin{theorem}[Boundary lower bound over all complex competitors]
\label{I:B:thm:all-k-boundary-lower}
For every fixed integer $k\ge1$,
\begin{equation}
 \liminf_{r\uparrow1}\delta_k(H_{f_{r,k}})
          \ge E_{k-1}=\sqrt{1-\Gamma_{k-1}^2}.
 \label{I:B:eq:all-k-boundary-lower}
\end{equation}
\end{theorem}

\begin{proof}
Let
\[
 v_{r,k}=\sqrt q\sum_{n\ge0}r^n e_{a^n b^k},\qquad
 u_{r,0}=\sqrt q\sum_{n\ge0}r^n e_{a^n}.
\]
The target's singular decomposition gives
$\Bnorm{v_{r,k}}=\Bnorm{u_{r,0}}=1$ and
$u_{r,0}^*H_{f_{r,k}}=v_{r,k}^*$. If $P_G$ projects onto
$\mathcal S_G$, then $G(I-P_G)=0$, and therefore
\begin{equation}
 \Bnorm{H_{f_{r,k}}-G}
   \ge\sqrt{1-\Bnorm{P_Gv_{r,k}}^2}.
 \label{I:B:eq:boundary-kernel-lower}
\end{equation}
The coordinates of $x\in\mathcal S_G$ satisfy
$x[a^n b^k]=\ell B^kA^nx$. Hence
\begin{equation}
 \Bnorm{P_Gv_{r,k}}
 =\Bnorm{\sqrt q\,\ell B^k(I-rA)^{-1}}
 \le\Bnorm{\ell B^{k-1}},
 \label{I:B:eq:boundary-overlap-bound}
\end{equation}
where the final inequality is \eqref{I:B:eq:resolvent-defect-contraction}.

Consider any sequence $r_j\uparrow1$ and arbitrary rank-at-most-$k$
competitors $G_j$. Zero-rank terms have zero overlap. For all remaining
terms, choose orthonormal coordinates on their row spaces and pad by
zero coordinates to dimension $k$. The padded tuples still satisfy
\eqref{I:B:eq:boundary-passivity}. Passing to a subsequence gives
$A_j\to A_0$, $B_j\to B_0$, $\ell_j\to\ell_0$.

If $I-A_0$ is invertible, the resolvent in
\eqref{I:B:eq:boundary-overlap-bound}'s middle expression stays bounded,
so that overlap tends to zero. Otherwise
$\mathcal E=\ker(I-A_0)$ is nonzero and reduces the contraction $A_0$.
Indeed $A_0x=x$ and contraction imply $A_0^*x=x$.
The limiting passive inequality forces $B_0\mathcal E=0$ and
$\ell_0\mathcal E=0$. In the orthogonal decomposition
$\mathcal E\oplus\mathcal F$ one can consequently write
\[
 B_0=\begin{pmatrix}0&C\\0&D\end{pmatrix},\qquad
 \ell_0=(0,l),\qquad \dim\mathcal F\le k-1.
\]
The $\mathcal F$-column part of passivity gives
$D^*D+l^*l\preceq I_{\mathcal F}$. If $k\ge2$,
$\ell_0B_0^{k-1}=(0,lD^{k-1})$ has norm at most $\Gamma_{k-1}$.
If $k=1$, $\mathcal F=\{0\}$ and $\ell_0=0$, giving the same conclusion
with $\Gamma_0=0$.

Thus every subsequential limiting overlap is at most $\Gamma_{k-1}$.
A compactness contradiction makes this uniform over all competitors:
\[
 \limsup_{r\uparrow1}\sup_{\substack{G\text{ bounded word-Hankel}\\
                                  \operatorname{rank}G\le k}}
       \Bnorm{P_Gv_{r,k}}\le\Gamma_{k-1}.
\]
Insert this bound into \eqref{I:B:eq:boundary-kernel-lower} and use
Lemma~\ref{I:B:lem:one-letter-projection}. No existence assumption for a
best word-Hankel competitor is required.
\end{proof}


For $k=2$, the passive problem has one state:
\[
 \Gamma_1=\max_{|d|^2+|l|^2\le1}|ld|=\frac12.
\]
The theorem therefore already gives
$\liminf_{r\uparrow1}\delta_2(H_{f_{r,2}})\ge\sqrt3/2$ without any
Hardy coefficient theorem.

\section{An explicit upper construction}
\label{I:B:sec:lift}

The lower reduction isolates a scalar one-letter approximation problem.
The construction below gives the matching direction without an abstract
extension theorem. It also controls the error before taking the boundary
limit.

\begin{theorem}[A finite-rank lift with an exact orthogonal split]
\label{I:B:thm:upper-lift}
Fix an integer $k\ge1$, put $m=k-1$, and let $0<r<1$,
$q=1-r^2$, and $s=\sqrt q$. Let $h=(h_j)_{j\ge0}$ have a bounded
scalar Hankel operator
\[
 H_h[i,j]=h_{i+j},\qquad i,j\ge0,
 \qquad \operatorname{rank}H_h\le m.
\]
For $m=0$ take $h=0$. Define the word series
\begin{equation}
 g_{r,h}(a^n b^j)=q r^n h_{j-1}
 \quad(n\ge0,\ j\ge1),\qquad
 g_{r,h}(w)=0\quad\text{otherwise}.
 \label{I:B:eq:lift-series}
\end{equation}
Then $H_{g_{r,h}}$ is bounded, has rank at most $m+1=k$, and
\begin{equation}
 \boxed{\displaystyle
 \Bnorm{H_{f_{r,k}}-H_{g_{r,h}}}^2
 \le \Bnorm{z^m-h}_{\BHtwo}^2
       +q\Bnorm{H_{z^m}-H_h}^2.}
 \label{I:B:eq:lift-finite-error}
\end{equation}
Consequently, if
\[
 E_m=\inf_{\substack{H_h\text{ bounded scalar Hankel}\\
                      \operatorname{rank}H_h\le m}}
        \Bnorm{z^m-h}_{\BHtwo},
\]
then
\begin{equation}
 \limsup_{r\uparrow1}\delta_k(H_{f_{r,k}})\le E_{k-1}.
 \label{I:B:eq:lift-limsup}
\end{equation}
All scalar coefficients and competing realizations may be complex.
\end{theorem}

\begin{proof}
If $h\ne0$, choose a scalar realization
$h_j=\alpha^*D^j\beta$ of dimension $n\le m$. The block matrices
\begin{equation}
 A_a=\begin{pmatrix}r&0\\0&0_n\end{pmatrix},\qquad
 A_b=\begin{pmatrix}0&\alpha^*\\0&D\end{pmatrix},\qquad
 \alpha_g^*=q(1,0_n),\qquad
 \beta_g=\begin{pmatrix}0\\\beta\end{pmatrix}
 \label{I:B:eq:lift-realization}
\end{equation}
give a realization of \eqref{I:B:eq:lift-series} in the convention
$g(a_1\cdots a_t)=\alpha_g^*A_{a_1}\cdots A_{a_t}\beta_g$.
Indeed $A_bA_a=0$, so a letter $a$ after a letter $b$ forces the
coefficient to vanish; the empty coefficient also vanishes. Direct block
multiplication gives the stated coefficients on $a^*b^+$. All additional
letter matrices are zero if the alphabet is larger. The case $h=0$ gives
the zero series. Thus the realization dimension is at most $m+1$.

We prove boundedness and the error estimate directly in the infinite
operator norm. Write $\ell^2(\mathbb N_0)$ for the scalar coefficient
space. Define two isometries from that space to $\BHH$ by
\begin{equation}
 V_r e_j=s\sum_{n\ge0}r^n e_{a^n b^{j+1}},\qquad
 W e_j=e_{b^j},\qquad j\ge0,
 \label{I:B:eq:lift-isometries}
\end{equation}
and define the unit vector
\[
 u_r=s\sum_{n\ge0}r^n e_{a^n}.
\]
The geometric normalization gives $V_r^*V_r=W^*W=I$, and
$u_r\perp\operatorname{ran}V_r$. Boundedness of $H_h$ implies
$h=H_he_0\in\ell^2(\mathbb N_0)$. With $\overline h$ denoting
coefficientwise conjugation, exact entry matching yields
\begin{equation}
 H_{g_{r,h}}
   =u_r(V_r\overline h)^*+sV_rH_hW^*.
 \label{I:B:eq:lift-exact-decomposition}
\end{equation}
The conjugation in the first term is necessary. At row $a^n$ and column
$a^t b^j$, that term equals $q r^{n+t}h_{j-1}$, rather than its
conjugate. At row $a^n b^j$ with $j\ge1$ and column $b^\ell$, the
second term equals $q r^n h_{j+\ell-1}$. These account for every
nonzero entry. In particular \eqref{I:B:eq:lift-exact-decomposition} proves
boundedness, so the finite realization gives a bounded word-Hankel
operator of rank at most $k$.

For the scalar polynomial $z^m$, its coefficient vector is $e_m$, and
the same identity reads
\begin{equation}
 H_{f_{r,k}}=u_r(V_re_m)^*+sV_rH_{z^m}W^*.
 \label{I:B:eq:lift-target-decomposition}
\end{equation}
Subtracting \eqref{I:B:eq:lift-exact-decomposition} from
\eqref{I:B:eq:lift-target-decomposition}, the two summands have orthogonal
output ranges. For each input $x\in\BHH$ this gives
\begin{align*}
 \Bnorm{(H_{f_{r,k}}-H_{g_{r,h}})x}^2
 &=\bigl|\langle V_r(e_m-\overline h),x\rangle\bigr|^2
   +q\Bnorm{V_r(H_{z^m}-H_h)W^*x}^2\\
 &\le\left(\Bnorm{e_m-\overline h}_2^2
           +q\Bnorm{H_{z^m}-H_h}^2\right)\Bnorm{x}^2.
\end{align*}
Since $e_m$ has real coefficients,
$\Bnorm{e_m-\overline h}_2=\Bnorm{e_m-h}_2
=\Bnorm{z^m-h}_{\BHtwo}$. This proves
\eqref{I:B:eq:lift-finite-error}.

For every fixed admissible $h$, the second term on the right of
\eqref{I:B:eq:lift-finite-error} tends to zero as $r\uparrow1$. First choose
$h$ within any positive tolerance of $E_m$, next take the boundary limit,
and finally let the tolerance decrease. This proves
\eqref{I:B:eq:lift-limsup} without assuming attainment of $E_m$.
\end{proof}


\subsection{A real two-state competitor and a sharper finite-parameter bound}
\label{I:B:subsec:lift-k2}

For $k=2$, the scalar one-state problem is explicit. Set
\begin{equation}
 y=\frac1{\sqrt2},\qquad c=y(1-y^2)=\frac1{2\sqrt2},\qquad
 h_j=cy^j.
 \label{I:B:eq:lift-k2-parameters}
\end{equation}
The lifted series is
\begin{equation}
 g_r(a^n b^j)=q c r^n y^{j-1}
 \quad(n\ge0,\ j\ge1),\qquad g_r(w)=0\quad\text{otherwise}.
 \label{I:B:eq:lift-k2-series}
\end{equation}
It has the real two-state realization
\begin{equation}
 A_a=\begin{pmatrix}r&0\\0&0\end{pmatrix},\qquad
 A_b=\begin{pmatrix}0&1\\0&y\end{pmatrix},\qquad
 \alpha_g^*=qc(1,0),\qquad\beta_g=(0,1)^{\mathsf T}.
 \label{I:B:eq:lift-k2-realization}
\end{equation}
The following stronger bound shows that this particular competitor
approaches its limiting error at order $q^2$.

\begin{proposition}[A quantitative two-state upper bound]
\label{I:B:prop:lift-k2-sharp}
For $g_r$ in \eqref{I:B:eq:lift-k2-series} and $0<q\le1/8$,
\begin{equation}
 \boxed{\displaystyle
 \frac34+\frac{q^2}{6}
 \le\Bnorm{H_{f_{r,2}}-H_{g_r}}^2
 \le\frac34+\frac{q^2}{6}+\frac43q^3.}
 \label{I:B:eq:lift-k2-sharp-bound}
\end{equation}
In particular,
\begin{equation}
 \delta_2(H_{f_{r,2}})^2
 \le\frac34+\frac{q^2}{6}+\frac43q^3,
 \qquad
 \limsup_{r\uparrow1}\delta_2(H_{f_{r,2}})\le\frac{\sqrt3}{2}.
 \label{I:B:eq:lift-k2-limsup}
\end{equation}
The lower inequality in \eqref{I:B:eq:lift-k2-sharp-bound} concerns this
explicit competitor only, rather than the structured optimum.
\end{proposition}

\begin{proof}
Use the orthonormal vectors
\[
 u=s\sum_{n\ge0}r^n e_{a^n},\qquad
 v_j=s\sum_{n\ge0}r^n e_{a^n b^j}\quad(j\ge1),
\]
and put
\[
 v_y=\sum_{j\ge1}y^{j-1}v_j,\qquad
 j_y=\sum_{\ell\ge0}y^\ell e_{b^\ell},\qquad e=e_\Beps.
\]
Their relevant norms and inner products are
\[
 \Bnorm{v_y}^2=\Bnorm{j_y}^2=(1-y^2)^{-1},\qquad
 \langle v_2,v_y\rangle=y.
\]
The exact target and competitor decompositions become
\begin{align*}
 H_{f_{r,2}}&=uv_2^*+s(v_2e^*+v_1e_b^*),\\
 H_{g_r}&=cu v_y^*+scv_yj_y^*.
\end{align*}
Define
\[
 d=v_2-cv_y,\qquad
 R=v_2e^*+v_1e_b^*-cv_yj_y^*,\qquad A=\Bnorm{d}^2.
\]
For the parameters \eqref{I:B:eq:lift-k2-parameters},
\[
 A=1-2cy+\frac{c^2}{1-y^2}=\frac34.
\]
The output vector $u$ is orthogonal to the range of $R$. Therefore, for
$E=H_{f_{r,2}}-H_{g_r}$,
\begin{equation}
 E=ud^*+sR,\qquad M:=E^*E=dd^*+qR^*R.
 \label{I:B:eq:lift-k2-orthogonal-error}
\end{equation}
Moreover,
\begin{equation}
 \Bnorm R\le1+\frac{c}{1-y^2}=1+y,
 \qquad \Bnorm R^2\le(1+y)^2<3.
 \label{I:B:eq:lift-k2-Rnorm}
\end{equation}
Here the first two summands in $R$ form a partial isometry of norm one.

The leading error vector has the useful cancellations
\begin{align*}
 e^*d&=0,& e_b^*d&=-sc,\\
 j_y^*d&=s\left(y^2-\frac{cy}{1-y^2}\right)=0.
\end{align*}
Consequently
\begin{equation}
 Rd=-scv_1.
 \label{I:B:eq:lift-k2-cancellation}
\end{equation}
Let $z=d/\sqrt A$ and decompose $M$ with respect to
$\operatorname{span}\{z\}\oplus z^\perp$. Its scalar upper-left entry
is
\begin{equation}
 a=\langle z,Mz\rangle
   =A+q\Bnorm{Rz}^2
   =\frac34+\frac{q^2c^2}{A}
   =\frac34+\frac{q^2}{6}.
 \label{I:B:eq:lift-k2-Rayleigh}
\end{equation}
This already proves the lower bound for the explicit competitor.
Let $B$ be the compression of $M$ to $z^\perp$, and let
$v=(I-zz^*)Mz$ be its off-diagonal vector. Since $dd^*$ vanishes on
$z^\perp$, \eqref{I:B:eq:lift-k2-Rnorm} and
\eqref{I:B:eq:lift-k2-cancellation} give
\[
 \Bnorm B\le q\Bnorm R^2<3q,
 \qquad
 \Bnorm v^2\le q^2\Bnorm R^2\Bnorm{Rz}^2
             <\frac{q^3}{2}.
\]
For $q\le1/8$ the spectral gap in this block decomposition obeys
\[
 a-\Bnorm B\ge\frac34-3q\ge\frac38.
\]
The elementary positive-block estimate
\[
 \lambda_{\max}(M)
 \le a+\frac{\Bnorm v^2}{a-\Bnorm B}
\]
therefore yields the upper inequality in
\eqref{I:B:eq:lift-k2-sharp-bound}. For completeness, this block estimate
follows from the Schur complement when the largest eigenvalue is above
$a$: if it is $\lambda$, then
$\lambda-a\le\Bnorm v^2/(\lambda-\Bnorm B)
\le\Bnorm v^2/(a-\Bnorm B)$. If $v=0$, the estimate is immediate.
\end{proof}


Thus the explicit upper competitor has the expansion
\[
 \Bnorm{H_{f_{r,2}}-H_{g_r}}
 =\frac{\sqrt3}{2}+\frac{q^2}{6\sqrt3}+O(q^3).
\]
This expansion is not asserted for the optimal structured error; the
matching boundary value follows by combining
\eqref{I:B:eq:lift-k2-limsup} with the independently proved lower bound.

\section{The exact Hardy coefficient constant}
\label{I:B:sec:hardy}

All scalar functions and approximation classes in this section are complex.
The circle carries normalized measure $d\theta/(2\pi)$.  For $0<p<\infty$,
the Hardy space $H^p$ consists of analytic functions on the unit disc for which
\[
 \Bnorm{F}_{H^p}
 =\sup_{0<\rho<1}\left(\int_0^{2\pi}
       |F(\rho e^{i\theta})|^p\,\frac{d\theta}{2\pi}\right)^{1/p}<\infty.
\]
We write $\BHtwo=H^2$ and identify it isometrically with the square-summable
Taylor coefficient sequences.  For $m\ge0$ define the coefficient-functional
norm
\begin{equation}
 C(m,p)=\sup\bigl\{|[z^m]F|:\Bnorm{F}_{H^p}=1\bigr\}.
 \label{I:B:eq:hardy-C-definition}
\end{equation}
This is the convention of Brevig--Saksman~\cite{brevig2020}; in particular,
it is not a norm of a finite Hankel block.  Point evaluation gives $C(0,p)=1$.

\subsection{The scalar rational problem and finite model spaces}

Let $\mathcal R_m$ be the set of analytic functions $h(z)=\sum_{j\ge0}h_jz^j$
whose coefficient Hankel $(h_{i+j})_{i,j\ge0}$ is bounded and has rank at most
$m$.  Set $\mathcal R_0=\{0\}$.  Equivalently, for $m\ge1$ this is the class
of stable rational functions
\begin{equation}
 h=\frac{N}{D},\qquad \deg N<m,\quad \deg D\le m,\quad D(0)=1,
 \label{I:B:eq:hardy-rational-class}
\end{equation}
after cancellation, with $D$ zero-free on the closed disc.  Polynomial parts
are included: a polynomial of degree $m-1$ can have coefficient-Hankel rank
$m$.  Omitting these parts would change the approximation class.  Define
\begin{equation}
 E_m=\inf_{h\in\mathcal R_m}\Bnorm{z^m-h}_{\BHtwo}.
 \label{I:B:eq:hardy-E-definition}
\end{equation}

For a finite Blaschke product $\Theta$, let
$K_\Theta=\BHtwo\ominus\Theta\BHtwo$.  Its degree is its number of zeros
inside the disc, counted with multiplicity.  A product of degree $l$ can be
written, up to a unimodular constant, as
\[
 \Theta=\frac{D_l^*}{D},\qquad
 D_l^*(z)=z^l\overline{D(1/\bar z)},\quad \deg D\le l,
\]
where $D$ is zero-free on the closed disc.  The reciprocal operation uses
the specified order $l$, even when $\deg D<l$.

\begin{lemma}[Classical finite-model-space formulation]
\label{I:B:lem:hardy-model}
The class $\mathcal R_m$ is the union of the spaces $K_\Theta$ for finite
Blaschke products of degree at most $m$.  If
$\Theta(z)=\sum_{j\ge0}\theta_jz^j$, then
\begin{equation}
 \Bnorm{z^m-P_{K_\Theta}z^m}_{\BHtwo}^2
 =\sum_{j=0}^m|\theta_j|^2.
 \label{I:B:eq:hardy-model-error}
\end{equation}
Consequently,
\begin{equation}
 E_m^2=\inf_{\deg\Theta\le m}\sum_{j=0}^m|\theta_j|^2.
 \label{I:B:eq:hardy-Blaschke-variational}
\end{equation}
\end{lemma}
\begin{proof}
For degree $l$, the standard rational description is
\[
 K_\Theta=\left\{\frac{N}{D}:\deg N<l\right\}.
\]
It follows, for example, from the $l$-element Takenaka--Malmquist orthonormal
basis associated with the zeros of $\Theta$: the common denominator is $D$,
the numerators have degree less than $l$, and dimension is $l$.  This
space is invariant under the scalar backward shift.  Hence the Hankel
columns of each member span a space of dimension at most $l$.  Conversely,
the stable rational realization of a bounded finite-rank scalar Hankel
has the form \eqref{I:B:eq:hardy-rational-class}, at its minimal order; its
denominator supplies such a finite Blaschke product.  This proves the union
assertion, including denominators of degree below the specified order.

Multiplication by $\Theta$ is an isometry on $\BHtwo$, so the orthogonal
projection onto $\Theta\BHtwo$ is
\[
 z^m-P_{K_\Theta}z^m
 =\Theta P_+(\overline\Theta z^m)
 =\Theta\sum_{j=0}^m\overline{\theta_j}z^{m-j}.
\]
Taking the squared norm proves \eqref{I:B:eq:hardy-model-error}, and minimizing
over the union proves \eqref{I:B:eq:hardy-Blaschke-variational}.
\end{proof}


This variational formulation belongs to classical rational approximation;
see also Theorem~1 and the monomial specialization in
Baratchart et al.~\cite{rational2026}.  We give the argument to fix the
rank and polynomial-part conventions precisely.

\subsection{A sharp coefficient-energy identity}

\begin{theorem}[Exact identification with the Hardy coefficient norm]
\label{I:B:thm:hardy-bridge}
For every integer $m\ge0$,
\begin{equation}
 \boxed{\displaystyle E_m=\frac1{C(m,2/3)}.}
 \label{I:B:eq:hardy-exact-bridge}
\end{equation}
The lower bound $E_m\ge1/C(m,2/3)$ is self-contained. Equality for $m\le3$ follows from the published low-degree outer extremals in~\cite{brevig2020}. For general $m$, the matching upper construction uses the outer-extremal theorem of Brevig--S\o vik~\cite{outer2026}.
\end{theorem}
\begin{proof}
The case $m=0$ is immediate, so assume $m\ge1$ and abbreviate
$C=C(m,2/3)$.

\emph{Lower bound, by a direct H\"older argument.}
First take $\Theta$ of degree exactly $m$.  Multiplying by a unimodular
constant does not affect \eqref{I:B:eq:hardy-model-error}, so write
\[
 \Theta=\frac{D_m^*}{D},\qquad
 D_m^*(z)=z^m\overline{D(1/\bar z)}.
\]
Let $b=P_{\le m}\Theta=\sum_{j=0}^m\theta_jz^j$.  Since Taylor truncation
does not affect the coefficient of degree $m$ after multiplication by an
analytic polynomial,
\begin{equation}
 [z^m](bD^2)=[z^m](\Theta D^2)
 =[z^m](D_m^*D)=\Bnorm{D}_{\BHtwo}^2.
 \label{I:B:eq:hardy-key-coefficient}
\end{equation}
The definition of $C$ and H\"older's inequality therefore give
\begin{equation}
 \Bnorm{D}_{\BHtwo}^2
 \le C\Bnorm{bD^2}_{H^{2/3}}
 \le C\Bnorm{b}_{\BHtwo}\Bnorm{D}_{\BHtwo}^2.
 \label{I:B:eq:hardy-holder-bound}
\end{equation}
For clarity, the second step is precisely
\[
 \left(\int |b|^{2/3}|D|^{4/3}\right)^{3/2}
 \le\left(\int |b|^2\right)^{1/2}\int |D|^2,
\]
using H\"older exponents $3$ and $3/2$ and normalized circle measure.
Division by $\Bnorm{D}_{\BHtwo}^2>0$ proves
$\Bnorm{b}_{\BHtwo}\ge1/C$.

If $\deg\Theta=l<m$, append any finite Blaschke factor $\Psi$ of degree
$m-l$.  The inclusion $K_\Theta\subset K_{\Theta\Psi}$ shows that the
projection error for $\Theta$ is at least that for the degree-$m$ product
$\Theta\Psi$.  Thus the same lower bound holds for every degree at most
$m$.  Lemma~\ref{I:B:lem:hardy-model} gives $E_m\ge1/C$.  This part of the
proof uses neither stationarity of a rational minimizer nor existence of
such a minimizer.

\emph{Matching upper bound, from an outer coefficient extremal.}
Coefficient extremals exist by a normal-family argument.  Theorem~1 and
Corollary~3 of Brevig--S\o vik~\cite{outer2026}, combined with the
polynomial structure in Brevig--Saksman~\cite{brevig2020}, imply that an
extremal for \eqref{I:B:eq:hardy-C-definition} at $p=2/3$ is $F=P^3$, where
$P$ is an outer polynomial of exact degree $m$.  Normalize its phase and
norm so that
\begin{equation}
 \Bnorm{P}_{\BHtwo}=1,\qquad [z^m]P^3=C>0.
 \label{I:B:eq:hardy-extremal-normalization}
\end{equation}
For every polynomial $Q$ one has
$\Bnorm{Q^3}_{H^{2/3}}=\Bnorm{Q}_{\BHtwo}^3$.  Hence $P$ maximizes
$\operatorname{Re}[z^m]Q^3$ over the entire unit sphere of the
$(m+1)$-dimensional polynomial Hilbert space; the variations do not need
to remain outer.

Write $P(z)=\sum_{j=0}^mp_jz^j$.  Differentiation with respect to the real
and imaginary parts of all $p_j$, with the unit-sphere constraint, yields
\[
 [z^{m-j}]P^2=\lambda\overline{p_j}\qquad(0\le j\le m).
\]
Multiplying by $p_j$ and summing shows
$\lambda=[z^m]P^3=C$, since $\sum_j|p_j|^2=1$.
Thus, for the reciprocal of order $m$,
\begin{equation}
 P^2=C P_m^*+O(z^{m+1}),\qquad
 P_m^*(z)=z^m\overline{P(1/\bar z)}.
 \label{I:B:eq:hardy-EL}
\end{equation}
The ratio $\Theta=P_m^*/P$ is a finite Blaschke product of degree at most
$m$, after cancelling any common factors on the unit circle.  Outerness
ensures $P(0)\ne0$, so division of \eqref{I:B:eq:hardy-EL} as a Taylor series
at zero gives
\[
 \Theta=P/C+O(z^{m+1}).
\]
Consequently,
\[
 \sum_{j=0}^m|\theta_j|^2
 =\frac{\Bnorm{P}_{\BHtwo}^2}{C^2}=\frac1{C^2}.
\]
The projection $h=P_{K_\Theta}z^m$ belongs to $\mathcal R_m$ and attains
this error.  Therefore $E_m\le1/C$, completing the proof.
\end{proof}


\begin{corollary}[A real optimizer may be chosen]
\label{I:B:cor:hardy-real}
The infimum defining $E_m$ is unchanged if $h$ is required to have real
Taylor coefficients.  In the upper construction one may choose $P$ with
nonnegative real coefficients.
\end{corollary}
\begin{proof}
Take a normalized complex extremal polynomial $P=\sum_jp_jz^j$ as above
and put $Q=\sum_j|p_j|z^j$.  Then $\Bnorm{Q}_{\BHtwo}=1$ and
\[
 [z^m]Q^3\ge\bigl|[z^m]P^3\bigr|=C(m,2/3).
\]
On the other hand, $\Bnorm{Q^3}_{H^{2/3}}=1$, so the coefficient inequality
forces equality.  For $m\le3$, the explicit nonnegative outer extremals in~\cite{brevig2020} give this step directly. For general $m$, the outer-extremal theorem makes $Q$ outer.
The corresponding $\Theta=Q_m^*/Q$ has real coefficients, and the formula
in Lemma~\ref{I:B:lem:hardy-model} makes $P_{K_\Theta}z^m$ real as well.
This real function attains the complex lower bound.
\end{proof}


\section{Completion of the theorem and algebraic construction}
The lower-bound and lift arguments may each be carried out over either scalar field: replace the passive tuples, invariant subspaces and scalar minimizers by real ones for the real problem. They give $\lim\delta_k^{\mathbb F}=E_{k-1}^{\mathbb F}$ independently of the outer-extremal input. In particular, the lower-bound section proves $\liminf\delta_k^{\BC}\ge E_{k-1}$ over the full complex comparison class. The lift section proves $\limsup\delta_k^{\BC}\le E_{k-1}$. The Hardy section identifies $E_{k-1}$ with the reciprocal coefficient constant. This proves the complex assertion in Theorem~\ref{I:B:thm:main}.

The Hardy extremal polynomial may be taken to have nonnegative real coefficients, as proved in the Hardy section. Its finite model-space projection of $z^m$ then has real coefficients, and the lifted word series is real. Therefore the same upper bound holds for $\delta_k^{\BR}$. Since $\delta_k^{\BC}\le\delta_k^{\BR}$, both limiting constants coincide. This proves the full theorem.

\begin{corollary}[Finite algebraic description of the limiting constant]
For each fixed $m\ge0$,
\[
C(m,2/3)=
\max_{\substack{p_0,\ldots,p_m\ge0\\\sum_{j=0}^m p_j^2=1}}
\sum_{i+j+\ell=m}p_i p_j p_\ell.
\]
Consequently $C(m,2/3)$ and its reciprocal are real algebraic numbers. Real quantifier elimination can in principle compute them and an algebraic optimal polynomial. This is exact decidability, not a polynomial-time claim.
\end{corollary}
\begin{proof}
Every polynomial cube in the displayed unit sphere has $H^{2/3}$ norm one, so its coefficient is at most $C(m,2/3)$. The outer-extremal theorem and the nonnegative-coefficient reduction supply a polynomial attaining equality. The feasible set is compact semialgebraic over $\mathbb Q$, and the objective is polynomial. Its maximum and algebraic witnesses can be obtained by real-algebraic decision and witness extraction.
\end{proof}


For a normalized extremal $P(z)=\sum_{j=0}^m p_jz^j$, let
$P^{\sharp}(z)=z^m\overline{P(1/\bar z)}$, with the reversal order fixed at $m$. The Hardy proof gives
\[
\Theta=P^{\sharp}/P,\qquad
h(z)=z^m-\frac{(P^{\sharp}(z))^2}{C(m,2/3)P(z)}.
\]
This $h$ is the projection of $z^m$ onto $K_\Theta$, so it has one-letter Hankel rank at most $m$. Apparent boundary factors cancel. A minimal stable realization of $h$, followed by the explicit word lift, gives a family with at most $m+1=k$ states and limiting error exactly $1/C(m,2/3)$. This construction optimizes the smaller scalar boundary problem; it does not solve the original word-Hankel optimum for each $r$.

For any fixed attained scalar minimizer $h$, the orthogonal lift bound gives
\[
\delta_k^{\BC}(H_{f_{r,k}})^2
\le E_{k-1}^2+q\Bnorm{H_{z^{k-1}}-H_h}^2.
\]
Since $E_{k-1}>0$, this immediately gives a one-sided $E_{k-1}+O(q)$ upper estimate. The next section strengthens the construction to an $O(q^2)$ upper excess for every fixed $k$. A matching finite-$r$ lower rate for the optimum is not proved.

\section{A quadratic approach to the boundary for the lifted competitors}
\label{I:B:sec:curvature}

The orthogonal split in Theorem~\ref{I:B:thm:upper-lift} contains more
information than its direct triangle estimate. For every fixed scalar
competitor, the squared error of its lift has no term of order $q$.
This is a statement about the explicit lifted construction; it does not
give an asymptotic expansion of the globally optimal structured error.

\begin{proposition}[Curvature of a fixed scalar lift]
\label{I:B:prop:lift-curvature}
Let $h$ satisfy the assumptions of Theorem~\ref{I:B:thm:upper-lift}, put
$m=k-1$, and write $d=z^m-h$ for its scalar coefficient residual.
Assume
\begin{equation}
 A=\Bnorm d_{\BHtwo}^2>0,\qquad
 L=\Bnorm{H_d}^2,\qquad
 T=\Bnorm{H_dS_+\overline d}_2^2,
 \label{I:B:eq:curvature-parameters}
\end{equation}
where $S_+e_j=e_{j+1}$ is the scalar forward shift and
$\overline d$ denotes coefficientwise conjugation. Let
$E_r=H_{f_{r,k}}-H_{g_{r,h}}$. If $qL\le A/2$, then
\begin{equation}
 \boxed{\displaystyle
 A+\frac{q^2T}{A}
 \le\Bnorm{E_r}^2
 \le A+\frac{q^2T}{A}+\frac{2q^3LT}{A^2}.}
 \label{I:B:eq:curvature-finite-bound}
\end{equation}
Consequently, for each fixed admissible $h$ with $A>0$,
\begin{align}
 \Bnorm{E_r}^2&=A+\frac{T}{A}q^2+O(q^3),
 \label{I:B:eq:curvature-squared-expansion}\\
 \Bnorm{E_r}&=\sqrt A+\frac{T}{2A^{3/2}}q^2+O(q^3)
 \qquad(q\downarrow0).
 \label{I:B:eq:curvature-norm-expansion}
\end{align}
The constants may depend on the fixed scalar residual through $A,L,T$.
\end{proposition}

\begin{proof}
Use the isometries $V_r,W$ and the unit vector $u_r$ from
\eqref{I:B:eq:lift-isometries}. Define $v_d=V_r\overline d$. The exact
error decomposition is
\[
 E_r=u_rv_d^*+\sqrt q\,V_rH_dW^*,
 \qquad u_r\perp\operatorname{ran}V_r.
\]
In particular,
\begin{equation}
 M:=E_r^*E_r=v_dv_d^*+qWH_d^*H_dW^*.
 \label{I:B:eq:curvature-error-square}
\end{equation}
The overlap between the two embeddings is especially simple:
\begin{equation}
 W^*V_re_j=\sqrt q\,e_{j+1},\qquad
 W^*v_d=\sqrt q\,S_+\overline d.
 \label{I:B:eq:curvature-overlap}
\end{equation}
Indeed the only term in $V_re_j$ supported on a pure power of $b$ is
its $n=0$ term, at the word $b^{j+1}$.

Let $z=v_d/\sqrt A$, so $\Bnorm z=1$, and decompose $M$ relative to
$\operatorname{span}\{z\}\oplus z^\perp$. By
\eqref{I:B:eq:curvature-overlap}, its upper-left scalar entry is
\begin{equation}
 a=\langle z,Mz\rangle
   =A+q\Bnorm{H_dW^*z}^2
   =A+\frac{q^2T}{A}.
 \label{I:B:eq:curvature-rayleigh}
\end{equation}
This is a Rayleigh quotient of $M$ and proves the lower inequality in
\eqref{I:B:eq:curvature-finite-bound} for the explicit error $E_r$.

Let $B_0$ be the compression of $M$ to $z^\perp$, and let
$v=(I-zz^*)Mz$ be the off-diagonal vector. The rank-one term
$v_dv_d^*$ vanishes on $z^\perp$. Since $W$ is an isometry,
\begin{equation}
 \Bnorm{B_0}\le qL.
 \label{I:B:eq:curvature-complement}
\end{equation}
Moreover, using \eqref{I:B:eq:curvature-error-square} and
\eqref{I:B:eq:curvature-overlap},
\begin{align}
 \Bnorm v^2
 &\le q^2\Bnorm{WH_d^*H_dW^*z}^2\\
 &\le q^2L\Bnorm{H_dW^*z}^2
  =\frac{q^3LT}{A}.
 \label{I:B:eq:curvature-offdiagonal}
\end{align}
The assumed inequality $qL\le A/2$ gives
$a-\Bnorm{B_0}\ge A/2$. Applying the positive-block estimate used in
Proposition~\ref{I:B:prop:lift-k2-sharp} yields
\[
 \Bnorm M\le a+\frac{\Bnorm v^2}{a-\Bnorm{B_0}}
 \le a+\frac{2q^3LT}{A^2},
\]
which is the upper inequality in
\eqref{I:B:eq:curvature-finite-bound}. The two expansions follow immediately;
in taking the square root, the next Taylor correction from the $q^2$
term is of order $q^4$ and is absorbed by the displayed $O(q^3)$.
\end{proof}


\subsection{An autocorrelation formula for Hardy extremal lifts}

For the scalar residual constructed from a Hardy coefficient extremal,
the curvature coefficient can be written directly in terms of its outer
polynomial. This makes the cancellation and its quantitative size
explicit.

\begin{corollary}[Curvature from the outer extremal polynomial]
\label{I:B:cor:hardy-curvature}
Let $m=k-1$, let $C=C(m,2/3)$, and choose an outer coefficient-extremal
polynomial
\[
 P(z)=\sum_{j=0}^{m}p_jz^j,\qquad
 \Bnorm P_{\BHtwo}=1,\qquad [z^m]P^3=C>0.
\]
Use reversal order $m$ in
$P^*(z)=z^m\overline{P(1/\overline z)}$ and let $h$ be the scalar
approximant supplied by the Hardy construction, whose residual is
\begin{equation}
 d=z^m-h=\frac{(P^*)^2}{CP}.
 \label{I:B:eq:curvature-hardy-residual}
\end{equation}
Put $J=\Bnorm P_{H^4}^4$ and $L=\Bnorm{H_d}^2$. Then
\begin{equation}
 A=C^{-2},\qquad
 T=\frac{J-1}{2C^4},\qquad
 \frac TA=\frac{J-1}{2C^2}.
 \label{I:B:eq:curvature-hardy-coefficients}
\end{equation}
For $qL\le1/(2C^2)$, the error of the corresponding explicit word lift
satisfies
\begin{equation}
 \boxed{\displaystyle
 \frac1{C^2}+\frac{J-1}{2C^2}q^2
 \le\Bnorm{E_r}^2
 \le\frac1{C^2}+\frac{J-1}{2C^2}q^2
           +L(J-1)q^3.}
 \label{I:B:eq:curvature-hardy-finite}
\end{equation}
In particular,
\begin{equation}
 \Bnorm{E_r}=\frac1C+\frac{J-1}{4C}q^2+O(q^3),
 \qquad
 \delta_k(H_{f_{r,k}})\le E_{k-1}+O(q^2).
 \label{I:B:eq:curvature-hardy-norm}
\end{equation}
The expansion in \eqref{I:B:eq:curvature-hardy-norm} is for this chosen
lift; the final inequality is only an upper bound for the optimum.
\end{corollary}

\begin{proof}
On the unit circle, $|P^*|=|P|$. Hence
\eqref{I:B:eq:curvature-hardy-residual} gives
\begin{equation}
 |d|^2=\frac{|P|^2}{C^2}\quad\text{almost everywhere}.
 \label{I:B:eq:curvature-boundary-modulus}
\end{equation}
Any boundary zeros of $P$ cancel in the rational expression for $d$;
they do not alter this identity or its coefficient consequences.
Integrating \eqref{I:B:eq:curvature-boundary-modulus} proves $A=C^{-2}$.

The $i$th scalar coefficient of $H_dS_+\overline d$ is
\begin{equation}
 (H_dS_+\overline d)_i
   =\sum_{j\ge0}d_{i+j+1}\overline{d_j},\qquad i\ge0.
 \label{I:B:eq:curvature-autocorrelation-d}
\end{equation}
These sums converge absolutely by Cauchy--Schwarz. They are the positive
Fourier coefficients of $|d|^2$. Applying
\eqref{I:B:eq:curvature-boundary-modulus} therefore gives
\[
 (H_dS_+\overline d)_{i-1}
   =\frac1{C^2}\sum_{j=0}^{m-i}p_{j+i}\overline{p_j}
 \quad(1\le i\le m),
\]
and these coefficients vanish for $i>m$. Consequently,
\begin{equation}
 T=\frac1{C^4}\sum_{i=1}^{m}
       \left|\sum_{j=0}^{m-i}p_{j+i}\overline{p_j}\right|^2.
 \label{I:B:eq:curvature-autocorrelation-P}
\end{equation}
Parseval's identity applied to the trigonometric polynomial $|P|^2$
shows that
\[
 \Bnorm P_{H^4}^4
 =\left(\sum_{j=0}^{m}|p_j|^2\right)^2
   +2\sum_{i=1}^{m}
       \left|\sum_{j=0}^{m-i}p_{j+i}\overline{p_j}\right|^2.
\]
The squared norm of $P$ is one, proving
\eqref{I:B:eq:curvature-hardy-coefficients}. Substitute these identities into
Proposition~\ref{I:B:prop:lift-curvature} to obtain
\eqref{I:B:eq:curvature-hardy-finite} and the norm expansion.
Finally, the scalar Hardy construction has
$E_{k-1}=1/C$, and its word lift has at most $k$ states. It is therefore
an admissible competitor for $\delta_k$, proving the upper bound in
\eqref{I:B:eq:curvature-hardy-norm}.
\end{proof}


\section{Scope and principal conclusions}
Let $\Sigma$ be finite, $d=|\Sigma|\ge2$, and let
\[
H_f[u,v]=f(uv),\qquad u,v\in\Sigma^*.
\]
Every norm of a word-Hankel operator below is its operator norm on $\FHH$ with counting measure. For bounded $H_f$ define
\[
\delta_k^{\FC}(H_f)=\inf\{\Fnorm{H_f-H_g}:H_g\text{ bounded word-Hankel},\ \Frank H_g\le k\}.
\]
For real $f$, define $\delta_k^{\FR}$ by requiring real coefficients of $g$. A bare $\delta_k$ means the complex optimum. The standard realization convention is
\[
f(a_1\cdots a_t)=\alpha^*A_{a_1}\cdots A_{a_t}\beta.
\]
Joint mean-square stability means $\rho(\Phi_A)<1$ for $\Phi_A(X)=\sum_a A_aXA_a^*$. Separate letter spectral radii do not replace this hypothesis. On rational input this condition is verifiable by solving $Y-\sum_aA_a^*YA_a=I$ and checking $Y\succ0$. A coefficient transpose $H^{\mathsf T}$ means $\mathcal JH^*\mathcal J$, with $\mathcal J$ coordinatewise conjugation; it has the same norm and singular values as $H$.

The finite-parameter family problems use the binary alphabet $\Sigma=\{a,b\}$. The synthesis theorem later allows any finite alphabet. Let
\[
q=1-r^2\in(0,1),\qquad f_{r,k}\text{ as in \eqref{I:eq:binary-source}}.
\]
Write $D_k^{\mathbb F}(q)=\delta_k^{\mathbb F}(H_{f_{\sqrt{1-q},k}})^2$.
The new finite-parameter and first-correction conclusions are:
\begin{enumerate}[leftmargin=*]
\item \textbf{Exact one-state global optimum.} Put
\[
\mathcal A_q(t)=(1-q)^2t^2+(q^2+2q-1)(t-q).
\]
Then $D_1^{\FC}(q)=D_1^{\FR}(q)$ is the unique selected root
\[
\boxed{\mathcal A_q(t)^2=4q^3(t-q)(1-2q+qt),\quad
\frac{2q}{1+q}<t<1,\quad\mathcal A_q(t)>0.}
\]
There is a single real-analytic optimal value branch throughout $(0,1)$. The complete proof and all minimizers appear below.
\item \textbf{Global two-state first correction.} For each field separately,
\[
\boxed{D_2^{\mathbb F}(q)=\frac34-\frac14q+O(q^2),
\qquad \mathbb F\in\{\FR,\FC\}.}
\]
This is not an assertion that $D_2^{\FR}(q)=D_2^{\FC}(q)$ at finite $q$. It implies that their difference is nonnegative and $O(q^2)$. The remainder constant and a numerical neighborhood are not explicitly evaluated; their existence follows from the analytic proof.
\item \textbf{Quantitative additive global synthesis.} For supplied stable rational input, an explicit finite enumeration produces a stable rational rank-at-most-$k$ competitor and rational $L,U$ with
\[
0\le L\le\delta_k(H_f)\le\Fnorm{H_f-H_g}\le U,\qquad U-L\le\varepsilon.
\]
Its candidate count is explicit and exponential in $2dk^2+4k$. Uniform moment and size bounds on the input give a nontrivial class on which the scheme is polynomial in input dimension, input bit length and $1/\varepsilon$ for fixed $d,k$. It is a theoretical enumeration scheme, not a practical fast minimizer or a polynomial algorithm in all parameters.
\end{enumerate}

The finite-parameter arguments use the boundary theorem as input and establish separate optimality statements. Classical completion, finite nets, rational Stein solves and analytic implicit-function arguments supply the supporting tools.

\section{Canonical facts used in the synthesis argument}
These are background facts. We include short proofs to make the construction's comparison class explicit.

\begin{lemma}[Stable observable canonical coordinates]\label{I:F:lem:canonical-input}
Every bounded word-Hankel $G$ of finite rank $j$ has a minimal $j$-state realization $g(w)=\ell T_w\beta$ such that
\[
\sum_a T_a^*T_a+\ell^*\ell=I,\qquad \Fnorm\ell\le1,
\qquad\Fnorm\beta\le\Fnorm G.
\]
Its completely positive transition map is strictly stable. Padding to a larger dimension preserves the displayed equality as an inequality.
\end{lemma}
\begin{proof}
Let $R_ae_u=e_{ua}$ be the right creation isometries and $\mathcal S=\Fran G$. The word-Hankel identity $R_a^*G=GL_a$, with $L_ae_v=e_{av}$, makes $\mathcal S$ invariant under $R_a^*$. Put $T_a=R_a^*|_{\mathcal S}$, $\beta=Ge_\Feps$ and $\ell x=x[\Feps]$. Then $G e_w=T_w\beta$ and $g(w)=\ell T_w\beta$. Partitioning words by last letter proves the energy identity. The observation map is inclusion: $\ell T_wx=x[w]$.

For a vector in this finite-dimensional space,
\[
\sum_{|w|=t}\Fnorm{T_wx}^2=\sum_{|v|\ge t}|x[v]|^2\longrightarrow0.
\]
Summing over a finite orthonormal basis gives uniform decay of the positive transition powers and hence a strictly smaller than one spectral radius. Reachability follows from the columns of $G$ spanning $\mathcal S$; observability from inclusion. The remaining bounds are immediate.
\end{proof}


The structured optimum is attained: a norm-bounded minimizing sequence has a weak-operator convergent subsequence; the Hankel entry equalities and all rank-$(k+1)$ minors pass to the limit, and the operator norm is lower semicontinuous. This existence fact is used in the radial comparison, not claimed as a new algorithm.

For a stable finite error realization $e(w)=cD_wb$, its rational or algebraic Gramians satisfy
\[
P=bb^*+\sum_aD_aPD_a^*,\qquad
Q=c^*c+\sum_aD_a^*QD_a.
\]
The nonzero eigenvalues of $PQ$ are the squared singular values of $H_e$, including multiplicities. They are real and nonnegative, even if $PQ$ is not Hermitian. This factorization evaluates the infinite operator norm without a word-horizon truncation.

\section{The global one-state optimum at every finite parameter}
\label{I:F:sec:F1}

Let $0<q<1$, $r=\sqrt{1-q}$, and let $f=f_{r,1}$ be the binary
series $f(a^n b)=qr^n$, zero elsewhere. All norms in this section are
ordinary counting-measure word-Hankel operator norms. Write
$D_1^{\FC}(q)=\delta_1^{\FC}(H_f)^2$ and similarly for real competitors.
The result below determines the global optimum, rather than an upper
bound within a prescribed family of amplitudes.

\begin{theorem}[Exact finite-parameter solution of F1]
\label{I:F:thm:F1-exact}
For every $0<q<1$, there is a unique $u=u(q)$ in
$q<u<\sqrt q$ satisfying
\begin{equation}
 q^2(1-u^2)(u-q)^2=(1-q)^3u^2(q-u^2).
 \label{I:F:eq:F1-u-quartic}
\end{equation}
The exact real and complex optima coincide and equal
\begin{equation}
 \boxed{\displaystyle
 D_1^{\FC}(q)=D_1^{\FR}(q)=t(q)
 =\frac{q^2+(1-2q)u(q)^2}{q(1-u(q)^2)}.}
 \label{I:F:eq:F1-value}
\end{equation}
They satisfy $2q/(1+q)<t(q)<1$. The selected functions $u(q)$ and
$t(q)$ are real analytic on $(0,1)$, with no interior change of optimal
branch.

An explicit real optimizer is $g_0(w)=c a^{\#_a(w)}y^{\#_b(w)}$, where
\begin{equation}
 a=\frac{1-u}{r},\quad
 y=\sqrt{\frac{u(1-u)}{1-q}},\quad
 m=\frac{u-q}{1-q},\quad
 \tau=\frac{u}{y(1+u)},\quad c=m\tau.
 \label{I:F:eq:F1-optimizer-parameters}
\end{equation}
These parameters obey $a,y,m,c>0$ and $a^2+y^2=1-m<1$.
All complex optimizers are precisely
\begin{equation}
 g_\theta(w)=c a^{\#_a(w)}y^{\#_b(w)}
             e^{i\theta(\#_b(w)-1)},\qquad
 \theta\in\FR/(2\pi\mathbb Z).
 \label{I:F:eq:F1-all-optimizers}
\end{equation}
The real optimizers are the two cases $\theta=0,\pi$.
\end{theorem}

The equation in $u$ has degree at most four; an exceptional vanishing of
its leading coefficient does not affect the unique interval-selected
analytic branch. A quartic directly in the optimum is given after the
proof.

\subsection{Geometric rank-one competitors and an exact residual lower bound}

Every nonzero bounded rank-one word-Hankel series has a one-state
realization
\[
 g(w)=c\lambda_w,\qquad c\ne0,\qquad
 \lambda_w=\prod_{j=1}^{|w|}\lambda_{w_j},\qquad
 |\lambda_a|^2+|\lambda_b|^2<1.
\]
Indeed minimal rank-one realization matrices are scalars, and boundedness
of the empty column forces the geometric coefficients to be square
summable. Conversely the displayed strict inequality gives a bounded
rank-one Hankel. Put
\[
 x=\lambda_a,\quad z=|\lambda_b|^2,\quad
 m=1-|x|^2-z>0,\quad
 p=\sqrt m\sum_w\lambda_w e_w.
\]
Then $\Fnorm p=1$ and $H_g=(c/m)pp^{\mathsf T}$; its input row space is
$\operatorname{span}\{\overline p\}$, with coefficientwise conjugation.

Write $s=\sqrt q$ and use the real unit vectors
\[
 u_0=s\sum_{n\ge0}r^n e_{a^n},\qquad
 v=s\sum_{n\ge0}r^n e_{a^n b},\qquad e=e_\Feps.
\]
Here $u_0\perp v$, $v\perp e$, and $\langle u_0,e\rangle=s$. The
target has the exact factorization
\begin{equation}
 H_f=u_0v^*+sve^*,\qquad H_f^*H_f=vv^*+qee^*.
 \label{I:F:eq:F1-target}
\end{equation}
Define
\begin{equation}
 A=\frac{qm}{|1-rx|^2},\qquad B=Az.
 \label{I:F:eq:F1-overlaps}
\end{equation}
The squared overlaps of $p$ with $u_0,v,e$ are $A,B,m$; since $v$ is
real, the squared overlap of $\overline p$ with $v$ is also $B$.

Let $P$ project onto $\overline p$. Every amplitude therefore satisfies
\begin{equation}
 \Fnorm{H_f-H_g}^2\ge R^2:=\Fnorm{H_f(I-P)}^2.
 \label{I:F:eq:F1-input-lower}
\end{equation}
The nonzero eigenvalues determining $R^2$ are those of the two-by-two
Hermitian matrix with diagonal entries $1-B,q(1-m)$ and off-diagonal
modulus $\sqrt{qBm}$. Its trace is $1+q-B-qm$ and determinant is
$q(1-B-m)$. For every $q<t<1$, its scalar Schur complement gives
\begin{equation}
 R^2\le t
 \quad\Longleftrightarrow\quad
 B\ge(1-t)\left(1+\frac{qm}{t-q}\right).
 \label{I:F:eq:F1-right-feasibility}
\end{equation}
The complementary diagonal is $t-q+qm>0$, so the equivalence has no
singular-denominator qualification. Rank-one interlacing also gives
$R^2\ge q$.

For fixed $|x|,z,m$, the left side of
\eqref{I:F:eq:F1-right-feasibility} is maximized by taking $x=|x|\ge0$,
because $|1-rx|\ge1-r|x|$. The phase of $\lambda_b$ has no effect.
Thus optimization of this lower bound over \emph{all complex}
competitors reduces exactly to $x=a\ge0$, $z\ge0$ and
$m=1-a^2-z>0$. A matching real completion, rather than an assumed
real-parameter restriction, will establish equality of the original
complex and real optima.

\subsection{A unique scalar branch solves the residual optimization}

For $q<t<1$, set
\[
 \lambda=\frac{1-t}{t-q}>0,\qquad
 d=(1-ra)^2,\qquad c_a=1-a^2.
\]
Condition~\eqref{I:F:eq:F1-right-feasibility} becomes
\begin{equation}
 q(c_a-z)(z-\lambda d)\ge(1-t)d.
 \label{I:F:eq:F1-concave-z}
\end{equation}
Positive feasibility requires $c_a>\lambda d$. The quadratic on the
left is strictly concave in $z$ and uniquely maximized at
\begin{equation}
 z=\frac{c_a+\lambda d}{2},\qquad
 m=\frac{c_a-\lambda d}{2}.
 \label{I:F:eq:F1-z-vertex}
\end{equation}
The maximum is $q(c_a-\lambda d)^2/4$. Therefore feasibility is equivalent
to a positive value satisfying
\begin{equation}
 \frac{\sqrt q}{2}
 \max_a\left\{\frac{1-a^2}{1-ra}-\lambda(1-ra)\right\}
 \ge\sqrt{1-t}.
 \label{I:F:eq:F1-residual-scalar}
\end{equation}

Set $u=1-ra\in(1-r,1]$. The expression to maximize is
\begin{equation}
 F_\lambda(u)
 =\frac{2-(1+(1-q)\lambda)u-q/u}{1-q}.
 \label{I:F:eq:F1-concave-u}
\end{equation}
It is strictly concave, with unconstrained maximum at
\[
 u_\lambda=\sqrt{\frac{q}{1+(1-q)\lambda}}.
\]
If $u_\lambda\le1-r$, its constrained maximum would occur at $1-r$
and equal $-\lambda(1-r)<0$, precluding feasibility. Also
$u_\lambda<\sqrt q<1$. A positive maximal value is equivalent to
$0<\lambda<1/q$; in this case $q<u_\lambda<\sqrt q$, so it lies
strictly inside the allowed interval. The maximum is
\[
 \max F_\lambda=
 \frac{2-2\sqrt{q(1+(1-q)\lambda)}}{1-q}.
\]
Using $t=(1+q\lambda)/(1+\lambda)$, the threshold condition is exactly
\begin{equation}
 \frac{q}{(1-q)^2}
 \left(1-\sqrt{q(1+(1-q)\lambda)}\right)^2
 \ge\frac{(1-q)\lambda}{1+\lambda}.
 \label{I:F:eq:F1-monotone-lambda}
\end{equation}
On $(0,1/q)$ its left side strictly decreases from a positive value to
zero, while its right side strictly increases from zero to a positive
value. There is consequently exactly one equality. Since $t$ strictly
decreases with $\lambda$, that equality is the minimal feasible residual
threshold.

At equality, set $u=u_\lambda$. Then
\begin{equation}
 \lambda=\frac{q-u^2}{(1-q)u^2},\qquad
 z=\frac{u(1-u)}{1-q},\qquad
 m=\frac{u-q}{1-q},
 \label{I:F:eq:F1-stationary-parameters}
\end{equation}
and the threshold is \eqref{I:F:eq:F1-value}. Its complement has the two
equivalent expressions
\begin{equation}
 1-t=\frac{q(u-q)^2}{(1-q)^2u^2}
     =\frac{(1-q)(q-u^2)}{q(1-u^2)}.
 \label{I:F:eq:F1-equality-u}
\end{equation}
Equating them gives \eqref{I:F:eq:F1-u-quartic}, without introducing an
extraneous sign. Conversely any root in $(q,\sqrt q)$ satisfies the
positive equality. The first expression in \eqref{I:F:eq:F1-equality-u}
strictly increases from zero in $u$, while the second strictly decreases
to zero. This proves the asserted root uniqueness directly. Their
difference has positive $u$ derivative, so the implicit function theorem
also proves analyticity on $(0,1)$.

We have thus proved that \emph{every} bounded complex rank-one competitor
has squared error at least the number $t(q)$ in
\eqref{I:F:eq:F1-value}. The zero competitor has error one, and cannot be
optimal since $t(q)<1$.

\subsection{Strict output slack and a genuine Hankel completion}

Take the positive-real geometric parameters in
\eqref{I:F:eq:F1-stationary-parameters}, with $y=\sqrt z$, and let $p$ be
the corresponding real unit vector. The output projection residual
$L^2=\Fnorm{(I-pp^*)H_f}^2$ is the larger eigenvalue of
\[
 \begin{pmatrix}
 1-A&-A\sqrt{qz}\\
 -A\sqrt{qz}&q(1-Az)
 \end{pmatrix}.
\]
For $t>q$, its Schur threshold is
\begin{equation}
 L^2\le t
 \quad\Longleftrightarrow\quad
 A\ge(1-t)\left(1+\frac{qAz}{t-q}\right)
 =(1-t)+\lambda qAz.
 \label{I:F:eq:F1-left-feasibility}
\end{equation}
Equation~\eqref{I:F:eq:F1-equality-u} gives $1-t=Am$. Hence the final
condition is $1-\lambda qz\ge m$. Substituting
\eqref{I:F:eq:F1-stationary-parameters}, its strict version is equivalent to
\[
 qu^2+(1-q)u-q^2>0.
\]
This polynomial is strictly increasing on $u>0$ and at $u=q$ equals
$q(1-q)^2>0$. Thus
\begin{equation}
 \Fnorm{H_f(I-pp^*)}^2=t,\qquad
 \Fnorm{(I-pp^*)H_f}^2<t.
 \label{I:F:eq:F1-two-corners}
\end{equation}
The strict second inequality supplies the required sequentially compatible
completion. In particular, we do not assert that $H_fpp^*$ itself is
Hankel.

\begin{lemma}[Scalar one-corner norm completion]
\label{I:F:lem:F1-completion}
Let $H$ be bounded and let $p$ be a unit vector used as the selected
one-dimensional input and output space. If
\[
 \Fnorm{H(I-pp^*)}\le\gamma,
 \qquad\Fnorm{(I-pp^*)H}<\gamma,
 \quad\gamma>0,
\]
there is an amplitude $\tau$ such that
$\Fnorm{H-\tau pp^*}\le\gamma$. If the first displayed norm equals
$\gamma$, this amplitude is unique. Real data admit a real amplitude.
\end{lemma}

\begin{proof}
This is a scalar contraction-completion result; we provide its elementary
Schur proof rather than relying on an external completion theorem (compare
the classical mechanism discussed in~\cite{parrott2025}). Write
\[
 H=\begin{pmatrix}h&b\\c&D\end{pmatrix}
\]
relative to $\operatorname{span}\{p\}\oplus p^\perp$. Since
$\Fnorm D<\gamma$, the complement block
\[
 \mathcal L=\begin{pmatrix}\gamma I&D\\D^*&\gamma I\end{pmatrix}
\]
is strictly positive. Put $t=\gamma^2$ and $x=h-\tau$.
The norm inequality for $M=\left(\begin{smallmatrix}x&b\\c&D\end{smallmatrix}\right)$
is equivalent to positivity of
$\left(\begin{smallmatrix}\gamma I&M\\M^*&\gamma I\end{smallmatrix}\right)$.
Its Schur complement with respect to $\mathcal L$ is the scalar block
\[
 \begin{pmatrix}
 \gamma\{1-b(tI-D^*D)^{-1}b^*\}
 &x+bD^*(tI-DD^*)^{-1}c\\
 \overline x+c^*(tI-DD^*)^{-1}Db^*
 &\gamma\{1-c^*(tI-DD^*)^{-1}c\}
 \end{pmatrix}.
\]
The given column and row norm bounds make both diagonal entries
nonnegative. Choose $x$ to make the off-diagonal zero. This proves
existence. If the column norm equals $\gamma$, its scalar Schur
complement, the first diagonal entry, is zero; positivity then forces the
off-diagonal to vanish, proving uniqueness. All formulas preserve real
data.
\end{proof}


\subsection{Explicit amplitude and classification of the minimizers}

Apply Lemma~\ref{I:F:lem:F1-completion} with $\gamma=\sqrt t$ and
\eqref{I:F:eq:F1-two-corners}. To identify its unique amplitude, set
\[
 z_0=(tI-H_f^*H_f)^{-1}p.
\]
This inverse exists because $q<t<1$ and
$\operatorname{spec}(H_f^*H_f)=\{0,q,1\}$. The right residual secular
equation gives $p^*z_0=0$; equivalently it follows by multiplying the
characteristic equation associated with
\eqref{I:F:eq:F1-right-feasibility}. Hence
\[
 H_f^*H_fz_0=t z_0-p,
 \qquad\Fnorm{H_fz_0}^2=t\Fnorm{z_0}^2.
\]
Every completion of norm at most $\sqrt t$ preserves this saturated input
direction. Taking its scalar cross term against $p$ gives
\begin{equation}
 \tau=-\frac1{p^*H_fz_0}.
 \label{I:F:eq:F1-amplitude-necessary}
\end{equation}
All quantities are real here. Put
$\alpha=\langle p,u_0\rangle=\sqrt A$ and
$\beta=\langle p,v\rangle=\alpha y$. The target factorization gives
\[
 p^*H_fz_0
 =\frac{\alpha\beta}{t-1}
  +\frac{\sqrt q\,\beta\sqrt m}{t-q}
 =-\frac{\alpha\beta(1-u\lambda)}{1-t}.
\]
The denominator in \eqref{I:F:eq:F1-amplitude-necessary} is nonzero because
\[
 1-u\lambda=\frac{(u-q)(1+u)}{(1-q)u}
             =\frac{m(1+u)}u>0.
\]
Finally $1-t=Am$, so
\[
 \tau=\frac{1-t}{Ay(1-u\lambda)}
       =\frac{u}{y(1+u)}.
\]
The completion lemma proves that this amplitude has error at most
$\sqrt t$, and the universal residual lower bound proves equality.
Since $pp^*=m(\lambda_w)_w(\lambda_w)_w^{\mathsf T}$ for these real
parameters, this is exactly the coefficient series in
\eqref{I:F:eq:F1-optimizer-parameters}. A real competitor has thus attained
the lower bound over all complex competitors, proving
$D_1^{\FR}=D_1^{\FC}=t$.

For completeness, every optimizer must saturate the inequalities in the
residual optimization. The strict phase bound for $|1-rx|$, together with
the strict quadratic maxima, forces $\lambda_a=a>0$ and the unique
modulus $|\lambda_b|=y$. Let
$U_\theta e_w=e^{i\theta\#_b(w)}e_w$. Because every target word has
exactly one $b$,
\[
 U_\theta H_fU_\theta=e^{i\theta}H_f.
\]
Consequently
$H_{g_\theta}=e^{-i\theta}U_\theta H_{g_0}U_\theta$ satisfies
\[
 H_f-H_{g_\theta}
 =e^{-i\theta}U_\theta(H_f-H_{g_0})U_\theta.
\]
Both side factors are unitary, so the errors have the same norm.
Conversely, a phase-$\theta$ competitor is reduced to the real geometric
pair by $G\mapsto e^{i\theta}U_\theta^*GU_\theta^*$ before applying
Lemma~\ref{I:F:lem:F1-completion}. This gives
\eqref{I:F:eq:F1-all-optimizers}. Strict output slack makes
the amplitude unique for each fixed geometric pair, so these exhaust the
optimizers. Since $g_\theta(\Feps)=ce^{-i\theta}\ne0$, real
coefficients force $\theta=0$ or $\pi$; both choices are indeed real.
This completes the proof of Theorem~\ref{I:F:thm:F1-exact}.

\subsection{A quartic directly in the optimum}

\begin{theorem}[Direct polynomial branch selector]
\label{I:F:thm:F1-t-quartic}
For $0<q<1$ put
\[
 B_q=q^2+2q-1,\qquad d_q(t)=1-2q+qt,\qquad
 A_q(t)=(1-q)^2t^2+B_q(t-q).
\]
Then $D_1^{\FC}(q)=D_1^{\FR}(q)$ is the unique real $t$ satisfying
\begin{equation}
 \boxed{\displaystyle A_q(t)^2=4q^3(t-q)d_q(t),\qquad
 \frac{2q}{1+q}<t<1,\quad A_q(t)>0.}
 \label{I:F:eq:F1-direct-quartic}
\end{equation}
The optimizer is reconstructed by
$u=\sqrt{q(t-q)/d_q(t)}$ and
\eqref{I:F:eq:F1-optimizer-parameters}.
\end{theorem}

\begin{proof}
On the stated interval $d_q(t)>0$, and
\[
 u^2=\frac{q(t-q)}{d_q(t)}
\]
lies strictly between $q^2$ and $q$. This is the inverse of
\eqref{I:F:eq:F1-value}. With this substitution,
\eqref{I:F:eq:F1-u-quartic} is equivalently
\[
 (u-q)^2=
 \frac{(1-q)^2(t-q)(1-t)}{d_q(t)}.
\]
Expanding the square and using the expression for $u^2$ gives
\[
 2q\,d_q(t)u=A_q(t).
\]
Its square is \eqref{I:F:eq:F1-direct-quartic}; the strict condition
$A_q(t)>0$ chooses the positive sign and recovers the preceding equality.
Thus the interval and sign conditions make the polynomial formulation
equivalent to the unique $u$ branch already proved. No additional
quantifier-elimination or local-optimization selection is required.
\end{proof}


\subsection{Short certified intervals without global elimination}

The selected scalar branch also provides an elementary certification
procedure. Let $[u_-,u_+]$ be the bracket after $N$ bisection steps from
$[q,\sqrt q]$, using the sign of the polynomial in
\eqref{I:F:eq:F1-u-quartic} moved to the left. Its sign agrees with the
strictly increasing difference in \eqref{I:F:eq:F1-equality-u}; it is
negative at $q$ and positive at $\sqrt q$. Thus bisection is certified
without searching over realizations.

Regarding the expression in \eqref{I:F:eq:F1-value} as $t_q(u)$,
\[
 t_q'(u)=\frac{2u(1-q)^2}{q(1-u^2)^2}\le\frac2{\sqrt q}
 \qquad(q<u<\sqrt q).
\]
Since $u_+-u_-=(\sqrt q-q)2^{-N}$, monotonicity gives certified
squared-error bounds
\[
 t_q(u_-)\le D_1(q)\le t_q(u_+),\qquad
 t_q(u_+)-t_q(u_-)\le2(1-\sqrt q)2^{-N}\le2^{1-N}.
\]
Their positive square roots bound $\delta_1$, with gap at most
$2^{(1-N)/2}$. In particular $N\ge1+2\log_2(1/\epsilon)$ suffices
for a norm interval of width at most $\epsilon$. This is a uniform
count of degree-four sign evaluations, rather than a quantifier-elimination
procedure.

For rational $q$ of binary length $B$, the bisection endpoints lie in the
quadratic field $\mathbb Q(\sqrt q)$. After clearing denominators, each
polynomial sign is the sign of $A+B'\sqrt{pd}$ with integers of bit length
$O(B+N)$ when $q=p/d$ in lowest terms. Signs reduce to integer signs and,
when necessary, comparison of $A^2$ with $(B')^2pd$. Consequently elementary
schoolbook arithmetic gives a polynomial bit bound, for example
$O(N(B+N)^2)$ with a common-denominator representation. The exact optimizer
itself is specified by the degree-at-most-four selected root and the
explicit radicals in \eqref{I:F:eq:F1-optimizer-parameters}; this avoids any
numerical perturbation of the optimal realization.

\section{A two-state branch that improves the first boundary correction}
\label{I:F:sec:f2}

The global two-state first correction is a different problem from the
one-state optimum. The previously available fixed lift has squared error
$3/4+q^2/6+O(q^3)$. The construction below proves that allowing the
competitor to depend on $q$ improves the error already at order $q$.
The matching global lower bound is proved separately by the two-dimensional passive-envelope argument.

\begin{theorem}[An improving real two-state branch]
\label{I:F:thm:f2-improvement}
For $0<q<1/2$, set
\[
 r=\sqrt{1-q},\qquad a=\sqrt{1-q-2q^2},\qquad
 c=\frac{q(1+2q)}2,
\]
and define the standard-order realization
\begin{equation}
 A_a=\begin{pmatrix}a&0\\0&0\end{pmatrix},\qquad
 A_b=\frac1{\sqrt2}\begin{pmatrix}q&1\\-q&1\end{pmatrix},\qquad
 \alpha_g^*=c(1,0),\qquad \beta_g=(-q,1)^{\mathsf T}.
 \label{I:F:eq:f2-branch}
\end{equation}
Let $g_q(w)=\alpha_g^*A_w\beta_g$, with zero transition matrices for any
additional alphabet letters. This is a real, jointly stable realization
of dimension two. For the target $f_{r,2}(a^n b^2)=qr^n$,
\begin{equation}
 \boxed{\displaystyle
 \Fnorm{H_{f_{r,2}}-H_{g_q}}^2
 =\frac34-\frac14q+O(q^2)\qquad(q\downarrow0).}
 \label{I:F:eq:f2-branch-expansion}
\end{equation}
Consequently
\begin{equation}
 D_2^{\FC}(q)\le D_2^{\FR}(q)
 \le\frac34-\frac14q+O(q^2).
 \label{I:F:eq:f2-optimum-upper}
\end{equation}
In particular the fixed lift approaching $3/4$ from above is not
asymptotically optimal at first order.
\end{theorem}

\begin{proof}
We establish joint stability, give exact error Gramians, and differentiate
the isolated leading squared singular value.

\emph{Joint stability.}
Direct multiplication yields
\begin{equation}
 A_a^*A_a+A_b^*A_b
   =\operatorname{diag}(1-q-q^2,1).
 \label{I:F:eq:f2-passivity}
\end{equation}
Let $D_0=\operatorname{diag}(q+q^2,0)$. The matrix
$D_0+A_a^*D_0A_a+A_b^*D_0A_b$ is positive definite. Indeed, its kernel
would be contained in $\ker D_0=\operatorname{span}\{(0,1)^{\mathsf T}\}$, but
$A_b(0,1)^{\mathsf T}=(1,1)^{\mathsf T}/\sqrt2$ has nonzero first coordinate.
Telescoping \eqref{I:F:eq:f2-passivity} therefore gives
\[
 \sum_{|w|=2}A_w^*A_w\prec I.
\]
The column completely positive map is a strict contraction after two
steps, and its adjoint has the same spectral radius. Thus the system is
mean-square stable. Its word-Hankel is bounded of rank at most two.

\emph{Competitor Gramians.}
Let $R_g,Q_g$ be its exact reachability and observation Gramians. Their
Stein equations have the solutions
\begin{equation}
 R_g=\begin{pmatrix}
 \dfrac{1+q^2}{q(1+q)}&\dfrac{1-q}{1+q}\\[4pt]
 \dfrac{1-q}{1+q}&\dfrac{2+q+2q^2+q^3}{1+q}
 \end{pmatrix},\qquad
 Q_g=\frac{q(1+2q)^2}{4(1+q)}I_2.
 \label{I:F:eq:f2-own-gramians}
\end{equation}
These identities can be verified by substitution in
$R_g=\beta_g\beta_g^*+\sum_j A_jR_gA_j^*$ and
$Q_g=\alpha_g\alpha_g^*+\sum_j A_j^*Q_gA_j$.
For example, writing the entries of $R_g$ as $u,v,w$, the off-diagonal
and lower-right equations give
$v=(1-q)/(1+q)$ and $w=2+q^2u-2qv$; the remaining equation gives
$u=(1+q^2)/(q(1+q))$. Joint stability ensures uniqueness.

\emph{Target and cross Gramians.}
Use the target realization
\[
 F_a=\operatorname{diag}(r,0,0),\qquad
 F_b=\begin{pmatrix}0&1&0\\0&0&1\\0&0&0\end{pmatrix},\qquad
 \alpha_f^*=q(1,0,0),\qquad \beta_f=(0,0,1)^{\mathsf T}.
\]
Its Gramians are
$R_f=\operatorname{diag}(q^{-1},1,1)$ and $Q_f=qI_3$.
Put $\Delta=1-ra>0$. Define $X,Y\in\FC^{3\times2}$ by their rows:
\begin{align}
 X_0&=\left(\frac{1+q+q^2-q^3}{2\Delta},
                 \frac{1-q+q^2+q^3}{2}\right),\notag\\
 X_1&=\frac1{\sqrt2}(1-q^2,1+q^2),\qquad X_2=(-q,1),
 \label{I:F:eq:f2-cross-reach}\\
 Y_0&=\frac{qc}{\Delta}(1,0),\qquad
 Y_1=Y_0A_b,\qquad Y_2=Y_0A_b^2.
 \label{I:F:eq:f2-cross-observe}
\end{align}
The triangular form of $F_b$ gives directly
\[
 X=\beta_f\beta_g^*+\sum_jF_jXA_j^*,\qquad
 Y=\alpha_f\alpha_g^*+\sum_jF_j^*YA_j.
\]
For example, the first equation gives successively
$X_2=\beta_g^*$, $X_1=\beta_g^*A_b^*$, and
$X_0=\beta_g^*(A_b^*)^2\operatorname{diag}(\Delta^{-1},1)$,
which is \eqref{I:F:eq:f2-cross-reach}.

For the composite error realization with output
$(\alpha_f^*,-\alpha_g^*)$ and input $(\beta_f,\beta_g)$, the exact
Gramians are therefore
\begin{equation}
 R=\begin{pmatrix}R_f&X\\X^*&R_g\end{pmatrix},\qquad
 Q=\begin{pmatrix}Q_f&-Y\\-Y^*&Q_g\end{pmatrix}.
 \label{I:F:eq:f2-error-gramians}
\end{equation}
Its infinite Hankel squared norm is the largest eigenvalue of
$R^{1/2}QR^{1/2}$, equivalently the largest eigenvalue of $RQ$.
Equations \eqref{I:F:eq:f2-own-gramians}--\eqref{I:F:eq:f2-error-gramians}
are thus a finite exact algebraic description of this chosen error;
there is no word-horizon cutoff.

\emph{Leading eigenvalue perturbation.}
Set $\widehat R=qR$ and $\widehat Q=Q/q$, so
$\widehat R\widehat Q=RQ$. Both extend analytically to $q=0$.
In particular $\Delta=q+q^2+O(q^4)$, removing the apparent
singularities in $qX$ and $Y/q$.
Index the five coordinates by $(f_0,f_1,f_2,g_0,g_1)$ and put
$y=1/\sqrt2$. The expansions are
\[
 \widehat R=R_0+qR_1+O(q^2),\qquad
 \widehat Q=Q_0+qQ_1+O(q^2),
\]
where
\begin{equation}
 R_0=\begin{pmatrix}
 1&0&0&1/2&0\\0&0&0&0&0\\0&0&0&0&0\\
 1/2&0&0&1&0\\0&0&0&0&0
 \end{pmatrix},\quad
 R_1=\begin{pmatrix}
 0&0&0&0&1/2\\0&1&0&y&y\\0&0&1&0&1\\
 0&y&0&-1&1\\1/2&y&1&1&2
 \end{pmatrix},
 \label{I:F:eq:f2-R-expansion}
\end{equation}
\begin{equation}
 Q_0=\begin{pmatrix}
 1&0&0&-1/2&0\\0&1&0&0&-y/2\\0&0&1&0&-1/4\\
 -1/2&0&0&1/4&0\\0&-y/2&-1/4&0&1/4
 \end{pmatrix},\quad
 Q_1=\begin{pmatrix}
 0&0&0&-1/2&0\\0&0&0&-y/2&-y/2\\0&0&0&1/4&-1/2\\
 -1/2&-y/2&1/4&3/4&0\\0&-y/2&-1/2&0&3/4
 \end{pmatrix}.
 \label{I:F:eq:f2-Q-expansion}
\end{equation}
The matrix $Z_0=R_0Q_0$ has only one nonzero row:
$Z_0[0,0]=3/4$, $Z_0[0,3]=-3/8$.
Its sole nonzero eigenvalue $3/4$ is simple. Right and left eigenvectors,
normalized by $w^*v=1$, are
\[
 v=e_0,\qquad w=e_0-\tfrac12 e_3.
\]
The analytic eigenvalue branch consequently has derivative
\begin{equation}
 \lambda'(0)=w^*(R_1Q_0+R_0Q_1)v=-\frac14.
 \label{I:F:eq:f2-eigen-derivative}
\end{equation}
More explicitly, the two terms are $-1/4$ and zero, respectively.
All other eigenvalues tend to zero, so this branch is the largest
squared Hankel singular value for all sufficiently small positive $q$.
Analyticity supplies a bounded $O(q^2)$ remainder and proves
\eqref{I:F:eq:f2-branch-expansion}. Admissibility of the real two-state
construction proves \eqref{I:F:eq:f2-optimum-upper}.
\end{proof}


\begin{remark}[What this does and does not determine]
The theorem refutes first-order optimality of the earlier fixed lift,
and refutes any proposed lower bound of the form
$D_2(q)\ge3/4+o(q)$ near zero. Taken alone it is an upper-bound result. The separate passive-envelope
lower bound, proved in this manuscript, establishes
$D_2^{\FC}(q)\ge3/4-q/4-O(q^2)$ against all complex competitors.
Combining the two proves the common first correction for the real and
complex optima. Local numerical searches played no role in the global
lower-bound certificate; higher coefficients and exact finite-$q$
rank-two optima are not determined here.
\end{remark}

\section{The global complex two-state first correction}
\label{I:F:sec:f2-global-lower}

We now match the real upper branch of
Theorem~\ref{I:F:thm:f2-improvement} by a lower bound against every bounded
complex word-Hankel competitor of rank at most two.  The central step is
an asymptotic maximization over all passive complex two-dimensional
systems.  The maximization is global; no restriction to a chosen real
matrix family is imposed.  The proof does not establish an exact
finite-$q$ classification of the optimum.

\subsection{From arbitrary two-state competitors to a passive matrix problem}

Write $q=1-r^2$, with $0<q<1$, and let
\begin{equation}
 u_r=\sqrt q\sum_{n\ge0}r^n e_{a^n},\qquad
 v_r=\sqrt q\sum_{n\ge0}r^n e_{a^n b^2}.
 \label{I:F:eq:f2-test-vectors}
\end{equation}
Both vectors have norm one in $\FHH$.  For the target
$X_f=H_{f_{r,2}}^{\mathsf T}$, direct substitution gives
$X_fu_r=v_r$.

Let $H_g$ be any bounded complex word-Hankel operator of rank at most
two, and let $\mathcal S=\operatorname{ran}H_g^{\mathsf T}$.
This is a finite-dimensional backward-shift invariant subspace of
$\FHH$.  In orthonormal canonical state coordinates, denote the
restrictions of $L_a^*$ and $L_b^*$ by $A$ and $B$, and denote
empty-word evaluation by the row $\ell$.  If necessary, pad the system
by zero coordinates to dimension two.  The canonical energy identity
implies
\begin{equation}
 A^*A+B^*B+\ell^*\ell\preceq I_2.
 \label{I:F:eq:f2-passive-class}
\end{equation}
Additional alphabet letters only add positive terms to the omitted part
of this identity.  Coordinate evaluation on $a^n b^2$ is the row
$\ell B^2A^n$.  Consequently, if $P_{\mathcal S}$ is orthogonal
projection onto $\mathcal S$,
\begin{equation}
 \Fnorm{P_{\mathcal S}v_r}^2
   =q\Fnorm{\ell B^2(I-rA)^{-1}}^2.
 \label{I:F:eq:f2-projection-passive}
\end{equation}
The resolvent exists because \eqref{I:F:eq:f2-passive-class} implies
$\Fnorm{A}\le1$ and $r<1$.  Since the range of $H_g^{\mathsf T}$ is
contained in $\mathcal S$,
\begin{equation}
 \Fnorm{H_{f_{r,2}}-H_g}^2
 \ge \Fnorm{(I-P_{\mathcal S})v_r}^2
 =1-q\Fnorm{\ell B^2(I-rA)^{-1}}^2.
 \label{I:F:eq:f2-competitor-lower}
\end{equation}
Thus it suffices to understand the universal passive-tuple maximum
\begin{equation}
 \mathcal P(q)=
 \max_{A^*A+B^*B+\ell^*\ell\preceq I_2}
 q\Fnorm{\ell B^2(I-rA)^{-1}}^2,
 \qquad r=\sqrt{1-q}.
 \label{I:F:eq:f2-P-definition}
\end{equation}
The maximum exists: the feasible set is compact and the resolvent is
continuous there for every fixed $r<1$.

\subsection{Exact elimination of the waiting transition}

The variable $A$ can be eliminated without relaxation.  For $0\le t\le1$,
define the compact envelope
\begin{equation}
 h(t)=\max\left\{
  |\ell B^2e_1|^2:
  V=\begin{pmatrix}B\\\ell\end{pmatrix}\in\FC^{3\times2},\quad
  V^*V\preceq I_2,\quad\Fnorm{Ve_1}^2=t
 \right\}.
 \label{I:F:eq:f2-h-definition}
\end{equation}
Here and below $e_1=(1,0)^{\mathsf T}$ is a state coordinate vector.

\begin{lemma}[Exact scalar reduction]
\label{I:F:lem:f2-A-elimination}
For every $0<q<1$,
\begin{equation}
 \mathcal P(q)=
 \max_{0\le t\le1}
 \frac{q\,h(t)}{\bigl(1-\sqrt{1-q}\sqrt{1-t}\bigr)^2}.
 \label{I:F:eq:f2-scalar-reduction}
\end{equation}
Moreover $0\le h(t)\le t$.
\end{lemma}

\begin{proof}
For fixed $A,B,\ell$, change variables in the row operator norm by
$x=(I-rA)y$.  This gives
\[
 \Fnorm{\ell B^2(I-rA)^{-1}}^2
 =\sup_{y\ne0}\frac{|\ell B^2y|^2}{\Fnorm{(I-rA)y}^2}.
\]
Normalize $\Fnorm{y}=1$ and set
$D=I-B^*B-\ell^*\ell\succeq0$ and
$t=\Fnorm{By}^2+|\ell y|^2=1-y^*Dy$.
The constraint $A^*A\preceq D$ yields
\begin{equation}
 \Fnorm{(I-rA)y}\ge1-r\Fnorm{Ay}
                 \ge1-r\sqrt{1-t}.
 \label{I:F:eq:f2-denominator-lower}
\end{equation}
A unitary state change sends $y$ to $e_1$ and preserves all constraints.
Taking suprema gives the upper direction in
\eqref{I:F:eq:f2-scalar-reduction}.

Conversely, fix $B,\ell,y$ satisfying the stated conditions.  If $t<1$,
define the rank-one matrix
\begin{equation}
 A=\frac{y\,y^*D}{\sqrt{1-t}}.
 \label{I:F:eq:f2-A-completion}
\end{equation}
Then $Ay=\sqrt{1-t}\,y$ and
\[
 A^*A=\frac{Dyy^*D}{y^*Dy}\preceq D.
\]
The last inequality is Cauchy--Schwarz in the positive semidefinite
$D$ metric.  Equality holds throughout
\eqref{I:F:eq:f2-denominator-lower}.  If $t=1$, take $A=0$ instead.
This proves attainability of the scalar denominator and hence
\eqref{I:F:eq:f2-scalar-reduction}.
Finally $\Fnorm{B},\Fnorm{\ell}\le1$, so
$|\ell B^2e_1|\le\Fnorm{Be_1}\le\sqrt t$.
\end{proof}


\subsection{The global compact envelope near zero}

\begin{lemma}[Two-term expansion of the complex envelope]
\label{I:F:lem:f2-h-expansion}
The envelope \eqref{I:F:eq:f2-h-definition} satisfies
\begin{equation}
 h(t)=\frac14t+\frac14t^2+O(t^3)
 \qquad(t\downarrow0).
 \label{I:F:eq:f2-h-expansion}
\end{equation}
The remainder applies to the global maximum over all complex
contractions in \eqref{I:F:eq:f2-h-definition}.
\end{lemma}

\begin{proof}
We first parametrize the entire feasible set, then remove the phase
symmetries at its limiting maximizing set.  The resulting maximum is a
nondegenerate analytic critical point in seven real variables.

\emph{An exact compact parametrization.}
Put $t=s^2$, initially allowing either sign of the real parameter $s$.
Write the first column of $V$ as $u=sz$, where
$z\in\FC^3$ and $\Fnorm{z}=1$, and write the second column as $v$.
For $|s|<1$, the Schur complement of $I-V^*V$ gives the exact equivalence
\begin{equation}
 V^*V\preceq I
 \quad\Longleftrightarrow\quad
 \Fnorm{v}^2+\frac{s^2}{1-s^2}|z^*v|^2\le1.
 \label{I:F:eq:f2-column-ellipsoid}
\end{equation}
Thus every feasible second column is uniquely of the form
\begin{equation}
 v=W_s(z)\nu,\qquad \Fnorm{\nu}\le1,\qquad
 W_s(z)=I_3-\bigl(1-\sqrt{1-s^2}\bigr)zz^*.
 \label{I:F:eq:f2-W-parametrization}
\end{equation}
The matrix $W_s(z)$ is analytic and even in $s$ near zero.  In these
coordinates,
\[
 B=\begin{pmatrix}sz_1&v_1\\sz_2&v_2\end{pmatrix},\qquad
 \ell=(sz_3,v_3),
\]
and direct multiplication gives
\begin{align}
 \ell B^2e_1&=s\,\mathcal J(s,z,\nu),\notag\\
 \mathcal J(s,z,\nu)
 &=z_2v_2v_3+s z_2(z_3v_1+z_1v_3)+s^2z_3z_1^2.
 \label{I:F:eq:f2-J-objective}
\end{align}
Consequently
\begin{equation}
 h(s^2)=s^2\mathcal M(s)^2,\qquad
 \mathcal M(s)=\max_{\Fnorm{z}=1,\,\Fnorm{\nu}\le1}
                  |\mathcal J(s,z,\nu)|.
 \label{I:F:eq:f2-M-definition}
\end{equation}

At $s=0$,
\[
 |\mathcal J(0,z,\nu)|=|z_2\nu_2\nu_3|\le\frac12.
\]
Equality is possible exactly on the phase torus
\begin{equation}
 z_1=z_3=\nu_1=0,\qquad
 |z_2|=1,\qquad |\nu_2|=|\nu_3|=2^{-1/2}.
 \label{I:F:eq:f2-maximizing-torus}
\end{equation}
The parameter space is compact, and the objective converges uniformly
as $s\to0$.  Hence every global maximizing pair for sufficiently small
$|s|$ lies arbitrarily close to this torus.

\emph{All complex phase degeneracies can be removed.}
The original optimization is invariant under
\begin{equation}
 B\longmapsto e^{i\theta}U^*BU,\qquad
 \ell\longmapsto e^{i\phi}\ell U,\qquad
 U=\operatorname{diag}(1,e^{i\psi}).
 \label{I:F:eq:f2-phase-action}
\end{equation}
The scalar $\ell B^2e_1$ is multiplied only by
$e^{i(\phi+2\theta)}$.  On the compact parameters, the same action is
\[
 z\longmapsto Dz,\qquad
 \nu\longmapsto e^{i\psi}D\nu,\qquad
 D=\operatorname{diag}\bigl(e^{i\theta},e^{i(\theta-\psi)},e^{i\phi}\bigr).
\]
Indeed \eqref{I:F:eq:f2-W-parametrization} is equivariant under diagonal
unitaries.  The phases of $(z_2,\nu_2,\nu_3)$ change respectively by
$(\theta-\psi,\theta,\phi+\psi)$, an invertible linear transformation
of the three phase parameters.  Since those three coordinates are
nonzero near \eqref{I:F:eq:f2-maximizing-torus}, they may all be fixed to be
positive real in a smooth local gauge.  No transverse complex direction
is excluded by this choice.

A global maximizer also satisfies $\Fnorm{\nu}=1$ for all sufficiently
small $|s|$.  To see this, use radial coordinates
$\nu=\rho\widehat\nu$, $\Fnorm{\widehat\nu}=1$.
At the limiting torus the derivative of
$|\mathcal J|$ in the increasing radial direction is $1$ at $\rho=1$.
It remains positive in a neighborhood of that compact torus.  Global
maximizers approach the torus, so a maximizing point with $\rho<1$
would admit a feasible radial increase of its objective.

\emph{Seven real coordinates and a strictly negative Hessian.}
After fixing the phases and the radial constraint, choose local
coordinates
\begin{equation}
 \begin{gathered}
 z_1,z_3,\nu_1\in\FC,\qquad \xi\in\FR,\qquad
 z_2=\sqrt{1-|z_1|^2-|z_3|^2},\\
 \nu_2=\sqrt{1-|\nu_1|^2}\cos(\pi/4+\xi),\qquad
 \nu_3=\sqrt{1-|\nu_1|^2}\sin(\pi/4+\xi).
 \end{gathered}
 \label{I:F:eq:f2-seven-coordinates}
\end{equation}
Let $x\in\FR^7$ list, in order,
$\Re z_1,\Im z_1,\Re z_3,\Im z_3,\Re\nu_1,\Im\nu_1,\xi$.
The gauge sends the entire limiting maximizing torus to $x=0$.
The objective
$F(s,x)=|\mathcal J(s,z(x),\nu(x))|$ is real analytic near $(0,0)$,
because $\mathcal J(0,0)=1/2\ne0$.
At $s=0$ its expansion is
\begin{equation}
 F(0,x)=\frac12-\frac{|z_1|^2+|z_3|^2}{4}
                  -\frac{|\nu_1|^2}{2}-\xi^2+O(\Fnorm{x}^4).
 \label{I:F:eq:f2-transverse-Hessian}
\end{equation}
Thus
\begin{equation}
 \partial_x^2F(0,0)=
 \operatorname{diag}\!\left(-\frac12,-\frac12,-\frac12,-\frac12,
                           -1,-1,-2\right)\prec0.
 \label{I:F:eq:f2-Hessian-matrix}
\end{equation}
The full expansion through total degree two is
\begin{align}
 F(s,x)=\frac12-\frac14s^2+\frac{s\,\Re z_1}{\sqrt2}
          -\frac{|z_1|^2+|z_3|^2}{4}
          -\frac{|\nu_1|^2}{2}-\xi^2
          +O\bigl((|s|+\Fnorm{x})^3\bigr).
 \label{I:F:eq:f2-full-quadratic}
\end{align}
For completeness, the pure $s^2$ term follows from the exact identity
$\mathcal J(s,0)=\sqrt{1-s^2}/2$.  In
\eqref{I:F:eq:f2-J-objective}, the only mixed term of total degree two
contributing to the modulus is $s\Re z_1/\sqrt2$.

The analytic implicit function theorem applied to
$\partial_xF(s,x)=0$ gives a unique analytic critical branch $x=x(s)$
near zero.  The negative definite Hessian persists in a fixed small
coordinate neighborhood, where the objective is strictly concave.
Every global maximizer lies in that neighborhood after the phase gauge,
by the compact localization already proved.  It must therefore be this
critical branch.  This justifies applying the local analytic expansion
to the \emph{global} maximum $\mathcal M(s)$.
Differentiating the critical equations using
\eqref{I:F:eq:f2-Hessian-matrix}--\eqref{I:F:eq:f2-full-quadratic} gives
\begin{equation}
 x'(0)=(\sqrt2,0,0,0,0,0,0).
 \label{I:F:eq:f2-critical-derivative}
\end{equation}
Substituting into \eqref{I:F:eq:f2-full-quadratic} yields
\begin{equation}
 \mathcal M(s)=\frac12+\frac14s^2+O(s^3).
 \label{I:F:eq:f2-M-preliminary}
\end{equation}

\emph{Physical evenness improves the remainder.}
The feasible physical matrices with first-column squared norm $s^2$
are identical for $s$ and $-s$.  Equivalently, replacing
$(s,z,\nu)$ by $(-s,-z,\nu)$ preserves $v$ and changes
$\mathcal J$ only by its sign.  Therefore
$\mathcal M(-s)=\mathcal M(s)$.  We have proved that the maximizing
value is analytic near zero, so all its odd Taylor coefficients vanish.
Equation \eqref{I:F:eq:f2-M-preliminary} sharpens to
\begin{equation}
 \mathcal M(s)=\frac12+\frac14s^2+O(s^4).
 \label{I:F:eq:f2-M-even}
\end{equation}
Finally \eqref{I:F:eq:f2-M-definition} gives
$h(t)=t(1/2+t/4+O(t^2))^2$, which is
\eqref{I:F:eq:f2-h-expansion}.
\end{proof}


\subsection{Uniform scalar localization and the matching lower bound}

\begin{theorem}[Global passive-tuple asymptotic]
\label{I:F:thm:f2-passive-asymptotic}
The maximum over all complex passive tuples in
\eqref{I:F:eq:f2-P-definition} satisfies
\begin{equation}
 \mathcal P(q)=\frac14+\frac14q+O(q^2)
 \qquad(q\downarrow0).
 \label{I:F:eq:f2-P-expansion}
\end{equation}
\end{theorem}

\begin{proof}
At the trial point $t=q$, the denominator in
\eqref{I:F:eq:f2-scalar-reduction} equals $q^2$ exactly.  Lemma~
\ref{I:F:lem:f2-h-expansion} therefore gives the lower bound
\begin{equation}
 \mathcal P(q)\ge\frac{h(q)}q
              =\frac14+\frac14q+O(q^2).
 \label{I:F:eq:f2-P-trial}
\end{equation}
We next localize every global scalar maximizer.  For $0\le q,t\le1$,
\[
 1-\sqrt{(1-q)(1-t)}
 \ge\frac{q+t-qt}{2}\ge\frac{q+t}{4}.
\]
Together with $h(t)\le t$, this gives, for $x=t/q$,
\begin{equation}
 \frac{q h(t)}{(1-\sqrt{(1-q)(1-t)})^2}
 \le\frac{16qt}{(q+t)^2}
 =\frac{16x}{(1+x)^2}.
 \label{I:F:eq:f2-scalar-coarse-localization}
\end{equation}
For $x\notin I=[1/100,100]$ the last expression is less than $1/5$.
The trial value \eqref{I:F:eq:f2-P-trial} exceeds $1/5$ for all sufficiently
small $q$.  Thus every global maximizing point has $x\in I$.
In particular $t=qx\to0$ uniformly on all possible maximizers, which
permits the global envelope expansion to be used uniformly.

On the fixed compact interval $I$, Taylor expansion gives
\begin{align}
 1-\sqrt{(1-q)(1-qx)}
 &=\frac q2(1+x)+\frac{q^2}{8}(x-1)^2+O_I(q^3),\notag\\
 \frac{q h(qx)}{(1-\sqrt{(1-q)(1-qx)})^2}
 &=f_0(x)+qf_1(x)+O_I(q^2),
 \label{I:F:eq:f2-uniform-scalar-expansion}
\end{align}
where
\begin{equation}
 f_0(x)=\frac{x}{(1+x)^2},\qquad
 f_1(x)=\frac{x}{(1+x)^2}
       \left(x-\frac{(x-1)^2}{2(1+x)}\right).
 \label{I:F:eq:f2-f0-f1}
\end{equation}
The remainder is uniform in $x\in I$: the numerator expansion comes
from \eqref{I:F:eq:f2-h-expansion}, and the leading denominator is bounded
away from zero after division by $q^2$.

The exact identity
\[
 \frac14-f_0(x)=\frac{(x-1)^2}{4(1+x)^2}
\]
gives a quadratic maximum at $x=1$.  On $I$ there is $c_I>0$ such that
$f_0(x)\le1/4-c_I(x-1)^2$.  Also $f_1(1)=1/4$ and
$|f_1(x)-f_1(1)|\le L_I|x-1|$ for a finite Lipschitz constant $L_I$.
Completing the square yields uniformly on $I$
\begin{align*}
 f_0(x)+qf_1(x)
 &\le\frac14+\frac q4-c_I(x-1)^2+qL_I|x-1|\\
 &\le\frac14+\frac q4+\frac{L_I^2}{4c_I}q^2.
\end{align*}
Combining this with \eqref{I:F:eq:f2-uniform-scalar-expansion} gives the
matching upper bound.  Equation \eqref{I:F:eq:f2-P-trial} then proves
\eqref{I:F:eq:f2-P-expansion}.
\end{proof}


\begin{theorem}[The two-state optimum to first boundary order]
\label{I:F:thm:f2-global-first-correction}
For each coefficient field $\mathbb F\in\{\FC,\FR\}$, the squared
structured two-state optimum for $f_{\sqrt{1-q},2}$ satisfies
\begin{equation}
 \boxed{\displaystyle
 D_2^{\mathbb F}(q)=\frac34-\frac14q+O(q^2)
 \qquad(q\downarrow0).}
 \label{I:F:eq:f2-global-first-correction}
\end{equation}
In particular
$0\le D_2^{\FR}(q)-D_2^{\FC}(q)=O(q^2)$.
This asserts equality of the displayed asymptotic coefficients, not
finite-$q$ equality of the real and complex optima.
\end{theorem}

\begin{proof}
Equations \eqref{I:F:eq:f2-competitor-lower} and
\eqref{I:F:eq:f2-P-expansion} apply to every bounded complex rank-at-most-two
word-Hankel competitor.  Taking the infimum yields
\[
 D_2^{\FC}(q)\ge\frac34-\frac14q-O(q^2).
\]
Since the complex competitor class contains the real class,
$D_2^{\FC}(q)\le D_2^{\FR}(q)$.  The independent real construction of
Theorem~\ref{I:F:thm:f2-improvement}, specifically
\eqref{I:F:eq:f2-optimum-upper}, gives
$D_2^{\FR}(q)\le3/4-q/4+O(q^2)$.  These inequalities squeeze each field's
optimum to \eqref{I:F:eq:f2-global-first-correction} and give the stated
$O(q^2)$ bound on their possible gap.
\end{proof}


The analytic proof supplies a neighborhood and bounded remainder
constants, but no numerical threshold $q_0$ or explicit numerical
remainder constant is provided here.  In particular,
\eqref{I:F:eq:f2-P-expansion} does not prove the stronger proposed exact
inequality $4r^2q\Fnorm{\ell B^2(I-rA)^{-1}}^2\le1$ for every sufficiently
large $r<1$.  Nor does \eqref{I:F:eq:f2-global-first-correction} determine the
full function $D_2^{\mathbb F}(q)$ away from the boundary, its second
correction coefficient, or exact real--complex equality at positive
$q$.

\section{A quantitative global synthesis scheme by radial damping}
\label{I:F:sec:g1-damping}

This section gives an explicit solution of the additive-certification task G1
for supplied stable realizations over the complex rationals. It uses a finite
net, rather than quantifier elimination. Its cost is exponential in the
reduced parameter count, and its certificate ledger may have the same size.
Neither a new short optimality certificate nor publication priority is claimed.

Let
\[
 f(w)=cA_wb,\qquad \Phi_A(X)=\sum_{a\in\Sigma}A_aXA_a^*,
 \qquad \rho(\Phi_A)<1,
\]
be an $n$-state input realization with complex rational entries. The norm
and comparison class are the ordinary counting-measure word-Hankel norm
and all bounded complex word-Hankel competitors of rank at most $k$.
Define the input moment
\begin{equation}
 J_2(f)^2=\sum_{w\in\Sigma^*}(|w|+1)|w|^2|f(w)|^2.
 \label{I:F:eq:g1-input-moment}
\end{equation}
This is finite under the supplied joint stability assumption.

\begin{theorem}[Certified additive synthesis from complete rational input]
\label{I:F:thm:g1-net}
Fix $1\le k<n$ and rational $\varepsilon>0$. From the input one can compute
rational bounds $M\ge\max\{1,\Fnorm{H_f}\}$ and
$J\ge\max\{1,J_2(f)\}$. Put
\begin{align}
 e&=\min\{\varepsilon,1\},& B&=2M,&
 \gamma&=1-\frac{e}{3J},& \mu&=1-\gamma^2,\nonumber\\
 \eta&=\frac{e\mu^2}{15B},&
 v&=\left\lceil\sqrt{2dk^2}\right\rceil,&
 p_0&=2dk^2+4k.
 \label{I:F:eq:g1-parameters}
\end{align}
An explicit enumeration of at most
\begin{equation}
 N\le\left(3+\frac{2v}{\eta}\right)^{p_0}
   \le\left(3+\frac{240vMJ^2}{e^3}\right)^{p_0}
 \label{I:F:eq:g1-net-size}
\end{equation}
stable rational $k$-state candidates returns a competitor $g$ and rational
$L,U$ satisfying
\begin{equation}
 0\le L\le\delta_k(H_f)\le\Fnorm{H_f-H_g}\le U,
 \qquad U-L\le\varepsilon.
 \label{I:F:eq:g1-global-gap}
\end{equation}
Every candidate has a joint column-contraction margin $\mu$. Let
\[
 Z=M+B/\mu,\qquad
 p_{\rm tr}=\left\lceil
 \frac{18Z\max\{1,\lceil\log_2(n+k)\rceil\}}e\right\rceil.
\]
Construction and independent verification cost
$N\,\operatorname{poly}(n+k,d,H,p_{\rm tr},\log(Z/e))$ bit operations,
where $H$ bounds the rational input and grid coefficient heights after a
common denominator is chosen. A complete interval ledger has the same
size bound. Thus for fixed $d,k$ and polynomially bounded numerical $M,J$,
this is a polynomial-time additive approximation scheme in input length,
$n$, and $1/e$. No polynomial dependence on $d,k$ or $\log(1/e)$ is asserted.
\end{theorem}

The parameter restriction applies to a known nearoptimal damped competitor,
not to the original optimum. The proof below quantifies its loss using the
input alone, thereby avoiding a restricted-comparison-class gap.

\subsection{Input moments and radial comparison}

Write $R=(I-\Phi_A)^{-1}$ and $S_0=bb^*$. Four successive rational Stein
solves give $M_j=R^{j+1}S_0$, $0\le j\le3$:
\[
 M_0=S_0+\Phi_A(M_0),\qquad
 M_j=M_{j-1}+\Phi_A(M_j)\quad(1\le j\le3).
\]
All are positive semidefinite. Counting word splittings and using the scalar
generating functions gives
\begin{align}
 \Fnorm{H_f}_{\rm HS}^2&=cM_1c^*,\label{I:F:eq:g1-HS-input}\\
 J_2(f)^2&=c\bigl(2\Phi_A(M_2)+6\Phi_A^2(M_3)\bigr)c^*.
 \label{I:F:eq:g1-moment-Stein}
\end{align}
Indeed the coefficient energy at length $t$ is
$c\Phi_A^t(S_0)c^*$, each coefficient occurs $t+1$ times in its Hankel,
and
\[
 \sum_{t\ge0}t^2(t+1)x^t
 =\frac{2x}{(1-x)^3}+\frac{6x^2}{(1-x)^4}.
\]
Integer ceilings of the square roots in
\eqref{I:F:eq:g1-HS-input}--\eqref{I:F:eq:g1-moment-Stein}, bounded below by one,
supply rational $M,J$. These calculations use no approximation of
$\delta_k$.

For $0<\gamma<1$, let $D_\gamma e_w=\gamma^{|w|}e_w$ and
$h_\gamma(w)=\gamma^{|w|}h(w)$. Entrywise,
\begin{equation}
 H_{h_\gamma}=D_\gamma H_hD_\gamma.
 \label{I:F:eq:g1-radial-Hankel}
\end{equation}
This preserves the rank upper bound, and its realization is obtained by
replacing every letter matrix by $\gamma A_a$. For the input,
\begin{equation}
 \Fnorm{H_f-H_{f_\gamma}}
 \le\Fnorm{H_f-H_{f_\gamma}}_{\rm HS}
 \le(1-\gamma)J_2(f),
 \label{I:F:eq:g1-radial-bias}
\end{equation}
because $1-\gamma^t\le t(1-\gamma)$.

Let $G=H_h$ attain $\delta_k(H_f)$. Since zero is admissible,
$\Fnorm G\le2\Fnorm{H_f}\le2M$. The canonical observable realization on
$\operatorname{ran}G$, using right backward shifts, has parameters
$T_a,\ell,\beta$ satisfying
\[
 \sum_aT_a^*T_a+\ell^*\ell=I,
 \qquad\Fnorm\ell\le1,
 \qquad\Fnorm\beta=\Fnorm{Ge_\varepsilon}\le B.
\]
Padding to $k$ states changes the equality to an inequality if needed.
Therefore its damped realization belongs to
\begin{equation}
 \sum_a\widehat T_a^*\widehat T_a\preceq\gamma^2I,
 \qquad\Fnorm\ell\le1,
 \qquad\Fnorm\beta\le B.
 \label{I:F:eq:g1-compact-class}
\end{equation}
Using that $D_\gamma$ is a contraction gives
\begin{align}
 \Fnorm{H_f-H_{h_\gamma}}
 &\le\Fnorm{H_f-H_{f_\gamma}}
       +\Fnorm{D_\gamma(H_f-G)D_\gamma}\nonumber\\
 &\le\delta_k(H_f)+(1-\gamma)J_2(f)
 \le\delta_k(H_f)+e/3.
 \label{I:F:eq:g1-nearoptimal-damped}
\end{align}
The original optimizer may have an arbitrarily poor stability margin; no
such margin enters this argument.

\subsection{An infinite-Hankel Lipschitz modulus}

Consider two parameter triples $(A,c,b)$ and $(A',c',b')$ satisfying
\eqref{I:F:eq:g1-compact-class}. Suppose
\[
 \Fnorm{\begin{bmatrix}A_a-A'_a\end{bmatrix}^{\rm col}}\le\eta,
 \qquad\Fnorm{c-c'}\le\eta,
 \qquad\Fnorm{b-b'}\le B\eta.
\]
Joint column contractivity gives
$\sum_{|w|=t}\Fnorm{A_wx}^2\le\gamma^{2t}\Fnorm x^2$.
Telescoping the product $A_w-A'_w$ and applying this energy inequality on
both sides of each differing letter gives, at each word length $t$,
\[
 \left(\sum_{|w|=t}|cA_wb-c'A'_wb'|^2\right)^{1/2}
 \le2B\eta\gamma^t+Bt\eta\gamma^{t-1}.
\]
For $t=0$ the last term is zero. The $t$ product-replacement terms are
summed by Minkowski in the length-$t$ coefficient space. A second
Minkowski inequality with weights $t+1$ yields
\begin{align}
 \Fnorm{H_h-H_{h'}}
 &\le\Fnorm{H_h-H_{h'}}_{\rm HS}\nonumber\\
 &\le\frac{2B\eta}{\mu}
       +\frac{B\eta\sqrt{2+4\gamma^2}}{\mu^2}
 <\frac{5B\eta}{\mu^2}=e/3.
 \label{I:F:eq:g1-Lipschitz}
\end{align}
Here
\[
 \sum_{t\ge0}(t+1)\gamma^{2t}=\mu^{-2},\qquad
 \sum_{t\ge1}(t+1)t^2\gamma^{2t-2}
 =(2+4\gamma^2)\mu^{-4}.
\]
This supplies the summability needed for the infinite counting-measure
operator norm, rather than transferring a worst-word bound without a tail
argument.

\subsection{A rational stable net}

There are $N_A=2dk^2$ real coordinates in the complex vertical letter stack
and $2k$ in each boundary vector. Set
$v_A=\lceil\sqrt{N_A}\rceil$ and $v_c=\lceil\sqrt{2k}\rceil$.
Use the following finite lattices, including ceiling endpoints.
\begin{enumerate}[leftmargin=*]
\item For each real letter coordinate in $[-\gamma,\gamma]$, use step
$\eta/v_A$. Retain a raw stack only if
$\Fnorm{A^{\rm raw}}\le\gamma+\eta/2$, checked by the rational Hermitian
inequality $(\gamma+\eta/2)^2I-(A^{\rm raw})^*A^{\rm raw}\succeq0$.
Scale it by $\gamma/(\gamma+\eta/2)$.
\item For each real readout coordinate in $[-1,1]$, use step $\eta/v_c$.
Retain $\Fnorm{c^{\rm raw}}\le1+\eta/2$ and divide by $1+\eta/2$.
\item For each real initial-state coordinate in $[-B,B]$, use step
$B\eta/v_c$. Retain $\Fnorm{b^{\rm raw}}\le B(1+\eta/2)$ and divide by
$1+\eta/2$.
\end{enumerate}
Every retained output is rational and satisfies
\eqref{I:F:eq:g1-compact-class}. Every target triple in that class is covered:
rounding contributes at most $\eta/2$, $\eta/2$, $B\eta/2$ respectively,
and final scaling contributes at most the same amount. Thus the total
errors meet \eqref{I:F:eq:g1-Lipschitz}. The number of choices per real
coordinate is at most $3+2v_A/\eta$, proving the first inequality in
\eqref{I:F:eq:g1-net-size}. Since $e\le1$, $J\ge1$,
\[
 \mu\ge e/(2J),\qquad \eta\ge e^3/(120MJ^2),
\]
which proves the second count. The known column bound makes every candidate
mean-square stable, with CP spectral radius at most $\gamma^2$.

\subsection{Finite rational norm certificates and the global lower bound}

For each candidate form the stable block realization of $f-g_i$, of
dimension $m_0=n+k$, and solve its rational Stein equations
\[
 P_i=b_ib_i^*+\sum_aD_{i,a}P_iD_{i,a}^*,\qquad
 Q_i=c_i^*c_i+\sum_aD_{i,a}^*Q_iD_{i,a}.
\]
Its squared Hankel norm is the largest eigenvalue of $P_iQ_i$; all these
eigenvalues are nonnegative. No matrix square root is needed for the
following interval computation. Let $p=p_{\rm tr}$ and
$a_i=\operatorname{tr}((P_iQ_i)^p)$. If $a_i=0$, the error norm is zero.
Otherwise
\begin{equation}
 (a_i/m_0)^{1/(2p)}\le\Fnorm{H_f-H_{g_i}}
                         \le a_i^{1/(2p)}.
 \label{I:F:eq:g1-trace-interval}
\end{equation}
The uniform norm bound $Z=M+B/\mu$ follows from
$\Fnorm{H_{g_i}}\le\Fnorm{H_{g_i}}_{\rm HS}\le B/\mu$.
The width of the exact endpoints in \eqref{I:F:eq:g1-trace-interval} is at most
\[
 Z\bigl(m_0^{1/(2p)}-1\bigr)
 \le\frac{Z\log m_0}{p}\le e/18,
\]
using $p\ge\log m_0$ and $e^x-1\le2x$ for $0\le x\le1/2$.
Enclose the lower root from below and the upper root from above, each
within $e/18$, by rational bisection on $[0,2Z]$. Its comparisons are
inequalities between rational $2p$th powers. The resulting interval
$[a_i^-,u_i]$ contains the norm and has width at most $e/6$.

Choose a candidate minimizing $u_i$, let $U$ be that endpoint, and set
$L=\max\{0,U-5e/6\}$. The existence of the damped competitor
\eqref{I:F:eq:g1-nearoptimal-damped}, its net approximation
\eqref{I:F:eq:g1-Lipschitz}, and the norm-interval width imply
\[
 U\le\delta_k(H_f)+e/3+e/3+e/6
   =\delta_k(H_f)+5e/6.
\]
The chosen candidate is admissible and has actual error at most $U$.
This proves \eqref{I:F:eq:g1-global-gap}.

The verifier may rerun the finite net or inspect a full ledger containing
the exact Stein solutions, rational trace powers and root brackets. It
checks the interval widths and the minimum-upper-endpoint choice. The
ledger and its verification may have $N$ entries. A short certificate is
not asserted.

\subsection{Complexity qualifications}

For rational input, vectorized Stein systems have order $(n+k)^2$ and
rational coefficients. Fraction-free exact linear algebra, rational matrix
powers to exponent $p_{\rm tr}$, and scalar rational-power bisection have
polynomial bit cost in the parameters stated in Theorem~\ref{I:F:thm:g1-net}.
All input/grid rational denominators can be combined explicitly; their
heights are polynomial in input bit length, $\log M$, $\log J$,
$\log(1/e)$, $d$ and $k$. Also
$p_{\rm tr}=O(MJ\log(n+k)/e^2)$. Thus the full precision and verification
costs are included in the finite bound.

A checkable positive class fixes $d,k$ and bounds the input moments $M,J$
by constants or by specified polynomials in input dimension. Membership is
checked through the input Stein computations; it is not defined by an easy
optimum. On this class the construction is polynomial in $n$, rational
input length and $1/\varepsilon$. This does not establish a strict
complexity improvement over every optimized QE formulation or earlier
algorithm; that comparison and scientific priority need a separate audit.

For individually encoded algebraic entries, the compositum field degree
can be exponential in the total encoding length. A polynomial bit claim
therefore needs a supplied common number field and its degree, height and
embedding costs as parameters. The rational subclass avoids this issue.
Finally, this scheme bypasses finite-block completion: it supplies no
quantitative convergence modulus for the original $\beta_{k,L}$ hierarchy.

\section{The resulting theorem and its scope}
Let $\Sigma=\{a,b\}$, let $H_f[u,v]=f(uv)$ on $\ell^2(\Sigma^*)$, and put
\[
 q=1-r^2\in(0,1),\qquad
 f_{r,d}\text{ denotes \eqref{I:eq:binary-source} with terminal index }d.
\]
For a field $\mathbb F\in\{\RR,\RC\}$ define
\[
 D_{n,d}^{\mathbb F}(q)=
 \inf_{\substack{H_g\text{ bounded word-Hankel over }\mathbb F\\
                 \Rrank H_g\le n}}
       \Rnorm{H_{f_{r,d}}-H_g}^{2}.
\]
Write $D_k^{\mathbb F}=D_{k,k}^{\mathbb F}$. Ordinary matrix rank is
used throughout; the comparison class contains every bounded structured
Hankel of that rank. The letter matrices are required to be jointly
mean-square stable only for the explicit constructions, not as an
extra uniform-margin restriction on the optimum.

For integers $p\ge j\ge1$, let
\[
 \gamma_{j,p}^{\mathbb F}
 =\max_{D^*D+l^*l\preceq I_j}\Rnorm{lD^p},\qquad
 (E_{j,p}^{\mathbb F})^2=1-(\gamma_{j,p}^{\mathbb F})^2.
\]
Equivalently, $E_{j,p}^{\mathbb F}$ is the distance in scalar $H^2$
from $z^p$ to stable rational coefficient sequences of Hankel rank at
most $j$, over the indicated field. Let $\mathcal M_{j,p}^{\mathbb F}$
contain the lossless triples $(D,l,x)$ with
\[
 D^*D+l^*l=I_j,\quad \Rnorm x=1,\quad
 lD^p x=\gamma_{j,p}^{\mathbb F}>0.
\]
Every such $D$ is strictly stable and $x$ is cyclic. For one such triple,
write $\gamma=\gamma_{j,p}^{\mathbb F}$ and define
\begin{align*}
 a_i&=lD^i x\quad(0\le i\le p),\\
 S&=\sum_{i=0}^{p-1}a_i\,x^*D^{p-i}x,
 &v&=\sum_{i=0}^{p-2}a_iD^{p-2-i}x,\\
 K&=\RHerm\!\left(\sum_{i=0}^{p}D^{p-i}x\,lD^i\right),
 &\RHerm T&=(T+T^*)/2.
\end{align*}
The empty sum defining $v$ is zero. The response matrix $K$ is positive
definite. Define the finite variational coefficient
\begin{equation}
 b_{n,d}^{\mathbb F}=
 \max_{(D,l,x)\in\mathcal M_{n-1,d-1}^{\mathbb F}}
 \left\{|a_{d-2}|^2-\gamma\Re S+
                     \gamma v^*K^{-1}v\right\}.
 \label{I:R:eq:main-response-coefficient}
\end{equation}
This maximization is part of the answer: the general theorem does not
assert a closed-form evaluation or an efficient optimizer for every
budget.

\begin{theorem}[First-order structured approximation at every fixed budget]
\label{I:R:thm:main-general}
For every fixed pair of integers $2\le n\le d$, over each field separately,
\begin{equation}
 \boxed{\displaystyle
 D_{n,d}^{\mathbb F}(q)=
 (E_{n-1,d-1}^{\mathbb F})^2-b_{n,d}^{\mathbb F}q+o(q),
 \qquad q\downarrow0.}
 \label{I:R:eq:main-general}
\end{equation}
The maximum in \eqref{I:R:eq:main-response-coefficient} exists. No uniqueness
or nondegenerate Hessian hypothesis is imposed on its scalar extremizers.
There is a stable $n$-state construction with squared error
$(E_{n-1,d-1}^{\mathbb F})^2-b_{n,d}^{\mathbb F}q+O(q^2)$.
The global lower remainder is $o(q)$ in general.
On the diagonal $n=d=k$, the leading constants and the response
coefficients agree over the two fields. Consequently
$D_k^{\RR}(q)-D_k^{\RC}(q)=o(q)$.
The general diagonal field comparison uses the external scalar-extremal
classification identified in Section~\ref{I:R:subsec:real-diagonal-orbits};
the explicit cases through $k=4$ use the published low-degree extremals.
\end{theorem}

The proof has two distinct parts. Section~\ref{I:R:sec:all-budget-response}
proves the unrestricted complex passive-envelope response and its
real analogue. Section~\ref{I:R:sec:v4-matching} constructs a genuine stable
word-Hankel attaining that response to first order. A scalar compactness
argument then gives \eqref{I:R:eq:main-general}. The independent-budget
version changes the fast state dimension to $j=n-1$ and the coefficient
index to $p=d-1$ in both proofs.

\begin{corollary}[Explicit consequences]
\label{I:R:cor:main-explicit}
For each field separately,
\[
 \boxed{D_3^{\mathbb F}(q)=\frac9{16}-\frac17q+O(q^2).}
\]
At four states,
\[
 \boxed{D_4^{\mathbb F}(q)=
 \frac{1103-33\sqrt{33}}{2048}
 -\frac{855+103\sqrt{33}}{15872}q+o(q).}
\]
More generally, for every fixed $d\ge2$, put
$a_d=(d-1)^{d-1}/d^d$. Then
\[
 \boxed{D_{2,d}^{\mathbb F}(q)=
       1-a_d-\frac{a_d}{(d-1)^2}q+O(q^2).}
\]
These common expansions do not assert finite-$q$ equality between the
real and complex optima.
\end{corollary}

On the diagonal $n=d=k$, the earlier Hardy identification is
\[
 E_{k-1,k-1}=\frac1{C(k-1,2/3)},\qquad
 C(m,p)=\sup_{\Rnorm F_{H^p}=1}|[z^m]F|.
\]
The Hardy constants and outer-extremal structure belong to the cited
sources~\cite{brevig2020,outer2026}. Their use here is as scalar input to the
passive response and structured-completion arguments. The value $9/16$
at three states follows from $C(2,2/3)=4/3$; the new calculation here is the finite-parameter coefficient and its global matching proof.

\subsection{How the rank-one benchmark constrains the extension}
The exact one-state optimum is the unique $\tau\in(q,1)$ satisfying
\[
 \boxed{1-\tau=
 \frac{q}{\left(1+\sqrt{q\tau/(\tau-q)}\right)^2}.}
\]
Section~\ref{I:R:sec:full-channel} derives this monotone radical formulation
from the complete target singular decomposition and relates it to the
selected quartic. The dominant-vector bound alone would instead give
$1-q/(1+\sqrt q)^2$. It misses the true error at order $q^{5/2}$.
The higher-state first-order result therefore needs the separate
structured upper construction; sharpness of a projection envelope is
insufficient as a general finite-$q$ principle.

\subsection{Canonical row spaces and the global projection bound}
Let $G=H_g$ be any bounded rank-at-most-$n$ word-Hankel, and regard
its coefficient rows as columns. Their span is
$S=\operatorname{ran}G^{\mathsf T}$, where coefficient transpose
preserves the operator norm. If $L_c e_w=e_{cw}$, the identity
$L_c^*(G^{\mathsf T}e_u)=G^{\mathsf T}e_{uc}$ makes $S$ invariant
under $L_a^*,L_b^*$. In orthonormal coordinates on $S$, let
$A=L_a^*|_S$, $B=L_b^*|_S$, and $\ell x=x[\varepsilon]$.
Partitioning words by their first letter gives
\[
 A^*A+B^*B+\ell^*\ell=I_S.
\]
The observation map $Ox=(\ell T_{\operatorname{rev}(w)}x)_w$
is the isometric inclusion of $S$. Moreover
\[
 \sum_{|w|=h}\Rnorm{T_{\operatorname{rev}(w)}x}^2
       =\sum_{|v|\ge h}|x[v]|^2\longrightarrow0.
\]
Finite dimensionality makes this decay uniform on the unit sphere,
so the tuple is jointly mean-square stable. Padding smaller spaces
preserves the energy equality as an inequality.

For the target, put
\[
 u_r=\sqrt q\sum_{j\ge0}r^je_{a^j},\qquad
 v_r=\sqrt q\sum_{j\ge0}r^je_{a^jb^d}.
\]
Both vectors are unit, and $H_{f_{r,d}}^{\mathsf T}u_r=v_r$.
Since $G^{\mathsf T}u_r\in S$,
\[
 \Rnorm{H_{f_{r,d}}-G}^2\ge1-\Rnorm{P_Sv_r}^2,
 \qquad
 \Rnorm{P_Sv_r}^2
   =q\Rnorm{\ell B^d(I-rA)^{-1}}^2.
\]
Thus maximizing the right overlap over all passive $n$-state tuples
is a valid lower bound against the entire comparison class. This
argument imposes no prescribed competitor stability margin.


% Begin all_budget_response.tex

\section{An all-budget first variation of the passive envelope}
\label{I:R:sec:all-budget-response}

Let $k\ge2$, put $m=k-1$, and let $q=1-r^2\in(0,1)$.  All state
spaces in this section are complex Euclidean spaces.  Define
\begin{equation}
 \mathcal P_k(q)=
 \max_{A^*A+B^*B+\ell^*\ell\preceq I_k}
 q\Rnorm{\ell B^k(I-rA)^{-1}}^2
 \label{I:R:eq:Pk-definition}
\end{equation}
and
\begin{equation}
 h_k(t)=\max\left\{|\ell B^ke_1|^2:
 V=\begin{pmatrix}B\\ \ell\end{pmatrix},\quad
 V^*V\preceq I_k,\quad \Rnorm{Ve_1}^2=t\right\},
 \qquad 0\le t\le1.
 \label{I:R:eq:hk-definition}
\end{equation}
Both maxima are over compact sets; the inverse in
\eqref{I:R:eq:Pk-definition} exists because $\Rnorm A\le1$ and $r<1$.
The intrinsic scalar boundary constant is
\begin{equation}
 \Gamma_m=\max_{D^*D+l^*l\preceq I_m}\Rnorm{lD^m},
 \qquad 0<\Gamma_m<1.
 \label{I:R:eq:Gamma-intrinsic}
\end{equation}
Its classical observation-space interpretation gives
$E_m^2=1-\Gamma_m^2$, where $E_m$ is the distance from $z^m$ to scalar
stable rational series of Hankel rank at most $m$.  The separate
identification $E_m=1/C(m,2/3)$ uses the polynomial coefficient-extremal
structure of Brevig--Saksman~\cite{brevig2020} and the outer-extremal theorem
of Brevig--S\o vik~\cite{outer2026}.  The argument below uses only
\eqref{I:R:eq:Gamma-intrinsic} and classical finite model spaces.

This section generalizes the passive projection mechanism.
Section~\ref{I:R:sec:v4-matching} then constructs a word-Hankel
approximant attaining its first-order response.

\subsection{Exact elimination of the slow transition}

\begin{proposition}[Exact scalar reduction in every dimension]
\label{I:R:prop:allk-eliminate-A}
For every $k\ge1$ and $0<q<1$,
\begin{equation}
 \boxed{\displaystyle
 \mathcal P_k(q)=\max_{0\le t\le1}
 \frac{q h_k(t)}{(1-r\sqrt{1-t})^2}.}
 \label{I:R:eq:allk-elimination}
\end{equation}
\end{proposition}
\begin{proof}
For fixed $B,\ell$, put $S=I-B^*B-\ell^*\ell\succeq0$.  The row
operator norm can be written as the supremum of
$|\ell B^ky|/\Rnorm{(I-rA)y}$ over unit vectors $y$.  Set
$t=\Rnorm{By}^2+|\ell y|^2$, so $y^*Sy=1-t$.  Every feasible $A$ satisfies
\[
 \Rnorm{(I-rA)y}\ge1-r\Rnorm{Ay}\ge1-r\sqrt{1-t}.
\]
If $t<1$, choose
\[
 A=\frac{y(y^*S)}{\sqrt{1-t}}.
\]
Then $Ay=\sqrt{1-t}\,y$, and the positive-semidefinite
Cauchy--Schwarz inequality gives
\[
 A^*A=\frac{Syy^*S}{y^*Sy}\preceq S.
\]
Thus the denominator bound is attained.  If $t=1$, choose $A=0$.
Conjugating the state coordinates to send $y$ to $e_1$ proves both
inequalities in \eqref{I:R:eq:allk-elimination}.
\end{proof}


\subsection{Cyclicity and quantitative control of the passivity defect}

We first prove the structural fact responsible for normal coercivity.
For integers $n\ge j\ge0$, write
\[
 \gamma_{j,n}=\max\{\Rnorm{lD^n}:D^*D+l^*l\preceq I_j\},
 \qquad \gamma_{0,n}=0.
\]

\begin{lemma}[Strict improvement with the scalar state dimension]
\label{I:R:lem:strict-dimension-improvement}
Fix $n\ge1$.  Then
$\gamma_{0,n}<\gamma_{1,n}<\cdots<\gamma_{n,n}$.
\end{lemma}
\begin{proof}
Identify $H^2$ with its square-summable coefficient sequences and let
$S^*$ denote the backward shift.  The observation-map formulation is
\begin{equation}
 \gamma_{j,n}=
 \sup_{\substack{K\subset H^2,\ S^*K\subset K\\ \dim K\le j}}
 \Rnorm{P_K z^n}.
 \label{I:R:eq:gamma-model-space}
\end{equation}
Indeed, a passive pair $(D,l)$ has a contractive observation map
$Ox=(lD^ax)_{a\ge0}$, with finite-dimensional invariant range $K$.
Consequently
$|lD^nx|\le\Rnorm{P_Kz^n}\Rnorm{Ox}\le\Rnorm{P_Kz^n}\Rnorm x$.
Conversely, on a finite-dimensional invariant $K$, take
$D=S^*|_K$ and $lx=x[0]$.  Then $D^*D+l^*l=I_K$ and
$\Rnorm{lD^n}=\Rnorm{P_Kz^n}$.  This proves
\eqref{I:R:eq:gamma-model-space}.  Compactness of passive parameter sets
also shows that its supremum is attained.

The classical finite-model-space description writes
$K=K_\Theta=H^2\ominus\Theta H^2$ for a finite Blaschke product
$\Theta$.  If $\dim K\le j\le n$, then $z^n\notin K$, because the
backward-shift cyclic span of $z^n$ has dimension $n+1$.  Thus
$f=z^n-P_Kz^n$ is a nonzero member of $\Theta H^2$.  Reproducing kernels
span $H^2$, so there is a point $a$ in the disk with
$\langle f,\Theta k_a\rangle\ne0$.  If $b_a$ is the elementary
Blaschke factor at $a$, then
\[
 K_{\Theta b_a}=K_\Theta\mathbin{\oplus}\Theta\operatorname{span}\{k_a\}
\]
has one additional dimension and strictly increases the projection
norm of $z^n$.  Apply this construction to an attained $j$-state maximum
for each $j<n$.
\end{proof}


\begin{lemma}[Every scalar boundary maximizer is cyclic and lossless]
\label{I:R:lem:cyclic-lossless-maximizer}
Suppose $D^*D+l^*l\preceq I_m$, $\Rnorm x=1$, and
$|lD^mx|=\Gamma_m$.  Then $x$ is cyclic for $D$,
$D^*D+l^*l=I_m$, and $\rho(D)<1$.
\end{lemma}
\begin{proof}
The observation map $O$ is contractive and its range is backward-shift
invariant of dimension at most $m$.  Hence
\[
 \Gamma_m=|lD^mx|\le\Gamma_m\Rnorm{Ox}\le\Gamma_m,
\]
which forces $\Rnorm{Ox}=1$.  If the cyclic span of $x$ under $D$ had
dimension below $m$, restricting to that span would contradict
Lemma~\ref{I:R:lem:strict-dimension-improvement}, with $n=m$.
Put $T=I-D^*D-l^*l\succeq0$.  The finite-horizon energy identity is
\[
 1=\sum_{j=0}^{N-1}|lD^jx|^2
   +\sum_{j=0}^{N-1}\Rnorm{T^{1/2}D^jx}^2+\Rnorm{D^Nx}^2.
\]
Since the first sum tends to $\Rnorm{Ox}^2=1$, all the defect terms
vanish and $D^Nx\to0$.  Cyclicity gives $T=0$ and $D^N\to0$ on the
whole state space.  The latter is equivalent to $\rho(D)<1$ in finite
dimension.
\end{proof}


Define the compact normalized boundary set
\begin{equation}
 \mathcal M_m=\{(D,l,x):D^*D+l^*l=I_m,\quad \Rnorm x=1,
                    \quad lD^mx=\Gamma_m>0\}.
 \label{I:R:eq:boundary-maximizer-set}
\end{equation}
It is nonempty by attainment and Lemma~\ref{I:R:lem:cyclic-lossless-maximizer}.
Every member has an invertible cyclic matrix
$[x,Dx,\ldots,D^{m-1}x]$.  Compactness therefore bounds its least
singular value uniformly away from zero.  This also proves
$0<\Gamma_m<1$: positivity follows from a scalar stable realization,
and equality to one would put $z^m$ in an invariant space of dimension
at most $m$.

\begin{lemma}[Quantitative passivity-defect control]
\label{I:R:lem:quantitative-defect-control}
There are $C_m,\varepsilon_m>0$ such that every passive $(D,l)$ with
$\Gamma_m^2-\Rnorm{lD^m}^2\le\varepsilon_m$ satisfies
\begin{equation}
 \operatorname{tr}(I-D^*D-l^*l)
 \le C_m\bigl(\Gamma_m^2-\Rnorm{lD^m}^2\bigr).
 \label{I:R:eq:defect-gap-control}
\end{equation}
\end{lemma}
\begin{proof}
Put $a=lD^m$, $\mu=\Rnorm a$, and $x=a^*/\mu$.
By compactness and Lemma~\ref{I:R:lem:cyclic-lossless-maximizer}, when
$\mu$ is sufficiently close to $\Gamma_m$ the matrix
$K_x=[x,Dx,\ldots,D^{m-1}x]$ has least singular value at least a fixed
$\kappa_m>0$.  With $T=I-D^*D-l^*l$, the observation-space comparison
and telescoping give
\begin{align*}
 \mu^2&\le\Gamma_m^2\Rnorm{Ox}^2,\\
 \sum_{j=0}^{m-1}\Rnorm{T^{1/2}D^jx}^2
 &\le1-\Rnorm{Ox}^2
 \le\frac{\Gamma_m^2-\mu^2}{\Gamma_m^2}.
\end{align*}
The sum on the left is
$\Rnorm{T^{1/2}K_x}_{\mathrm F}^2\ge\kappa_m^2\operatorname{tr}T$.
This proves \eqref{I:R:eq:defect-gap-control}.
\end{proof}


\subsection{An explicit positive normal-response matrix}

For $(D,l,x)\in\mathcal M_m$, abbreviate $\Gamma=\Gamma_m$ and set
\begin{align}
 a_j&=lD^jx\quad(0\le j\le m),\qquad a_m=\Gamma,
 \label{I:R:eq:response-a}\\
 S&=\sum_{j=0}^{m-1}a_j\,x^*D^{m-j}x,
 &v&=\sum_{j=0}^{m-2}a_jD^{m-2-j}x,
 \label{I:R:eq:response-S-v}\\
 K&=\operatorname{Herm}\left(\sum_{j=0}^{m}D^{m-j}xlD^j\right),
 &\operatorname{Herm}(Z)&=\frac{Z+Z^*}{2}.
 \label{I:R:eq:response-K}
\end{align}
An empty sum is zero, so $v=0$ when $m=1$.

\begin{lemma}[Strict positivity of the response matrix]
\label{I:R:lem:response-K-positive}
For every $(D,l,x)\in\mathcal M_m$,
\begin{equation}
 \boxed{\displaystyle
 K\succeq\Gamma\sum_{j=0}^{m-1}D^jxx^*D^{*j}\succ0.}
 \label{I:R:eq:response-K-coercivity}
\end{equation}
Consequently $K^{-1}$ is uniformly bounded on $\mathcal M_m$.
\end{lemma}
\begin{proof}
Fix $Y\succeq0$ and perturb the lossless isometry
$U=\begin{pmatrix}D\\l\end{pmatrix}$ to
$U_\epsilon=U(I-\epsilon Y)^{1/2}$.  Its components satisfy
\[
 l_\epsilon D_\epsilon^mx
 =\Gamma-\frac\epsilon2
   \sum_{j=0}^{m}lD^jYD^{m-j}x+O(\epsilon^2).
\]
Its passivity defect is exactly $\epsilon Y$.  The observation-space
inequality and the first $m$ terms of the energy identity give
\[
 |l_\epsilon D_\epsilon^mx|^2
 \le\Gamma^2\left(1-\epsilon
   \sum_{j=0}^{m-1}x^*D^{*j}YD^jx+O(\epsilon^2)\right).
\]
Comparison of first derivatives yields
\[
 \operatorname{tr}(KY)\ge\Gamma\operatorname{tr}\left(
     \sum_{j=0}^{m-1}D^jxx^*D^{*j}Y\right).
\]
This holds for every $Y\succeq0$, proving the first inequality in
\eqref{I:R:eq:response-K-coercivity}.  Its right side is positive definite
by cyclicity.  Uniform boundedness of the inverse follows from
compactness.
\end{proof}


\begin{theorem}[All-budget first variation of the passive envelope]
\label{I:R:thm:all-budget-response}
For every fixed $k\ge2$, put $m=k-1$ and define
\begin{equation}
 b_k=\max_{(D,l,x)\in\mathcal M_m}
 \left\{|a_{m-1}|^2-\Gamma_m\operatorname{Re}S
                  +\Gamma_m v^*K^{-1}v\right\}.
 \label{I:R:eq:bk-response-formula}
\end{equation}
Then
\begin{equation}
 \boxed{\displaystyle
 h_k(t)=\Gamma_m^2t+b_kt^2+o(t^2),\qquad
 \mathcal P_k(q)=\Gamma_m^2+b_kq+o(q).}
 \label{I:R:eq:all-budget-envelope-expansion}
\end{equation}
The maximum in \eqref{I:R:eq:bk-response-formula} exists.  No uniqueness or
nondegenerate Hessian assumption on the scalar extremizers is needed.
In particular, for the binary target family, the projection lower bound
implies
\begin{equation}
 D_k(q)\ge E_{k-1}^2-b_kq+o(q).
 \label{I:R:eq:all-budget-D-lower}
\end{equation}
\end{theorem}

\begin{proof}
\emph{Exact first-column coordinates.}
Put $t=s^2$, with $s$ real near zero.  Every contraction $V$ in
\eqref{I:R:eq:hk-definition} can be written
\begin{equation}
 V=[sz,W_s(z)N],\quad \Rnorm z=1,\quad \Rnorm N\le1,\qquad
 W_s(z)=I-(1-\sqrt{1-s^2})zz^*.
 \label{I:R:eq:first-column-chart}
\end{equation}
Indeed, the Schur complement of $I-V^*V$ says that the remaining columns
$V_F$ satisfy
\[
 V_F^*\left(I+\frac{s^2}{1-s^2}zz^*\right)V_F\preceq I;
\]
its inverse positive square root is $W_s(z)$.  Partition
\[
 z=\begin{pmatrix}z_0\\z_F\\z_o\end{pmatrix},\qquad
 N=\begin{pmatrix}c\\D\\l\end{pmatrix},
 \qquad z_F\in\RC^m,
\]
where $c$ is a row.  Define $J(s,z,N)=\ell B^{m+1}e_1/s$ by its
continuous polynomial extension at $s=0$, and
$M(s)=\max_{z,N}|J(s,z,N)|$.  Then
\begin{equation}
 h_k(s^2)=s^2M(s)^2,\qquad J(0,z,N)=lD^mz_F,\qquad M(0)=\Gamma.
 \label{I:R:eq:M-and-J}
\end{equation}
At every leading maximizer, Lemma~\ref{I:R:lem:cyclic-lossless-maximizer}
forces $z_0=z_o=0$, $\Rnorm{z_F}=1$, $D^*D+l^*l=I_m$, and $c=0$.
Thus the leading maximizers are precisely the embedded scalar boundary
set, up to an output phase.

\emph{Uniform localization of the normal directions.}
Expansion of the matrix product gives
\begin{align}
 J(s,z,N)&=lD^mz_F+sJ_1(z,N)+O(s^2),
 \label{I:R:eq:J-first-order}\\
 J_1(z,N)&=z_ocD^{m-1}z_F+z_0lD^{m-1}z_F
       +\sum_{j=0}^{m-2}(lD^jz_F)(cD^{m-2-j}z_F).
 \label{I:R:eq:J1-explicit}
\end{align}
The remainder is uniform on the compact parameter domain.  Put
$\delta=\Gamma^2-|lD^mz_F|^2$.  Near the leading maximizers,
\[
 |z_0|^2+|z_o|^2\le\frac\delta{\Gamma^2},\qquad
 \Rnorm c^2\le\operatorname{tr}(I-D^*D-l^*l)\le C\delta
\]
by Lemma~\ref{I:R:lem:quantitative-defect-control}.  Therefore
$|J_1|\le C\sqrt\delta$ there.  Away from a fixed neighborhood of the
leading maximizers, $\delta$ is bounded below, so the same estimate holds
globally after enlarging $C$.  It follows that
\begin{equation}
 |J(s,z,N)|\le\Gamma-\frac\delta{2\Gamma}
                         +C|s|\sqrt\delta+Cs^2.
 \label{I:R:eq:normal-gap-upper}
\end{equation}
A fixed leading maximizer gives $M(s)\ge\Gamma-Cs^2$.
Completing the square in \eqref{I:R:eq:normal-gap-upper} proves
\begin{equation}
 M(s)=\Gamma+O(s^2),\qquad
 \delta=O(s^2)\quad\hbox{at every global maximizer.}
 \label{I:R:eq:normal-localization}
\end{equation}
In particular, all maximizing $z_0,z_o,c$ are $O(s)$ and their remaining
passivity defect is $O(s^2)$.

\emph{Polar coordinates and the complex phase.}
At any maximizing point write
\[
 z_0=s\lambda,\quad z_o=s\zeta,\quad c=sw,\quad
 z_F=\sqrt{1-s^2(|\lambda|^2+|\zeta|^2)}\,x,\qquad \Rnorm x=1.
\]
Let $Y_{\rm raw}=\begin{pmatrix}D_{\rm raw}\\l_{\rm raw}\end{pmatrix}$
denote the lower block of $N$, and put
$T=I-Y_{\rm raw}^*Y_{\rm raw}$.  Then
$T\succeq c^*c$, $\operatorname{tr}T=O(s^2)$, and the polar
factorization is
\begin{equation}
 Y_{\rm raw}=U(I-T)^{1/2},\quad U^*U=I_m,\quad
 T=s^2(w^*w+F),\qquad F\succeq0.
 \label{I:R:eq:normal-polar}
\end{equation}
The variables $\lambda,\zeta,w,F$ are uniformly bounded at maximizing
points.  Write $U=\begin{pmatrix}D\\l\end{pmatrix}$.  The slow-coordinate
unitary $\operatorname{diag}(e^{i\psi},I_m)$ preserves feasibility and
rotates both $J$ and $lD^mx$ by $e^{i\psi}$.  It therefore permits the
normalization $a=lD^mx>0$.  More explicitly, it changes
$x$ to $e^{i\psi}x$, $c$ to $e^{-i\psi}c$, $z_o$ to
$e^{i\psi}z_o$, and leaves $z_0,D,l$ unchanged.  Thus this normalization
discards no complex direction.  The base triples $(D,l,x)$ approach
$\mathcal M_m$, and always satisfy $a\le\Gamma$.

\emph{The normal quadratic form.}
Expand $(I-T)^{1/2}$, $W_s(z)$, and the finite product $\ell B^{m+1}e_1/s$.
Uniformly over the localized parameters,
\begin{equation}
 |J(s)|=a+s^2Q+O(s^3).
 \label{I:R:eq:J-normal-expansion}
\end{equation}
At a boundary triple in $\mathcal M_m$, the quadratic coefficient is
\begin{align}
 Q={}&-\frac12\operatorname{Re}S
   +\operatorname{Re}(\lambda a_{m-1})-\frac\Gamma2|\lambda|^2
   +\operatorname{Re}(wv)-\frac12wKw^*\nonumber\\
   &\hspace{18mm}-\frac\Gamma2|\zeta|^2
                       -\frac12\operatorname{tr}(KF).
 \label{I:R:eq:normal-quadratic-Q}
\end{align}
Off the boundary set the same formula uses the continuously varying
phase-normalized coefficients and replaces $\Gamma$ in the two
normalization costs by $a$.

For completeness, each contribution in
\eqref{I:R:eq:normal-quadratic-Q} can be read directly from the block
product.  Normalizing $z_F$ contributes
$-\Gamma(|\lambda|^2+|\zeta|^2)/2$.  The slow diagonal entry
$s^2\lambda$ contributes $\lambda a_{m-1}$.  A single slow-to-fast return
through the row $sw$ contributes
\[
 \sum_{j=0}^{m-2}(lD^jx)(wD^{m-2-j}x)=wv.
\]
Right contraction by $(I-s^2(w^*w+F))^{1/2}$ contributes
\[
 -\frac12\operatorname{Re}\sum_{j=0}^{m}
 lD^j(w^*w+F)D^{m-j}x
 =-\frac12wKw^*-\frac12\operatorname{tr}(KF).
\]
The first-column Schur factor contributes the perturbation
$D\mapsto D-s^2xx^*D/2$, hence $-S/2$.  Finally the slow readout is
$s^2\zeta$; using it requires a return to the slow state and first affects
$J$ at order $s^3$.  These observations prove the displayed expansion,
including its empty-sum case $m=1$.

Lemma~\ref{I:R:lem:response-K-positive} allows the quadratic form to be
maximized explicitly:
\begin{equation}
 \begin{gathered}
 F=0,\quad\zeta=0,\quad
 \lambda=\frac{\overline{a_{m-1}}}{\Gamma},\quad w=v^*K^{-1},\\
 Q_{\max}=-\frac12\operatorname{Re}S
       +\frac{|a_{m-1}|^2}{2\Gamma}+\frac12v^*K^{-1}v.
 \end{gathered}
 \label{I:R:eq:normal-quadratic-max}
\end{equation}
For the global upper bound, every sequence of maximizers has a
subsequence whose base triples approach $\mathcal M_m$.  Their base
contribution $a-\Gamma$ is nonpositive, all quadratic coefficients
converge continuously, and $K$ remains uniformly positive.  Therefore
\[
 \limsup_{s\to0}\frac{M(s)-\Gamma}{s^2}
 \le\max_{\mathcal M_m}Q_{\max}.
\]
For the reverse bound choose a boundary triple maximizing
$Q_{\max}$, take the parameters in
\eqref{I:R:eq:normal-quadratic-max}, and use the exact square roots in
\eqref{I:R:eq:first-column-chart} and \eqref{I:R:eq:normal-polar}.
They form a feasible family for all sufficiently small $|s|$ and give
the reverse liminf.  Consequently
\[
 M(s)=\Gamma+s^2\max_{\mathcal M_m}Q_{\max}+o(s^2).
\]
Equation~\eqref{I:R:eq:M-and-J} proves the first expansion in
\eqref{I:R:eq:all-budget-envelope-expansion}, with the coefficient
\eqref{I:R:eq:bk-response-formula}.

\emph{The final scalar maximization.}
The elementary estimate $h_k(t)\le t$ and
\[
 1-\sqrt{(1-q)(1-t)}\ge\frac{q+t}{4}
\]
show that the objective in \eqref{I:R:eq:allk-elimination} is at most
$16qt/(q+t)^2$.  The trial point $t=q$ has value tending to
$\Gamma^2>0$, so all maximizers have $x=t/q$ in a fixed compact
subinterval of $(0,\infty)$ when $q$ is small.  Uniformly on such
intervals,
\begin{align*}
 \frac{q h_k(qx)}{(1-\sqrt{(1-q)(1-qx)})^2}
 ={}&\frac{4\Gamma^2x}{(1+x)^2}\\
 &+q\left(\frac{4b_kx^2}{(1+x)^2}
        -\frac{2\Gamma^2x(x-1)^2}{(1+x)^3}\right)+o(q).
\end{align*}
The leading term has its unique maximum $\Gamma^2$ at $x=1$.
Thus all maximizing $x$ tend to one.  The leading term is never above
$\Gamma^2$ and the bracket tends to $b_k$ along every maximizing
sequence.  Taking $x=1$ gives the matching lower bound, proving the
second expansion.  The established projection inequality
$D_k(q)\ge1-\mathcal P_k(q)$ proves
\eqref{I:R:eq:all-budget-D-lower}.
\end{proof}


\subsection{Independent state budget and terminal delay}

The state budget and the number of terminal letters need not coincide.
This gives a larger class of targets without another nondegeneracy
assumption.

\begin{theorem}[Independent budget and delay]
\label{I:R:thm:independent-budget-delay-response}
Let $n\ge2$ and $d\ge n$.  Put $j=n-1$ and $p=d-1$, and define
\begin{align*}
 \mathcal P_{n,d}(q)&=
 \max_{A^*A+B^*B+\ell^*\ell\preceq I_n}
 q\Rnorm{\ell B^d(I-rA)^{-1}}^2,\\
 h_{n,d}(t)&=\max\{|\ell B^de_1|^2:
                  V=\begin{pmatrix}B\\\ell\end{pmatrix}
                  \text{ is a contraction},\ \Rnorm{Ve_1}^2=t\},\\
 \Gamma_{j,p}&=\max_{D^*D+l^*l\preceq I_j}\Rnorm{lD^p}.
\end{align*}
Let
\[
 \mathcal M_{j,p}=\{(D,l,x):D^*D+l^*l=I_j,\ \Rnorm x=1,
                                     \ lD^px=\Gamma_{j,p}>0\}.
\]
For a triple in this set put $\Gamma=\Gamma_{j,p}$ and
\begin{align*}
 a_i&=lD^ix\quad(0\le i\le p),\\
 S_p&=\sum_{i=0}^{p-1}a_i x^*D^{p-i}x,
 &v_p&=\sum_{i=0}^{p-2}a_iD^{p-2-i}x,\\
 K_p&=\operatorname{Herm}\sum_{i=0}^{p}D^{p-i}xlD^i,
 &b_{n,d}&=\max_{\mathcal M_{j,p}}
    \bigl\{|a_{p-1}|^2-\Gamma\operatorname{Re}S_p
                           +\Gamma v_p^*K_p^{-1}v_p\bigr\}.
\end{align*}
Then $\mathcal M_{j,p}$ is compact and nonempty,
\begin{equation}
 K_p\succeq\Gamma\sum_{i=0}^{j-1}D^ixx^*D^{*i}\succ0,
 \label{I:R:eq:independent-delay-K-positive}
\end{equation}
and
\begin{equation}
 \boxed{\displaystyle
 h_{n,d}(t)=\Gamma_{j,p}^2t+b_{n,d}t^2+o(t^2),\qquad
 \mathcal P_{n,d}(q)=\Gamma_{j,p}^2+b_{n,d}q+o(q).}
 \label{I:R:eq:independent-delay-expansion}
\end{equation}
For the target supported on $a^*b^d$, its squared rank-$n$ error therefore
satisfies
\[
 D_{n,d}(q)\ge1-\Gamma_{j,p}^2-b_{n,d}q+o(q).
\]
\end{theorem}
\begin{proof}
Proposition~\ref{I:R:prop:allk-eliminate-A} concerns neither the exponent
of $B$ nor its state dimension, so it gives exactly
\[
 \mathcal P_{n,d}(q)=\max_t
       \frac{q h_{n,d}(t)}{(1-r\sqrt{1-t})^2}.
\]
Since $p\ge j$, Lemma~\ref{I:R:lem:strict-dimension-improvement} applies to
the coefficient $z^p$ at every dimension up to $j$.  Its observation-map
proof implies that a maximizing unit $x$ for $lD^px$ is cyclic in the
$j$-dimensional state space.  The energy identity then forces
$D^*D+l^*l=I_j$ and strict stability exactly as in
Lemma~\ref{I:R:lem:cyclic-lossless-maximizer}.  Quantitative defect control
uses the invertible matrix $[x,Dx,\ldots,D^{j-1}x]$, so
\eqref{I:R:eq:defect-gap-control} holds with $lD^p$ in place of $lD^m$.

The exact first-column chart has fast state dimension $j$, and the
normalized product is now $J=\ell B^{p+1}e_1/s$.  Its leading term is
$lD^pz_F$; its first derivative is
\[
 z_ocD^{p-1}z_F+z_0lD^{p-1}z_F
   +\sum_{i=0}^{p-2}(lD^iz_F)(cD^{p-2-i}z_F).
\]
The same quantitative defect estimate controls this derivative by the
square root of the leading gap.  Consequently every maximizing normal
coordinate is $O(s)$ and the remaining defect is $O(s^2)$.
The polar coordinates \eqref{I:R:eq:normal-polar} therefore apply with
$j$ fast states.  Expanding the product gives the quadratic form
\eqref{I:R:eq:normal-quadratic-Q} with $a_{m-1},S,v,K$ replaced by
$a_{p-1},S_p,v_p,K_p$.

To justify the inverse, perturb the lossless fast isometry by
$(I-\epsilon Y)^{1/2}$.  Differentiating its coefficient $lD^px$ gives
$K_p$, while telescoping just the first $j$ energy terms gives the right
side of \eqref{I:R:eq:independent-delay-K-positive}.  Cyclicity makes that
matrix positive definite.  Thus the normal quadratic form can again be
maximized by a positive matrix inverse.  Compactness over
$\mathcal M_{j,p}$ and the nonpositive tangential base gap give the
matching upper and lower limits, proving the expansion of $h_{n,d}$.
The scalar maximization in $t/q$ is identical to that in
Theorem~\ref{I:R:thm:all-budget-response}, with leading maximum at $t/q=1$.
This proves the expansion of $\mathcal P_{n,d}$.  The target's projection
inequality gives the final lower bound.
\end{proof}


\subsection{Low-order response checks}

\begin{proposition}[The two-state response]
\label{I:R:prop:two-state-response-check}
For $k=2$, formula \eqref{I:R:eq:bk-response-formula} gives $b_2=1/4$.
\end{proposition}
\begin{proof}
Here $m=1$, and all boundary maximizers are equivalent by scalar phases
to $D=l=1/\sqrt2$, $x=1$, $\Gamma=1/2$.  Then
$a_0=1/\sqrt2$, $S=1/2$, $v=0$, and $K=1$.
Thus $b_2=1/2-1/4=1/4$.
\end{proof}


\begin{proposition}[Response of the degree-two scalar extremizer]
\label{I:R:prop:degree-two-response-check}
The scalar extremizer with denominator $1+z+z^2/2$ has
$\Gamma_2^2=7/16$ and normal-response coefficient $1/7$.
If all scalar degree-two boundary maximizers are equivalent to this
extremizer under the scalar symmetries, then $b_3=1/7$.
\end{proposition}
\begin{proof}
Use the observable companion coordinates
\[
 T=\begin{pmatrix}0&1\\-1/2&-1\end{pmatrix},\quad
 l_0=(1,0),\quad
 Q=\begin{pmatrix}8/5&4/5\\4/5&12/5\end{pmatrix},
 \qquad Q-T^*QT=l_0^*l_0.
\]
The canonical Euclidean realization is
$D=Q^{1/2}TQ^{-1/2}$, $l=l_0Q^{-1/2}$.  Equivalently, compute in the
$Q$ inner product, in which
\[
 \Gamma=\frac{\sqrt7}{4},\qquad
 x_{\rm raw}=\frac{(-1,-3)^{\mathsf T}}{2\sqrt7},\quad
 a_0=-\frac1{2\sqrt7},\quad a_1=-\frac3{2\sqrt7},\quad a_2=\Gamma.
\]
Direct rational algebra gives
\[
 S=\frac{23\sqrt7}{196},\qquad
 v_{\rm raw}=\frac{(1,3)^{\mathsf T}}{28},\qquad
 K_{\rm raw}=\begin{pmatrix}
 2\sqrt7/7&-\sqrt7/7\\-\sqrt7/28&13\sqrt7/28
 \end{pmatrix}.
\]
The last matrix is $Q$-self-adjoint and represents the same positive
response operator.  In particular,
\[
 v_{\rm raw}^*QK_{\rm raw}^{-1}v_{\rm raw}=\frac{3\sqrt7}{196}.
\]
The three terms in \eqref{I:R:eq:bk-response-formula} are therefore
\[
 \frac9{28}-\frac{23}{112}+\frac3{112}=\frac17.
\]
The value for this extremizer follows without any uniqueness assertion.
The last statement follows by taking the maximum in
\eqref{I:R:eq:bk-response-formula}.
\end{proof}


\subsection{The exact rank-one benchmark and the role of matching}

The case $k=1$ is singular because its limiting passive-overlap constant is zero (the limiting squared approximation error is one).
For a single-column contraction $V=(B,\ell)^{\mathsf T}$,
\[
 h_1(t)=\max_{|B|^2+|\ell|^2=t}|\ell B|^2=\frac{t^2}{4}.
\]
Proposition~\ref{I:R:prop:allk-eliminate-A} then gives the exact envelope
\begin{equation}
 \mathcal P_1(q)=\frac{q}{(1+\sqrt q)^2},\qquad
 t_{\max}=\frac{2\sqrt q}{1+\sqrt q}.
 \label{I:R:eq:rankone-envelope-benchmark}
\end{equation}
Thus
\[
 1-\mathcal P_1(q)
 =1-q+2q^{3/2}-3q^2+4q^{5/2}+O(q^3).
\]
The exact rank-one Hankel optimum instead has the expansion
\[
 D_1(q)=1-q+2q^{3/2}-3q^2+5q^{5/2}+O(q^3).
\]
The envelope is sharp through $q^2$ in this benchmark and has a strict
gap at order $q^{5/2}$.  Accordingly, the all-budget response theorem
supplies the global lower mechanism. The matching, $q$-dependent
sequential realization needed to identify the first correction of
$D_k$ is supplied in Section~\ref{I:R:sec:v4-matching}.

All arguments above also hold over the real field, restricting
$\mathcal M_m$ and the normal variables accordingly.  In
Lemma~\ref{I:R:lem:strict-dimension-improvement}, a real point $a\in(-1,1)$
can be chosen because a nonzero real analytic residual cannot vanish
on that entire interval.  The scalar Hardy extremal can be chosen real,
so the leading $\Gamma_m$ agrees in the two fields.
Section~\ref{I:R:subsec:real-diagonal-orbits} proves equality of the
diagonal response maxima by classifying the phases of all complex
scalar extremizers. For independent budget and delay, the two
field-specific response maxima remain distinct definitions.

% End all_budget_response.tex


% Begin matching.tex

\section{A structured realization matching the envelope to first order}
\label{I:R:sec:v4-matching}

The projection estimate becomes useful for optimal approximation only
when a genuine Hankel realization attains it to the required order.
The following construction supplies this step in every fixed dimension.
It applies over the complex field; real input data give a real
competitor.  It does not assert exact finite-parameter completion.

\begin{theorem}[Structured first-order matching]\label{I:R:thm:v4-matching}
Fix integers $n\ge2$ and $d\ge n$, and put $m=n-1$.  Let
$D\in\mathbb C^{m\times m}$ and $l\in\mathbb C^{1\times m}$ satisfy
\[
 D^*D+l^*l=I_m,\qquad \rho(D)<1,\qquad
 \Gamma=\|lD^{d-1}\|\in(0,1).
\]
Write $x=(lD^{d-1})^*/\Gamma$.  Fix any
$\lambda\in\mathbb C$ and $w\in\mathbb C^{1\times m}$.
For sufficiently small positive $s$, set $q=s^2$, $r=\sqrt{1-q}$ and
\begin{align}
 F_s&=(I_m-s^2w^*w)^{1/2},&
 N_s&=\begin{pmatrix}sw\\DF_s\\lF_s\end{pmatrix},&
 z_s&=\begin{pmatrix}s\lambda\\
       \sqrt{1-s^2|\lambda|^2}\,x\\0\end{pmatrix},\label{I:R:eq:v4-match-trial}\\
 W_s&=I_{n+1}-(1-r)z_sz_s^*,&
 V_s&=[\,sz_s,\ W_sN_s\,]
       =\begin{pmatrix}B_s\\l_s\end{pmatrix},&
 v_s&=\begin{pmatrix}r\\-sN_s^*z_s\end{pmatrix},\qquad
 A_s=e_0v_s^*.\label{I:R:eq:v4-match-tuple}
\end{align}
Define the row and the boundary vector
\begin{equation}
 \omega_s=s\,l_sB_s^d(I-rA_s)^{-1},\qquad
 \beta_s=s^{-1}(I-rA_s)\omega_s^*.
 \label{I:R:eq:v4-match-beta}
\end{equation}
For a word $u$, let $T_u$ denote the standard product of the letters
$T_a=A_s,T_b=B_s$, and set
\begin{equation}
 g_s(u)=l_sT_{\operatorname{rev}(u)}\beta_s.
 \label{I:R:eq:v4-match-series}
\end{equation}
This is an $n$-state word series: in the standard, unreversed convention
its readout is $\beta_s^{\mathsf T}$, its letters are
$A_s^{\mathsf T},B_s^{\mathsf T}$, and its boundary is
$l_s^{\mathsf T}$.  For all sufficiently small $s>0$ it is jointly
mean-square stable, and
\begin{equation}
 \boxed{\displaystyle
 \|H_{f_{r,d}}-H_{g_s}\|^2
     =1-\|\omega_s\|^2+O(q^2).}
 \label{I:R:eq:v4-match-conclusion}
\end{equation}
In particular, if a trial direction has
$\|\omega_s\|^2=\Gamma^2+cq+O(q^2)$, it gives a genuine structured
upper bound $1-\Gamma^2-cq+O(q^2)$.  Constants and neighborhoods may
depend on the fixed data.  The same conclusion holds for any other
analytic branch having the exact energy identity and the block orders
used in the proof below.
\end{theorem}

\begin{proof}
We keep ordinary coefficient transpose ${}^{\mathsf T}$ distinct from
Hermitian adjoint ${}^*$.  This distinction matters for complex data.

\emph{Passivity and the cancellation.}
Direct calculation gives $N_s^*N_s=I_m$, $\|z_s\|=1$ and
\[
 I_n-V_s^*V_s=v_sv_s^*,\qquad
 A_s^*A_s+B_s^*B_s+l_s^*l_s=I_n.
\]
The zero in the last component of $z_s$ is exact.  On
$\mathbb Ce_0\oplus\mathbb C^m$, the construction therefore has
\begin{align}
 B_s&=\begin{pmatrix}O(s^2)&O(s)\\
                sx+O(s^3)&D+O(s^2)\end{pmatrix},&
 l_s&=(0,\ l+O(s^2)),\label{I:R:eq:v4-match-blocks}\\
 A_s&=\begin{pmatrix}r&-sx^*D+O(s^3)\\0&0\end{pmatrix},&
 (I-rA_s)^{-1}
 &=\begin{pmatrix}q^{-1}&-rsq^{-1}x^*D+O(s)\\0&I_m\end{pmatrix}.
 \label{I:R:eq:v4-match-resolvent}
\end{align}
For fixed $d$, multiplication of the blocks gives
\[
 l_sB_s^de_0=s\,lD^{d-1}x+O(s^3)=s\Gamma+O(s^3),\qquad
 l_sB_s^d|_{\mathbb C^m}=lD^d+O(s^2).
\]
The identity $\Gamma x^*=lD^{d-1}$ cancels the order-$s$ fast
component of the overlap:
\begin{equation}
 \omega_s=(\Gamma+O(s^2),\
 s\{lD^d-\Gamma x^*D\}+O(s^3))
          =(\Gamma+O(q),O(s^3)).
 \label{I:R:eq:v4-match-cancellation}
\end{equation}
Consequently $\beta_s=(s\Gamma+O(s^3),O(s^2))$.

\emph{Joint stability by finite observability.}
Let $\Phi_s(X)=A_sXA_s^*+B_sXB_s^*$. The exact energy identity gives
$\Phi_s^*(I)=I-l_s^*l_s$. Since $\rho(D)<1$ and
$D^*D+l^*l=I_m$, the pair $(l,D)$ is observable: if all $lD^jx$
vanish, telescoping gives $\|D^t x\|=\|x\|$ for every $t$, hence $x=0$.
Thus the $m$ rows $lD^j$, $0\le j<m$, have full column rank.
Furthermore,
\[
 l_sB_s^dA_s=(l_sB_s^de_0)v_s^*
           =(s\Gamma+O(s^3))v_s^*.
\]
For small $s>0$ its scalar coefficient is nonzero. The $m+1$ rows
$v_s^*,l_s,l_sB_s,\ldots,l_sB_s^{m-1}$ converge to
$e_0^*,(0,l),(0,lD),\ldots,(0,lD^{m-1})$, a full-rank system.
Consequently a finite observation Gramian
\[
 Q_{s,L}=\sum_{|u|\le L}T_u^*l_s^*l_sT_u\succ0,
 \qquad L=d+1,
\]
is positive definite. Telescoping the energy identity yields
$(\Phi_s^*)^{L+1}(I)=I-Q_{s,L}\preceq(1-\eta_s)I$ with
$\eta_s=\lambda_{\min}(Q_{s,L})>0$. Positivity now implies
$\rho(\Phi_s)^{L+1}\le1-\eta_s<1$.
The transposed realization has the same stability spectral radius.

\emph{A reachability Gram estimate.}
Let
\[
 P_s=\sum_u T_{\operatorname{rev}(u)}\beta_s\beta_s^*
                    T_{\operatorname{rev}(u)}^*
      =\begin{pmatrix}p&z^*\\z&X\end{pmatrix}.
\]
It exists by joint stability and solves
$P_s=\beta_s\beta_s^*+A_sP_sA_s^*+B_sP_sB_s^*$.
The blocks above yield, with uniform implicit constants,
\begin{align*}
 qp&=q\Gamma^2+O(q^2)+O(s\|z\|)+O(q\|X\|)+O(q^2p),\\
 \|z\|&=O(s^3(1+p))+O(q\|z\|)+O(s\|X\|),\\
 \|X\|&=O(qp)+O(s\|z\|)+O(q^2).
\end{align*}
The last line follows by inverting the uniformly stable fast Stein
map.  Substitution in the middle line and absorption give
$\|z\|=O(s^3(1+p))$ and $\|X\|=O(qp+q^2)$.
The scalar equation becomes
$p=\Gamma^2+O(q(1+p))$.  Absorbing the last small multiple of $p$
proves
\begin{equation}
 (P_s)_{00}=\Gamma^2+O(q),\qquad
 (P_s)_{0F}=O(s^3),\qquad (P_s)_{FF}=O(q),\qquad
 \|P_s-\omega_s^*\omega_s\|=O(q).
 \label{I:R:eq:v4-match-gram}
\end{equation}

\emph{Exact interpolation and a fixed spectral gap.}
Define the reversed observation isometry and reachability operator by
\[
 (\mathcal O_s c)[u]=l_sT_{\operatorname{rev}(u)}c,\qquad
 \mathcal R_s e_u=T_{\operatorname{rev}(u)}\beta_s.
\]
The energy identity and joint stability imply
$\mathcal O_s^*\mathcal O_s=I_n$.
Exact coefficient multiplication gives
$H_{g_s}^{\mathsf T}=\mathcal O_s\mathcal R_s$ and
$\mathcal R_s\mathcal R_s^*=P_s$.
Let $S_s=\operatorname{ran}\mathcal O_s$, and introduce the real
orthonormal target vectors
\[
 u_r=s\sum_{j\ge0}r^j e_{a^j},\qquad
 u_h=s\sum_{j\ge0}r^j e_{a^jb^h}\quad(1\le h\le d),
 \qquad v_r=u_d.
\]
The target decomposes exactly as
\[
 H_{f_{r,d}}=u_rv_r^*+
       s\sum_{j=0}^{d-1}u_{d-j}e_{b^j}^*.
\]
By \eqref{I:R:eq:v4-match-beta},
\[
 \mathcal R_su_r=s(I-rA_s)^{-1}\beta_s=\omega_s^*,\qquad
 \mathcal O_s^*v_r=\omega_s^*,\qquad
 H_{g_s}^{\mathsf T}u_r=P_{S_s}v_r.
\]
Set $T=H_{f_{r,d}}^{\mathsf T}-H_{g_s}^{\mathsf T}$.
Then $Tu_r=(I-P_{S_s})v_r$ and
\begin{align*}
 T^*Tu_r
 &=H_{f_{r,d}}(I-P_{S_s})v_r\\
 &=a_su_r-s\sum_{j=0}^{d-1}u_{d-j}
                  [l_sB_s^j\omega_s^*],\qquad
 a_s=1-\|\omega_s\|^2.
\end{align*}
The first equality uses that the residual is orthogonal to the range
of $H_{g_s}^{\mathsf T}$.  The target is real, so its Hermitian adjoint
in this expression is its coefficient transpose.  For $j\ge1$,
$l_sB_s^je_0=O(s)$, while $l_se_0=0$.
Together with \eqref{I:R:eq:v4-match-cancellation}, this shows that the
component of $T^*Tu_r$ perpendicular to $u_r$ has norm $O(q)$.

Finally the isometry and \eqref{I:R:eq:v4-match-gram} give exactly
\[
 \|H_{g_s}^{\mathsf T}P_{u_r^\perp}\|^2
 =\|\mathcal R_sP_{u_r^\perp}\mathcal R_s^*\|
 =\|P_s-\omega_s^*\omega_s\|=O(q).
\]
Also $\|H_{f_{r,d}}^{\mathsf T}P_{u_r^\perp}\|=s$.
Hence the compression of $T^*T$ to $u_r^\perp$ has norm $O(q)$.
In this orthogonal decomposition,
\[
 T^*T=\begin{pmatrix}a_s&O(q)\\O(q)&O(q)\end{pmatrix},\qquad
 a_s\longrightarrow1-\Gamma^2>0.
\]
The scalar Schur-complement estimate across this fixed gap gives
$\|T\|^2=a_s+O(q^2)$.  Ordinary coefficient transpose preserves
operator norm, proving \eqref{I:R:eq:v4-match-conclusion}.
\end{proof}


\begin{corollary}[Matching the three-state correction]
\label{I:R:cor:v4-three-state-upper}
For the three-state, three-delay problem there is an explicit real
stable competitor with squared error
\[
 \frac9{16}-\frac17q+O(q^2).
\]
Combined with the independently established complex global envelope
$\mathcal P_{3,3}(q)=7/16+q/7+O(q^2)$, this gives, for each field
separately,
\[
 D_3^{\mathbb F}(q)=\frac9{16}-\frac17q+O(q^2),
 \qquad\mathbb F\in\{\mathbb R,\mathbb C\}.
\]
\end{corollary}

\begin{proof}
Take
\[
 D=\begin{pmatrix}-4/7&\sqrt5/7\\-5\sqrt5/14&-3/7\end{pmatrix},
 \qquad l=(-\sqrt7/14,\sqrt{35}/7).
\]
One checks $D^*D+l^*l=I_2$, $\sigma(D)=\{(-1+i)/2,(-1-i)/2\}$,
and $lD^2=(\sqrt7/4,0)$, so $\Gamma^2=7/16$ and $x=e_1$.
The explicit trial from the envelope computation takes
$N_s=\exp(sK)N_0$, $N_0=\begin{pmatrix}0_{1\times2}\\D\\l\end{pmatrix}$,
with the only nonzero upper-triangular entries
$K_{01}=2/7$, $K_{02}=\sqrt5/35$ and $K^{\mathsf T}=-K$, and
\[
 z_s=\frac{(-6s/7,1,0,0)^{\mathsf T}}
                   {\sqrt{1+36s^2/49}}.
\]
It has all the exact identities and parity orders in the theorem's
proof.  Its direct envelope expansion is
$\|\omega_s\|^2=7/16+q/7+O(q^2)$.
Apply \eqref{I:R:eq:v4-match-conclusion}, then sandwich the real upper
bound with the lower bound over all complex competitors.
This does not assert real--complex equality at finite $q$.
\end{proof}


% End matching.tex


\begin{proof}[Proof of Theorem~\ref{I:R:thm:main-general}]
Use the independent-budget/delay envelope theorem with fast dimension
$j=n-1$ and coefficient index $p=d-1$. It gives the global lower bound
$(E_{j,p}^{\mathbb F})^2-b_{n,d}^{\mathbb F}q+o(q)$.
Choose a boundary triple attaining the response maximum in
\eqref{I:R:eq:main-response-coefficient}. Cauchy--Schwarz equality implies
$x=(lD^p)^*/\gamma$. In Theorem~\ref{I:R:thm:v4-matching}, take
\[
 \lambda=\overline{a_{p-1}}/\gamma,\qquad w=v^*K^{-1}.
\]
These are precisely the maximizing normal directions of the envelope.
The evenness needed for the remainder follows from the exact trial.
Put $E=\operatorname{diag}(-1,I_m)$ and
$\widetilde E=\operatorname{diag}(-1,I_m,1)$. The definitions give
$N_{-s}=\widetilde E N_s$, $z_{-s}=\widetilde E z_s$,
$A_{-s}=EA_sE$, $B_{-s}=EB_sE$, and $l_{-s}=l_sE$.
It follows that $\omega_{-s}=-\omega_sE$: its slow coordinate
$J_s=\omega_se_0$ is even, and its fast coordinates are odd.
The apparent division by $s$ in $J_s$ is removable by
\eqref{I:R:eq:v4-match-cancellation}. Hence $J_s$ and
$\|\omega_s\|^2$ are real analytic and even. Their third-order
remainder is therefore $O(s^4)=O(q^2)$.
The fixed analytic trial has
$\Rnorm{\omega_s}^2=\gamma^2+b_{n,d}^{\mathbb F}q+O(q^2)$:
its slow component is the divided coefficient $J_s$, and the fast
component is $O(q^{3/2})$. The matching theorem therefore gives the
upper bound $(E_{j,p}^{\mathbb F})^2-b_{n,d}^{\mathbb F}q+O(q^2)$.
The two estimates prove the stated expansion in each field separately.
\end{proof}



% Begin three_state.tex

\section{The three-state envelope and its sharp first correction}
\label{I:R:sec:three-state}

The next state budget admits an unconditional complex envelope expansion.
The global Hessian calculation below improves the general first-response
remainder from $o(q)$ to $O(q^2)$.  This section establishes the envelope
and the resulting Hankel lower bound. Combined with
Theorem~\ref{I:R:thm:v4-matching}, it proves the three-state equality in
Corollary~\ref{I:R:cor:main-explicit}.

\begin{theorem}[Three-state passive envelope]\label{I:R:thm:three-envelope}
Let
\[
 h_3(t)=\max\{|\ell B^3e_1|^2:
 V=[B;\ell]\in\RC^{4\times3},\ V^*V\preceq I_3,
 \ \|Ve_1\|^2=t\}.
\]
As $t\downarrow0$ and $q\downarrow0$, respectively,
\begin{align}
 h_3(t)&=\frac7{16}t+\frac17t^2+O(t^3),
 \label{I:R:eq:three-h}\\
 \mathcal P_3(q)&:=\sup_{A^*A+B^*B+\ell^*\ell\preceq I_3}
 q\|\ell B^3(I-\sqrt{1-q}A)^{-1}\|^2
 =\frac7{16}+\frac17q+O(q^2).
 \label{I:R:eq:three-P}
\end{align}
Both expansions hold over the full complex comparison class.  Consequently,
for the binary target $f_{r,3}(a^jb^3)=(1-r^2)r^j$ and $q=1-r^2$,
\[
 D_3^{\RC}(q)\ge\frac9{16}-\frac17q-Cq^2
\]
for some $C<\infty$ and all sufficiently small positive $q$.
The same lower bound holds for real competitors.
\end{theorem}

\subsection{The entire boundary maximizing set}

Put $t=s^2$.  The exact contraction chart is
\[
 V=[sz,W_s(z)N],\quad \|z\|=1,\quad
 W_s(z)=I-(1-\sqrt{1-s^2})zz^*,\quad N^*N\preceq I_2.
\]
For fixed $s,z$, the amplitude $J_s=\ell B^2z_B$ is a holomorphic
polynomial in $N$, and $\ell B^3e_1=sJ_s$.  Applying the maximum-modulus
principle to the two singular-value variables of $N$ shows that its maximum
can be attained on $N^*N=I_2$.  Thus the exact maximization takes place on
\[
 X=S^7\times\operatorname{St}_{\RC}(4,2),\qquad
 h_3(s^2)=s^2\max_X|J_s|^2.
\]
This compact real-analytic manifold has dimension $7+12=19$.
The reduction asserts equality of maximum values; it does not assert that
every maximizer in the original contraction ball is isometric.

At $s=0$, write $z=(z_0,z_F,z_o)$ and $N=[c;D;l]$, with
$z_F\in\RC^2$. Then
\[
 J_0=lD^2z_F,\qquad
 \max|J_0|^2=\Gamma_2^2=1-E_2^2=\frac7{16}.
\]
Here $E_2=3/4$ is the established scalar rational-approximation constant.
The cyclic-lossless boundary lemma gives, at every maximum,
\[
 z_0=z_o=0,\quad \|z_F\|=1,\quad c=0,\quad D^*D+l^*l=I_2.
\]
For clarity, its rigidity has a short direct proof in this dimension.
A nonobservable pair reduces to one state, where
$\max |ld^2|^2=\max_{0\le x\le1}(1-x)x^2=4/27<7/16$.
For an observable pair let $G\succ0$ be its observation Gramian, so
$G\preceq I$.  Whitening cannot decrease $\|lD^2\|$; equality at the
maximum forces $G w^*=w^*$ for $w=lD^2$.
With $Q=I-G$, passivity gives $D^*QD\preceq Q$, making
$\ker Q$ invariant under $D$.
The maximizing vector $w^*$ is cyclic for the whitened pair, and the
whitened and original dynamics agree on $\ker Q$.
It follows that $Q=0$.  Losslessness then forces $c=0$.

There is also no unclassified scalar extremizer at this budget.
Write a minimizing degree-two model-space denominator as $P$ and put
$b=P_{\le2}(P^*/P)$.  Equality in the Hardy coefficient and H\"older
bounds makes $bP^2$ a coefficient extremal in $H^{2/3}$.
The Brevig--Saksman equality case identifies it, up to phases, with the
cube of
\[
 P_\theta(z)=\frac23\left(1+e^{i\theta}z+
                 \frac12 e^{2i\theta}z^2\right).
\]
Thus $bP^2=\lambda P_\theta^3$.
The right side has degree six, forcing $\deg b=\deg P=2$.
Its two roots are distinct; since $P^2$ divides $P_\theta^3$, necessarily
$P$ is proportional to $P_\theta$.
Lower-dimensional model spaces have strictly smaller monomial projection,
so no lower-degree case is omitted.
The canonical pair is consequently unique up to unitary state coordinates,
circle rotation, and readout phase.

A convenient representative is
\begin{equation}
 z_*=\begin{pmatrix}0\\1\\0\\0\end{pmatrix},\qquad
 N_*=
 \begin{pmatrix}
 0&0\\
 -4/7&\sqrt5/7\\
 -5\sqrt5/14&-3/7\\
 -\sqrt7/14&\sqrt{35}/7
 \end{pmatrix}.
 \label{I:R:eq:three-base}
\end{equation}
Its middle block $D$ and last row $l$ satisfy
$N_*^*N_*=I_2$ and $lD^2=(\sqrt7/4,0)$.
The nonzero lower-left entry of $D$ verifies cyclicity directly.
All boundary maxima form the orbit of \eqref{I:R:eq:three-base} under
\[
 B\mapsto e^{i\theta}U^*BU,\qquad
 \ell\mapsto e^{i\phi}\ell U,\qquad
 U=\operatorname{diag}(1,U_2),\quad U_2\in U(2).
\]
This orbit has dimension six.  Its continuous stabilizer is trivial:
fixing $z_*$ makes $U_2$ diagonal; the nonzero diagonal and off-diagonal
entries of $D$ force the letter and basis phases to vanish; the nonzero
entries of $l$ then force the readout phase to vanish.

\subsection{Exact response in rational coordinates}

The all-budget normal-response formula provides a compact exact computation
of the coefficient.  In a lossless canonical system, with
$x=(lD^2)^*/\Gamma$, put
\[
 a_j=lD^jx,\quad
 S=a_0x^*D^2x+a_1x^*Dx,\quad v=a_0x,
\quad
 K=\operatorname{Herm}\!\left(
 D^2xl+DxlD+xlD^2\right).
\]
The coefficient is
\begin{equation}
 b_3=|a_1|^2-\Gamma\operatorname{Re}S+
                  \Gamma v^*K^{-1}v.
 \label{I:R:eq:three-response}
\end{equation}
It suffices to use rational observable coordinates with their Gram metric:
\[
 T=\begin{pmatrix}0&1\\-1/2&-1\end{pmatrix},\quad
 l_0=(1,0),\quad
 Q=\begin{pmatrix}8/5&4/5\\4/5&12/5\end{pmatrix},\quad
 Q-T^*QT=l_0^*l_0.
\]
The corresponding canonical quantities are represented in the $Q$ inner
product by
\[
 \Gamma=\frac{\sqrt7}{4},\quad
 x=\frac{(-1,-3)^T}{2\sqrt7},\quad
 a_0=-\frac1{2\sqrt7},\quad a_1=-\frac3{2\sqrt7},
\]
\[
 S=\frac{23\sqrt7}{196},\qquad
 v=\frac{(1,3)^T}{28},\qquad
 K_Q=\begin{pmatrix}
 2\sqrt7/7&-\sqrt7/7\\
 -\sqrt7/28&13\sqrt7/28
 \end{pmatrix}.
\]
Here $K_Q$ represents the operator $K$ and is self-adjoint for the metric
$Q$. Direct exact algebra gives
\[
 v^*QK_Q^{-1}v=\frac{3\sqrt7}{196},\qquad
 b_3=\frac9{28}-\frac{23}{112}+\frac3{112}=\frac17.
\]
Uniqueness of the scalar boundary orbit makes this the global response,
rather than the response of only one chosen scalar extremizer.

\subsection{Complex transverse curvature and the stronger remainder}

The following exact certificate verifies every complex tangent direction.
Use seven real sphere coordinates
\[
 v=(x_1+iy_1,\ iy_2,\ x_3+iy_3,\ x_4+iy_4)^T,
 \qquad z(s)=\frac{z_*+sv}{\sqrt{1+s^2\|v\|^2}},
\]
and sixteen real coordinates for $N(s)=e^{isH}N_*$, $H=H^*$.
Let $h_0,\ldots,h_3$ be its diagonal entries and list its upper off-diagonal
entries lexicographically as
\[
 h_4+ih_5,\ h_6+ih_7,\ h_8+ih_9,\
 h_{10}+ih_{11},\ h_{12}+ih_{13},\ h_{14}+ih_{15}.
\]
For the resulting 23-vector $\xi$, direct multiplication gives
\[
 |J_s(z(s),N(s))|^2=\frac7{16}+
 s^2\left(a+g^T\xi+\frac12\xi^TH\xi\right)+O(s^3),
\]
where
\[
 a=-\frac{23}{112},\qquad
 g^T\xi=-\frac34x_1-\frac17h_5-
             \frac{5\sqrt5}{56}h_7-\frac{\sqrt7}{56}h_9.
\]
The real symmetric matrix $H$ has characteristic polynomial
\begin{align*}
 \det(\lambda I-H)
 ={}&\frac{\lambda^{10}(8\lambda+7)^4
                 (16\lambda^2+42\lambda+25)^2}{368293445632}\\
 &\quad\cdot(10976\lambda^2+56620\lambda+53525)
       (32\lambda^3+308\lambda^2+655\lambda+350).
\end{align*}
Every nonzero eigenvalue is negative: all roots are real by symmetry,
and each displayed factor has positive coefficients and nonzero constant.
Thus $\operatorname{rank}H=13$.
The 23-coordinate chart has four redundant directions, since the map
$H\mapsto e^{iH}N_*$ has the four-dimensional stabilizer $U(2)$ on the
orthogonal complement of $\operatorname{ran}N_*$.
The physical Hessian therefore has a six-dimensional kernel on the
19-dimensional manifold.  These are precisely the six symmetry directions
already exhibited; the transverse Hessian is strictly negative.

For an additional direct check, a solution of $H\xi=-g$ has only
\[
 x_1=-6/7,\qquad h_5=-2/7,\qquad h_7=-\sqrt5/35
\]
nonzero, and its optimized quadratic value is
\[
 a+\tfrac12g^T\xi=-23/112+39/112=1/7.
\]
The displayed chart specifies $H$ explicitly; its characteristic polynomial,
the response equation and the optimized value can be checked with exact
algebraic arithmetic.

Compactness forces every maximizing sequence for $s\to0$ toward the
single boundary orbit. On a local analytic slice transverse to that orbit,
the implicit-function theorem supplies a unique analytic maximizing branch.
The exact symmetry makes its value independent of the chosen orbit point,
so it is the global maximum for all sufficiently small $|s|$.
The physical class depends only on $s^2$; hence this analytic value is even.
Its expansion is therefore
\[
 \max_X|J_s|^2=\frac7{16}+\frac17s^2+O(s^4),
\]
which proves \eqref{I:R:eq:three-h}.

An exactly feasible real branch attaining this response is
\[
 z(s)=\frac{e_2-(6s/7)e_1}{\sqrt{1+36s^2/49}},\qquad
 N(s)=e^{sL}N_*,
\]
where $L$ is real skew-symmetric with
$L_{12}=2/7$, $L_{13}=\sqrt5/35$ and all other upper entries zero.
Both $\|z(s)\|=1$ and $N(s)^TN(s)=I_2$ hold exactly.
This branch certifies the real lower side of the envelope expansion;
the full complex Hessian certifies its global upper side.

Finally, exact elimination of $A$ gives
\[
 \mathcal P_3(q)=\max_{0\le t\le1}
 \frac{q h_3(t)}{(1-\sqrt{(1-q)(1-t)})^2}.
\]
The bound $h_3(t)\le t$ localizes all maximizing $t/q$ to a compact
subinterval of $(0,\infty)$.  The leading rescaled objective is
$7x/[4(1+x)^2]$, with a unique nondegenerate maximum at $x=1$.
Substitution of \eqref{I:R:eq:three-h} proves \eqref{I:R:eq:three-P}.
The analytic scalar critical branch additionally satisfies
$t_*(q)=q+(32/49)q^2+O(q^3)$.

% End three_state.tex


% Begin four_state.tex

\section{An exact four-state first correction}
\label{I:R:sec:four-state-correction}

The all-state response formula can be evaluated exactly at the
degree-three scalar extremizer.  Its classification follows from
equality in the Hardy-space bridge, rather than from a new scalar
Hessian calculation.  The scalar extremal and its uniqueness are the
external results of Brevig--Saksman~\cite[Theorem 2]{brevig2020}; the response computation below is
an additional finite algebraic calculation.

\subsection{Classification through equality in the Hardy bridge}

\begin{lemma}[Transfer of scalar-extremizer classification]
\label{I:R:lem:Hardy-classification-transfer}
Let $m\ge1$.  Suppose every normalized extremal for the $m$th
coefficient functional on $H^{2/3}$ is an outer polynomial cube, and
let its collection of cube roots be $\mathcal P_m$.  The minimizing
degree-$m$ finite Blaschke products in the rational approximation of
$z^m$ are precisely
\[
 \Theta=P^*/P,\qquad P\in\mathcal P_m,
\]
up to unimodular constants and the normalization of $P$.
Consequently, uniqueness of the coefficient extremal up to rotation
gives uniqueness of the boundary passive extremizer up to unitary
state coordinates and the corresponding scalar phases.
\end{lemma}
\begin{proof}
Write $C=C(m,2/3)$.  The established Hardy bridge gives
$E_m=1/C$.  Every minimizing model space has exact dimension $m$:
otherwise adjoining a reproducing-kernel direction with nonzero
overlap with the residual strictly improves its fit.  Write its
Blaschke product as $\Theta=D^*/D$, with $D$ an outer polynomial of
degree at most $m$, and set $b=P_{\le m}\Theta$.
The coefficient identity and H\"older inequality used in the bridge
are
\begin{equation}
 \Rnorm D_2^2=[bD^2]_m
 \le C\Rnorm{bD^2}_{2/3}
 \le C\Rnorm b_2\Rnorm D_2^2.
 \label{I:R:eq:classification-Holder-chain}
\end{equation}
At a minimizing space $\Rnorm b_2=E_m=1/C$, so equality holds
throughout.  In particular normalized $bD^2$ is a coefficient
extremal.  It is outer by the stated hypothesis, hence $b$ is outer
as well.  Equality in H\"older gives
$|b|/\Rnorm b_2=|D|/\Rnorm D_2$ almost everywhere on the circle.
Uniqueness of an outer function with its boundary modulus therefore
gives $b=cD$ for a nonzero scalar $c$.  It follows that the normalized
cube of $D$ is a coefficient extremal: $D$ is a scalar multiple of
some $P\in\mathcal P_m$.  Conversely the stationarity identity
$P^2=C P^*+O(z^{m+1})$ in the upper half of the Hardy bridge makes
$P^*/P$ a minimizing Blaschke product.  This proves the asserted
classification.

For a lossless observable boundary tuple $(D,l,x)$, its observation
map is an isometry onto the minimizing model space and
$O^*e_m=(lD^m)^*=\Gamma_m x$.  Thus the model space fixes the tuple
up to a unitary state change and fixes the unit cyclic vector by this
display.  A rotation of the scalar variable only produces the
corresponding phase changes in the tuple.
\end{proof}


The hypothesis is supplied in general by the inspected outer-extremal
theorem and polynomial structure.  For $m=2,3$, the complete
Brevig--Saksman classifications also state directly that the
extremizers are the displayed outer cubes, unique up to rotation.
The response
\begin{equation}
 b(D,l,x)=|a_{m-1}|^2-\Gamma_m\Re S
                         +\Gamma_m v^*K^{-1}v,
 \label{I:R:eq:four-state-response-recalled}
\end{equation}
with the notation of the all-state theorem, is invariant under these
state unitaries and phases.  For example, under
$D\mapsto e^{i\theta}D$, $x\mapsto e^{-im\theta}x$, one has
$a_j\mapsto e^{i(j-m)\theta}a_j$, $S\mapsto S$,
$v\mapsto e^{-i(m+2)\theta}v$, and $K\mapsto K$.
Thus evaluating one real representative determines the global
response for $m=3$.

\subsection{An exact algebraic certificate}

Set
\begin{equation}
 \alpha=\frac{\sqrt{15-\sqrt{33}}}{8},\qquad
 \beta=6\alpha-\frac{64}{3}\alpha^3,
 \qquad\gamma=2-8\alpha^2.
 \label{I:R:eq:four-state-alpha}
\end{equation}
Then $64\alpha^4-30\alpha^2+3=0$, $0<\alpha^2<1/4$, and
$P(z)=1+\beta z+\gamma z^2+\alpha z^3$ is the
Brevig--Saksman extremal denominator.  Equivalently
$\beta=\sqrt{3+\sqrt{33}/3}/2$ and
$\gamma=(1+\sqrt{33})/8$.

Use the observable companion coordinates
\[
 T=\begin{pmatrix}0&1&0\\0&0&1\\-\alpha&-\gamma&-\beta\end{pmatrix},
 \qquad l_0=(1,0,0).
\]
An exact solution of $Q-T^{\mathsf T}QT=l_0^{\mathsf T}l_0$ is
\begin{equation}
 Q=\frac1{333}
 \begin{pmatrix}
 342+768\alpha^2&1116\alpha-4224\alpha^3&528-2304\alpha^2\\
 1116\alpha-4224\alpha^3&798-576\alpha^2&1980\alpha-6272\alpha^3\\
 528-2304\alpha^2&1980\alpha-6272\alpha^3&858-192\alpha^2
 \end{pmatrix}.
 \label{I:R:eq:four-state-Q}
\end{equation}
Its leading principal minors are
\[
 \frac{38}{37}+\frac{256}{111}\alpha^2,\quad
 \frac{244}{111}+\frac{896}{333}\alpha^2,\quad
 \frac{968}{333}+\frac{11008}{999}\alpha^2,
\]
all positive.  Since $(T,l_0)$ is observable, the positive Stein
identity gives strict stability.  Whitening yields the Euclidean
lossless pair $D=Q^{1/2}TQ^{-1/2}$, $l=l_0Q^{-1/2}$.

The following calculation stays in companion coordinates and the
$Q$ inner product.  Put
\[
 x_0=Q^{-1}(l_0T^3)^{\mathsf T}
 =\begin{pmatrix}
 -\alpha^3\\21\alpha^2/8-9/16\\27\alpha^3/4-15\alpha/8
 \end{pmatrix},\qquad
 G=l_0T^3x_0=\frac{45}{64}-\frac{33}{32}\alpha^2.
\]
Then $\Gamma_3=\sqrt G$, while the normalized cyclic vector has raw
coordinates $x=x_0/\sqrt G$.
To evaluate \eqref{I:R:eq:four-state-response-recalled} without any
additional square roots, define
\begin{align*}
 c_j&=l_0T^jx_0,\\
 S_0&=\sum_{j=0}^{2}c_jx_0^{\mathsf T}QT^{3-j}x_0,\\
 v_0&=c_0Tx_0+c_1x_0,\\
 M_0&=\sum_{j=0}^{3}T^{3-j}x_0l_0T^j,\\
 K_0&=\tfrac12(M_0+Q^{-1}M_0^{\mathsf T}Q).
\end{align*}
Here $K=K_0/\sqrt G$, $v=v_0/G$, and $S=S_0/G^{3/2}$ in raw
coordinates.  The symmetric matrix $H_0=QK_0$ simplifies to
\begin{equation}
 H_0=\begin{pmatrix}
 21/32+7\alpha^2/16&\alpha/2-3\alpha^3&5/8-9\alpha^2/4\\
 \alpha/2-3\alpha^3&7/4-5\alpha^2/2&3\alpha-20\alpha^3/3\\
 5/8-9\alpha^2/4&3\alpha-20\alpha^3/3&3-19\alpha^2/3
 \end{pmatrix}.
 \label{I:R:eq:four-state-H0}
\end{equation}
Its leading principal minors are
\[
 \frac{21}{32}+\frac7{16}\alpha^2,\qquad
 \frac{1287}{1024}-\frac{915}{512}\alpha^2,\qquad
 \frac{153}{64}-\frac{189}{32}\alpha^2.
\]
They are positive using $0<\alpha^2<1/4$.  Thus the response matrix
is positive definite directly in this example.

Reduction of the following scalar products modulo
$64\alpha^4-30\alpha^2+3$ gives
\begin{align}
 c_2^2&=\frac{6075}{32768}-\frac{7767}{16384}\alpha^2,\notag\\
 S_0&=\frac{6669}{65536}-\frac{1809}{32768}\alpha^2,\notag\\
 v_0^{\mathsf T}QK_0^{-1}v_0
 &=\frac{171}{65536}+\frac{5625}{32768}\alpha^2.
 \label{I:R:eq:four-state-scalar-certificate}
\end{align}
Therefore the exact response is
\begin{equation}
 b_4=\frac{c_2^2-S_0+v_0^{\mathsf T}QK_0^{-1}v_0}{G}
    =\frac{75}{496}-\frac{103}{248}\alpha^2
    =\boxed{\frac{855+103\sqrt{33}}{15872}}.
 \label{I:R:eq:four-state-b4}
\end{equation}
The Stein equation, response matrix and coefficient admit direct exact
arithmetic checks in $\mathbb Q(\alpha)$ using the displayed identities.
Numerically $b_4=0.09114730044061346723567\ldots$; this decimal is
not used to obtain the coefficient.

\begin{theorem}[Four-state first boundary correction]
\label{I:R:thm:four-state-first-correction}
For $f_{r,4}(a^j b^4)=qr^j$, zero elsewhere, the squared optimum for
each field satisfies
\begin{equation}
 \boxed{\displaystyle
 D_4^{\mathbb F}(q)
  =\frac{1103-33\sqrt{33}}{2048}
    -\frac{855+103\sqrt{33}}{15872}\,q+o(q),
 \qquad\mathbb F\in\{\mathbb R,\RC\}.}
 \label{I:R:eq:four-state-optimum}
\end{equation}
\end{theorem}
\begin{proof}
Lemma~\ref{I:R:lem:Hardy-classification-transfer} and the external
degree-three classification identify every complex boundary
extremizer with the real representative used above, modulo symmetries
preserving the response.  The all-state normal-response lower bound
therefore has global coefficient \eqref{I:R:eq:four-state-b4}.
The real structured first-order matching construction uses this real
lossless tuple and its maximizing real normal directions.  It gives
the same coefficient in a rank-at-most-four competitor, with
$O(q^2)$ error relative to its analytic overlap branch.  Since
$\Gamma_3^2=G=(945+33\sqrt{33})/2048$, the complex lower and real
upper bounds give \eqref{I:R:eq:four-state-optimum} for each field.
The global lower remainder established here is $o(q)$; a global
$O(q^2)$ remainder is not inferred without a further tangential
nondegeneracy argument.
\end{proof}


% End four_state.tex


% Begin real_diagonal_lemma.tex

\section{Real representatives of every diagonal boundary extremizer}
\label{I:R:subsec:real-diagonal-orbits}

The diagonal case admits a further field comparison.  This argument
uses the exact scalar Hardy bridge and the strict increase of the
coefficient constants supplied by the October 2026 outer-extremal
result~\cite{outer2026}; it is not asserted for independent budget and delay.

\begin{lemma}[Positive coefficients and affine phases]
\label{I:R:lem:positive-Hardy-extremizers}
Fix $m\ge1$ and let $F=P^3$ be any normalized extremal for the
$m$-th coefficient on $H^{2/3}$, with $\Rnorm P_2=1$ and $P$ outer of
exact degree $m$.  There exist a polynomial
$Q(z)=\sum_{j=0}^m q_jz^j$ with every $q_j>0$ and real phases
$\alpha,\theta$ such that
\[
 P(z)=e^{i\alpha}Q(e^{i\theta}z).
\]
\end{lemma}
\begin{proof}
Write $P(z)=\sum_{j=0}^m p_jz^j$ and put
$Q(z)=\sum_{j=0}^m|p_j|z^j$.  Then $\Rnorm Q_2=1$ and
\[
 C(m,2/3)=|[P^3]_m|\le[Q^3]_m\le C(m,2/3).
\]
Thus $Q^3$ is also a coefficient extremal.  Its support
$S=\{j:q_j>0\}$ contains $0,m$: $P$ is outer of exact degree $m$.
If $q_j=0$, varying that real coefficient and renormalizing the
Euclidean norm shows $[Q^2]_{m-j}=0$; otherwise a positive first-order
variation would increase $[Q^3]_m$.  Consequently
\[
 i,h\in S,\quad i+h\le m\quad\Longrightarrow\quad m-i-h\in S.
\]
Taking $h=0$ shows $m-S=S$.  Applying this reflection once more gives
closure under addition when the sum is at most $m$.  Applying the rule
to $h$ and $m-i$ gives $i-h\in S$ whenever $i\ge h$ belong to $S$.
Let $g$ be the smallest positive element of $S$.  Repeated subtraction
shows every element of $S$ is divisible by $g$, and addition shows all
multiples up to $m$ occur.  Hence $g\mid m$ and
$Q(z)=R(z^g)$ with $\deg R=m/g$.  If $g>1$, then
\[
 C(m,2/3)=[Q^3]_m=[R^3]_{m/g}
 \le C(m/g,2/3)<C(m,2/3),
\]
a contradiction to the strict increase of coefficient constants.
Therefore $g=1$ and every $q_j$ is positive.

Equality in the first triangle inequality now says that every monomial
$p_ip_hp_k$ with $i+h+k=m$ has the same argument.  Put
$p_j=q_je^{i\alpha_j}$.  Comparing the nonzero terms with index triples
$(0,j,m-j)$ and $(1,j-1,m-j)$ gives
\[
 \alpha_j-\alpha_{j-1}=\alpha_1-\alpha_0\pmod{2\pi}
 \qquad(1\le j\le m).
\]
Thus $\alpha_j=\alpha_0+j\theta$ modulo $2\pi$, proving the claim.
\end{proof}


\begin{proposition}[All diagonal boundary models have real representatives]
\label{I:R:prop:real-diagonal-boundary-models}
Every maximizing complex passive triple in $\mathcal M_m$ is equivalent
to a real triple through state unitaries and phases of the scalar
transition and output.  In particular,
\[
 b_{k,k}^{\RR}=b_{k,k}^{\RC}\qquad(k=m+1\ge2).
\]
Consequently the real and complex structured optima have the same
first-order diagonal expansion, without an assertion of equality at
finite $q$.
\end{proposition}
\begin{proof}
A maximizing passive pair is lossless, stable and observable by
Lemma~\ref{I:R:lem:cyclic-lossless-maximizer}.  Its observation map is an
isometry onto an $m$-dimensional model space $K_\Theta$ minimizing the
error of approximation to $z^m$.  Write
$\Theta=e^{i\phi}D_0^*/D_0$, with $D_0$ outer of degree at most $m$
and reversal of order $m$, and set $b=P_{\le m}\Theta$.
The Hardy bridge gives $\Rnorm b_2=1/C(m,2/3)$, while its exact
coefficient identity gives
\[
 \Rnorm{D_0}_2^2=|[bD_0^2]_m|
 \le C(m,2/3)\Rnorm{bD_0^2}_{2/3}
 \le C(m,2/3)\Rnorm b_2\Rnorm{D_0}_2^2
 =\Rnorm{D_0}_2^2.
\]
Therefore $bD_0^2$, after normalization and a phase, is a Hardy
coefficient extremal and hence outer.  Since $D_0$ is outer, $b$ is
outer as well.  Equality in H\"older gives
$|b|^2=c|D_0|^2$ almost everywhere on the circle.  Uniqueness of outer
functions with prescribed boundary modulus then gives $b=\mu D_0$.
It follows that the normalized $D_0^3$ is a coefficient extremal.
Lemma~\ref{I:R:lem:positive-Hardy-extremizers} shows that $D_0$ has real
positive coefficients after a scalar phase and circle rotation.
Thus $K_\Theta$ is a circle rotation of a model space invariant under
coefficient conjugation.  Choosing a real orthonormal basis of that
space gives real canonical matrices and readout.  Every original
maximizing passive pair is unitarily equivalent to its canonical
model, and its maximizing $x$ equals $(lD^m)^*/\Gamma_m$.
This proves the stated orbit classification.

For clarity, the response formula is invariant under
\[
 (D,l,x)\mapsto(U^*DU,lU,U^*x),\quad
 (D,l,x)\mapsto(D,e^{i\varphi}l,e^{-i\varphi}x),\quad
 (D,l,x)\mapsto(e^{i\theta}D,l,e^{-im\theta}x).
\]
The first operation conjugates $K$ and transforms $v$ by $U^*$.
The second fixes $a_j,S,K$ and multiplies $v$ by a phase.
For the third, $a_j$ acquires $e^{i(j-m)\theta}$,
$S,K$ are unchanged, and $v$ acquires $e^{-i(m+2)\theta}$.
Every term in the response coefficient is unchanged.  Since every
complex boundary triple has a real representative, maximizing over
complex triples gives the same value as maximizing over real triples.
The structured first-order theorem supplies the final conclusion.
\end{proof}


% End real_diagonal_lemma.tex


% Begin full_channel.tex

\section{All-state, all-delay projection bounds retaining every singular channel}
\label{I:R:sec:full-channel}

The dominant-vector argument controls only one singular direction of the target.
We now retain all of them and eliminate the waiting transition exactly. This
gives a finite-dimensional characterization of an invariant-subspace relaxation
for every state budget. Structured Hankel completion remains an additional step.

Fix integers $n,m\ge1$, $0<q<1$, and $r=\sqrt{1-q}$. On the binary alphabet
$\{a,b\}$ let
\[
 f_{r,m}(a^j b^m)=qr^j\quad(j\ge0),\qquad
 f_{r,m}(w)=0\text{ otherwise},\qquad
 D_{n,m}(q)=\delta_n^{\RC}(H_{f_{r,m}})^2.
\]
All norms here use counting measure on words. With $e_w$ the coordinate vectors,
put
\[
 v=\sqrt q\sum_{j\ge0}r^j e_{a^j b^m},\qquad
 e_j=e_{b^j}\quad(0\le j<m),\qquad
 \mathcal U=(v,e_0,\ldots,e_{m-1}).
\]
These columns are orthonormal, and direct grouping of the target's rows gives
\begin{equation}
 H_{f_{r,m}}^*H_{f_{r,m}}
   =vv^*+q\sum_{j=0}^{m-1}e_je_j^*
   =\mathcal U C\mathcal U^*,\qquad C=\operatorname{diag}(1,qI_m).
 \label{I:R:eq:fc-singular}
\end{equation}
Indeed the pure-$a$ rows give the unit singular channel; rows $a^i b^h$,
$1\le h\le m$, give the other $m$ orthogonal channels.

Let $L_a,L_b$ be the left creation shifts and define
\begin{equation}
 R_{n,m}(q)=\inf_{\substack{\dim S\le n\\L_a^*S,L_b^*S\subseteq S}}
          \Rnorm{H_{f_{r,m}}(I-P_S)}^2.
 \label{I:R:eq:fc-relaxation}
\end{equation}
For every bounded rank-at-most-$n$ word-Hankel $G$, the row space
$S=\operatorname{ran}G^*$ is invariant under these backward shifts and
$G(I-P_S)=0$. Consequently
\begin{equation}
 D_{n,m}(q)\ge R_{n,m}(q).
 \label{I:R:eq:fc-relaxation-bound}
\end{equation}
The definition allows invariant spaces that need not be cyclic row spaces of a
word-Hankel operator. Equality in \eqref{I:R:eq:fc-relaxation-bound} therefore
requires a separate argument.

\subsection{An exact residual matrix for the relaxation}

Use standard matrix products $T_{a_1\cdots a_j}=T_{a_1}\cdots T_{a_j}$, with
$T_a=A,T_b=B$. For a passive tuple
\begin{equation}
 A^*A+B^*B+\ell^*\ell\preceq I_n,
 \label{I:R:eq:fc-passive}
\end{equation}
the appropriate observation convention is
\[
 (\mathcal Ox)[w]=\ell T_{\operatorname{rev}(w)}x.
\]
The reversal is essential for left-backward-shift coordinates: the coordinate
at $a^j b^m$ is $\ell B^m A^j x$. Telescoping
\eqref{I:R:eq:fc-passive} over complete word levels proves
$\Rnorm{\mathcal O}\le1$, and
$L_c^*\mathcal O=\mathcal OT_c$ for $c=a,b$.
Define
\begin{equation}
 w_0=\sqrt q\,\ell B^m(I-rA)^{-1},\qquad
 Z=\begin{pmatrix}\ell\\\ell B\\\vdots\\\ell B^{m-1}\end{pmatrix},
 \qquad W=\begin{pmatrix}w_0\\Z\end{pmatrix},\qquad
 E=C-C^{1/2}WW^*C^{1/2}.
 \label{I:R:eq:fc-E}
\end{equation}
Here $W=\mathcal U^*\mathcal O$. The inverse exists uniformly over the passive
class because $r<1$ and $\Rnorm{A}\le1$.

\begin{theorem}[Exact passive representation of the projection relaxation]
\label{I:R:thm:fc-residual}
For every $n,m\ge1$ and $0<q<1$,
\begin{equation}
 R_{n,m}(q)=\min_{A^*A+B^*B+\ell^*\ell\preceq I_n}
                  \lambda_{\max}(E).
 \label{I:R:eq:fc-passive-min}
\end{equation}
\end{theorem}
\begin{proof}
For an invariant $S$ in \eqref{I:R:eq:fc-relaxation}, take orthonormal state
coordinates, set $A=L_a^*|_S$, $B=L_b^*|_S$, and let $\ell$ evaluate the empty
coordinate. The observation map is the isometric inclusion of $S$ and the
energy identity gives \eqref{I:R:eq:fc-passive}, with equality. Formula
\eqref{I:R:eq:fc-singular} then yields
\[
 \Rnorm{H_{f_{r,m}}(I-P_S)}^2
 =\lambda_{\max}\!\left(C^{1/2}\mathcal U^*(I-P_S)\mathcal U C^{1/2}\right)
 =\lambda_{\max}(E).
\]
Conversely, for any passive tuple, $K=\operatorname{ran}\mathcal O$ is a
backward-invariant space of dimension at most $n$, and
$\mathcal O\mathcal O^*\preceq P_K$. Thus its matrix $E$ dominates the exact
residual matrix for $K$ in the positive-semidefinite order. Its largest
eigenvalue is at least $R_{n,m}(q)$. The two comparisons prove equality of the
infima. The passive class is compact and its objective continuous, so the
minimum exists. Smaller state spaces are included by zero padding.
\end{proof}


\subsection{Threshold reduction and exact elimination of the waiting matrix}

For a threshold $\tau>q$, set
\begin{equation}
 Q_\tau=I_n+\frac{q}{\tau-q}Z^*Z\succ0.
 \label{I:R:eq:fc-Q}
\end{equation}
The lower-right block of $\tau I-E$ is
$(\tau-q)I_m+qZZ^*\succ0$. Its Schur complement, followed by the Woodbury
identity, gives the exact equivalence
\begin{equation}
 E\preceq\tau I_{m+1}
 \quad\Longleftrightarrow\quad
 w_0Q_\tau^{-1}w_0^*\ge1-\tau.
 \label{I:R:eq:fc-threshold}
\end{equation}
This includes every subdominant target channel through $Z^*Z$.

\begin{lemma}[A weighted trust-region elimination]
\label{I:R:lem:fc-trust}
Fix $B,\ell$ with $D=I-B^*B-\ell^*\ell\succeq0$, and let $Q\succ0$.
For a unit vector $y$, put $d=y^*Dy$ and
\begin{equation}
 \mathcal T_Q(y,d)=\min_{\Rnorm{z}^2\le d}(y-rz)^*Q(y-rz).
 \label{I:R:eq:fc-trust-value}
\end{equation}
Then
\begin{align}
 &\max_{A^*A\preceq D}
 q\ell B^m(I-rA)^{-1}Q^{-1}(I-rA^*)^{-1}(B^*)^m\ell^*
 \nonumber\\[-2pt]
 &\hspace{30mm}=\max_{\Rnorm{y}=1}
       \frac{q|\ell B^m y|^2}{\mathcal T_Q(y,y^*Dy)}.
 \label{I:R:eq:fc-eliminate-A}
\end{align}
If $d=0$, then $\mathcal T_Q(y,0)=y^*Qy$. If $d>0$, its unique minimizer is
\begin{equation}
 z_\mu=r(r^2Q+\mu I)^{-1}Qy,
 \qquad \Rnorm{z_\mu}^2=d,
 \label{I:R:eq:fc-secular}
\end{equation}
where the last equation has a unique positive solution $\mu$.
\end{lemma}
\begin{proof}
As $A$ ranges over $A^*A\preceq D$, its value $Ay$ ranges over exactly the ball
$\Rnorm{z}\le\sqrt d$. Necessity follows from $\Rnorm{Ay}^2\le y^*Dy$. For
$d>0$, the rank-one choice
\[
 A=zy^*D/d
\]
has $Ay=z$ and $A^*A\preceq D$, since
$Dyy^*D\preceq dD$ and $\Rnorm{z}^2\le d$. The case $d=0$ is immediate.

For a fixed $A$, the weighted row norm on the left of
\eqref{I:R:eq:fc-eliminate-A} is
\[
 \max_{x\ne0}\frac{q|\ell B^m(I-rA)^{-1}x|^2}{x^*Qx}
 =\max_{y\ne0}\frac{q|\ell B^m y|^2}
                          {((I-rA)y)^*Q((I-rA)y)}.
\]
Homogeneity permits $\Rnorm{y}=1$. Maximizing over $A,y$ and using the exact
ball of attainable $Ay$ gives \eqref{I:R:eq:fc-eliminate-A}. The denominator is
strictly positive because $r<1$ and $d\le1$.

The trust-region objective is strictly convex. Its unconstrained minimizer is
$y/r$, whose norm $1/r$ exceeds the feasible radius $\sqrt d\le1$.
The constraint is active, and the multiplier equations give
\eqref{I:R:eq:fc-secular}. Diagonalizing $Q$ shows that
$\Rnorm{z_\mu}^2$ decreases strictly from $1/r^2$ to zero as $\mu$ increases
from zero to infinity. This proves existence, uniqueness and sufficiency.
\end{proof}


Combining the preceding statements gives an unconditional all-state threshold
formula. Define
\begin{equation}
 \Psi_{n,m}(q,\tau)=
 \max_{\substack{B^*B+\ell^*\ell\preceq I_n\\\Rnorm{y}=1}}
 \frac{q|\ell B^m y|^2}{\mathcal T_{Q_\tau}(y,y^*Dy)}.
 \label{I:R:eq:fc-Psi}
\end{equation}
Here $Z,Q_\tau,D$ depend only on $B,\ell$. Then, for every $\tau>q$,
\begin{equation}
 \boxed{R_{n,m}(q)\le\tau
       \quad\Longleftrightarrow\quad
       \Psi_{n,m}(q,\tau)\ge1-\tau.}
 \label{I:R:eq:fc-general-threshold}
\end{equation}
The remaining optimization is compact. Each objective evaluation requires one
positive scalar secular root, in place of an optimization over $A$.
Equation \eqref{I:R:eq:fc-general-threshold} characterizes the relaxation at
thresholds above $q$; it does not assert structured attainability, or determine
values below $q$ when the state budget is large.

\subsection{The exact rank-one benchmark is recovered}

For $n=m=1$, write $\kappa=q/(\tau-q)$ and $K=1+\kappa$. Phases may be aligned
so that $a=|A|\ge0$ minimizes $|1-rA|$. With $x=|B|^2$ and $z=|\ell|^2$,
the overlap in \eqref{I:R:eq:fc-threshold} is
\[
 \frac{qzx}{(1-ra)^2(1+\kappa z)},\qquad a^2+x+z\le1.
\]
At fixed $a$, the available budget $t=1-a^2$ is fully used. Direct
differentiation gives
\[
 z_* =\frac{t}{1+\sqrt{1+\kappa t}},\qquad
 \max_{0\le z\le t}\frac{z(t-z)}{1+\kappa z}
       =\frac{t^2}{(1+\sqrt{1+\kappa t})^2}.
\]
The remaining maximization is elementary. Equivalently maximize
\[
 F(a)=\frac{\sqrt{1+\kappa(1-a^2)}-1}{1-ra},\qquad0\le a\le1.
\]
Its interior stationary equation is
\[
 \sqrt{1+\kappa(1-a^2)}+\frac{\kappa a}{r}=1+\kappa.
\]
Solving with this unsquared sign condition gives its unique maximum at
$a_*=r\sqrt K/(\sqrt K+\sqrt q)$, where
$F(a_*)=\kappa/(1+\sqrt{qK})$. Consequently
\begin{equation}
 \boxed{\Psi_{1,1}(q,\tau)
      =\frac{q}{\left(1+\sqrt{q\tau/(\tau-q)}\right)^2}.}
 \label{I:R:eq:fc-scalar-Psi}
\end{equation}

\begin{theorem}[Monotone radical representation of the exact F1 value]
\label{I:R:thm:fc-F1-radical}
The exact rank-one value $D_1(q)=D_{1,1}(q)$ is the unique root $\tau\in(q,1)$
of
\begin{equation}
 \boxed{1-\tau
       =\frac{q}{\left(1+\sqrt{q\tau/(\tau-q)}\right)^2}.}
 \label{I:R:eq:fc-F1-root}
\end{equation}
It equals the real optimum as well.
\end{theorem}
\begin{proof}
The right side of \eqref{I:R:eq:fc-F1-root} increases continuously from zero as
$\tau$ increases from $q$, while $1-\tau$ decreases from $1-q$ to zero.
There is exactly one crossing in $(q,1)$. Equations
\eqref{I:R:eq:fc-general-threshold}--\eqref{I:R:eq:fc-scalar-Psi} identify it as the
projection optimum $R_{1,1}(q)$. The previously proved rank-one output
completion theorem supplies a real rank-one Hankel attaining this relaxation,
so its value is the exact real and complex F1 optimum.

For an algebraic comparison with the selected quartic, put
\[
 \mathcal L_q(\tau)=(1+q)(\tau-q)-(1-q)\tau^2,\qquad
 \mathcal A_q(\tau)=(1-q)^2\tau^2+(q^2+2q-1)(\tau-q).
\]
The radical equation implies
$\mathcal L_q(\tau)=2(\tau-q)\sqrt{q(1-\tau)}>0$, and direct expansion gives
\begin{align*}
 &\mathcal A_q(\tau)^2-4q^3(\tau-q)(1-2q+q\tau)\\
 &\qquad=(1-q)^2\bigl[\mathcal L_q(\tau)^2
                        -4q(1-\tau)(\tau-q)^2\bigr].
\end{align*}
Thus \eqref{I:R:eq:fc-F1-root} provides a monotone, unambiguous representation of
the already selected quartic branch.
\end{proof}


This benchmark also detects the limitation of keeping only the dominant
singular direction. In scalar dimension its local envelope is
$h_{1,1}(t)=t^2/4$, and exact maximization gives
\[
 P_{1,1}(q)=\frac{q}{(1+\sqrt q)^2},\qquad
 t_* =\frac{2\sqrt q}{1+\sqrt q}.
\]
The corresponding lower bound $L_{\mathrm{dom}}=1-P_{1,1}$ agrees with the
exact F1 value through $q^2$, but
\begin{align*}
 D_1(q)&=1-q+2q^{3/2}-3q^2+5q^{5/2}-8q^3+O(q^{7/2}),\\
 L_{\mathrm{dom}}(q)&=1-q+2q^{3/2}-3q^2+4q^{5/2}-5q^3+O(q^{7/2}),\\
 D_1(q)-L_{\mathrm{dom}}(q)&=q^{5/2}-3q^3+O(q^{7/2}).
\end{align*}
The full-channel threshold therefore passes an exact benchmark that the
dominant-vector envelope fails at higher order. Its equality with the Hankel
optimum has been established here only for F1; higher-state matching
constructions remain separate problems.

% End full_channel.tex


% Begin terminal_languages.tex

\section{Arbitrary terminal languages: an exact boundary reduction}
\label{I:R:sec:terminal-languages}

The waiting-letter mechanism extends beyond a monomial suffix.  The
appropriate scalar approximation problem is replaced by simultaneous
approximation of all terminal residuals in one common invariant space.
This section proves that extension, without a uniqueness or a
nondegeneracy assumption on its terminal extremizers.

Let $\Delta$ be a finite alphabet, let $a\notin\Delta$, and write
$\Sigma=\{a\}\sqcup\Delta$.  Let
$\tau:\Delta^*\to\RC$ have bounded counting-measure word-Hankel operator
$H_\tau$.  In particular $\tau\in\ell^2(\Delta^*)$; a polynomial
terminal language always satisfies this hypothesis.  For
$q=1-r^2\in(0,1)$ define
\begin{equation}
 f_{r,\tau}(a^j w)=qr^j\tau(w)
 \quad(j\ge0,\ w\in\Delta^*),
 \qquad f_{r,\tau}=0\text{ on the remaining words}.
 \label{I:R:eq:terminal-target}
\end{equation}
The parsing in \eqref{I:R:eq:terminal-target} is unique.  Fix a state budget
$n\ge1$ and put $\tau_j(v)=\tau(jv)$ for $j\in\Delta$.
Let $\mathfrak K_{n-1}$ be the collection of subspaces of
$\ell^2(\Delta^*)$ of dimension at most $n-1$ invariant under every
left backward shift.  Define
\begin{equation}
 E_{n-1}(\tau)^2=
 \inf_{K\in\mathfrak K_{n-1}}
       \sum_{j\in\Delta}\Rnorm{\tau_j-P_K\tau_j}_2^2.
 \label{I:R:eq:terminal-E}
\end{equation}
The same $K$ is used for every $j$; separate residual fits do not give
the stated state budget.  When $n=1$, use $K=\{0\}$.

\subsection{Exact elimination and dimension loss}

We use the canonical backward-shift convention
$B_w=B_{w_s}\cdots B_{w_1}$ for $w=w_1\cdots w_s$, and
$B_\varepsilon=I$.  Thus coordinate evaluation is $\ell B_wx$.
For a passive tuple
\begin{equation}
 A^*A+\sum_{j\in\Delta}B_j^*B_j+\ell^*\ell\preceq I_n,
 \qquad V=\begin{pmatrix}(B_j)_{j\in\Delta}\\\ell\end{pmatrix},
 \label{I:R:eq:terminal-passive}
\end{equation}
set
\[
 T_\tau(B,\ell)=\sum_{w\in\Delta^*}
          \overline{\tau(w)}\,\ell B_w,
 \qquad
 h_{\tau,n}(t)=
 \max_{\substack{V^*V\preceq I_n\\\Rnorm{Ve_1}^2=t}}
                  |T_\tau(B,\ell)e_1|^2.
\]
The observation map $Ox=(\ell B_wx)_w$ is contractive by telescoping
the passive energy inequality.  In fact
\begin{equation}
 \sum_w|\ell B_wy|^2
 \le |\ell y|^2+\sum_j\Rnorm{B_jy}^2=\Rnorm{Vy}^2.
 \label{I:R:eq:terminal-refined-energy}
\end{equation}
Consequently the series defining $T_\tau$ converges in operator norm
uniformly over all passive tuples, and
$0\le h_{\tau,n}(t)\le\Rnorm\tau_2^2t$.  Compactness then gives the
maxima in this definition and below.

\begin{lemma}[Universal waiting-transition elimination]
\label{I:R:lem:terminal-elimination}
For every $q\in(0,1)$,
\begin{align}
 P_{\tau,n}(q)
 &:=\max_{\eqref{I:R:eq:terminal-passive}}
       q\Rnorm{T_\tau(B,\ell)(I-rA)^{-1}}^2\notag\\
 &=\max_{0\le t\le1}
       \frac{q h_{\tau,n}(t)}{(1-r\sqrt{1-t})^2}.
 \label{I:R:eq:terminal-elimination}
\end{align}
\end{lemma}
\begin{proof}
For a unit vector $y$, put $D=I-V^*V$ and
$t=\Rnorm{Vy}^2=1-y^*Dy$.  The substitution $x=(I-rA)y$ in the row
operator norm and the passive inequality give
\[
 \Rnorm{(I-rA)y}\ge1-r\Rnorm{Ay}\ge1-r\sqrt{1-t}.
\]
A unitary state change sends $y$ to $e_1$ and preserves the terminal
row and constraints.  This proves the upper direction of
\eqref{I:R:eq:terminal-elimination}.  Conversely, for fixed $V,y$ and
$t<1$, choose
\[
 A=\frac{yy^*D}{\sqrt{1-t}}.
\]
Then $Ay=\sqrt{1-t}\,y$ and
$A^*A=Dyy^*D/(y^*Dy)\preceq D$, by Cauchy--Schwarz in the $D$
metric.  Both denominator inequalities are equalities.  When $t=1$
choose $A=0$.  This proves exact attainability of the scalar
denominator.
\end{proof}


An exact factorization makes the localization transparent.  For
$\rho(x,y)=|x-y|/|1-xy|$ on the real unit interval, put
\begin{equation}
 K_q(t)=\frac{qt}{(1-\sqrt{(1-q)(1-t)})^2}
       =1-\rho\bigl(\sqrt{1-q},\sqrt{1-t}\bigr)^2.
 \label{I:R:eq:terminal-kernel}
\end{equation}
Thus $0\le K_q(t)\le1$ and $K_q(q)=1$.  This is just
$(1-r\sqrt{1-t})^2-qt=(r-\sqrt{1-t})^2$.

\begin{lemma}[The residual dimension-drop coefficient]
\label{I:R:lem:terminal-dimension-drop}
As $t\downarrow0$,
\begin{equation}
 h_{\tau,n}(t)=a_{\tau,n}t+o(t),
 \qquad
 a_{\tau,n}=|\tau(\varepsilon)|^2+
       \max_{K\in\mathfrak K_{n-1}}
             \sum_{j\in\Delta}\Rnorm{P_K\tau_j}_2^2.
 \label{I:R:eq:terminal-leading-coefficient}
\end{equation}
The maximum is attained, and $P_{\tau,n}(q)\to a_{\tau,n}$.
\end{lemma}
\begin{proof}
Write $t=s^2$ and let $N=|\Delta|n+1$.  Every feasible $V$ has the
exact compact parametrization
\begin{equation}
 V=[sz,W_s(z)Z],\qquad
 \Rnorm z=1,\quad Z^*Z\preceq I_{n-1},\qquad
 W_s(z)=I_N-(1-\sqrt{1-s^2})zz^*.
 \label{I:R:eq:terminal-parametrization}
\end{equation}
Indeed this is the Schur complement of $I-V^*V$.  For a polynomial
$\tau$, the quotient $T_\tau e_1/s$ has an analytic extension to
$s=0$.  For general $\tau\in\ell^2$, the tail bound from
\eqref{I:R:eq:terminal-refined-energy}, divided by $|s|$, is at most
the terminal tail's $\ell^2$ norm.  Polynomial approximation therefore
gives a uniform continuous extension on the compact parameter domain.

At $s=0$ write
\[
 B_j(0)=\begin{pmatrix}0&b_j\\0&C_j\end{pmatrix},\qquad
 \ell(0)=(0,l),\qquad z_j=(\alpha_j,\xi_j),
\]
where $z_j$ is the first-column derivative block for $B_j$, and let
$z_o$ be its scalar output block.  The limiting divided terminal row is
\begin{equation}
 \overline{\tau(\varepsilon)}z_o+
 \sum_j R_j(C,l)\xi_j,
 \qquad
 R_j(C,l)=\sum_v\overline{\tau(jv)}\,lC_v.
 \label{I:R:eq:terminal-divided-limit}
\end{equation}
The $\alpha_j$ do not contribute.  Maximizing over unit $z$ gives
$|\tau(\varepsilon)|^2+\sum_j\Rnorm{R_j(C,l)}^2$.
The remaining columns in \eqref{I:R:eq:terminal-parametrization} imply
$\sum_j C_j^*C_j+l^*l\preceq I_{n-1}$.  Conversely every such reduced
tuple is feasible by choosing $b_j=0$.  Uniform convergence of the
compact objective proves
\[
 a_{\tau,n}=|\tau(\varepsilon)|^2+
 \max_{\sum C_j^*C_j+l^*l\preceq I_{n-1}}
                    \sum_j\Rnorm{R_j(C,l)}^2.
\]
For $n=1$ the reduced state space is zero and the sum is zero.

This passive maximum equals the invariant-space maximum in
\eqref{I:R:eq:terminal-leading-coefficient}.  A canonical observation map
on $K$ is its isometric inclusion, so $R_j$ has norm
$\Rnorm{P_K\tau_j}$.  Conversely a passive reduced tuple has a
contractive observation map $O$ whose finite-dimensional range $K$ is
backward invariant.  Since $OO^*\preceq P_K$,
$\Rnorm{O^*\tau_j}\le\Rnorm{P_K\tau_j}$.  This gives the reverse
inequality after summing.  A maximizing passive tuple exists by
compactness and uniform convergence; its range attains the
invariant-space maximum as well.

Finally write \eqref{I:R:eq:terminal-elimination} as
$P_{\tau,n}(q)=\max_t K_q(t)h_{\tau,n}(t)/t$.
The trial $t=q$ tends to $a_{\tau,n}$.  On $t\ge\eta>0$ the kernel
$K_q(t)$ tends uniformly to zero, whereas on $0<t<\eta$ it is at most
one and $h(t)/t\to a_{\tau,n}$.  Splitting at $\eta$ and then sending
$\eta$ to zero proves the upper limit.  This also covers
$a_{\tau,n}=0$.
\end{proof}


\subsection{The exact boundary theorem and an explicit lift}

\begin{theorem}[Simultaneous terminal approximation]
\label{I:R:thm:terminal-boundary}
For every $n\ge1$ and terminal series satisfying the stated boundedness
hypothesis,
\begin{equation}
 \boxed{\displaystyle
 \lim_{r\uparrow1}\delta_n^{\RC}(H_{f_{r,\tau}})^2
       =E_{n-1}(\tau)^2.}
 \label{I:R:eq:terminal-boundary-theorem}
\end{equation}
The infimum in \eqref{I:R:eq:terminal-E} is attained.  For real $\tau$, the
same statement holds separately for real competitors and real
invariant spaces.  Equality of the two fields' constants is not
assumed.
\end{theorem}
\begin{proof}
The unit vector
$u_r=\sqrt q\sum_{m\ge0}r^m e_{a^m}$ is mapped by
$H_{f_{r,\tau}}^{\mathsf T}$ to
\[
 v_r=\sqrt q\sum_{m,w}r^m\tau(w)e_{a^m w},\qquad
 \Rnorm{v_r}=\Rnorm\tau_2.
\]
For a bounded rank-at-most-$n$ competitor, its transposed range $S$ is
backward invariant.  Canonical orthonormal state coordinates yield
\eqref{I:R:eq:terminal-passive} and
\[
 \Rnorm{P_Sv_r}^2
   =q\Rnorm{T_\tau(B,\ell)(I-rA)^{-1}}^2.
\]
Orthogonal projection thus gives
\[
 \delta_n^{\RC}(H_{f_{r,\tau}})^2
   \ge\Rnorm\tau_2^2-P_{\tau,n}(q).
\]
Lemma~\ref{I:R:lem:terminal-dimension-drop} proves the lower boundary
bound in \eqref{I:R:eq:terminal-boundary-theorem}, because
$\Rnorm\tau_2^2=|\tau(\varepsilon)|^2+\sum_j\Rnorm{\tau_j}_2^2$.

For the upper bound choose a maximizing space $K$ in
Lemma~\ref{I:R:lem:terminal-dimension-drop}.  Put
\[
 h(\varepsilon)=\tau(\varepsilon),\qquad
 h(jv)=(P_K\tau_j)(v),\qquad e=\tau-h.
\]
All residuals $x_j=P_K\tau_j$ share $K$.  With $C_j$ the canonical
backward shifts on $K$ and $l$ evaluation at the empty word, the usual
forward realization of the lifted series
$g_r(a^m w)=qr^m h(w)$ uses at most $1+\dim K\le n$ states:
\begin{equation}
 A_a=\begin{pmatrix}r&0\\0&0\end{pmatrix},\qquad
 A_j=\begin{pmatrix}0&x_j^{\mathsf T}\\0&C_j^{\mathsf T}\end{pmatrix},
 \quad\alpha^*=q(1,0),\quad
 \beta=\begin{pmatrix}\tau(\varepsilon)\\l^{\mathsf T}\end{pmatrix}.
 \label{I:R:eq:terminal-lift}
\end{equation}
These are coefficient transposes, not Hermitian adjoints.  Products
give $h(jv)=lC_vx_j$ with the reversed canonical word order.  A letter
$a$ following a terminal letter makes the product zero, as required.
The common terminal tuple is jointly strictly stable: in its
finite-dimensional shift-invariant space, finite-horizon tail energies
tend uniformly to zero.  Thus its lifted Hankel is bounded.

Here is also a direct infinite-operator check.  Let $U_r$ be the
isometry from $\ell^2(\Delta^+)$ given by
$U_re_w=\sqrt q\sum_mr^m e_{a^m w}$, let $W$ embed
$\ell^2(\Delta^*)$ as words without $a$, and let
$K_e[w,v]=e(wv)$ with $w\ne\varepsilon$.  Since
$e(\varepsilon)=0$, direct entry matching gives
\begin{equation}
 H_{f_{r,\tau}}-H_{g_r}
   =u_r(U_r\overline e)^*+\sqrt q\,U_rK_eW^*.
 \label{I:R:eq:terminal-error-decomposition}
\end{equation}
The two output ranges are orthogonal.  Also $K_e$ is bounded: the
terminal Hankel of $\tau$ is bounded by hypothesis, and that of $h$ is
bounded by its stable finite realization.  Consequently
\[
 \Rnorm{H_{f_{r,\tau}}-H_{g_r}}^2
   \le\Rnorm e_2^2+q\Rnorm{K_e}^2
   =E_{n-1}(\tau)^2+q\Rnorm{K_e}^2.
\]
This proves the upper boundary bound and attainment.  The argument
over real coefficients uses orthogonal state changes and real spaces
throughout.
\end{proof}


\begin{proposition}[A fixed optimal lift has no positive linear error]
\label{I:R:prop:terminal-fixed-lift}
For the fixed lift above, suppose $E=\Rnorm e_2>0$ and put
$M=\Rnorm{K_e}$.  Then
\begin{equation}
 \Rnorm{H_{f_{r,\tau}}-H_{g_r}}^2
 =E^2+\frac{\Rnorm{K_e\overline e}^2}{E^2}q^2+O(q^3).
 \label{I:R:eq:terminal-fixed-lift-expansion}
\end{equation}
If $E=0$, the lift equals the target and its error is zero.
\end{proposition}
\begin{proof}
Put $v=U_r\overline e$ and $R=K_e^*K_e$.  Equation
\eqref{I:R:eq:terminal-error-decomposition} gives
$H_{\rm err}^*H_{\rm err}=vv^*+qWRW^*$, while
$\Rnorm v=E$ and $W^*v=\sqrt q\,\overline e$.
In the split $\operatorname{span}\{v\}\oplus v^\perp$, the
upper-left perturbation is
$a_q=q^2\Rnorm{K_e\overline e}^2/E^2$, the off-diagonal block has
norm at most $q^{3/2}M^2$, and the lower block has norm at most $qM^2$.
For $qM^2\le E^2/2$, the Schur complement of the isolated top
eigenvalue yields
\[
 E^2+a_q\le\Rnorm{H_{\rm err}}^2
 \le E^2+a_q+2q^3M^4/E^2.
\]
This proves \eqref{I:R:eq:terminal-fixed-lift-expansion}.  Complex
coefficient conjugation in $K_e\overline e$ is essential.
\end{proof}


This proposition is an upper construction.  A $q$-dependent lift can
have a negative linear correction, as the original two-state example
demonstrates; no matching first-order lower coefficient is inferred
from the fixed lift.

\subsection{An exact mixed-terminal example}

Take $\Delta=\{b\}$, $n=2$, and let
$\tau(b^2)=\beta$, $\tau(b^3)=\gamma$, all other coefficients zero,
where $\beta,\gamma\ge0$ and $\beta^2+\gamma^2=1$.
The reduced tuple in Lemma~\ref{I:R:lem:terminal-dimension-drop} is scalar,
so phase alignment gives
\[
 a_{\tau,2}=\max_{0\le\rho\le1}
          (1-\rho^2)\rho^2(\beta+\gamma\rho)^2.
\]
For $\beta,\gamma>0$ its unique maximizing point satisfies
\[
 \beta(1-2\rho^2)+\gamma\rho(2-3\rho^2)=0.
\]
Uniqueness follows from strict concavity of the logarithm of the
positive objective on $(0,1)$.  When
$\beta=\gamma=1/\sqrt2$,
\begin{gather}
 \rho=\frac{1+\sqrt{13}}6,\qquad
 a_{\tau,2}=\frac{587+143\sqrt{13}}{2916},\notag\\
 \boxed{\displaystyle
 \lim_{r\uparrow1}\delta_2(H_{f_{r,\tau}})^2
 =\frac{2329-143\sqrt{13}}{2916}
 =0.62188140178623097774\ldots .}
 \label{I:R:eq:terminal-mixed-example}
\end{gather}
A real scalar reduced tuple attains the complex maximum here, so the
real and complex limits agree for this example.  This is a boundary
value for a nonmonomial terminal language, not an exact finite-$r$
optimization statement.

\subsection{A sharp two-state correction for every terminal degree}
\label{I:R:subsec:all-terminal-degrees}

The same complex phase calculation also gives a family of sharp
first-order results with a fixed two-state budget and arbitrary
terminal degree.  This extension changes the target family rather than
the state budget.

\begin{theorem}[All terminal degrees with two states]
\label{I:R:thm:two-state-all-degrees}
Fix an integer $d\ge2$ and let
$f_{r,d}(a^j b^d)=qr^j$, zero elsewhere.  Put
\[
 a_d=\frac{(d-1)^{d-1}}{d^d}.
\]
For each field separately,
\begin{equation}
 \boxed{\displaystyle
 \delta_2^{\mathbb F}(H_{f_{r,d}})^2
   =1-a_d-\frac{a_d}{(d-1)^2}q+O(q^2),
 \qquad\mathbb F\in\{\mathbb R,\RC\}.}
 \label{I:R:eq:two-state-all-degrees}
\end{equation}
The constants may depend on fixed $d$.  In particular $d=2$ gives
$3/4-q/4+O(q^2)$, whereas $d=3$ gives
$23/27-q/27+O(q^2)$.
\end{theorem}
\begin{proof}
We first prove the lower bound globally over complex tuples.  Write
\[
 h_{2,d}(t)=\max_{\substack{[B;\ell]^*[B;\ell]\preceq I_2\\
                       \Rnorm{[B;\ell]e_1}^2=t}}
                         |\ell B^d e_1|^2.
\]
In the exact parametrization
$[B;\ell]=[sz,W_s(z)\nu]$ put
$\mathcal J_d=\ell B^de_1/s$.  At $s=0$,
\[
 \mathcal J_d=z_2\nu_3\nu_2^{d-1},\qquad
 \max|\mathcal J_d|=M_0=lc^{d-1}=\sqrt{a_d},\qquad
 c=\sqrt{\frac{d-1}{d}},\quad l=d^{-1/2}.
\]
Equality holds precisely when $z_1=z_3=\nu_1=0$,
$|z_2|=1$, $|\nu_2|=c$, and $|\nu_3|=l$.
The exact symmetries
$B\mapsto e^{i\theta}U^*BU$, $\ell\mapsto e^{i\phi}\ell U$,
$U=\operatorname{diag}(1,e^{i\psi})$, change the phases of
$(z_2,\nu_2,\nu_3)$ by
$(\theta-\psi,\theta,\phi+\psi)$.  This invertible phase map fixes
$z_2,\nu_2,\nu_3$ positive near the maximizing torus without discarding
any complex transverse direction.  The outward
radial derivative in $\nu$ is $dM_0>0$, so maximizers have
$\Rnorm\nu=1$ for small $|s|$.

Let $\theta_0=\arccos c$.  In the resulting seven real coordinates,
take $z_1,z_3,\nu_1\in\RC$, $\xi\in\mathbb R$, and
\[
 z_2=\sqrt{1-|z_1|^2-|z_3|^2},\quad
 \nu_2=\sqrt{1-|\nu_1|^2}\cos(\theta_0+\xi),\quad
 \nu_3=\sqrt{1-|\nu_1|^2}\sin(\theta_0+\xi).
\]
The analytic objective $F_d=|\mathcal J_d|$ has the expansion
\begin{align}
 F_d(s,x)=M_0
 &-\frac{d-1}{2}M_0s^2
   +\frac{M_0}{c}s\Re z_1
   +\frac{(d-2)M_0}{c^2}s\Re\nu_1\notag\\
 &-\frac{M_0}{2}(|z_1|^2+|z_3|^2)
   -\frac{dM_0}{2}|\nu_1|^2-dM_0\xi^2
   +O((|s|+\Rnorm x)^3).
 \label{I:R:eq:degree-d-Hessian}
\end{align}
For clarity, the pure $s^2$ term follows from the exact value
$F_d(s,0)=M_0(1-s^2)^{(d-1)/2}$.  The mixed $z_1$ term is the path
with one extra waiting step in the first state.  There are $d-2$
paths producing the mixed $\nu_1$ term: one return from the second
state to the first, followed immediately by re-entry.  The angular
second logarithmic derivative of
$\cos^{d-1}\theta\sin\theta$ at $\theta_0$ is $-2d$.
These observations give every coefficient in
\eqref{I:R:eq:degree-d-Hessian}.

Its transverse Hessian is negative definite.  Compact localization,
the complete phase gauge, and the analytic implicit function theorem
therefore identify its local maximum with the global maximum for
small $|s|$.  The optimizing first derivatives are
\[
 \Re z_1=s/c+O(s^2),\qquad
 \Re\nu_1=\frac{d-2}{dc^2}s+O(s^2),
\]
with all other transverse coordinates $O(s^2)$.  The maximum modulus's
$s^2$ coefficient is
\[
 \frac{M_0}{2}\left[-(d-1)+\frac1{c^2}
                         +\frac{(d-2)^2}{dc^4}\right]
   =\frac{M_0}{2(d-1)^2}.
\]
Physical evenness removes the odd analytic coefficients.  Hence
\begin{equation}
 h_{2,d}(t)=a_dt+\frac{a_d}{(d-1)^2}t^2+O(t^3),\qquad
 P_{2,d}(q)=a_d+\frac{a_d}{(d-1)^2}q+O(q^2).
 \label{I:R:eq:two-state-all-degree-envelope}
\end{equation}
The second equality follows by scalar localization in
Lemma~\ref{I:R:lem:terminal-elimination}; its leading profile is
$4a_dx/(1+x)^2$ for $t=qx$, with a unique quadratic maximum at $x=1$.
Projection gives the claimed complex lower bound.

For completeness the matching real trial is explicit.  With
$s=\sqrt q$, set
\[
 \lambda=1/c,\quad w=\frac{d-2}{d-1},\quad
 F_s=\sqrt{1-s^2w^2},\quad
 N_s=\begin{pmatrix}sw\\cF_s\\lF_s\end{pmatrix},\quad
 z_s=\begin{pmatrix}s\lambda\\\sqrt{1-s^2\lambda^2}\\0\end{pmatrix}.
\]
For sufficiently small $s>0$, define
\[
 V_s=[sz_s,W_s(z_s)N_s]=\begin{pmatrix}B_s\\\ell_s\end{pmatrix},
 \quad v_s=\begin{pmatrix}r\\-sN_s^{\mathsf T}z_s\end{pmatrix},
 \quad A_s=e_1v_s^{\mathsf T}.
\]
Then $N_s^{\mathsf T}N_s=1$, $\Rnorm{z_s}=1$, and
$A_s^{\mathsf T}A_s+B_s^{\mathsf T}B_s+
\ell_s^{\mathsf T}\ell_s=I_2$ exactly.  Define
\[
 \omega_s=s\ell_sB_s^d(I-rA_s)^{-1},\qquad
 b_s=s^{-1}(I-rA_s)\omega_s^{\mathsf T},\qquad
 g_s(u)=\ell_sT_{\operatorname{rev}(u)}b_s.
\]
Here $T_u$ is the standard ordered product of $A_s,B_s$ along $u$.
This is a real two-state series.  The structured first-order matching
argument applies with terminal power $d$: its required cancellation is
\[
 l c^{d-1}=M_0,\qquad
 l c^d-M_0c=0.
\]
In particular
$\omega_s=(M_0+O(q),O(s^3))$, $b_s=(sM_0+O(s^3),O(q))$.
The exact energy identity and the block orders
\[
 B_s=\begin{pmatrix}O(q)&O(s)\\s+O(s^3)&c+O(q)\end{pmatrix},
 \qquad
 A_s=\begin{pmatrix}r&-sc+O(s^3)\\0&0\end{pmatrix}
\]
give joint stability for small positive $s$: the unique transition
eigenvalue near one is $1-q+O(q^2)$ and all other eigenvalues remain
strictly smaller in modulus.  The reachability Gram has slow entry
$M_0^2+O(q)$, off-diagonal $O(s^3)$ and fast entry $O(q)$.
As in the structured matching proof, interpolation on the dominant
input is exact, the error's squared dominant Rayleigh quotient is
$1-\Rnorm{\omega_s}^2$, its off-diagonal block is $O(q)$, and its
remaining block is $O(q)$.  The fixed spectral gap $1-a_d>0$ gives
\[
 \Rnorm{H_{f_{r,d}}-H_{g_s}}^2
       =1-\Rnorm{\omega_s}^2+O(q^2).
\]
The chosen $\lambda,w$ realize the quadratic maximum in
\eqref{I:R:eq:degree-d-Hessian}, so
$\Rnorm{\omega_s}^2=a_d+a_dq/(d-1)^2+O(q^2)$.
This matches the global complex lower bound by a real construction,
proving \eqref{I:R:eq:two-state-all-degrees} for each field.  Equality of
their full finite-$q$ optimum functions is not asserted.
\end{proof}


% End terminal_languages.tex

\section{Scalar input for the response examples}
\label{I:R:app:Hardy-bridge}
The examples use the single Hardy bridge proved in
Theorem~\ref{I:B:thm:hardy-bridge}; its proof is not repeated here.
The lower bound is a direct H\"older inequality, the equality for
$m\le3$ uses published low-degree extremals, and equality in general
uses the external outer-extremal theorem~\cite{outer2026}.

\section{Balancing and certified dyadic recertification}
\label{I:U:sec:dyadic-recertification}
A spectral-radius margin by itself does not supply the resolvent estimate
needed for a spectral truncation bound. Balancing remains useful, provided
the reduced candidate is recertified. The following statement separates
these two steps.

\begin{theorem}[Balancing and dyadic recertification]
\label{I:U:thm:dyadic-recertification}
\begin{enumerate}[label=(\roman*)]
\item A minimal mean-square-stable realization with positive definite
Gramians admits a balancing transformation
\[
 S^{-1}P_fS^{-*}=S^*Q_fS=\Sigma_H.
\]
A leading-block truncation whose mean-square stability has been verified
is a valid candidate of the stated dimension. Its error can be certified
by Theorem~\ref{I:U:thm:exact}; no additional spectral-tail formula is asserted
here.
\item Let $g=(T_a,b,c)$ have real algebraic coefficients and be
mean-square stable. Choose a positive definite algebraic $X$ with
$X-\sum_aT_aXT_a^*=I$. For perturbed matrices
$\widehat T_a=T_a+E_a$, define
\[
 \theta=\left\|\sum_a(T_aXE_a^*+E_aXT_a^*+E_aXE_a^*)\right\|.
\]
If a certified bound gives $\theta<1$, then
\begin{equation}
 X-\sum_a\widehat T_aX\widehat T_a^*\succeq(1-\theta)I,
 \qquad
 \rho(\Phi_{\widehat T})\le1-\frac{1-\theta}{\|X\|}<1.
 \label{I:U:eq:dyadic-verified-margin}
\end{equation}
An elementary sufficient upper bound for $\theta$ is
\[
 \sum_a\left(2\|T_aX^{1/2}\|\|E_aX^{1/2}\|
                         +\|E_aX^{1/2}\|^2\right).
\]
\item For every positive rational $\varepsilon$, an algorithm can
increase dyadic precision, round all entries of $T_a,b,c$, and return a
stable dyadic realization $\widehat g$ with
$\|H_g-H_{\widehat g}\|\le\varepsilon$. Each candidate is accepted
only after verifying \eqref{I:U:eq:dyadic-verified-margin} and a joint
Gramian error certificate. This search terminates. This statement gives
no uniform polynomial bound on the required precision or search time.
\end{enumerate}
\end{theorem}
\begin{proof}
For (i), diagonalize
$P_f^{1/2}Q_fP_f^{1/2}=U\Sigma_H^2U^*$ and take
$S=P_f^{1/2}U\Sigma_H^{-1/2}$. Direct substitution proves both
identities. Stability of a proposed truncation is an explicit condition,
not a consequence asserted by this proof.

For (ii), subtract the perturbation from the Stein equality. The
operator-norm bound gives the first inequality in
\eqref{I:U:eq:dyadic-verified-margin}. Since $X\succeq I$ and
$I\succeq X/\|X\|$, the positive map $\Phi_{\widehat T}$ contracts
the order-unit norm induced by $X$ by at most the displayed factor.
This bounds its spectral radius. Factoring each cross term through
$X^{1/2}$ proves the sufficient bound.

For (iii), choose dyadic approximations converging entrywise to the
algebraic input. Their $\theta$ tends to zero, so all sufficiently fine
approximations are stable. Form the joint error realization
\[
 B_a=\Udiag(T_a,\widehat T_a),\qquad
 \beta=(b^\top,\widehat b^\top)^\top,\qquad
 \gamma=(c^\top,-\widehat c^\top)^\top.
\]
Its Gramians $P_E,Q_E$ are obtained from the two joint Stein equations.
The block diagonal matrix $\Udiag(X,X)$ supplies a strict Lyapunov
inequality in a neighborhood of the unperturbed tuple. Thus the joint
Stein inverses and $\operatorname{tr}(P_EQ_E)$ depend continuously on
the perturbation there. At zero perturbation the error series vanishes,
so this trace is zero. Consequently it eventually becomes strictly
smaller than $\varepsilon^2$. Algebraic arithmetic can certify this
strict inequality and the Lyapunov margin, and
\[
 \|H_g-H_{\widehat g}\|^2
 \le\|H_g-H_{\widehat g}\|_{\rm HS}^2
 =\operatorname{tr}(P_EQ_E)<\varepsilon^2.
\]
Enumerating increasing precisions with these tests therefore terminates.
\end{proof}


\paragraph{Why a radius margin is not a resolvent norm.}
For $A_t=\left(\begin{smallmatrix}0&t\\0&0\end{smallmatrix}\right)$,
the positive map $\Phi(X)=A_tXA_t^*$ has $\rho(\Phi)=0$, but
\[
 (I-\Phi)^{-1}(I)=I+A_tA_t^*
       =\Udiag(1+t^2,1).
\]
Its induced operator norm is at least $1+t^2$, although the inverse
spectral-radius margin equals one. Thus a bound by
$\max(\gamma^{-1},\gamma_k^{-1})$ is not a general resolvent estimate.
This counterexample to that proof step does not by itself disprove its
separate balanced-truncation inequality. That inequality remains
unestablished here. A uniform spectral-tail guarantee and a complete efficient synthesis theorem remain separate questions.


\subsection{Quantitative truncation with two verified Stein witnesses}
\label{I:U:sec:quantitative-truncation}
The following result supplies one explicit answer to the quantitative
truncation target. Its constants are allowed to deteriorate when the
source approaches the obstruction of Theorem~\ref{I:U:thm:separation}.
It uses full matrix inequalities, not a spectral-radius margin.

\begin{theorem}[Conditional square-root bound for balanced truncation]
\label{I:U:thm:quantitative-truncation}
Let a minimal stable realization $(A_a,b,c)$ be balanced with
$P=Q=\Sigma=\Udiag(\Sigma_1,\Sigma_2)\succ0$. Retain the first $k$
coordinates and put $B_a=A_{a,11}$, $g(w)=c_1^*B_wb_1$.
Assume the tuple $B$ is independently certified mean-square stable.
Suppose positive definite matrices $X,Y$ and $0\le r_X,r_Y<1$ satisfy
\[
 \sum_a A_aXA_a^*\preceq r_X^2X,\qquad
 \sum_a A_a^*YA_a\preceq r_Y^2Y.
\]
Write $\kappa_X=\|X\|/\lambda_{\min}(X)$ and similarly for $Y$.
Then
\begin{equation}
 \|H_f-H_g\|\le\sqrt{\sigma_1\Utr\Sigma_2}
 \left(\frac{\sqrt{\kappa_X}}{1-r_X}
             +\frac{\sqrt{\kappa_Y}}{1-r_Y}\right).
 \label{I:U:eq:quantitative-truncation}
\end{equation}
In particular $\Utr\Sigma_2\le(n-k)\sigma_{k+1}$ gives an explicit
square-root rate with all other displayed parameters fixed.
\end{theorem}
\begin{proof}
Let $E:\mathbb F^k\to\mathbb F^n$ be the coordinate inclusion and
let $C,O$ and $C_B,O_B$ denote reachability and observation maps.
The leading Stein blocks show $P_B,Q_B\preceq\Sigma_1$: for example,
iterate $\Sigma_1\succeq b_1b_1^*+\sum_aB_a\Sigma_1B_a^*$ and
use stability of $B$. Hence $\|O\|=\sqrt{\sigma_1}$ and
$\|C_B\|\le\sqrt{\sigma_1}$.

For $z_w=A_wb-EB_wb_1$, the tree recursion is
\[
 z_\varepsilon=(0,b_2)^\top,\qquad
 z_{aw}=A_az_w+(0,A_{a,21}B_wb_1)^\top.
\]
On the Hilbert sum of state vectors indexed by words, with each state
norm $\|z\|_Y$, the map $(z_w)\mapsto(A_az_w)_{aw}$ has norm at
most $r_Y$. Root and nonroot forcing coordinates are disjoint. Their
squared Euclidean sum is
\[
 \|b_2\|^2+\sum_a\Utr(A_{a,21}P_BA_{a,21}^*)
 \le\Utr\Sigma_2,
\]
by the lower reachability Stein block. The Neumann series on the tree
and equivalence of the state norms therefore give
\[
 \|C-EC_B\|_{\rm HS}
 \le\frac{\sqrt{\kappa_Y\Utr\Sigma_2}}{1-r_Y}.
\]
Apply the identical argument to $(A_a^*,c)$ with witness $X$, reversing
word order, to obtain
$\|O-O_BE^*\|_{\rm HS}\le
\sqrt{\kappa_X\Utr\Sigma_2}/(1-r_X)$. Finally,
\[
 H_f-H_g=O(C-EC_B)+(O-O_BE^*)EC_B.
\]
The operator norm is at most the displayed Hilbert--Schmidt bounds
times $\sqrt{\sigma_1}$, proving the assertion.
\end{proof}

Both witnesses exist for each stable finite tuple, but their quantitative
constants need not be uniform. A witness for the full tuple does not
automatically give a witness for its leading block. For example,
\[
 A=\begin{pmatrix}2&1\\-4&-2\end{pmatrix},\qquad A^2=0,
 \qquad X=I+AA^*,\qquad X-AXA^*=I,
\]
whereas $A_{11}=2$ is unstable. This refutes the proposed inheritance
step. Balanced truncation and singular-value automata have established
antecedents \citep{sva2019}; the theorem above records the precise
conditional operator-norm estimate used here, without a novelty claim.

\begin{theorem}[A sufficient dyadic precision chosen before recertification]
\label{I:U:thm:explicit-dyadic}
Let an $n$-state tuple $(A_a,b,c)$ have witnesses as above. Choose
$r\in(0,1)$ and $\kappa\ge1$ with
$r_X,r_Y\le r$, $\kappa_X,\kappa_Y\le\kappa$, and put $a=1-r$.
Round every scalar entry to a dyadic number with absolute error at most
$h=2^{-b_0}$. For complex entries, replace these scalar-error bounds
by bounds on their complex moduli. Define
\[
 D_b=\frac{2\sqrt{\kappa n}}a+
                  \frac{2\kappa\sqrt d\,n\|b\|}{a^2},\qquad
 D_c=\frac{2\sqrt{\kappa n}}a+
                  \frac{2\kappa\sqrt d\,n\|c\|}{a^2},
\]
\[
 K= D_c\frac{\sqrt\kappa\|b\|}a
       +\frac{2\sqrt\kappa(\|c\|+1)}a D_b.
\]
The choices
\begin{equation}
 h\le\min\left\{n^{-1/2},
       \frac{a}{2\sqrt\kappa\sqrt d\,n},\frac\epsilon K\right\}
 \label{I:U:eq:explicit-dyadic}
\end{equation}
guarantee mean-square stability of the rounded tuple and
$\|H_f-H_{\widehat f}\|\le\epsilon$.
If $K=0$, the last condition is omitted. Rational inputs and rational
upper bounds for these constants give polynomial bit cost in the input
length and $b_0$ for rounding and exact joint-Gramian recertification.
\end{theorem}
\begin{proof}
In either witness norm the tree-operator perturbation has norm at most
$\sqrt\kappa\sqrt d\,nh\le a/2$. Thus both rounded tree operators
have norm at most $1-a/2$, proving stability. The resolvent identity,
$\|\delta b\|,\|\delta c\|\le\sqrt n h$, and the original and
rounded resolvent bounds $a^{-1}$ and $2a^{-1}$ give
\[
 \|\widehat C-C\|_{\rm HS}\le D_bh,\qquad
 \|\widehat O-O\|_{\rm HS}\le D_ch.
\]
Indeed, the two terms for reachability are the changed root vector
and the changed tree operator applied to the original reachability
sequence; their bounds are exactly the two terms of $D_bh$.
The dual argument gives $D_ch$. Also
$\|C\|\le\sqrt\kappa\|b\|/a$ and
$\|\widehat O\|\le2\sqrt\kappa(\|c\|+1)/a$.
Use $\widehat H-H=(\widehat O-O)C+\widehat O(\widehat C-C)$.
This proves the bound $Kh$. Rational Stein systems and their joint
error systems have polynomial dimension; exact linear algebra and
positive-semidefinite tests have polynomial bit cost. This is a
bound for specified witnesses, not for finding a well-conditioned
balancing transform or a uniformly well-conditioned truncation.
\end{proof}


\subsection{A second truncation certificate from positive Stein resolvents}
\label{I:U:sec:resolvent-truncation}
The extended proposal supplies a useful alternative to the tree estimate
above. Its positive-resolvent argument is valid after correcting the
block inequalities. This subsection retains that argument and exploits
the commuting reference Gramians; it does not supply a uniform AAK factor.

\begin{lemma}[Hermitian norm of a positive Stein resolvent]
\label{I:U:lem:positive-resolvent}
Let $\Phi(Z)=\sum_a T_aZT_a^*$ have spectral radius less than one,
and put $\mathcal R=(I-\Phi)^{-1}$ and $X=\mathcal R(I)$.
For every Hermitian $M$,
\[
 -\|M\|X\preceq\mathcal R(M)\preceq\|M\|X,
 \qquad \|\mathcal R(M)\|\le\|X\|\,\|M\|.
\]
Thus the induced spectral norm on the Hermitian subspace is exactly
$\|X\|$, with equality at $M=I$.
\end{lemma}
\begin{proof}
Apply the positive map $\sum_{j\ge0}\Phi^j$ to
$-\|M\|I\preceq M\preceq\|M\|I$. The order bounds imply the norm
bound, and $M=I$ attains it. This is the elementary Hermitian case of
the classical norm-at-identity property for positive maps
\citep{bourin2019}, not a new general operator-theory result.
\end{proof}


\begin{theorem}[Commuting-reference certificate for balanced truncation]
\label{I:U:thm:reference-truncation}
Let $(A_a,b,c)$ be minimal, mean-square stable, and balanced with
$P=Q=\Sigma=\Udiag(\Sigma_1,\Sigma_2)\succ0$, where $1\le k<n$
coordinates are retained. Set $B_a=A_{a,11}$ and
$g(w)=c_1^*B_wb_1$, and independently assume $B$ is mean-square
stable. Write $\delta=\|\Sigma_2\|=\sigma_{k+1}$ and
$\sigma_1=\|\Sigma\|$. Define the four Stein solutions
\[
 X_A-\sum_aA_aX_AA_a^*=I_n,\quad
 X_B-\sum_aB_aX_BB_a^*=I_k,
\]
\[
 Y_A-\sum_aA_a^*Y_AA_a=I_n,\quad
 Y_B-\sum_aB_a^*Y_BB_a=I_k,
\]
and $K_P=\max(\|X_A\|,\|X_B\|)$,
$K_Q=\max(\|Y_A\|,\|Y_B\|)$.
On the joint space with block sizes $(k,n-k,k)$ put
\[
 T_a=\Udiag(A_a,B_a),\quad \beta=(b_1,b_2,b_1)^\top,
 \quad\chi=(c_1,c_2,-c_1)^\top,
\]
\[
 P_0=\begin{pmatrix}\Sigma_1&0&\Sigma_1\\0&\Sigma_2&0\\
 \Sigma_1&0&\Sigma_1\end{pmatrix},\qquad
 Q_0=\begin{pmatrix}\Sigma_1&0&-\Sigma_1\\0&\Sigma_2&0\\
 -\Sigma_1&0&\Sigma_1\end{pmatrix}.
\]
Let
\[
 D_P=\beta\beta^*+\sum_aT_aP_0T_a^*-P_0,\qquad
 D_Q=\chi\chi^*+\sum_aT_a^*Q_0T_a-Q_0,
\]
and take any $e_P\ge K_P\|D_P\|$, $e_Q\ge K_Q\|D_Q\|$.
Then
\begin{equation}
 \|H_f-H_g\|^2\le
 \max\{(\delta+e_P)(\delta+e_Q),\,
                    2\sigma_1\max(e_P,e_Q)+e_Pe_Q\}.
 \label{I:U:eq:reference-truncation}
\end{equation}
In particular, define
\[
 M_P=\Big\|\sum_aA_{a,12}A_{a,12}^*\Big\|,\quad
 M_Q=\Big\|\sum_aA_{a,21}^*A_{a,21}\Big\|,
\]
\[
 c_P=K_P(2M_P+\sqrt{M_P}),\qquad
 c_Q=K_Q(2M_Q+\sqrt{M_Q}).
\]
One may take $e_P=c_P\delta$, $e_Q=c_Q\delta$, giving
\begin{equation}
 \|H_f-H_g\|^2\le
 \max\{\delta^2(1+c_P)(1+c_Q),\,
             2\sigma_1\delta\max(c_P,c_Q)+c_Pc_Q\delta^2\}.
 \label{I:U:eq:reference-square-root}
\end{equation}
This is an $O(\sqrt\delta)$ bound when the other displayed quantities
are bounded uniformly. It contains no explicit tail-dimension factor,
but its coupling and resolvent constants may depend on that dimension.
\end{theorem}
\begin{proof}
Let $P_E,Q_E$ be the exact joint error Gramians. Subtracting the
reference Stein equations gives
$P_E-P_0=(I-\Phi_T)^{-1}D_P$, and the dual identity for $Q_E$.
The identity forcings have solutions
$\Udiag(X_A,X_B)$ and $\Udiag(Y_A,Y_B)$, respectively.
Lemma~\ref{I:U:lem:positive-resolvent} therefore gives
\[
 0\preceq P_E\preceq P_0+e_PI=:\overline P,\qquad
 0\preceq Q_E\preceq Q_0+e_QI=:\overline Q.
\]
For positive semidefinite matrices the quantity
$\lambda_{\max}(P^{1/2}QP^{1/2})$ is separately monotone in $P,Q$:
use congruence for $Q$, and the equality of the nonzero spectra of
$P^{1/2}QP^{1/2}$ and $Q^{1/2}PQ^{1/2}$ for $P$.
Apply Theorem~\ref{I:U:thm:exact}. The reference matrices commute.
On the retained symmetric and antisymmetric subspaces their pairs
of eigenvalues are $(2\sigma_j,0)$ and $(0,2\sigma_j)$; on a tail
coordinate they are $(\sigma_j,\sigma_j)$.
The largest eigenvalue of $\overline P\overline Q$ is exactly the
right-hand side of~\eqref{I:U:eq:reference-truncation}.

For the explicit estimates put
$S=\sum_aA_{a,12}\Sigma_2A_{a,12}^*$ and
$U=\sum_aA_{a,22}\Sigma_2A_{a,12}^*$.
Direct multiplication of the balanced Stein equations gives
\[
 D_P=-\begin{pmatrix}0&0&S\\0&0&U\\S&U^*&S\end{pmatrix}.
\]
Here $\|S\|\le M_P\delta$, and operator Cauchy--Schwarz, together
with $\sum_aA_{a,22}\Sigma_2A_{a,22}^*\preceq\Sigma_2$, gives
$\|U\|\le\sqrt{M_P}\delta$.
Split $D_P$ into its $S$ and $U$ parts; their norms are at most
$2\|S\|$ and $\|U\|$. Hence
$\|D_P\|\le(2M_P+\sqrt{M_P})\delta$.
Apply the same calculation to $A_a^*,c$, followed by changing the
sign of the last joint block, to obtain the asserted $D_Q$ bound.
\end{proof}


\paragraph{Correct balanced block inequalities.}
The identities imply
\[
 b_2b_2^*+\sum_aA_{a,22}\Sigma_2A_{a,22}^*\preceq\Sigma_2,
 \qquad
 \Big\|\sum_aA_{a,12}\Sigma_2A_{a,12}^*\Big\|
                  \le\min\{\sigma_1,M_P\delta\}.
\]
They do not permit replacing the norm of a positive sum by the sum
of its spectral norms, nor multiplying the two independent bounds.
For a concrete balanced counterexample let
\[
 A_0=E_{12},\quad A_1=E_{23},\quad
 b=\sqrt t\,e_3,\quad c=\sqrt t\,e_1,\qquad 0<t<1,
\]
and retain $k=1$. All products of three letters vanish, and the
reachability and observation Gramians are both $tI_3$; in particular
the representation is stable and minimal. The retained tuple is zero.
Nevertheless,
\[
 \|b_2\|^2+\sum_a\|A_{a,22}\Sigma_2^{1/2}\|^2=2t>t,
 \qquad
 \Big\|\sum_aA_{a,12}\Sigma_2A_{a,12}^*\Big\|=t>t^2.
\]
The reference-Gramian theorem uses neither false strengthening.

\section{A rational full-operator upper certificate in the commuting example}\label{I:S:app:commupper}
Use $b=13/24$, $c=197/200$, and the rank-one series $g(b^n)=cb^n$, zero on all words containing $a$. To avoid overloading the letter notation, the scalar $b$ in this sentence is the coefficient of the alphabet letter $b$. Take the orthogonal basis
\[
 e_\Seps,\ e_a,\ e_b,\ e_{aa},\
 h=\sum_{n\ge2}b^ne_{b^n}.
\]
Its metric matrix is $D=\Sdiag(1,1,1,1,d_0)$, $d_0=b^4/(1-b^2)$. The kernel has coordinates $k=(1,0,b,0,1)^\top$. Put $H_5=H\oplus0$ for the four-square target in Theorem~\ref{I:S:thm:commutative}. The error bilinear-form matrix is
$L=DH_5-cDkk^\top D$. The operator vanishes on the orthogonal complement of this five-dimensional space. Thus its norm is less than $s=29/25$ exactly when both $sD+L$ and $sD-L$ are positive definite. Their leading principal minors are:
\begin{center}\small
\begin{tabular}{@{}ccc@{}}
\toprule
Order & $sD+L$ & $sD-L$\\\midrule
1 & $47/40$ & $229/200$\\
2 & $1269/500$ & $229/1250$\\
3 & $3481227/2000000$ & $1037897/4500000$\\
4 & $13775541/100000000$ & $16067927/450000000$\\
5 & $75291604330033/79511520000000000$ & $131840555454959/1073405520000000000$\\
\bottomrule
\end{tabular}
\end{center}
All are strictly positive. These entries are exact rational values, and the displayed inequalities require no numerical search.


\chapter{Probability Guarantees and Positive Realization}\label{sec:probability}

\section{From a linear representation to a probability law}
Finite Hankel rank describes the dimension of a linear realization
\citep{carlyle1971,fliess1974}. It
does not impose nonnegative transitions, almost-sure termination, or
robust predictions after conditioning on a rare prefix. These constraints
are distinct and must enter a probabilistic reduction theorem separately.
Section~\ref{sec:hankel} treats the operator approximation problem.
Here we ask when spectral information or an acquired dictionary can
certify a proper approximation of a finite-word law.

Let $\Sigma$ have $d\ge2$ letters and let $p$ be a probability law on
completed words, including the empty word. All logarithms in this section
are natural. The word-Hankel operator is $H_p[u,v]=p(uv)$ and
$t_p(T)=\Pr_p\{|W|>T\}$ for integer $T\ge0$. We assume
$p(\varepsilon)\ge\mu$ with $0<\mu<1$ when using the empty pivot in the spectral
bridge. Let $\mathcal Q_N^+$ be the proper nonnegative automata with
at most $N$ states, and write
\[
 \mathcal K(p)=\inf_{q\in\mathcal Q_1^+}D_{\rm KL}(p\|q).
\]
A one-state law consists of an independent stop/letter event distribution.
Its information projection is explicit: with $M=\mathbb E_p|W|$,
letter-frequency law $\pi$, and
$g(M)=(M+1)\log(M+1)-M\log M$,
\[
 \mathcal K(p)=g(M)+M H(\pi)-H(p)
\]
when the terms are finite, with the conventions in
Proposition~\ref{II:U:prop:iprojection}. This classical calculation fixes
the comparison class; the new issue is the sharp dependence on a small
second Hankel singular value.

\section{Sharp spectral-to-probability regimes}
\begin{theorem}[Three length regimes]
Fix $0<\mu<1$, $C>0$, $M_s>0$, $M_0>0$, and the alphabet size; use the indicated parameters in each class.
\begin{enumerate}
\item If $t_p(T)\le C\rho^{T+1}$ for all integers $T\ge0$, where
$0<\rho<1$, then
\[
 \sup_{\substack{p(\varepsilon)\ge\mu,\ t_p(T)\le C\rho^{T+1}\\
                   \sigma_2(H_p)\le\epsilon}}\mathcal K(p)
 =\Theta\bigl(\epsilon^\gamma\log^2(1/\epsilon)\bigr),\qquad
 \gamma=\frac{\log(1/\rho)}{\log(1/\rho)+\tfrac12\log d}.
\]
\item If $s>1$ and $\mathbb E|W|^s\le M_s$, the corresponding supremum is
\[
 \Theta\bigl((\log(1/\epsilon))^{1-s}\log\log(1/\epsilon)\bigr).
\]
\item Under only $\mathbb E|W|\le M_0$, that supremum equals
$g(M_0)+M_0\log d$ for every $\epsilon>0$, but equals zero at
$\epsilon=0$.
\end{enumerate}
The asymptotic bounds use sufficiently small
$0<\epsilon\le\mu/2$ and integer truncation lengths. Constants depend
on the fixed class parameters; they are not uniform as $s\downarrow1$.
\end{theorem}
The exact statements and proofs are
Theorems~\ref{II:U:thm:kl-exp}, \ref{II:U:thm:kl-moment} and
\ref{II:U:thm:kl-mean}.

The proof first converts the rank-one spectral defect into a proper law
$q$ through the nonvanishing empty pivot. For every integer $T\ge0$,
\[
 \operatorname{TV}(p,q)
 \le\frac{dC_\mu}{\mu}\epsilon\sqrt{\sum_{j=0}^Td^j}
       +\frac{t_p(T)}{\mu^2},\qquad
 C_\mu=(1+1/\mu)(1+2/\mu).
\]
The first term is a finite-head Cauchy--Schwarz estimate; the second uses
positivity and the source tail. Length-tail information then controls
entropy and the cross-entropy loss after smoothing the event law.
Explicit integer choices of $T$ and $0<\Delta\le1/2$ are given in the
proofs. Smoothing is a step in the bound, not a claim that an arbitrary
spectral truncation is positive.

Sharp lower bounds use laws concentrating a small probability on a long
length shell. For
$p_T=(1-\delta)\delta_\varepsilon+\delta\operatorname{Unif}(\Sigma^T)$,
$T\ge2$, the shell-basis singular decomposition gives
$\sigma_2(H_{p_T})=\delta d^{-T/2}$ exactly, while
$\mathcal K(p_T)=g(\delta T)-h_2(\delta)$.
The choices of $\delta$ corresponding to exponential and moment tails
produce the stated rates. Under a mean bound, a single long word with
balanced letter counts supplies the discontinuity. A separate argument
in Theorem~\ref{II:U:thm:kl-fixed-states} shows that allowing any fixed
number of positive states does not restore a universal spectral-to-KL
continuity theorem under the hypotheses stated there.

\section{A constructive certificate from a supplied dictionary}
Let $f_1,\ldots,f_N$ be probability laws, let $\alpha$ be a probability
row, and assume explicitly
\[
 p=\sum_i\alpha_i f_i.
\]
For nonnegative occupation variables $x_i,z_i,y_{aij}$, impose
\[
 x_j=\alpha_j+\sum_{a,i}y_{aij},\qquad
 x_i=z_i+\sum_{a,j}y_{aij},\qquad \sum_i x_i\le V.
\]
Define the residuals
$R_i=x_if_i-z_i\delta_\varepsilon-\sum_{a,j}y_{aij}L_af_j$.
The local defects must be weighted by how often the decoded process
visits the corresponding states.

\begin{theorem}[Occupation certificate]
For $x_i>0$, decode $T_a(i,j)=y_{aij}/x_i$ and
$\tau_i=z_i/x_i$; use an immediate stop for zero-occupation rows.
The resulting law $q$ started from $\alpha$ is proper and satisfies
\[
 \mathbb E_q|W|\le V-1,\qquad
 \operatorname{TV}(p,q)\le\frac12\sum_i\|R_i\|_1.
\]
If the source mixture is only accurate to $\epsilon_0$ in $\ell^1$,
add $\epsilon_0/2$ to the bound.
\end{theorem}
Theorem~\ref{II:U:thm:flows} proves properness and the telescoping residual
estimate. The budget $V$ controls expected visits; an unweighted maximum
local defect cannot replace the occupation-weighted total in general.

The finite-sample version fixes the dictionary before validation, or uses
an independent split. For each dictionary law, empirical bins up to an
integer horizon $L$ and one overflow bin give a simultaneous confidence
event. On this event a finite LP may optimize its flows using the same
histograms, since the event is uniform over feasible flows. The bound
$\operatorname{TV}(p,\widehat q)\le\widehat C_L/2$ certifies an actually
returned feasible model. This is a high-probability statement under the
sampling model, not an unconditional certificate from arbitrary traces.
The complete oracle inequality is Theorem~\ref{II:U:thm:flow-lp}.

\section{What positive acquisition requires}
The preceding theorem starts with a dictionary and its source-mixture
relation. Discovering a small suitable dictionary is a different problem.
The conditional spectral--conic bridge explicitly combines prefix-mass
bounds, spectral conditioning, a feasible visit budget, and a conic defect.
Its bound charges the whole product of these parameters; a small local
closure defect alone is insufficient.

The constructive results below separate several access models:
a supplied finite candidate pool; positive closure in dual state
coordinates; rank-two monotone endpoint acquisition; observable-coordinate
rounding; relatively separated prefix pools; and bounded-depth promise
classes. Each theorem states its own conditioning, positivity, sampling
and feasibility assumptions. None supplies unrestricted discovery from
uninterpreted observations.

A two-state proper law already prevents replacing latent positive states
by an arbitrary finite set of actual prefix residuals with exact closure.
The golden-ratio example makes this obstruction quantitative. Additional
counterexamples show that uniform Doeblin positivity does not imply
arbitrary hidden-state purification and that changes to a rare-prefix
posterior invalidate naive conditional smoothing bounds. These examples
are part of the theorem scope, rather than failures of an implementation.

\section{Word laws, conditional futures and filtering}
A total-variation guarantee for complete words does not directly become
a uniform next-event guarantee after every prefix. Conditioning divides
by prefix mass. Likewise, a projective contraction statement for a filter
state and a KL statement for the law of an entire word are different
claims. Projective contraction and stability of nonlinear filters have
a substantial prior literature~\citep{cohen,atar}. The final proofs in this section give the sharp projective-to-KL
envelope, finite-horizon word-law estimates, and the precise columnwise
limit hypothesis used for word laws. Intermittent contraction controls the
filter state under its own assumptions.

The open acquisition problem is to find, with quantified sample and bit
costs, a small dictionary satisfying the hypotheses that make the entire
certificate useful. A successful theorem must account for rare-prefix
sampling, adaptive dictionary selection and numerical feasibility. The
transfer-operator constructions in Section~\ref{sec:memory} assume
a supplied dynamical model.

\section{Reading guide}
The proofs below use the linear preliminaries and attainment result in
Subsections~\ref{I:U:sec:preliminaries} and~\ref{I:U:sec:realization}.
They begin with the three KL regimes and the fixed-state obstruction.
The occupation and positive-realization section contains all constructive
and conditional acquisition results. The final sections give projective
bounds and exact counterexamples. Numerical agreement is not used to
establish global optimality or bibliographic priority.
\label{II:II:main-end}

\section*{Detailed proofs and supporting results}

\section{Probability guarantees: rank one and fixed state budgets}\label{II:U:sec:probability}
This section concerns probability distributions on \emph{completed finite
words}. All logarithms are natural. Let $p$ be such a law, $d\ge2$, and
$p(\varepsilon)\ge\mu>0$. Define
\[
 \mathcal Q_1=\left\{q_{a,b}(w)=a\prod_{i=1}^{|w|}b_{w_i}:
 a>0,\ b_x\ge0,\ a+\sum_xb_x=1\right\},\quad
 \mathcal K(p)=\inf_{q\in\mathcal Q_1}D_{\rm KL}(p\|q).
\]
An impulse word-Hankel is a finite partial permutation of norm one.
Consequently $H_p$ is a norm-convergent sum of finite-rank operators and is
compact with $\|H_p\|\le1$. A nonzero normalized nonnegative rank-one
word-Hankel has a positive empty pivot and multiplicative coefficients;
its law is in $\mathcal Q_1$.

\subsection{The information projection is explicit}
Let $M=\mathbb E_p|W|$, $c_x=\mathbb E_p\#_x(W)$, and
$g(M)=(1+M)\log(1+M)-M\log M$, with $g(0)=0$.
\begin{proposition}[Classical one-state information projection]\label{II:U:prop:iprojection}
If $M<\infty$, the optimum is attained at
\[
 a_*=(1+M)^{-1},\qquad b_{*,x}=c_x/(1+M),
\]
and, for $M>0$ and $\pi_x=c_x/M$,
\[
 \mathcal K(p)=g(M)+M H(\pi)-H(p).
\]
The value is zero when $M=0$.
\end{proposition}
\begin{proof}
Cross-entropy is $-\log a-\sum_xc_x\log b_x$ on the event simplex.
The entropy maximization or Lagrange multiplier identity normalizes the
counts $(1,c_x)$ to obtain the displayed parameters. A zero $c_x$ means
that letter is absent almost surely, so it causes no support problem.
Finite mean implies $H(p)\le g(M)+M\log d<\infty$: first maximize
entropy at fixed length, then use the geometric maximum-entropy law at
fixed mean length. Subtracting $H(p)$ gives the formula.
\end{proof}

This proposition is an ingredient, not a novelty claim. The substantive
question is how small $\sigma_2(H_p)$ forces $\mathcal K(p)$ to be.

\subsection{From a spectral tail to a proper approximation}
Write $t_p(T)=\Pr_p\{|W|>T\}$ and $N_d(T)=\sum_{j=0}^Td^j$.
\begin{lemma}[Spectral-to-TV bridge]\label{II:U:lem:spectral-tv}
If $\epsilon=\sigma_2(H_p)\le\mu/2$, there is a proper
$q\in\mathcal Q_1$ such that, for every integer $T\ge0$,
\[
 \UTV(p,q)\le \frac{dC_\mu}{\mu}\epsilon\sqrt{N_d(T)}
                    +\frac{t_p(T)}{\mu^2},\qquad
 C_\mu=(1+1/\mu)(1+2/\mu).
\]
It can be chosen with stop probability $p(\varepsilon)$, hence with
$\Pr_q\{|W|>T\}\le(1-\mu)^{T+1}$.
\end{lemma}
\begin{proof}
Set $a=p(\varepsilon)$ and $c_x=p(x)/a$. Take a leading rank-one SVD
term $H_p=su v^*+E$ with $\|E\|=\epsilon$ and $s\le1$.
The empty pivot implies $|v_\varepsilon|\ge a-\epsilon\ge\mu/2$.
For $z_x=e_x-c_xe_\varepsilon$, the row
$h_x(v')=p(xv')-c_xp(v')$ satisfies $z_x^*H_pe_\varepsilon=0$.
Eliminating the rank-one term using this equality gives
$\|h_x\|_2\le\epsilon\|z_x\|(1+|v_\varepsilon|^{-1})
\le C_\mu\epsilon$.

On words of length at most $T$, sum Cauchy--Schwarz over the $d$ rows.
On longer words, positivity gives
\[
 \sum_x\sum_{|v'|>T}|h_x(v')|
 \le t_p(T+1)+\left(\sum_xc_x\right)t_p(T)
 \le t_p(T)/a.
\]
The last inequality uses $\sum_xp(x)\le1-a$. Put
$b_x=\sum_{v'}p(xv')$, so $\sum b_x=1-a$ and
$b_x-c_x=\sum h_x$. The defect in
$p=a\delta_\varepsilon+Kp$, where $(Kf)(xv')=b_xf(v')$, has
$\ell^1$ norm at most $2\sum_x\|h_x\|_1$.
The proper law $q=a\prod b_{w_i}$ solves that recurrence exactly.
Since $\|K\|_{1\to1}=1-a$, its inverse resolvent has norm at most
$1/a$. Divide the resulting $\ell^1$ bound by two and use $a\ge\mu$.
\end{proof}


\subsection{Three sharp regimes}
For fixed $C>0$, $0<\rho<1$, and $0<\mu<1$, let
\[
 \mathcal P_{\exp}=\{p:p(\varepsilon)\ge\mu,
       \ t_p(t)\le C\rho^{t+1}\text{ for every integer }t\ge0\},
 \quad \gamma=\frac{\log(1/\rho)}{\log(1/\rho)+\tfrac12\log d}.
\]
\begin{theorem}[Sharp exponential-tail rate]\label{II:U:thm:kl-exp}
As $\epsilon\downarrow0$,
\[
 \sup_{\substack{p\in\mathcal P_{\exp}\\\sigma_2(H_p)\le\epsilon}}
 \mathcal K(p)
 =\Theta\!\left(\epsilon^\gamma\log^2(1/\epsilon)\right).
\]
The asymptotic range may be chosen with $0<\epsilon\le\mu/2$ and with the TV upper bound $0<\Delta\le1/2$ used below; all truncation horizons are integers. The constants depend on $d,\mu,C,\rho$. The claim is for the best
normalized one-state law, not a particular spectral estimator.
\end{theorem}
\begin{proof}[Upper bound]
Restrict to sufficiently small $0<\epsilon\le\mu/2$ and put
$\lambda=\log(1/\rho)$ and
$T=\lfloor\log(1/\epsilon)/(\lambda+\tfrac12\log d)\rfloor\in\mathbb Z_{\ge0}$ in
Lemma~\ref{II:U:lem:spectral-tv}. Both terms are $O(\epsilon^\gamma)$.
Let $\Delta\le1/2$ be a positive such upper bound. Smooth the event
law $\pi=(a,b_1,\ldots,b_d)$ to
$\pi_\Delta=(1-\Delta)\pi+\Delta\operatorname{Unif}(d+1)$,
and call the resulting proper law $q_\Delta$.

Couple $W\sim p,V\sim q$ maximally, enlarging the mismatch event $E$
to probability $\Delta$ if necessary. Exponential tails imply
\[
 \mathbb E[|W|1_E]\le\Delta A_p(\Delta),\quad
 \mathbb E[|V|1_E]\le\Delta A_q(\Delta),\quad A_p,A_q=O(\log(1/\Delta)).
\]
For example, summing $\min(\Delta,C\rho^{t+1})$ proves the first
bound. With $A=\max(A_p,A_q)$, conditional entropy and the geometric
maximum-entropy bound imply
\[
 |H(p)-H(q)|\le h_2(\Delta)+\Delta\{g(A)+A\log d\}.
\]
All smoothed event probabilities are at least $\Delta/(d+1)$. The stop
term cancels when comparing two word cross-entropies, giving a difference
at most $\Delta(A_p+A_q)\log((d+1)/\Delta)$.
Finally the event-by-event renewal calculation yields
\[
 D_{\rm KL}(q\|q_\Delta)=a^{-1}D_{\rm KL}(\pi\|\pi_\Delta)
 \le-\mu^{-1}\log(1-\Delta).
\]
Adding the three estimates gives $D_{\rm KL}(p\|q_\Delta)
=O(\Delta\log^2(1/\Delta))$. Optimization can only reduce it.
\end{proof}

\begin{proof}[Lower bound]
For an integer $T\ge2$ take
$p_T=(1-\delta)\delta_\varepsilon+\delta\operatorname{Unif}(\Sigma^T)$.
In the orthonormal uniform shell basis of lengths $0,\ldots,T$, its
Hankel is $bJ+(1-\delta)e_0e_0^*$, where $J$ reverses the coordinates
and $b=\delta d^{-T/2}$. All orthogonal shell directions vanish.
The endpoint block is
$\left(\begin{smallmatrix}1-\delta&b\\b&0\end{smallmatrix}\right)$;
one singular value exceeds $b$, one is smaller, and the $T-1$ interior
singular values equal $b$. Hence $\sigma_2=b$ exactly.
Proposition~\ref{II:U:prop:iprojection} gives
\[
 \mathcal K(p_T)=g(\delta T)-h_2(\delta).
\]
Choose $\delta=c_0\rho^T$, $c_0=\min(C,1)$. The tail constraint holds,
and $p_T(\varepsilon)\ge\mu$ for large $T$.
Expanding at zero gives
$g(\delta T)-h_2(\delta)\sim\lambda c_0\rho^TT^2$.
Since $\epsilon_T=c_0(\rho/\sqrt d)^T$, this is the required order.
Successive $\epsilon_T$ have a fixed ratio, so selecting the first one
below a given tolerance extends the bound from a subsequence to all small
$\epsilon$.
\end{proof}


\begin{theorem}[Only a mean bound: exact discontinuity]\label{II:U:thm:kl-mean}
Fix $M_0>0$ and $0<\mu<1$. For every $\epsilon>0$,
\[
 \sup_{\substack{p(\varepsilon)\ge\mu,\ \mathbb E|W|\le M_0\\
                      \sigma_2(H_p)\le\epsilon}}\mathcal K(p)
 =g(M_0)+M_0\log d.
\]
At $\epsilon=0$ this supremum is zero.
\end{theorem}
\begin{proof}
The information-projection formula, $H(p)\ge0$, $H(\pi)\le\log d$,
and monotonicity of $g$ give the upper bound. Let $T$ run through
positive integer multiples of $d$ with $T\ge M_0/(1-\mu)$, let $w_T$ have equal symbol counts, and put
$p_T=(1-M_0/T)\delta_\varepsilon+(M_0/T)\delta_{w_T}$.
Its mean is $M_0$, its entropy is $h_2(M_0/T)$, and its expected-count
law is uniform. Removing the rank-one empty impulse shows
$\sigma_2(H_{p_T})\le M_0/T$, since an impulse Hankel has norm one.
Thus its optimal KL tends to the upper bound while satisfying any
positive spectral tolerance. If $\sigma_2=0$, positivity of the empty
pivot makes $p$ itself a proper multiplicative rank-one law.
\end{proof}


\begin{theorem}[Higher length moments]\label{II:U:thm:kl-moment}
Fix $s>1$, $M_s>0$, and $0<\mu<1$. Then
\[
 \sup_{\substack{p(\varepsilon)\ge\mu,\ \mathbb E|W|^s\le M_s\\
                      \sigma_2(H_p)\le\epsilon}}\mathcal K(p)
 =\Theta\!\left((\log(1/\epsilon))^{1-s}
                     \log\log(1/\epsilon)\right).
\]
This is an $\epsilon\downarrow0$ assertion; take $0<\epsilon\le\mu/2$ sufficiently small that the TV bound $\Delta$ in the proof is at most $1/2$. All truncation and shell lengths are integers. The constants are not uniform as $s\downarrow1$.
\end{theorem}
\begin{proof}
Markov's inequality and Lemma~\ref{II:U:lem:spectral-tv}, with
$T=\lfloor\log(1/\epsilon)/\log d\rfloor$, give a proper $q$ with
$\UTV(p,q)\le\Delta=O((\log(1/\epsilon))^{-s})$.
Its length is geometric with stop at least $\mu$, so its $s$th moment
is uniformly bounded. In the preceding coupling, H\"older gives
$\mathbb E[(|W|+|V|)1_E]=O(\Delta^{1-1/s})$.
The same entropy and smoothed cross-entropy argument gives KL
$O(\Delta^{1-1/s}\log(1/\Delta))$.

For the lower bound use the uniform shell source with $\delta=M_s/T^s$
and large integer $T$ with $M_s/T^s\le1-\mu$. Its $s$th moment is exactly $M_s$,
$\epsilon_T=M_sT^{-s}d^{-T/2}$, and
\[
 \mathcal K(p_T)=g(M_sT^{1-s})-h_2(M_sT^{-s})
 \sim M_s(s-1)T^{1-s}\log T.
\]
Here $\log(1/\epsilon_T)\sim T\log d/2$; consecutive spectral values
have ratios bounded away from zero. The asserted order follows.
\end{proof}


\subsection{Why the probability theorem does not contradict the obstruction}
The binary family in Subsection~\ref{I:U:sec:obstruction} is not the fixed-tail
probability class above. Here rank-one positivity and the nonvanishing empty
pivot give a special closure recurrence; the tail condition converts an
$\ell^2$ defect to $\ell^1$ and controls entropy. None of those steps gives a
general rank-$k$ proper realization with $k$ positive states. Nor does a
small completed-word TV guarantee uniform next-token error after every
prefix: conditioning divides by the prefix probability. Those are separate
positive-realization and conditional-prediction questions.
\subsection{A KL discontinuity at every fixed positive-state budget}
\label{II:U:sec:kl-fixed-states}
Let $\mathcal Q_N^+$ be the laws generated by proper probabilistic
automata with at most $N$ states. Concretely,
$q(w)=\alpha T_w\tau$, where $\alpha\ge0$, $\alpha\mathbf1=1$,
$T_a\ge0$, $\tau\ge0$, and
$\tau+\sum_aT_a\mathbf1=\mathbf1$; termination is almost sure from
$\alpha$. Zero transition weights and unused states are permitted.
Write $\mathcal K_N(p)=\inf_{q\in\mathcal Q_N^+}D_{\rm KL}(p\|q)$,
so that $\mathcal K_1=\mathcal K$.

\begin{theorem}[Mean-constrained KL discontinuity for fixed $N\ge2$]
\label{II:U:thm:kl-fixed-states}
Fix $d\ge2$, $M_0>0$, and $\mu\in(0,1)$. For an integer $m\ge1$,
put $K=d^m$ and $L=mK$, and concatenate the $K$ distinct words of
$\Sigma^m$, each exactly once and in any order, to form $w_L$.
Whenever $L>M_0/(1-\mu)$, define
\[
 \delta_L=\frac{M_0}{L},\qquad
 p_L=(1-\delta_L)\delta_\varepsilon+\delta_L\delta_{w_L}.
\]
Then $\mathbb E_{p_L}|W|=M_0$, $p_L(\varepsilon)\ge\mu$, and
$\sigma_{k+1}(H_{p_L})\le\sigma_2(H_{p_L})\le M_0/L$ for every
$k\ge1$. For every integer $N\ge2$,
\begin{align}
 M_0\log d-\frac{M_0\log N}{m}-h_2(\delta_L)
 &\le\mathcal K_N(p_L)\notag\\
 &\le M_0\log d+\frac{M_0}{L}
 \left[\log L-(L-1)\log(1-1/L)\right],
 \label{II:U:eq:kl-fixed-sandwich}
\end{align}
where $h_2(s)=-s\log s-(1-s)\log(1-s)$. Consequently, for each
\emph{fixed} $N\ge2$,
\[
 \lim_{m\to\infty}\mathcal K_N(p_L)=M_0\log d>0.
\]
For the same family, $\mathcal K_1(p_L)\to g(M_0)+M_0\log d$,
whereas $L+1$ states realize $p_L$ exactly.
\end{theorem}
\begin{proof}
The mass and mean assertions are immediate. The impulse Hankel
$H_{\delta_{w_L}}$ is a partial permutation of norm one. Subtracting
the rank-one operator $(1-\delta_L)H_{\delta_\varepsilon}$ therefore
gives $\sigma_2(H_{p_L})\le\delta_L$.

For any $q\in\mathcal Q_N^+$, pad its realization with unused states
if necessary, and let $T=\sum_aT_a$. Since $T\mathbf1\le\mathbf1$,
also $T^m\mathbf1\le\mathbf1$. For $B\in\Sigma^m$ put
$\psi(B)=\|T_B\|_{\infty\to\infty}
=\max_i(T_B\mathbf1)_i$. Nonnegativity yields
\[
 \sum_{B\in\Sigma^m}\psi(B)
 \le\sum_i\sum_{B\in\Sigma^m}(T_B\mathbf1)_i
 =\sum_i(T^m\mathbf1)_i\le N.
\]
If $w_L=B_1\cdots B_K$, submultiplicativity, $\|\alpha\|_1=1$,
and $\|\tau\|_\infty\le1$ imply
\[
 q(w_L)\le\prod_{r=1}^K\psi(B_r)
 \le\left(\frac NK\right)^K.
\]
The last inequality is arithmetic--geometric mean; it remains valid
when a factor vanishes. Thus
$-\log q(w_L)\ge L\log d-(L/m)\log N$.
Using $q(\varepsilon)\le1$ gives
\[
 D_{\rm KL}(p_L\|q)
 \ge-h_2(\delta_L)-\delta_L\log q(w_L),
\]
which proves the lower bound, including the case of infinite KL.

For the upper bound use the following two-state automaton:
\[
 \alpha=(1,0),\qquad
 \tau=\begin{pmatrix}1-M_0/L\\1/L\end{pmatrix},\qquad
 T_a=\begin{pmatrix}
 0&M_0/(Ld)\\0&(1-1/L)/d
 \end{pmatrix}\quad(a\in\Sigma).
\]
It is proper, has mean length $(M_0/L)L=M_0$, and satisfies
\[
 q_L(\varepsilon)=1-M_0/L,\qquad
 q_L(w)=\frac{M_0}{L^2}(1-1/L)^{L-1}d^{-L}
 \quad(|w|=L).
\]
Its divergence is exactly the upper expression in
\eqref{II:U:eq:kl-fixed-sandwich}. The two residual terms in that expression
vanish, as do $m^{-1}\log N$ and $h_2(\delta_L)$ for fixed $N$.

Each letter occurs $md^{m-1}=L/d$ times in $w_L$. Hence
Proposition~\ref{II:U:prop:iprojection} gives
$\mathcal K_1(p_L)=g(M_0)+M_0\log d-h_2(\delta_L)$.
Finally, a deterministic path of length $L$, with initial stop
probability $1-\delta_L$, initial continuation probability $\delta_L$,
and terminal stop probability one, realizes $p_L$ with $L+1$ states.
\end{proof}


\paragraph{Quantifiers and comparison classes.}
For every fixed $N\ge2$ and every positive spectral tolerance, the
mean-constrained worst-case value of $\mathcal K_N$ is at least
$M_0\log d$. At tolerance $\sigma_2(H_p)=0$, it is zero because
$p\in\mathcal Q_1\subset\mathcal Q_N^+$. The theorem identifies the
exact limiting value on the displayed family, not the exact worst-case
supremum over all sources at a positive tolerance. It also does not assert
that every source with small spectral tail has large KL error. The
quantifier that $N$ is fixed before $m\to\infty$ is essential.
Replacing $\delta_{w_L}$ by the uniform law on $\Sigma^L$ changes the
example: the same two-state $q_L$ then has divergence
$\frac{M_0}{L}[\log L-(L-1)\log(1-1/L)]\to0$.

\section{Positive realization through occupation flows}\label{II:U:sec:flows}
The second probability bridge starts from a finite dictionary rather than a
singular value. Let $f_1,\ldots,f_N$ be probability laws on finite words,
$\alpha$ a probability row vector, and initially assume
$p=\sum_i\alpha_i f_i$. Let $(L_af)(av)=f(v)$ and zero elsewhere.
For a visit budget $V\ge1$, choose nonnegative $x_i,z_i,y_{aij}$ with
\begin{equation}\label{II:U:eq:flow}
 x_j=\alpha_j+\sum_{a,i}y_{aij},\qquad
 x_i=z_i+\sum_{a,j}y_{aij},\qquad \sum_i x_i\le V.
\end{equation}
These imply $\sum z_i=1$ and $x\ge\alpha$. Define
\[
 R_i=x_if_i-z_i\delta_\varepsilon-\sum_{a,j}y_{aij}L_af_j,
 \qquad J_\infty=\sum_i\|R_i\|_1.
\]

\begin{theorem}[Properness and the occupation-weighted certificate]\label{II:U:thm:flows}
For $x_i>0$ set $T_a(i,j)=y_{aij}/x_i$ and $\tau_i=z_i/x_i$.
For $x_i=0$, take a zero transition row and $\tau_i=1$.
The automaton started at $\alpha$ terminates almost surely and defines
a normalized law $q$ with
\[
 \mathbb E_q|W|\le V-1,\qquad \UTV(p,q)\le J_\infty/2.
\]
If $\|p-\sum_i\alpha_if_i\|_1\le\epsilon_0$, add
$\epsilon_0/2$ to the TV bound.
\end{theorem}
\begin{proof}
Let $T=\sum_aT_a$. The constraints give $\tau+T\mathbf1=\mathbf1$
and $x=\alpha+xT$. Iteration shows
$\nu=\sum_{t\ge0}\alpha T^t\le x$. Survival probabilities
$\alpha T^t\mathbf1$ have finite sum, hence tend to zero. Thus
termination holds from $\alpha$ and
$1+\mathbb E|W|=\nu\mathbf1\le V$.

For $x_i>0$ write $r_i=R_i/x_i$, and let $q_i$ be the possibly
subprobability law starting in state $i$. The two recurrences imply
\[
 f_i-q_i=r_i+\sum_{a,j}T_a(i,j)L_a(f_j-q_j).
\]
After $t$ expansions and weighting by $\alpha$, the remainder has norm
at most $2\alpha T^t\mathbf1\to0$. Prepending is an $\ell^1$
isometry, so
$\|p-q\|_1\le\sum_i\nu_i\|r_i\|_1\le J_\infty$.
A zero-visit state has no incoming flow and $\alpha_i=0$, hence makes
no contribution. Unreachable circulation in $x-\nu$ is harmless.
The misspecification term follows by the triangle inequality.
\end{proof}


\paragraph{Sharpness of the visit factor.}
Take a state looping on $c$ with probability $\beta<1$, exiting on $a,b$
to an immediate-stop state with probabilities $\gamma$ each, and otherwise
stopping; let $0<2\gamma\le1-\beta$. Change the two source exit
probabilities to $\gamma+\eta,\gamma-\eta$, $0<\eta\le\gamma$.
Using the source future law and the terminal law as dictionary gives
$r_1=\eta(\delta_a-\delta_b)$, $r_2=0$, and
$\nu_1=(1-\beta)^{-1}$. The discrepancies occur on disjoint pairs
$c^na,c^nb$, so
$\UTV(p,q)=\eta/(1-\beta)=\nu_1\|r_1\|_1/2$.
The factor cannot be replaced by a uniform unweighted local error.

\subsection{A finite-sample linear program}\label{II:U:sec:flow-lp}
Fix $L\ge1$. For each dictionary law take $n_i$ independent samples and
record words up to length $L$ plus an overflow category. The dictionary
must be fixed before these samples, or chosen on an independent split.
Independence between the sample sets for different dictionary entries is
not needed. Let
\[
 K_L=1+\sum_{j=0}^Ld^j,\qquad
 b_i=\min\left\{2,\sqrt{\frac{2(K_L\log2+\log(N/\delta))}{n_i}}\right\}.
\]
A union bound over all sign vectors and all dictionary entries, with
Hoeffding's inequality, gives simultaneous histogram $\ell^1$ errors
at most $b_i$ with probability at least $1-\delta$.

Use hats for empirical probabilities and tails. For feasible flows put
\begin{align*}
 \widehat H_L&=\sum_i\sum_{|w|\le L}
  \left|x_i\widehat f_i(w)-z_i1_{w=\varepsilon}
             -\sum_{a,j}y_{aij}(L_a\widehat f_j)(w)\right|,\\
 \widehat U_L&=\sum_i x_i\widehat t_i(L)
                 +\sum_{a,i,j}y_{aij}\widehat t_j(L-1),\\
 S(x)&=\sum_i x_ib_i+\sum_{a,i,j}y_{aij}b_j
       =\sum_i(2x_i-\alpha_i)b_i,\\
 \widehat C_L&=\widehat H_L+\widehat U_L+S(x).
\end{align*}
Minimizing $\widehat C_L$ subject to \eqref{II:U:eq:flow} is a finite LP,
using the usual epigraph variables for absolute values.

\begin{theorem}[Uniform data-dependent certificate and defect oracle]\label{II:U:thm:flow-lp}
Assume $f_i$ are probability laws, $\alpha$ is a probability vector and $p=\sum_i\alpha_i f_i$. Fix the dictionary independently of the validation data and use the sampling scheme above. On one event of probability at least $1-\delta$, every feasible decoded
law satisfies $\UTV(p,q)\le\widehat C_L/2$. In particular the LP
solution $\widehat q$ satisfies
\[
 \UTV(p,\widehat q)\le
 \inf_{\text{feasible flows}}\left\{\tfrac12J_\infty
                     +\sum_i x_it_i(L)+S(x)\right\}.
\]
If $t_i(L)\le t_{\max}(L)$ and $b_i\le b$, this is at most
$\tfrac12\inf J_\infty+Vt_{\max}(L)+(2V-1)b$.
\end{theorem}
\begin{proof}
Define $J_L=H_L+U_L$ from the true probabilities. Tail majorization
gives $J_\infty\le J_L$. Head bins and the overflow bin are disjoint;
each shifted suffix histogram is a coarsening of the same sampled
histogram. Consequently
$|J_L-\widehat H_L-\widehat U_L|\le S(x)$ for \emph{all} feasible
flows on the simultaneous event. Theorem~\ref{II:U:thm:flows} proves the
first assertion, including for fitted flows chosen from these data.

For a comparison flow, optimality implies
$\widehat C_L^*\le J_L+2S(x)$. For nonnegative tail vectors $u,v$,
$\|u\|_1+\|v\|_1-\|u-v\|_1\le2\|u\|_1$.
Applying this with the source tail gives
$J_L\le J_\infty+2\sum_i x_it_i(L)$. Divide by two and infimize.
Finally $\sum x_i\le V$ and $\sum\alpha_i=1$ bound $S(x)$.
\end{proof}


The oracle compares occupation-weighted \emph{closure defects in the supplied
dictionary}, not the best arbitrary $N$-state TV or KL approximation.
Cancellations can make actual error smaller than the certificate. An
additive LP objective gap $\zeta$ costs at most $\zeta/2$ in TV if the
returned flow is feasible. Numerical feasibility itself must be checked or
repaired. The LP has size $O(dN^2+NK_L)$ up to constants; $K_L$ is
exponential in $L$. Under uniform exponential dictionary tails and
$n_i\ge n$, balancing terms gives a rate of order
\[
 Vn^{-\lambda/(2\lambda+\log d)}
       +V\sqrt{\log(N/\delta)/n},\qquad\lambda=\log(1/\rho),
\]
apart from the approximation defect and fixed constants. This is not a
dimension-free or minimax-optimal claim.

\subsection{A safe KL extension with an explicitly charged extra state}\label{II:U:sec:flow-smoothing}
For any proper $q$ with $\UTV(p,q)\le\Delta<1$, set
$u(w)=2^{-|w|-1}d^{-|w|}$ and
$q^+=(1-\Delta)q+\Delta u$. This mixture can be realized by adding
one probabilistic state and modifying the initial mixture. Define
$\Psi_p(\Delta)=\sum_{t\ge0}\min(\Delta,t_p(t))$.
\begin{proposition}[Smoothing a completed-word TV approximation]
For finite $\Psi_p(\Delta)$,
\[
 D_{\rm KL}(p\|q^+)\le-\log(1-\Delta)+\Delta\log2
                            +\log(2d)\Psi_p(\Delta).
\]
\end{proposition}
\begin{proof}
Split a common submeasure from $p$ and $q$, of mass $1-\Delta$,
enlarging the excess if their actual TV is smaller. Let the remaining
submeasure of $p$ be $p_1$ of mass $\Delta$ and the common one $p_0$.
The denominator $q^+$ dominates $(1-\Delta)p_0+\Delta u$.
The log-sum inequality bounds the common part by
$-(1-\Delta)\log(1-\Delta)$ and the excess by
$\Delta D_{\rm KL}(p_1/\Delta\|u)$.
Dropping the nonnegative entropy of $p_1/\Delta$ bounds the latter by
$\Delta\log2+\log(2d)\sum_w|w|p_1(w)$.
The last sum is at most $\Psi_p(\Delta)$ by summing tail probabilities.
Finally replace the common term by $-\log(1-\Delta)$.
\end{proof}

For exponential source tails this is $O(\Delta\log(1/\Delta))$.
The extra state changes the approximation class; it does not improve the
sharp one-state rate in Theorem~\ref{II:U:thm:kl-exp} within the same class.

\paragraph{Keeping the same positive-state count.}
Alternatively, for $0<\Delta<1$ smooth each state locally:
\[
 \tau_i^{(\Delta)}=(1-\Delta)\tau_i+\frac\Delta{d+1},\qquad
 T_a^{(\Delta)}=(1-\Delta)T_a+\frac\Delta{d+1}I_N.
\]
The rows remain stochastic and every state has stop probability at least
$\Delta/(d+1)$. The all-original and all-smoothing choices in the expanded
product give the sum of two nonnegative contributions:
\[
 q_\Delta(w)\ge(1-\Delta)^{|w|+1}q(w)
                   +\left(\frac\Delta{d+1}\right)^{|w|+1}.
\]
Applying log-sum to the common and excess submeasures as before gives,
when $\UTV(p,q)\le\Delta$,
\begin{align*}
 D_{\rm KL}(p\|q_\Delta)
 &\le-\log(1-\Delta)(1+\mathbb E_p|W|)\\
 &\quad+\Delta\log(d+1)+\Psi_p(\Delta)\log\frac{d+1}{\Delta}.
\end{align*}
For exponential tails this is $O(\Delta\log^2(1/\Delta))$.
Sharpness for a general fixed $N$ is not established by this construction.
Local smoothing can increase the original visit budget and changes the
transition topology; both effects must be accounted for.

\subsection{The unresolved acquisition problem}\label{II:U:sec:acquisition-obstruction}
For a prefix anchor $u_i$, a residual dictionary law is
$f_i(v)=p(u_iv)/P_p(u_i\Sigma^*)$. The denominator is a prefix mass,
not the completed-word value $p(u_i)$. Rare anchors may require many
samples; source tails do not provide uniform residual-tail constants
without accounting for that denominator. The supplied-dictionary theorem
does not discover a small nearly closed dictionary or bound its acquisition
cost. Linear rank and positive state size also differ: nonnegative rational
series need not have finite positive realizations. A rank bound alone cannot
replace the positive dictionary hypothesis.

\begin{proposition}[A classical positive-realization obstruction]\label{II:U:prop:positive-obstruction}
For $0<x<1$ and binary words, put $d(w)=\#_a(w)-\#_b(w)$ and
\[
 p_x(w)=(1-x)^2(x/2)^{|w|}(1+d(w))^2.
\]
This is a probability law of counting-Hankel rank exactly three, with
positive empty mass and an exponential tail, but no finite nonnegative
weighted-automaton realization.
\end{proposition}
\begin{proof}
Over all words of length $t$, the mean of $d(w)$ is zero and its second
moment is $t$. Therefore the length mass is
$(1-x)^2(t+1)x^t$, summing to one. The tail has an exponential bound
with any fixed rate strictly between $x$ and one.
Since $d(uv)=d(u)+d(v)$, expanding the square gives rank at most three.
Restriction to words $a^j$ gives a nondegenerate quadratic times a
geometric sequence and hence rank at least three. Its zero language is
$\{w:\#_b(w)=\#_a(w)+1\}$, which is nonregular: intersect with
$a^*b^*$ to obtain $\{a^nb^{n+1}:n\ge0\}$.
A finite automaton with nonnegative weights has regular nonzero support,
because a coefficient is positive exactly when a positive-weight path
exists. Its zero language is consequently regular as well, a contradiction.
\end{proof}

This is a classical type of signed-versus-positive realization obstruction
\citep{denis2008}; this is not a new impossibility claim.
It already has $\sigma_4(H_{p_x})=0$ and shows why the higher-rank bridge
cannot rely on a spectral tail alone to guarantee an exact finite positive
realization.
\subsection{Conditioned residual anchors and signed spectral closure}
\label{II:U:sec:anchored-closure}
The following bridge is conditional on a chosen dictionary and a positive
flow defect. It supplies the previously missing explicit hypotheses; it
does not derive positivity or a small dictionary from singular values.
All data in this subsection are real, and $p$ is a probability law.

\begin{definition}[Conditioned residual anchors]
\label{II:U:def:conditioned-anchors}
Fix $k\ge1$, $N\ge k$, and constants $\mu_{\min}>0$ and
$\kappa<\infty$. Choose prefix anchors
$\mathcal U=(u_1,\ldots,u_N)$ with $u_1=\varepsilon$ and
\[
 \mu_i=P_p(u_i\Sigma^*)=\sum_v p(u_iv)\ge\mu_{\min},\qquad
 f_i(v)=\frac{p(u_iv)}{\mu_i}.
\]
Each $f_i$ is a probability law, $f_1=p$, and, with
$E_{\mathcal U}=[e_{u_1},\ldots,e_{u_N}]$ and
$D_\mu=\Udiag(\mu_1,\ldots,\mu_N)$,
\[
 F_{\mathcal U}=[f_1,\ldots,f_N]
 =H_p^*E_{\mathcal U}D_\mu^{-1}.
\]
Write $H_p=U_k\Sigma_kV_k^*+E_k$ for a truncated SVD and put
$\epsilon_k=\|E_k\|=\sigma_{k+1}(H_p)$. The spectral coordinate
matrix must satisfy
\begin{equation}
 C_{\mathcal U}=V_k^*F_{\mathcal U}
 =\Sigma_kU_k^*E_{\mathcal U}D_\mu^{-1},\qquad
 \Urank C_{\mathcal U}=k,\qquad
 \|C_{\mathcal U}^{\dagger}\|_{2\to2}
 =\frac1{\sigma_k(C_{\mathcal U})}\le\kappa.
 \label{II:U:eq:anchor-coordinates}
\end{equation}
These conditions define
$\mathcal H_{\rm anchor}(k,N,\mu_{\min},\kappa)$.
\end{definition}

The orientation is forced by \eqref{I:U:eq:intertwining}: its adjoint is
$L_a^*H_p^*=H_p^*R_a$. Therefore
\begin{equation}
 f_i=\mu_i^{-1}H_p^*e_{u_i},\qquad
 L_a^*f_i=\mu_i^{-1}H_p^*R_ae_{u_i}
 =\mu_i^{-1}H_p^*e_{u_i a},\qquad
 (L_a^*f_i)(v)=f_i(av).
 \label{II:U:eq:anchor-orientation}
\end{equation}
The future subspace is the right singular subspace of $H_p$; using its
left singular vectors here would interchange the prefix and suffix roles.

\begin{lemma}[Conditioned signed linear closure]
\label{II:U:lem:signed-anchor-closure}
Under Definition~\ref{II:U:def:conditioned-anchors}, set
\[
 w_{a,i}^{\rm lin}=C_{\mathcal U}^{\dagger}V_k^*(L_a^*f_i).
\]
These coefficients, which need not be nonnegative, satisfy
\begin{equation}
 \|w_{a,i}^{\rm lin}\|_2\le\kappa,\qquad
 \left\|L_a^*f_i-F_{\mathcal U}w_{a,i}^{\rm lin}\right\|_2
 \le c_{\mathcal U}\epsilon_k,\qquad
 c_{\mathcal U}=\frac{1+\kappa\sqrt N}{\mu_{\min}}.
 \label{II:U:eq:signed-anchor-defect}
\end{equation}
This is linear closure, not affine or positive closure.
\end{lemma}
\begin{proof}
Put $P=V_kV_k^*$ and $g=L_a^*f_i$. Full row rank of
$C_{\mathcal U}$ gives $PF_{\mathcal U}w_{a,i}^{\rm lin}=Pg$, so
\[
 g-F_{\mathcal U}w_{a,i}^{\rm lin}
 =(I-P)g-(I-P)F_{\mathcal U}w_{a,i}^{\rm lin}.
\]
By \eqref{II:U:eq:anchor-orientation}, $\|(I-P)g\|_2\le
\epsilon_k/\mu_{\min}$. Each column of $(I-P)F_{\mathcal U}$
has the same bound, hence its operator norm is at most
$\sqrt N\epsilon_k/\mu_{\min}$.
Finally $\|g\|_2\le\|f_i\|_1=1$ yields
$\|w_{a,i}^{\rm lin}\|_2\le\kappa$. The triangle inequality proves
\eqref{II:U:eq:signed-anchor-defect}.
\end{proof}


\subsection{The conic hypothesis and the conditional higher-rank bridge}
\label{II:U:sec:conditional-cone-bridge}
Fix the anchored dictionary above and $\alpha=(1,0,\ldots,0)$.
For $V\ge1$, let $\mathcal F_N(V)$ consist of nonnegative flows
$(x,z,y)$ satisfying \eqref{II:U:eq:flow}; explicitly,
\[
 x_j=\alpha_j+\sum_{a,i}y_{aij},\qquad
 x_i=z_i+\sum_{a,j}y_{aij},\qquad \sum_i x_i\le V.
\]
This is a nonempty compact polytope. In particular,
$\sum_i z_i=1$ and $\sum_{a,i,j}y_{aij}=\sum_i x_i-1$.

\begin{definition}[Spectral-reference conic defect in $\ell^2$]
\label{II:U:def:cone-defect}
For a feasible flow, define
\begin{align}
 \eta_2(x,z,y)
 &=\sum_i|x_if_i(\varepsilon)-z_i|\notag\\
 &\quad+\sum_{a,i}
 \left\|F_{\mathcal U}
       (x_iw_{a,i}^{\rm lin}-y_{a,i,\cdot})\right\|_2,
 \label{II:U:eq:cone-defect}
\end{align}
where $y_{a,i,\cdot}=(y_{ai1},\ldots,y_{aiN})^\top$.
The condition $\mathcal H_{\rm cone}^{(2)}(V,\eta_{\rm cone})$
means that some flow in $\mathcal F_N(V)$ has
$\eta_2(x,z,y)\le\eta_{\rm cone}$.
Equivalently, the attained minimum
$\eta_2^*(\mathcal U,V)=\min_{\mathcal F_N(V)}\eta_2$ is at most
$\eta_{\rm cone}$. This hypothesis is separate from
$\mathcal H_{\rm anchor}$ and the spectral bound.
\end{definition}

\paragraph{SOCP versus LP.}
Let $G=F_{\mathcal U}^*F_{\mathcal U}$, so
$G_{jj'}=\langle f_j,f_{j'}\rangle_{\ell^2}$ and
$\|F_{\mathcal U}c\|_2=\|G^{1/2}c\|_2$.
With $G$, $w_{a,i}^{\rm lin}$, and $f_i(\varepsilon)$ fixed and
known, epigraph constraints
$\|G^{1/2}(x_iw_{a,i}^{\rm lin}-y_{a,i,\cdot})\|_2\le s_{ai}$,
together with linear epigraphs for the absolute values, give a finite
second-order cone program (SOCP). Singularity of $G$ is harmless.
An arbitrary infinite Hankel oracle does not automatically supply these
finite coefficients with certified accuracy. Computing or estimating the
SVD, prefix masses, Gram matrix, and feasible flows must be charged.
The LP of Theorem~\ref{II:U:thm:flow-lp} instead minimizes a confidence-adjusted
upper bound for the truncated $\ell^1$ closure certificate; it neither
computes nor minimizes \eqref{II:U:eq:cone-defect} directly.

\begin{theorem}[Conditional spectral--conic TV bridge]
\label{II:U:thm:conditional-cone-bridge}
Assume $p=\sum_i\alpha_i f_i$ with the anchored choice $f_1=p$ and $\alpha=e_1$, as well as $\mathcal H_{\rm anchor}(k,N,\mu_{\min},\kappa)$ and
$\mathcal H_{\rm cone}^{(2)}(V,\eta_{\rm cone})$ and choose a
witnessing feasible flow. Decode $\widehat q$ as in
Theorem~\ref{II:U:thm:flows}. It has at most $N$ states, is proper, and
$\mathbb E_{\widehat q}|W|\le V-1$. Define
$t_p(s)=P_p(|W|>s)$ for nonnegative integers $s$, and
$D_T=\sum_{\ell=0}^{T-1}d^\ell$ for $T\ge1$. Then
\begin{align}
 J_\infty
 &\le\sqrt{D_T}\bigl(dV c_{\mathcal U}\epsilon_k
                        +\eta_{\rm cone}\bigr)
       +\frac{2V-1}{\mu_{\min}}t_p(T-1),
       \label{II:U:eq:conditional-cone-j}\\
 \UTV(p,\widehat q)&\le\tfrac12J_\infty.
       \label{II:U:eq:conditional-cone-tv}
\end{align}
The conic defect in this theorem is the $\ell^2$ quantity
\eqref{II:U:eq:cone-defect}, not the total $\ell^1$ defect $J_\infty$.
\end{theorem}
\begin{proof}
For $R_i=x_if_i-z_i\delta_\varepsilon-\sum_{a,j}y_{aij}L_af_j$,
let $A=\sum_i|R_i(\varepsilon)|$. On the branch starting with $a$,
\[
 L_a^*R_i=x_i(L_a^*f_i-F_{\mathcal U}w_{a,i}^{\rm lin})
           +F_{\mathcal U}(x_iw_{a,i}^{\rm lin}-y_{a,i,\cdot}).
\]
Lemma~\ref{II:U:lem:signed-anchor-closure} and
Definition~\ref{II:U:def:cone-defect} give
\[
 \sum_{a,i}\|L_a^*R_i\|_2
 \le dV c_{\mathcal U}\epsilon_k+\eta_{\rm cone}-A.
\]
Cauchy--Schwarz on the $D_T$ suffixes of lengths $0$ through $T-1$
therefore bounds the contribution of $|w|\le T$ by
$\sqrt{D_T}(dV c_{\mathcal U}\epsilon_k+\eta_{\rm cone})$,
since $\sqrt{D_T}\ge1$ absorbs $A$.

The exact residual-tail relation and its useful upper bound are
\begin{equation}
 t_{f_i}(s)=\frac{\sum_{|v|>s}p(u_iv)}{\mu_i}
 \le\frac{t_p(s+|u_i|)}{\mu_i}
 \le\frac{t_p(s)}{\mu_{\min}}.
 \label{II:U:eq:conditional-anchor-tail}
\end{equation}
Positivity bounds the remaining contribution by
\begin{align*}
 \sum_i\sum_{|w|>T}|R_i(w)|
 &\le\sum_i x_it_{f_i}(T)
           +\sum_{a,i,j}y_{aij}t_{f_j}(T-1)\\
 &=\sum_i x_it_{f_i}(T)
           +\sum_j(x_j-\alpha_j)t_{f_j}(T-1)\\
 &\le\frac{Vt_p(T)+(V-1)t_p(T-1)}{\mu_{\min}}\\
 &\le\frac{2V-1}{\mu_{\min}}t_p(T-1).
\end{align*}
Adding head and tail proves \eqref{II:U:eq:conditional-cone-j}.
Since $f_1=p$, Theorem~\ref{II:U:thm:flows} applies with no initial-mixture
error and proves all remaining assertions.
\end{proof}


\paragraph{An alternative direct $\ell^1$ hypothesis.}
One can instead require a feasible flow with
\[
 \eta_1(x,z,y):=\sum_i|x_if_i(\varepsilon)-z_i|
       +\sum_{a,i}\left\|x_iL_a^*f_i-\sum_jy_{aij}f_j\right\|_1
       \le\eta_1^{\rm target}.
\]
The empty word and letter branches have disjoint supports, so
$\eta_1=J_\infty$ exactly and the TV bound is
$\eta_1^{\rm target}/2$, without spectral truncation or a separate
tail term. This directly certifying formulation must not be
identified with Definition~\ref{II:U:def:cone-defect}.
For a KL bound without increasing the number of states, apply the local
smoothing construction in Section~\ref{II:U:sec:flow-smoothing} with a
\emph{prescribed positive upper bound} $\Delta$ on the certified TV.
Its tail assumptions and possible increase in expected visits remain
part of the guarantee. A mean bound alone does not make KL vanish
uniformly, by Theorem~\ref{II:U:thm:kl-fixed-states}.

\subsection{Certified search and the unresolved acquisition theorem}
\label{II:U:sec:acquisition-contract}
An anchor-search heuristic may start with $\varepsilon$, fit a feasible
flow, and add a previously absent child $u_i a$ whose estimated branch
defect is large and whose prefix mass passes a declared threshold.
The exact local identity
\[
 L_a^*f_{u_i}=\frac{P_p(u_i a\Sigma^*)}{P_p(u_i\Sigma^*)}f_{u_i a}
\]
holds whenever the child mass is positive. It closes that branch if its
coefficient is used, but the resulting flows must still jointly satisfy
\eqref{II:U:eq:flow}. Neither the visit budget, the spectral conditioning,
nor all other branch defects are thereby controlled.

\paragraph{What a successful run certifies.}
Freeze the dictionary before the certification sample, or establish a
uniform concentration event covering its adaptive selection. On the event
in Theorem~\ref{II:U:thm:flow-lp}, an actually returned feasible flow with
$\widehat C_L\le2\varepsilon$ certifies
$\UTV(p,\widehat q)\le\varepsilon$. This is a probability-at-least
$1-\delta$ statement when empirical data are used, not an unconditional
deterministic guarantee. The LP may optimize its flows on those data
because its event is uniform over feasible flows for the fixed dictionary.
If successive selected dictionaries are certified on fresh data, allocate
failure probabilities $\delta_r$ with $\sum_r\delta_r\le\delta$.
Reusing a holdout to select dictionaries without such control is not
covered. Numerical flow feasibility and the achieved objective must also
be verified; the unknown optimum alone is not a certificate. This LP
certificate does not require spectral full rank or a condition-number
bound unless the spectral bridge is separately claimed.

\paragraph{The remaining acquisition problem.}
A discovery theorem must specify an independently checkable source class,
an input/access model, and an algorithm, then prove all of the following:
\begin{enumerate}[label=(A\arabic*)]
\item A dictionary of at most $N(\varepsilon)$ anchors is found in
bounded time, with explicitly controlled minimum prefix masses and
spectral conditioning, rather than assumed to exist.
\item A globally feasible flow is obtained with controlled visit budget
$V$ and conic defect. Their dependence on the accuracy must make the
\emph{whole} right side of \eqref{II:U:eq:conditional-cone-j} at most
$2\varepsilon$ for a stated choice of $T$. Merely requiring
$\eta_{\rm cone}\le\varepsilon$ is insufficient when $D_T$, $V$,
$\kappa$, or $\mu_{\min}^{-1}$ grow.
\item Anchor discovery, conditional sampling, spectral/Gram estimation,
finite precision, and optimization have explicit sample, bit, and runtime
costs, including rare-prefix and adaptive-selection effects.
\item The algorithm returns a verified model under the stated assumptions;
outside them, failure to find a certificate is reported as unresolved,
not as nonexistence of a positive approximation.
\end{enumerate}
The smallness of $\sigma_{k+1}(H_p)$ alone proves none of these
acquisition claims. Proposition~\ref{II:U:prop:positive-obstruction} rules out
an exact universal positive-realization implication even at zero spectral
tail. It does not rule out approximation with a growing state count:
conditioning $p$ on $|W|\le T$ gives a finite-support law at TV distance
$t_p(T)$, realizable by a finite prefix tree. This construction neither
provides a small state bound nor finite KL for unbounded-support $p$.
The transfer-operator result in Subsection~\ref{III:U:sec:oracle} concerns
already supplied contracting two-state homographic models and supplies
no missing anchor-discovery theorem.


\subsection{What can be certified over a finite candidate family}
\label{II:U:sec:finite-family-search}
Existence of a good hidden anchor tree does not prove that a local greedy
rule reaches it within the same number of nodes. A finite exhaustive
search has the following weaker, explicit guarantee.

\begin{theorem}[Finite-family search with a common validation event]
\label{II:U:thm:finite-family-search}
Fix $M$ candidate prefix laws $f_1=p,f_2,\ldots,f_M$ before validation,
and a finite family $\mathcal D$ of subsets containing index one, each
of size at most $N_*$. Assume conditional samples from these laws are
available. Each candidate dictionary is started with $\alpha=e_1$, so $p=\sum_i\alpha_i f_i$ holds exactly. For an integer horizon $L\ge1$, take $n_i$ independent samples per candidate
law, with no independence required between different candidate laws,
and use
\[
 b_i=\min\left\{2,\sqrt{
    \frac{2(K_L\log2+\log(M/\delta))}{n_i}}\right\},\qquad
 K_L=1+\sum_{j=0}^Ld^j.
\]
Solve the flow LP for every dictionary in $\mathcal D$ with budget $V$,
using the same histograms and the applicable $b_i$. With probability
at least $1-\delta$, every feasible returned model satisfies its
certificate $\UTV(p,q)\le\widehat C_L/2$ simultaneously.
If one candidate dictionary admits a flow satisfying
\begin{equation}
 J_\infty+2\sum_i x_it_i(L)+2S(x)\le2\varepsilon,
 \label{II:U:eq:finite-family-witness}
\end{equation}
then, with exact LP solves, the search finds a certified model of error
at most $\varepsilon$ after at most $|\mathcal D|$ solves.
Otherwise it may return \texttt{UNRESOLVED}. For all subsets of size
at most $N_*$, the number of solves is bounded by
$\sum_{j=0}^{N_*-1}\binom{M-1}{j}$, with terms beyond $M-1$ omitted.
\end{theorem}
\begin{proof}
Apply the histogram concentration argument of
Theorem~\ref{II:U:thm:flow-lp} once to the $M$ fixed candidate laws.
Its simultaneous event covers every feasible flow on every subset, so
subsequent dictionary selection does not require another union bound.
For the witnessing flow,
$\widehat C_L^*\le J_L+2S(x)
\le J_\infty+2\sum_i x_it_i(L)+2S(x)$.
Condition \eqref{II:U:eq:finite-family-witness} therefore forces an accepted
dictionary when it is examined. The finite subset count is immediate.
\end{proof}


This theorem neither discovers the candidate universe nor promises a
small search. Access to conditional samples, certification of prefix
masses, numerical LP feasibility, and spectral data are separate costs.
For numerical LP solves, a verified objective gap must be included in
\eqref{II:U:eq:finite-family-witness}. An existential witness-tree assumption
does not yield the proposed greedy guarantee. In particular, ordinary
column pivoting is not the strong rank-revealing QR algorithm of
Gu--Eisenstat \citep{gu1996}; moreover the inequality
$2/\sqrt{k(N_*-k)+1}\ge1$ is already false at $(k,N_*)=(3,6)$.

\subsection{Positive closure in the correct dual state coordinates}
\label{II:U:sec:dual-positive-closure}
For $p(w)=c^\top A_wb$, prefix residuals are represented by
\emph{dual} state vectors. The vectors $A_ub$ instead represent suffix
columns and cannot be substituted without changing the word order.

\begin{theorem}[A sufficient dual-coordinate positive closure condition]
\label{II:U:thm:dual-positive-closure}
Let $p$ have a real minimal mean-square-stable realization of dimension
$k$. Choose anchors with $u_1=\varepsilon$, masses
$\mu_i\ge\mu_{\min}>0$, and vectors
\[
 \ell_i=\mu_i^{-1}A_{u_i}^\top c,\qquad
 (\mathcal R_b\ell)(v)=\ell^\top A_vb.
\]
Assume the $\ell_i$ span $\mathbb R^k$. Suppose nonnegative numbers
$t_{aij}$ satisfy, for every $i,a$,
\begin{align}
 \sum_{a,j}t_{aij}&=1-f_i(\varepsilon),\qquad
 f_i(\varepsilon)\ge\gamma_{\rm stop}>0,\notag\\
 \left\|A_a^\top\ell_i-\sum_jt_{aij}\ell_j\right\|_2
 &\le\eta_{\rm pos}.\label{II:U:eq:dual-positive-hypothesis}
\end{align}
Then the conditioned-anchor hypothesis holds with
$\kappa=\|(V_k^*F_{\mathcal U})^\dagger\|<\infty$.
Moreover, there is a feasible flow with
\[
 V=\gamma_{\rm stop}^{-1},\qquad
 \eta_2\le dV\|\mathcal R_b\|\eta_{\rm pos},\qquad
 \|\mathcal R_b\|^2=\|P_p\|.
\]
Thus Theorem~\ref{II:U:thm:conditional-cone-bridge} applies with
$\epsilon_k=0$. This result presupposes the supplied spanning anchors
and the jointly normalized positive coefficients; it is not an
anchor-discovery theorem.
\end{theorem}
\begin{proof}
The correct identities are
$f_i=\mathcal R_b\ell_i$ and
$L_a^*f_i=\mathcal R_bA_a^\top\ell_i$.
Minimality makes $P_p\succ0$, so $\mathcal R_b$ is injective and
the spanning assumption makes $F_{\mathcal U}$ have rank $k$.
Its image is the right singular space of $H_p$; hence the stated
$\kappa$ is finite and the signed closure in
Lemma~\ref{II:U:lem:signed-anchor-closure} is exact.

Take $\alpha=e_1^\top$, put $T_a=(t_{aij})$ and $T=\sum_aT_a$. Its row sums are at most
$1-\gamma_{\rm stop}$. Therefore
$x=\alpha(I-T)^{-1}\ge0$ and $\sum_i x_i\le V$.
Set $y_{aij}=x_it_{aij}$ and $z_i=x_if_i(\varepsilon)$.
These satisfy both flow equations. Their stop defect is zero, and
the exact signed closure gives
\[
 \eta_2=\sum_{a,i}x_i
 \left\|L_a^*f_i-\sum_jt_{aij}f_j\right\|_2
 \le dV\|\mathcal R_b\|\eta_{\rm pos}.
\]
Finally $\|\mathcal R_b\ell\|_2^2=\ell^\top P_p\ell$ proves
the asserted operator norm. No assertion that this cone is simplicial
is needed; it may have more than $k$ generators.
\end{proof}


\subsection{Conditional predictions and the effect of smoothing}
\label{II:U:sec:conditional-event-bounds}
Let $p_u$ denote the next-event law on $\Sigma\cup\{\mathrm{STOP}\}$
after observing a prefix $u$. All KL quantities below use natural logs.

\begin{theorem}[Prefix-conditioned TV and a valid smoothing bound]
\label{II:U:thm:conditional-event-bounds}
Let $p,q$ be proper word laws with $\UTV(p,q)\le\varepsilon$.
For $s=P_p(u\Sigma^*)\ge\mu_{\rm pref}>\varepsilon$,
$P_q(u\Sigma^*)>0$ and
\begin{equation}
 \UTV(p_u,q_u)\le\frac{\varepsilon}{s}
                         \le\frac{\varepsilon}{\mu_{\rm pref}}.
 \label{II:U:eq:conditional-event-tv}
\end{equation}
Suppose $q$ is an $N$-state automaton with
$1+\mathbb E_q|W|\le V$, and apply the state-preserving local smoothing
of Section~\ref{II:U:sec:flow-smoothing} with parameter $0<\Delta<1$.
Put $e=\varepsilon+\Delta V$. If $\mu_{\rm pref}>e$, then
\begin{align}
 \UTV(p_u,(q_\Delta)_u)&\le e/\mu_{\rm pref},\notag\\
 D_{\rm KL}(p_u\|(q_\Delta)_u)
 &\le\min\left\{\log\frac{d+1}{\Delta},\;
 \frac{2(d+1)}{\Delta}\left(\frac e{\mu_{\rm pref}}\right)^2
 \right\}.\label{II:U:eq:conditional-event-kl}
\end{align}
The number of states remains $N$; no preservation of the visit budget
after smoothing is asserted.
\end{theorem}
\begin{proof}
Let $E=u\Sigma^*$, $t=Q(E)$, and
$A=\sum_{w\in E}|p(w)-q(w)|$. The complement has total variation
mass at least $|s-t|$, so $A+|s-t|\le2\varepsilon$.
Also $t\ge s-\varepsilon>0$. Consequently
\[
 \sum_{w\in E}\left|\frac{p(w)}s-\frac{q(w)}t\right|
 \le\frac{A+|s-t|}{s}\le\frac{2\varepsilon}s.
\]
Mapping a conditional suffix to its first letter or STOP cannot
increase TV; this proves \eqref{II:U:eq:conditional-event-tv}.
The sum of absolute differences of next-event masses is \emph{at most},
not necessarily equal to, the sum over all completed suffixes.

Couple the original and smoothed automata until a smoothing choice
occurs. For an original word $W$, there are $|W|+1$ event choices, so
$\UTV(q,q_\Delta)\le
\mathbb E[1-(1-\Delta)^{|W|+1}]\le\Delta V$.
Apply the first part to $p,q_\Delta$. Every next-event probability
under $q_\Delta$, for any posterior state mixture, is at least
$a=\Delta/(d+1)$. For probability vectors $r,z$ with $z_j\ge a$,
the inequality $\log t\le t-1$ gives
\[
 D_{\rm KL}(r\|z)\le\sum_j\frac{(r_j-z_j)^2}{z_j}
 \le\frac{2\UTV(r,z)^2}{a}.
\]
The last step uses that positive and negative differences each sum to
$\UTV(r,z)$. Also $D_{\rm KL}(r\|z)\le\log(1/a)$ by dropping
the nonnegative entropy of $r$. These prove the two KL bounds.
\end{proof}


\paragraph{Counterexample to prefix-independent smoothing error.}
Take $p=q$ on two states with initial row $(1/2,1/2)$,
\[
 T_a=\Udiag(t,0),\quad T_b=\Udiag(0,1/2),\quad
 \tau=(1-t,1/2)^\top,
 \qquad t=10^{-6},\quad\Delta=10^{-3}.
\]
After prefix $a$, the unsmoothed next-event law (STOP,$a$,$b$) is
$(1-t,t,0)$, and its prefix mass is $t/2>0$. After local smoothing,
the posterior weight of the first state is
\[
 \omega=\frac{(1-\Delta)t+\Delta/3}
                   {(1-\Delta)t+2\Delta/3}.
\]
The new next-event law is approximately
$(0.7499565235,0.0003338336,0.2497096429)$.
Its KL from the original law exceeds $0.2877$, whereas a proposed uniform smoothing bound at $\varepsilon=0$ is
$\Delta\log(3/\Delta)<0.0081$. The changed posterior after a rare
prefix cannot be ignored. This example also has original visit budget
less than two; it is not a failure caused by infinite expected length.


\subsection{Finite positive states need not be finite prefix anchors}
\label{II:U:sec:rank-two-obstruction}
The rank-two results below are conditional on a supplied, exact
two-dimensional model and a certified interval of probability laws.
They do not complete arbitrary-rank spectral learning. They also
distinguish latent generators, limiting residuals, and actual finite
prefix residuals; these are different admissible dictionaries.

\begin{lemma}[The aggregate of a minimal proper-word realization is stable]
\label{II:U:lem:proper-aggregate-stable}
Let $p(w)=c^\top A_wb$ be a minimal finite-dimensional real realization
of a probability law on finite words, and put $T=\sum_a A_a$.
Then $T^n\to0$ and $\rho(T)<1$.
\end{lemma}
\begin{proof}
Choose words $u_i,v_j$ so that the matrix $O$ with rows
$c^\top A_{u_i}$ and the matrix $C$ with columns $A_{v_j}b$ are
square and invertible; minimality guarantees these choices. For each
fixed pair $(i,j)$,
\[
 (OT^nC)_{ij}=\sum_{|w|=n}p(u_iwv_j)
 \le P_p\{|W|=n+|u_i|+|v_j|\}\longrightarrow0.
\]
Hence $T^n=O^{-1}(OT^nC)C^{-1}\to0$. In finite dimension this
implies that every eigenvalue has modulus strictly below one.
\end{proof}

Thus a Perron eigenvalue equal to one for the unnormalized aggregate
is incompatible with this convention. A normalized continuation
operator is a different object.

\begin{proposition}[A strictly positive two-state prefix obstruction]
\label{II:U:prop:golden-prefix-obstruction}
Set
\[
 A_0=\frac1{10}\begin{pmatrix}2&1\\1&1\end{pmatrix},\quad
 A_1=\frac1{10}\begin{pmatrix}1&1\\1&2\end{pmatrix},\quad
 \alpha=(1/2,1/2),\quad \tau=(1/2,1/2)^\top.
\]
The law $p(w)=\alpha A_w\tau$ is proper and has Hankel rank two.
Both letters map the positive orthant strictly into its interior,
but no finite dictionary of actual prefix residuals containing $p$
in its convex hull has exact positive closure.
\end{proposition}
\begin{proof}
Here $T\mathbf1=\frac12\mathbf1$, so termination is geometric.
The two-row reachability matrix with rows $\alpha,\alpha A_0$
and the two-column observability matrix with columns $\tau,A_0\tau$
both have determinant $-1/40$, proving rank two.
For this example write $f_t(v)=(t,1-t)A_v\tau$; its first coordinate
updates by
\[
 F_0(t)=\frac{1+t}{2+t},\qquad F_1(t)=\frac1{3-t}.
\]
On $[0,1]$ their images are $[1/2,2/3]$ and $[1/3,1/2]$.
With $r=(\sqrt5-1)/2$, both maps preserve $(1-r,r)$, and every
finite-prefix coordinate lies in that interval. The words $0^n$
and $1^n$ approach its two endpoints. For any finite dictionary let
$m<r$ be its maximum coordinate. Then
\[
 F_0(m)-m=\frac{1-m-m^2}{2+m}>0.
\]
The next residual is outside the dictionary's convex hull. The
coordinate map is affine and injective since
$f_t(0)=1/10+t/20$. Nonnegative representations of normalized laws
have coefficient sum one, so convex and conic closure give the same
obstruction after normalizing each nonzero shift.
\end{proof}


The same obstruction rules out $J_\infty=0$ for every finite feasible
occupation flow with this source in its initial mixture: restrict to
states with positive occupation. Zero-occupation states receive no
flow; zero defect would make the active dictionary invariant. For a
fixed finite visit budget the feasible flow set is compact, hence a
nonempty such problem has strictly positive minimum defect. No
uniform positive lower bound over increasing dictionaries is claimed.
The distinction between positive rational realization and finite
generation by residuals is classical \citep{denis2008}.

For example,
\[
 t_\varepsilon=\tfrac12,\quad t_0=\tfrac35,\quad t_{00}=\tfrac8{13},
 \qquad f_{00}=-\tfrac2{13}f_\varepsilon+\tfrac{15}{13}f_0.
\]
Thus linear QR residuals vanish after two independent anchors while
conic residuals need not. Also $2T$ has eigenvalues $1,1/5$, yet
$0^n$ and $1^n$ have different limiting conditional states. A gap for
the normalized aggregate does not give convergence of every letter
product to a common ray.

\begin{lemma}[A sufficient simplex-expanding exchange]
\label{II:U:lem:simplex-exchange}
Let $f_1,\ldots,f_R$ be linearly independent normalized laws, and let
another law satisfy $f=\sum_i\beta_i f_i$, where $\sum_i\beta_i=1$.
If $\beta_j>1$ for one index $j$ and $\beta_i\le0$ for every $i\ne j$,
replacing $f_j$ by $f$ preserves linear independence and strictly
enlarges the convex hull.
\end{lemma}
\begin{proof}
The identity
\[
 f_j=\beta_j^{-1}f+
       \sum_{i\ne j}(-\beta_i/\beta_j)f_i
\]
is a convex combination. Independence follows from $\beta_j\ne0$,
and $f$ was outside the old simplex by uniqueness of its coordinates.
\end{proof}

This is a sufficient exchange test, not a termination theorem.
In higher affine dimension an exterior point can have several
positive coordinates, so an eligible exchange need not exist.

\subsection{Rank-two monotone realization and finite-prefix approximation}
\label{II:U:sec:rank-two-monotone}
Choose two distinct normalized residual laws $f_0,f_1$ spanning the
rank-two residual space. In this subsection the convention is
\[
 f_t=(1-t)f_0+tf_1,\qquad
 L_1=\|f_1-f_0\|_1,\qquad \|f_t-f_s\|_1=L_1|t-s|.
\]
The notation differs from the first-coordinate convention used in
Proposition~\ref{II:U:prop:golden-prefix-obstruction}; changing orientation
does not change any law or certificate. There is a real realization
$f_t(v)=(1-t,t)A_v\tau$. By Lemma~\ref{II:U:lem:proper-aggregate-stable}
and normalization, $(I-\sum_a A_a)^{-1}\tau=\mathbf1$.
For $A_a=\left(\begin{smallmatrix}a_a&b_a\\c_a&d_a\end{smallmatrix}\right)$
define
\begin{align}
 w_a(t)&=(1-t,t)A_a\mathbf1,\qquad s(t)=f_t(\varepsilon)
                  =1-\sum_a w_a(t),\notag\\
 F_a(t)&=\frac{b_a+(d_a-b_a)t}{w_a(t)},\qquad
 F_a'(t)=\frac{\det A_a}{w_a(t)^2}.
 \label{II:U:eq:rank-two-projective}
\end{align}
When $w_a(t)>0$, $L_a^*f_t=w_a(t)f_{F_a(t)}$.

\begin{theorem}[Monotone rank-two limiting and prefix realizations]
\label{II:U:thm:rank-two-monotone}
Suppose a compact interval $I$ contains $t_0$ with $p=f_{t_0}$,
every $f_t$ for $t\in I$ is a probability law, and every letter satisfies
$w_a(t)>0$, $F_a(I)\subseteq I$, and
$0<F_a'(t)\le\bar\gamma_a<1$ on $I$.
Let $r_a$ be its unique fixed point and set
\[
 r_-=\min_a r_a=r_{a_-},\qquad
 r_+=\max_a r_a=r_{a_+}.
\]
Assume $r_-<t_0<r_+$ and
$s_*=\min_{[r_-,r_+]}s(t)>0$.
Then the closed convex hull of all prefix residuals is
$\operatorname{conv}(f_{r_-},f_{r_+})$. In this ordered dictionary,
the following matrices give an exact proper two-state realization:
\begin{align}
 \lambda_a^\pm&=\frac{F_a(r_\pm)-r_-}{r_+-r_-},\qquad
 B_a=\begin{pmatrix}
 w_a(r_-)(1-\lambda_a^-)&w_a(r_-)\lambda_a^-\\
 w_a(r_+)(1-\lambda_a^+)&w_a(r_+)\lambda_a^+
 \end{pmatrix},\notag\\
 \alpha&=\left(\frac{r_+-t_0}{r_+-r_-},
              \frac{t_0-r_-}{r_+-r_-}\right),\qquad
 \tau_B=(s(r_-),s(r_+))^\top.
 \label{II:U:eq:rank-two-limit-model}
\end{align}

For $L\ge1$, let
$l_L=F_{a_-}^L(t_0)$ and $h_L=F_{a_+}^L(t_0)$.
They are actual prefix coordinates, with
$r_-<l_L<t_0<h_L<r_+$.
Replace $r_-,r_+$ in the initial and stopping vectors by $l_L,h_L$,
and replace each successor $F_a(t)$ in the transition formula by
$\Pi_{[l_L,h_L]}F_a(t)$, for $t=l_L,h_L$. Denote the resulting
matrices by $\widehat B_{a,L}$, the law by $q_L$, and
$Q_L=\sum_a\widehat B_{a,L}$. Then
\[
 x_L=\alpha_L(I-Q_L)^{-1}\ge0,\qquad x_L\mathbf1\le s_*^{-1},
\]
and its feasible occupation flows have exact defect
\begin{equation}
 J_L^{\rm clip}
 =L_1\sum_{t\in\{l_L,h_L\}}x_{L,t}
        \sum_a w_a(t)\operatorname{dist}(F_a(t),[l_L,h_L]).
 \label{II:U:eq:rank-two-clip-defect}
\end{equation}
Here $J_L^{\rm clip}$ is an untruncated $J_\infty$ certificate;
it is not the finite-word objective $J_L$ of Section~\ref{II:U:sec:flow-lp}.
Writing $\delta_L=\max(l_L-r_-,r_+-h_L)$ gives
\begin{equation}
 \UTV(p,q_L)\le\tfrac12J_L^{\rm clip}
 \le\frac{L_1(1-s_*)}{2s_*}\delta_L.
 \label{II:U:eq:rank-two-clip-tv}
\end{equation}
No pointwise ordering of the different maps $F_a$ is required.
\end{theorem}
\begin{proof}
For an increasing contraction with fixed point $r_a$, a point below
$r_a$ moves strictly toward $r_a$ without crossing it, and a point
above does the same from the other side. Thus each map preserves
$[r_-,r_+]$. All reachable points lie in it; the repeated letters
$a_-^n,a_+^n$ approach its endpoints. Affinity and continuity into
$\ell^1$ prove the closed convex hull assertion.

All $\lambda_a^\pm$ lie in $[0,1]$. The rows sum to $w_a(r_\pm)$,
and the residual recurrences hold exactly. Alternatively, with
$S=\left(\begin{smallmatrix}1-r_-&r_-\\1-r_+&r_+\end{smallmatrix}\right)$,
$B_a=SA_aS^{-1}$ and $\tau_B=S\tau$, which proves equality of
every word coefficient. Each stop probability is at least $s_*$,
so the model is proper.

The same row normalization holds after clipping, with stop at least
$s_*$. Hence $\rho(Q_L)\le1-s_*$ and the displayed occupation vector
exists. Take $y_{aij}=x_i\widehat B_{a,L}(i,j)$ and
$z_i=x_i\tau_{L,i}$ to obtain the flow equations. The source is
exactly the mixture $\alpha_L(f_{l_L},f_{h_L})^\top$.
The stop residual is zero, and the residuals for different first
letters have disjoint support. Each has norm
$L_1w_a(t)|F_a(t)-\Pi_{[l_L,h_L]}F_a(t)|$, proving
\eqref{II:U:eq:rank-two-clip-defect}.
Invariance bounds every distance by $\delta_L$; also
$\sum_a w_a(t)\le1-s_*$ and $x_L\mathbf1\le1/s_*$.
Apply Theorem~\ref{II:U:thm:flows}.
\end{proof}


\begin{corollary}[Exact cross-ratio depth rates]
\label{II:U:cor:rank-two-cross-ratio}
For either extremal map let $r$ be its attracting fixed point and
$\gamma=F'(r)\in(0,1)$. If its second fixed point $\bar r$ is finite,
put $c=(t_0-r)/(t_0-\bar r)$. Then
\begin{equation}
 F^L(t_0)=\frac{r-c\gamma^L\bar r}{1-c\gamma^L},
 \qquad
 |F^L(t_0)-r|
 =\frac{|r-\bar r|\,|c|\,\gamma^L}{|1-c\gamma^L|}.
 \label{II:U:eq:rank-two-cross-ratio}
\end{equation}
If $F$ is affine, the formula is instead
$F^L(t_0)-r=\gamma^L(t_0-r)$.
Consequently $\delta_L\le C\gamma_*^L$, with
$\gamma_*=\max(F_{a_-}'(r_-),F_{a_+}'(r_+))<1$
and an explicit finite, source-dependent constant $C$.
\end{corollary}
\begin{proof}
The two fixed points are distinct because $F'(r)<1$.
For a nonaffine map the second fixed point is real and outside $I$,
since the contraction has only one fixed point in $I$.
The map $\psi(t)=(t-r)/(t-\bar r)$ conjugates $F$ to multiplication
by $F'(r)$, giving \eqref{II:U:eq:rank-two-cross-ratio}.
Since $t_0,r$ lie on the same side of $\bar r$, $c<1$ and
$1-c\gamma^L\ge\min(1,1-c)>0$. One may use
$C_r=|r-\bar r||c|/\min(1,1-c)$, or $C_r=|t_0-r|$
in the affine case, and take the larger endpoint constant.
\end{proof}

The closed interval here is a convex hull, not necessarily the IFS
attractor. A local multiplier $\gamma_*$ is justified by the exact
cross-ratio identity, not by treating it as a global Lipschitz constant.

\subsection{Both orientations: endpoint acquisition without word enumeration}
\label{II:U:sec:rank-two-envelope}
The even/odd powers of one decreasing letter do not generally find the
extrema created by mixed words. For instance, the proper rank-two model
\[
 A_+=\tfrac1{10}\begin{pmatrix}2&1\\1&1\end{pmatrix},\quad
 A_-=\tfrac1{20}\begin{pmatrix}1&2\\2&1\end{pmatrix},\quad
 \alpha=(1/2,1/2),\quad \tau=(11/20,13/20)^\top
\]
has first-coordinate updates
$F_+(t)=(1+t)/(2+t)$ and $F_-(t)=(2-t)/3$.
All powers of $F_-$ fix $1/2$, but
$F_-(F_+(1/2))=7/15$. The following procedure treats all letters.

\begin{theorem}[Rank-two endpoint iteration for either orientation]
\label{II:U:thm:rank-two-envelope}
Use the affine-law setup above. Suppose every $f_t$ for $t$ in a
compact interval $I$ is a probability law, $t_0\in I$,
$w_a(t)>0$, $F_a(I)\subseteq I$, and every $F_a$ is monotone
(increasing, decreasing, or constant) and $\gamma$-Lipschitz on $I$,
with $0<\gamma<1$. Assume $s(t)\ge s_*>0$ on $I$.
Initialize $K_0=\{t_0\}$ and iterate
\begin{equation}
 K_{n+1}=\operatorname{conv}
   \left(\{t_0\}\cup\bigcup_a F_a(K_n)\right)
       =[l_{n+1},h_{n+1}].
 \label{II:U:eq:rank-two-envelope}
\end{equation}
Then $K_n$ is the convex hull of coordinates of prefixes of length
at most $n$, and its endpoints have actual prefix witnesses.
The intervals increase to a unique compact fixed interval $K_*$
whose endpoints are limits of prefix residuals, with
\[
 d_H(K_n,K_*)\le\gamma^n\operatorname{diam}I.
\]
If the reachable residual space has rank two, $K_*$ is nondegenerate,
and $K_n$ is nondegenerate for all sufficiently large $n$.
The endpoint laws of $K_*$ give an exact two-state realization.
Clipping into any nondegenerate $K_n$ gives a proper two-state law with
\[
 \UTV(p,q_n)\le
 \frac{L_1(1-s_*)}{2s_*}\gamma^n\operatorname{diam}I.
\]
At each iteration only the $2d$ endpoint images and $t_0$ are needed.
Thus $n$ iterations require $O(dn)$ exact real-arithmetic evaluations
and comparisons; no bit-complexity bound is asserted.
\end{theorem}
\begin{proof}
A continuous monotone map sends an interval to the interval between
its endpoint images. Induction therefore identifies $K_n$ with the
convex hull of the finite set of coordinates at depths at most $n$;
each new extreme is attained by a previous extreme followed by one
letter, or by the empty prefix. The interval map in
\eqref{II:U:eq:rank-two-envelope} is monotone for inclusion and is a
$\gamma$-contraction in Hausdorff distance: each image map is, and
adjoining the same point and taking interval hulls cannot increase the
maximum endpoint discrepancy. The space of nonempty compact
subintervals of $I$ is complete. The contraction theorem gives a
unique $K_*$ and the error bound. Since $K_0\subseteq K_1$, induction
gives nesting and identifies $K_*$ as the closed convex hull of all
reachable coordinates.

Rank two rules out a singleton reachable hull. Its endpoint laws are
probabilities by continuity and are invariant under all letters.
The barycentric construction and clipping proof of
Theorem~\ref{II:U:thm:rank-two-monotone} now apply on $K_*$ and $K_n$.
Prefix witnesses can be retained as endpoint backpointers; there are
at most two new witnesses per round. Exact arithmetic and certified
input hypotheses, including the interval of nonnegative laws, are
assumptions of this construction.
\end{proof}

This is a constructive theorem for a declared rank-two interval class.
It does not discover that class from noisy Hankel data, certify its
positivity interval, or establish the missing arbitrary-rank acquisition
theorem. General contraction alone gives no fixed-$R$ simplex theorem
in higher dimension.

\subsection{An exact golden-ratio clipping certificate}
\label{II:U:sec:golden-constructive}
Return to the first-coordinate convention of
Proposition~\ref{II:U:prop:golden-prefix-obstruction}. In the ordered
limiting dictionary $(f_r,f_{1-r})$ the exact matrices are
\[
 B_0=\begin{pmatrix}(3+\sqrt5)/20&0\\1/5&(3-\sqrt5)/20\end{pmatrix},
 \qquad
 B_1=\begin{pmatrix}(3-\sqrt5)/20&1/5\\0&(3+\sqrt5)/20\end{pmatrix}.
\]
The initial vector remains $(1/2,1/2)$ and both stops are $1/2$.
The matrices equal $SA_aS^{-1}$ for
$S=\left(\begin{smallmatrix}r&1-r\\1-r&r\end{smallmatrix}\right)$.
The limiting laws are not finite-prefix residuals.

\begin{proposition}[Exact two-prefix defect on the golden-ratio source]
\label{II:U:prop:golden-clip}
For $L\ge1$ set
\[
 a=F_0^L(1/2),\quad
 \gamma=\frac{7-3\sqrt5}{2},\quad h=9-4\sqrt5,\quad z_L=h\gamma^L.
\]
In the ordered dictionary $(f_a,f_{1-a})$, clipping gives
\[
 \widehat B_{0,L}=
 \begin{pmatrix}
 (2+a)/10&0\\
 (-a^2+3a-1)/(10(2a-1))&(-a^2+4a-2)/(10(2a-1))
 \end{pmatrix}.
\]
\[
 \widehat B_{1,L}=P\widehat B_{0,L}P,\qquad
 P=\begin{pmatrix}0&1\\1&0\end{pmatrix}.
\]
These matrices are nonnegative, have total transition row sum $1/2$,
and yield occupation vector $x_L=(1,1)$ from $(1/2,1/2)$.
Their exact untruncated flow defect is
\begin{equation}
 J_L^{\rm clip}
 =\frac{1-a-a^2}{25}
 =\frac{z_L}{5(1+z_L)^2},\qquad
 \UTV(p,q_L)\le\frac{z_L}{10(1+z_L)^2}.
 \label{II:U:eq:golden-exact-defect}
\end{equation}
\end{proposition}
\begin{proof}
Here $3/5\le a<r$. The inward successors lie in $[1-a,a]$;
the two outward ones are clipped to their respective endpoints.
Their barycentric coefficients give the displayed matrices.
The aggregate has both row and column sums $1/2$, hence $x_L=(1,1)$.
If $f^{[1]},f^{[2]}$ are the laws from the two original pure states,
then $(1,-1)A_0=(1/10,0)$ and $(1,-1)A_1=(0,-1/10)$.
Their difference has positive mass $1/10$ on words starting with $0$
and negative mass $-1/10$ on words starting with $1$, with zero
empty-word difference. Thus
$\|f_t-f_s\|_1=|t-s|/5$.
The outward coordinate displacement is
$(1-a-a^2)/(2+a)$ and its letter probability is $(2+a)/10$.
The two occupied rows therefore give $(1-a-a^2)/25$.
The cross-ratio formula gives
$r-a=\sqrt5\,z_L/(1+z_L)$, so
$1-a-a^2=5z_L/(1+z_L)^2$.
Theorem~\ref{II:U:thm:flows} completes the proof.
\end{proof}

At $L=1$, $J_1^{\rm clip}=1/625$ and $\UTV(p,q_1)\le1/1250$.
These are certificate values, not assertions that the actual TV
distance or the optimized flow objective equals the displayed bound.


\subsection{Observable coordinates, explicit rounding, and conditional sampling}
\label{II:U:sec:rank-two-numerical}
Fix a nondegenerate two-prefix dictionary $(f_l,f_h)$ from either
rank-two construction, with $l<h$ and $l\le t_0\le h$.
The exact clipped model is denoted by $q$, with stop probabilities at
least $s_*>0$ and defect $J^{\rm clip}$. Choose a fixed measurable
suffix event $E\subseteq\Sigma^*$ such that
$c_h=f_h(E)>c_l=f_l(E)$, replacing $E$ by its complement if necessary.
Put $\Delta_E=c_h-c_l>0$. For any coordinate $t$,
\[
 \frac{t-l}{h-l}=\frac{f_t(E)-c_l}{c_h-c_l}.
\]
If $L_E=|f_{r_+}(E)-f_{r_-}(E)|$ denotes separation of the limiting
laws, then the normalization is
\begin{equation}
 \Delta_E=L_E\,\frac{h-l}{r_+-r_-},\qquad
 D_{l,h}:=\|f_h-f_l\|_1=L_1(h-l).
 \label{II:U:eq:observable-gap-correct}
\end{equation}
The interval width cannot be silently omitted. The transition
barycentric weights are the clipped ratios for the successor
residuals $f^{(u_l a)},f^{(u_h a)}$, and the initial weight is the
ratio for $p$. Thus these quantities are observable from the fixed
query set
\[
 \mathcal W=\{\varepsilon,u_l,u_h\}
       \cup\{u_l a,u_h a:a\in\Sigma\},\qquad M=|\mathcal W|\le3+2d.
\]
All these prefixes have positive mass under the interval hypotheses.

\begin{lemma}[Stopped-automaton perturbation by coupling]
\label{II:U:lem:stopped-coupling}
Let $q,\widetilde q$ be proper finite-state word laws with initial
vectors $\alpha,\widetilde\alpha$ on the same states. Let $K_i$ and
$\widetilde K_i$ be the state-$i$ probability vectors on the finite
outcomes $\{\mathrm{STOP}\}\cup(\Sigma\times\{1,\ldots,N\})$.
If the visit vector $\nu=\sum_{n\ge0}\alpha(\sum_a B_a)^n$ is finite, then
\[
 \UTV(q,\widetilde q)
 \le\UTV(\alpha,\widetilde\alpha)
       +\sum_i\nu_i\UTV(K_i,\widetilde K_i).
\]
\end{lemma}
\begin{proof}
Maximally couple initial states and, while the states agree, maximally
couple the next finite outcome. The words agree unless a mismatch
occurs. At any visit to $i$ before mismatch, its probability is at
most $\UTV(K_i,\widetilde K_i)$. Summing these probabilities is bounded
by the expected number $\nu_i$ of visits in the first automaton.
The coupling inequality gives the assertion. This proof uses the
original visits, not visits of the perturbed chain.
\end{proof}


\begin{theorem}[A deterministic rank-two estimation and rounding certificate]
\label{II:U:thm:rank-two-rounding}
Suppose every queried event probability $f^{(u)}(E)$, $u\in\mathcal W$,
has an estimate with absolute error at most $\epsilon_E<\Delta_E/4$.
At each of the two anchors suppose the estimated vector of stop and
letter probabilities is nonnegative, sums to one, and has $\ell^1$
error at most $\epsilon_m<s_*$. Define
\[
 e=\frac{4\epsilon_E}{\Delta_E},\qquad
 h_b=2^{-b},\quad b\in\mathbb N,\ b\ge1.
\]
Form initial and successor barycentric weights from the estimated
ratios, clipped to $[0,1]$. Multiply each estimated letter mass by
its two barycentric weights to obtain a preliminary transition row.
Round each of its $2d$ transition entries down to a multiple of $h_b$,
and assign the remaining row mass to STOP. Round the initial weight
of the second state down to the same grid and define the first as
its complement. The resulting two-state law $\widetilde q_b$ is proper
and satisfies
\begin{equation}
 \UTV(p,\widetilde q_b)\le
 \frac{J^{\rm clip}}2+e+h_b+
 \frac{\epsilon_m/2+e+2dh_b}{s_*}.
 \label{II:U:eq:rank-two-rounding}
\end{equation}
No perturbation of the original occupation weights is assumed.
\end{theorem}
\begin{proof}
Let $a=c_l,b'=c_h$ and let $z$ denote any queried event probability.
If $z\in[a,b']$, write $z=(1-t)a+tb'$ with $t\in[0,1]$.
The numerator of the ratio perturbation, after multiplication by the
estimated denominator, is
$\delta z-(1-t)\delta a-t\delta b'$, of absolute value at most
$2\epsilon_E$. The denominator is at least
$\Delta_E-2\epsilon_E$. Clipping cannot increase the discrepancy.
If $z<a$ or $z>b'$, the same bound follows by comparison with the
corresponding endpoint. Thus every clipped barycentric error,
including the initial error, is at most
$2\epsilon_E/(\Delta_E-2\epsilon_E)\le e$.

For a preliminary transition row, separate the error in stop/letter
masses from the error in splitting each letter between two states.
The former has TV at most $\epsilon_m/2$ and the latter at most
$\sum_a\widehat w_a e\le e$. Downward rounding transfers less than
$2dh_b$ mass from transitions to STOP, hence changes row TV by at most
$2dh_b$. The initial TV error is at most $e+h_b$.
Because the stop/letter vectors are normalized, their stop-coordinate
error is at most $\epsilon_m/2$. Downward rounding only increases the
stop mass, so every final stop is at least
$s_*-\epsilon_m/2>0$. Properness follows.
Apply Lemma~\ref{II:U:lem:stopped-coupling}, using the original visit bound
$\nu\mathbf1\le1/s_*$, and then
$\UTV(p,q)\le J^{\rm clip}/2$.
\end{proof}

The theorem concerns represented dyadic transition probabilities and
their mathematical law. Runtime sampling bias and unaccounted
floating-point errors in evaluating the estimated ratios are not
covered; verified bounds for them must be added. SVD-denoised estimates
are eligible only after their stated probability errors and normalized
mass vectors have been independently certified.

\begin{corollary}[One-sample conditional-frequency certificate]
\label{II:U:cor:rank-two-sampling}
Fix $0<\delta<1$, the two prefix words, and $E$ before drawing $n$ independent source
words, or use an independent validation split. Define
\[
 \nu_*=\min_{u\in\mathcal W}P_p(u\Sigma^*)>0.
\]
For each query use the empirical frequency of $E$ among suffixes of
sampled words beginning with that prefix. At the two anchors use the
empirical stop/next-letter histogram. If
$n\nu_*\ge8\log(3M/\delta)$, then with probability at least $1-\delta$
all prefix counts are at least $n\nu_*/2$ and the errors satisfy
\begin{align}
 \epsilon_E&\le
 \sqrt{\frac{\log(6M/\delta)}{n\nu_*}},\notag\\
 \epsilon_m&\le
 2\sqrt{\frac{(d+1)\log2+\log(6/\delta)}{n\nu_*}}.
 \label{II:U:eq:rank-two-sample-radii}
\end{align}
Whenever these radii meet the hypotheses of
Theorem~\ref{II:U:thm:rank-two-rounding}, its conclusion holds on this event.
\end{corollary}
\begin{proof}
Multiplicative Chernoff and a union bound make the probability of any
prefix count below half its expectation at most
$M\exp(-n\nu_*/8)\le\delta/3$.
Conditional on its count, the suffixes for a fixed prefix are iid from
that residual law. Hoeffding with a union bound over the $M$ event
queries gives the first radius with failure probability at most
$\delta/3$. For a histogram with $d+1$ categories, taking a union bound
over its $2^{d+1}$ sign vectors gives
$P(\|\widehat m-m\|_1>t)\le2^{d+1}\exp(-kt^2/2)$ at count $k$.
Apply this to the two anchor histograms and use $k\ge n\nu_*/2$ to
obtain the second radius with total failure probability at most
$\delta/3$. Dependence between different prefix queries causes no
problem for these union bounds.
\end{proof}

If a required prefix count is zero, the empirical construction returns
\texttt{UNRESOLVED}. A numerical certificate also requires verified
lower bounds for prefix masses and observable separation; unchecked
plug-in estimates of these lower bounds are not covered.

\paragraph{Depth selection and its scope.}
In Theorem~\ref{II:U:thm:rank-two-monotone}, compactness gives
\[
 w_*=\min_{\substack{a\in\Sigma\\t\in[r_-,r_+]}}w_a(t)>0,\qquad
 \mu=\min_{t\in[r_-,r_+]}\min(w_{a_-}(t),w_{a_+}(t))>0.
\]
All required queries, including the extra successor letter, satisfy
$\nu_*\ge w_*\mu^L$. Also
$\Delta_E^{(L)}\ge\Delta_E^{(1)}>0$.
Put $A=d+\log((d+1)/\delta)$,
$a=\log(1/\gamma_*)$, and $c=\log(1/\mu)$.
For sufficiently large $n/A$, the declared choice
\[
 L=\left\lfloor\frac{\log(n/A)}{2a+c}\right\rfloor\ge1
\]
and the preceding results give
\[
 \UTV(p,\widetilde q_{L,b})
 =O_{\mathfrak m}\left((A/n)^{a/(2a+c)}+d\,2^{-b}\right).
\]
The constants $\mathfrak m$ include the source contraction constants,
stop margin, event separation, and $w_*$. The count and small-error
conditions must still hold; they do for sufficiently large $n$ with
these fixed constants. This is a completed-word TV statement.
It assumes the extremal letters, event, and parameter bounds needed
for choosing $L$ are supplied or certified separately. Reusing the
same data for their adaptive choice needs a further uniform event.
It proves neither empirical singular-subspace stability nor a
depth-independent $n^{-1/2}$ boundary-estimation theorem, and it gives
no unqualified KL or uniform rare-prefix prediction guarantee.

\subsection{Restoring an anchored initial state}
\label{II:U:sec:rank-two-anchored}
\begin{proposition}[Three-state anchoring with correct occupation weights]
\label{II:U:prop:rank-two-anchored}
For any clipped two-prefix model above, adjoin a state representing
$p=f_{t_0}$, start there, keep its exact stop probability $s(t_0)$,
and send each letter to the clipped barycentric mixture in $[l,h]$.
Let $\beta$ be the total row of incoming probabilities into the two
old states and retain their aggregate matrix $Q$. The occupation
vector of the resulting proper three-state automaton is
\[
 x^{\rm anch}=(1,v),\qquad v=\beta(I-Q)^{-1}.
\]
If the invariant limiting interval is $[l_*,h_*]$, let
$\delta=\max(l-l_*,h_*-h)$. The exact defect is
\begin{align}
 J^{\rm anch}
 &=L_1\sum_a w_a(t_0)
       \operatorname{dist}(F_a(t_0),[l,h])\notag\\
 &\quad+L_1\sum_{t\in\{l,h\}}v_t
       \sum_a w_a(t)\operatorname{dist}(F_a(t),[l,h]),
 \label{II:U:eq:rank-two-anchored-defect}
\end{align}
and
$\UTV(p,q^{\rm anch})\le J^{\rm anch}/2
\le L_1(1-s_*)\delta/(2s_*)$.
\end{proposition}
\begin{proof}
The new state has no incoming transitions and is visited once.
The old states satisfy $v=\beta+vQ$. All three stops are at least
$s_*$, so their total expected visits are at most $1/s_*$.
The residual calculation is the same disjoint-first-letter calculation
as in \eqref{II:U:eq:rank-two-clip-defect}, including the new row.
Each displacement is at most $\delta$, giving the bound.
\end{proof}

In general $v\ne x_L$ and $v\ne x_L-\alpha_L$; neither identity may be
substituted without proof. On the golden-ratio source, for every
$L\ge1$ the initial successors $2/5,3/5$ are inside $[1-a,a]$,
so the new-row defect is zero. Symmetry gives
$\beta=(1/4,1/4)$ and $v=(1/2,1/2)$.
Consequently $J^{\rm anch}=J_L^{\rm clip}/2$ and
\[
 \UTV(p,q_L^{\rm anch})\le\frac{z_L}{20(1+z_L)^2}.
\]
The factor improvement is a certificate for this construction, not an
optimality claim. A pure-prefix two-state model uses a mixed initial
vector; imposing the empty-prefix initial state adds the third state.


\subsection{Observable coordinates and a finite positive certificate}
\label{II:U:sec:observable-positive}
The results below distinguish a supplied positive basis, selection from a
supplied candidate pool, and discovery of such a pool. Throughout this
continuation, matrix $\|\cdot\|_\infty$ is the maximum absolute row sum,
$\|\cdot\|_{\max}$ is the maximum absolute entry, and vector norms have
their usual meanings. All positive realizations are row-normalized:
$\alpha\mathbf1=1$ and $\sum_a A_a\mathbf1+\tau=\mathbf1$.

\begin{theorem}[Observable normalization and finite positivity certificate]
\label{II:U:thm:observable-certificate}
Let $p(w)=\alpha_*A_{*,w}\tau_*$ have a minimal $R$-state positive
realization with $\tau_*\ge s_*\mathbf1>0$. Fix distinct core prefixes
$u_i$ and suffixes $v_j$, $1\le i,j\le R$, for which
\[
 H_{ij}=p(u_iv_j),\quad (M_a)_{ij}=p(u_iav_j),\quad
 h_j=p(v_j),\quad t_i=p(u_i)
\]
give an invertible $H$. Put $B_a=M_aH^{-1}$, $\beta=hH^{-1}$,
$q_i=p(u_i\Sigma^*)$, and $D_q=\Udiag(q_i)$. Then $q_i>0$ and
\[
 q=(I-\textstyle\sum_aB_a)^{-1}t,\qquad
 \widetilde B_a=D_q^{-1}B_aD_q,\quad
 \widetilde\beta=\beta D_q,\quad \widetilde t=D_q^{-1}t
\]
satisfy $\widetilde\beta\mathbf1=1$ and
$\sum_a\widetilde B_a\mathbf1+\widetilde t=\mathbf1$.
For any invertible $V$ with $V\mathbf1=\mathbf1$, define
\[
 \alpha=\widetilde\beta V^{-1},\qquad
 A_a=V\widetilde B_aV^{-1},\qquad \tau=V\widetilde t.
\]
If $\alpha\ge0$, every $A_a\ge0$, and $\tau\ge s\mathbf1$ for some
$s>0$, these finitely many inequalities certify an exact proper
$R$-state realization of $p$.
\end{theorem}
\begin{proof}
Let $P_{i,:}=\alpha_*A_{*,u_i}$ and let column $j$ of $C$ be
$A_{*,v_j}\tau_*$. Since $H=PC$ is invertible, both factors are
invertible. Thus $B_a=PA_{*,a}P^{-1}$, $\beta=\alpha_*P^{-1}$,
and $t=P\tau_*$. Row normalization and the stopping margin give
$(I-\sum_aA_{*,a})^{-1}\tau_*=\mathbf1$, hence $q=P\mathbf1$.
Every row of $P$ is nonzero and nonnegative, proving $q_i>0$.
The normalized matrix $\widetilde P=D_q^{-1}P$ is row-stochastic;
the displayed identities follow by similarity. A second similarity
by $V$ preserves every coefficient. Moreover $V^{-1}\mathbf1=\mathbf1$,
so the resulting initial vector and transition/stop rows sum to one.
The nonnegative aggregate has row sums at most $1-s$; hence it is
transient, every state terminates, and the word law has mass one.
\end{proof}

The certificate eliminates an infinite family of word-positivity tests
\emph{after} $V$ is obtained. It does not find $V$. Observable coordinates
can be signed; they are not hidden-state posterior probabilities.

\subsection{Contracting cones and finite trailing-word dictionaries}
\label{II:U:sec:trailing-word}
\begin{theorem}[Conditional finite-prefix construction under projective contraction]
\label{II:U:thm:trailing-word}
Suppose a supplied finite-dimensional proper cone $\mathcal C$ carries
positive linear letter maps and a linear map $x\mapsto f_x$ to
nonnegative summable word functions. Let $\ell(x)=\sum_w f_x(w)$ be
strictly positive on $\mathcal C\setminus\{0\}$; normalized states have
$\ell(x)=1$. Assume letter shifts and normalized updates agree with this
representation, the initial law is $p=f_{x_0}$, and every normalized
state has stopping probability at least $s>0$. Each nonzero letter map
sends $\mathcal C\setminus\{0\}$ into its interior, has Hilbert diameter
at most $\Delta<\infty$, and has projective contraction coefficient at
most $\gamma\in(0,1)$. Then for every $L\ge1$ there is a proper automaton
$q_L$, using at most $N_d(L)=\sum_{j=0}^Ld^j$ actual-prefix residuals, with
\begin{equation}
 \UTV(p,q_L)\le\frac{1-s}{s}
       \bigl(e^{\Delta\gamma^{L-1}}-1\bigr).
 \label{II:U:eq:trailing-tv}
\end{equation}
The same bound holds after keeping one representative of each vertex of
the convex hull of the residuals of length at most $L$.
\end{theorem}
\begin{proof}
Use a state for every positive-mass prefix of length at most $L$, initialized at
$\varepsilon$. At state $u$, retain the source's stopping and letter
probabilities. For $|u|<L$, route letter $a$ to $ua$. For
$u=u_1\cdots u_L$, route it to $u_2\cdots u_La$. The latter source
state and the true successor are obtained by the \emph{same} sequence
of $L$ normalized letter maps, applied to $x_0$ and $x_{u_1}$.
The first map bounds their distance by $\Delta$, and the remaining
$L-1$ maps contract it to $D=\Delta\gamma^{L-1}$.

To pass to probability laws, take cone-order bounds
$m y\le_{\mathcal C}x\le_{\mathcal C}M y$. Positivity implies
$m f_y\le f_x\le M f_y$. Normalization gives $m\le1\le M$;
from $M/m\le e^D$ one gets $M\le e^D$ and hence
$\|f_x-f_y\|_1\le2(e^D-1)$, by summing the positive part.
Limits give the same bound when order extrema are not attained.
No coordinate $\ell^1$ bound on an arbitrary cone is used.

Each row defect is at most $2(1-s)(e^D-1)$. The constructed automaton
stops with probability at least $s$ at every visit, so its occupation
sum is at most $1/s$. The occupation-flow certificate gives
\eqref{II:U:eq:trailing-tv}. For pruning, first merge duplicate residuals,
retain the distinct vertices, and express the initial residual and each
routed successor as convex mixtures of retained vertices. The same row
defect and stopping estimates apply. Keeping the empty-prefix state
instead is always permitted within the original $N_d(L)$ budget.
\end{proof}

The Hilbert contraction estimate is the classical Birkhoff theorem
\citep{gaubert2015}. For fixed $s,\Delta,\gamma$ and $d\ge2$, the state bound is
$N(\epsilon)=O(\epsilon^{-\log d/\log(1/\gamma)})$:
use $e^z-1\le ze^\Delta$ for $0\le z\le\Delta$ and choose
$L=O(\log(1/\epsilon))$. This is exponential enumeration in depth,
not an efficient rank-$R$ discovery theorem. Exact transition queries,
cone data, and contraction bounds are inputs to this construction.

\subsection{Relative Doeblin anchors and a corrected simplex inclusion}
\label{II:U:sec:doeblin-sandwich}
Here $\Delta(Z)$ denotes the convex hull of the rows of a matrix $Z$.
\begin{lemma}[Inverse-entry bound for a nearly diagonal stochastic matrix]
\label{II:U:lem:inverse-entry}
If $\Pi\ge0$, $\Pi\mathbf1=\mathbf1$, and $\Pi_{ii}\ge1-\eta$
with $0\le\eta<1/2$, then $\Pi$ is invertible,
$\Pi^{-1}\mathbf1=\mathbf1$, and
\[
 (\Pi^{-1})_{ij}\ge-\frac{\eta}{1-2\eta}.
\]
\end{lemma}
\begin{proof}
For $\eta>0$, write $\Pi=(1-\eta)I+\eta K$ with $K$ row-stochastic.
Expand the inverse in powers of $-\eta K/(1-\eta)$. Each entry of
$K^n$ belongs to $[0,1]$. Dropping the even nonnegative terms bounds
every entry below by the sum of the odd negative coefficients,
$-\eta/(1-2\eta)$. The case $\eta=0$ is immediate.
\end{proof}


\begin{theorem}[Conditional Doeblin simplex inclusion and strict margins]
\label{II:U:thm:doeblin-sandwich}
Let a proper $R$-state law have $\tau_*\ge s_*\mathbf1>0$,
$\alpha_{*,j}\ge\mu$, and
\[
 A_{*,a,ij}\ge\mu w_{a,i},\qquad
 w_a=A_{*,a}\mathbf1,\qquad 0<\mu<1/R.
\]
Set $a_\mu=1-R\mu$, $J_\mu=a_\mu I+\mu\mathbf1\mathbf1^\top$,
and $V_*=\widetilde P^{-1}$ in Theorem~\ref{II:U:thm:observable-certificate}.
Suppose a supplied true anchor matrix has the form
$U=\Pi J_\mu V_*$, where $\Pi$ satisfies
Lemma~\ref{II:U:lem:inverse-entry}. Write
$c_{\max}=\max_j\sum_i\Pi_{ij}$ and choose a known upper bound
$C\ge c_{\max}$. Define
\[
 m_- =\frac{\eta}{1-2\eta},\qquad
 m_+=\frac{\mu}{a_\mu C}.
\]
If $m_-\le m_+$, every $m\in[m_-,m_+]$ gives an exact proper positive
realization through
\[
 D_m=(1+Rm)I-m\mathbf1\mathbf1^\top,\quad
 V=D_mU=WV_*,\quad W=D_m\Pi J_\mu.
\]
In particular,
$\Delta(J_\mu)\subseteq\Delta(W)\subseteq\Delta(I)$ and
$W\tau_*\ge s_*\mathbf1$. If $m_-<m<m_+$, put
$\delta_m=\min(m-m_-,m_+-m)$. If also
$w_{\min}=\min_{a,i}w_{a,i}>0$, all initial and transition entries are
at least
\begin{equation}
 \gamma_* =\frac{\delta_m}{1+Rm}\min(1,w_{\min})>0.
 \label{II:U:eq:sandwich-margin}
\end{equation}
One may use $C=1+(R-1)\eta$. A sufficient nonempty-interval condition is
\begin{equation}
 a_\mu(R-1)\eta^2+[1-(R-2)\mu]\eta-\mu\le0.
 \label{II:U:eq:sandwich-feasibility}
\end{equation}
\end{theorem}
\begin{proof}
Let $\widetilde W=D_m\Pi$. The inverse formula and row sums give
\[
 \widetilde W^{-1}=\frac{\Pi^{-1}+m\mathbf1\mathbf1^\top}{1+Rm}\ge0.
\]
It is row-stochastic. Since $J_\mu=\widetilde W^{-1}W$, this proves
the inner inclusion. For the outer inclusion one must retain column
sums $c_j$:
\[
 \widetilde W_{ij}=(1+Rm)\Pi_{ij}-mc_j,\qquad
 W_{ij}=a_\mu\widetilde W_{ij}+\mu
       \ge\mu-a_\mu m C\ge0.
\]
Also $W\mathbf1=\mathbf1$. For every row $z$ satisfying
$z_j\ge\mu(z\mathbf1)$,
\[
 (zJ_\mu^{-1})_j=\frac{z_j-\mu(z\mathbf1)}{a_\mu}\ge0.
\]
Consequently
\[
 \alpha=(\alpha_*J_\mu^{-1})\widetilde W^{-1},\qquad
 A_a=W(A_{*,a}J_\mu^{-1})\widetilde W^{-1},\qquad
 \tau=W\tau_*
\]
are nonnegative and normalized, with the asserted stopping margin.
If $m$ is strictly inside the interval, entries of $\widetilde W^{-1}$
are at least $\delta_m/(1+Rm)$ and $W_{ij}\ge a_\mu C\delta_m>0$.
The vector $\alpha_*J_\mu^{-1}$ has mass one. The row sums of
$W A_{*,a}J_\mu^{-1}$ equal $(Ww_a)_i\ge w_{\min}$, which proves
\eqref{II:U:eq:sandwich-margin}. Finally $c_j\le1+(R-1)\eta$; substituting
this bound into $m_-\le m_+$ gives \eqref{II:U:eq:sandwich-feasibility}.
\end{proof}

The margin in \eqref{II:U:eq:sandwich-margin} uses the stochasticity of $W$;
it does not require a cross-row letter-probability ratio. The hypotheses
concern relative anchors inside $\Delta(J_\mu)$, not nearly pure hidden
states in $\Delta(I)$. They are additional structural assumptions.

\subsection{Selection from a supplied relatively separable prefix pool}
\label{II:U:sec:relative-spa}
\begin{theorem}[Conditional relative SPA selection, with corrected noise budget]
\label{II:U:thm:relative-spa}
Let $G=J_\mu V_*$ be invertible, and let a supplied pool of $M$ true
prefix-coordinate rows form $F=HG$, where $H\ge0$ and
$H\mathbf1=\mathbf1$. Suppose there are distinct witness rows $i_r$
with $H_{i_r,r}\ge1-\eta_0$. Let
$\max_i\|\widehat F_{i,:}-F_{i,:}\|_2\le e_F$ and define
\begin{align*}
 \sigma&=\sigma_R(G),\qquad \kappa=\|G\|_2/\sigma,
 & C_{\rm SPA}&=1+80\kappa^2,\\
 e&=e_F+\sqrt2\eta_0\|G\|_2,
 & \theta_R&=\min\{(2\sqrt{R-1})^{-1},1/4\}.
\end{align*}
For $R\ge2$, if $e<\sigma\theta_R/C_{\rm SPA}$, ordinary SPA on
$\widehat F^\top$ returns $R$ indices in $O(MR^2)$ exact-real arithmetic
operations. After a permutation, their \emph{true} relative coordinates
satisfy
\begin{equation}
 H_{\widehat i_r,r}\ge1-\eta,\qquad
 \eta=\frac{e_F+C_{\rm SPA}e}{\sigma}.
 \label{II:U:eq:spa-relative-budget}
\end{equation}
If $\eta<1/2$ and \eqref{II:U:eq:sandwich-feasibility} is strict, these true
selected anchors satisfy Theorem~\ref{II:U:thm:doeblin-sandwich}.
\end{theorem}
\begin{proof}
Replace the witness rows of $F$ by the corresponding rows of $G$,
leaving other rows unchanged, to obtain an exactly separable matrix
$F^0$. The replacement changes each row by at most
$\sqrt2\eta_0\|G\|_2$, since a simplex vector $h$ with
$h_r\ge1-\eta_0$ obeys $\|h-e_r\|_2\le\sqrt2\eta_0$.
Apply the squared-Euclidean specialization of
\citet[Theorem 3 and Remark 2]{gillis2014} to $(F^0)^\top$ and
$\widehat F^\top$, using the larger factor $C_{\rm SPA}$.
Its hypotheses allow signed generator coordinates. The selected
\emph{observed} rows are within $C_{\rm SPA}e$ of distinct rows of $G$.
Their true rows are within $e_F+C_{\rm SPA}e$. Multiplication by
$G^{-1}$ bounds each relative-coordinate error by
$(e_F+C_{\rm SPA}e)/\sigma$, proving \eqref{II:U:eq:spa-relative-budget}.
\end{proof}

This is an application of a classical selection theorem, not a new SPA
algorithm. Its cost starts with the pool matrix already available.
The pool size, generation cost, prefix masses, and uniform estimation
error are separate inputs. The estimated rows $\widehat F_{\widehat i_r,:}$
are not the unknown true rows $F_{\widehat i_r,:}$; replacing the latter
in $V=D_mU$ requires an additional numerical error certificate.

\begin{proposition}[Uniform observable coordinates for a fixed prefix pool]
\label{II:U:prop:pool-estimation}
Use a fixed core with $|u_i|,|v_j|\le L_0$ and
$\sigma_R(H)\ge\sigma_0>0$, and prefixes in a fixed
pool $\mathcal P\subseteq\Sigma^{\le L_0}$ with masses at least $p_*>0$.
From $n$ independent words, observe the first $2L_0$ letters together
with whether and when the word terminates by that depth. Estimate
$H$, the core prefix masses $q$, the pool prefix masses $q_w$, and
$h_{w,j}=p(wv_j)$ by counts. Put
\[
 K=R^2+R+MR+M,\quad
 t=\sqrt{\log(2K/\delta)/(2n)},\quad h_0=\sqrt R/\sigma_0.
\]
If $Rh_0t\le1/2$ and $t\le p_*/2$, then with probability at least
$1-\delta$, all estimates
\[
 \widehat f_w=(\widehat h_w/\widehat q_w)\widehat H^{-1}D_{\widehat q}
\]
exist and have row $\ell^2$ error at most
\begin{equation}
 e_F\le\frac{8Rh_0t}{p_*}+(2Rh_0^2+h_0)t.
 \label{II:U:eq:pool-estimation}
\end{equation}
\end{proposition}
\begin{proof}
A union bound and Hoeffding give error at most $t$ for all $K$ entries.
The Neumann bound gives $\|\widehat H^{-1}\|_\infty\le2h_0$ and
$\|\widehat H^{-1}-H^{-1}\|_\infty\le2Rh_0^2t$.
Both conditional rows $h_w/q_w$ and $\widehat h_w/\widehat q_w$ are
nonnegative with mass at most one, because the $v_j$ are distinct.
Each coordinate differs by at most $4t/p_*$. Telescope
$(h_w/q_w)H^{-1}D_q$; the conditional-row error contributes at most
$8Rh_0t/p_*$ in $\ell^1$, and the remaining product error contributes
at most $(2Rh_0^2+h_0)t$. The $\ell^2$ norm is no larger.
\end{proof}

The uniform event permits choosing indices by SPA on the same pool.
It does not permit choosing the pool or core after observing these
data without an additional uniform argument or sample split.

\subsection{Pathwise stability and a complete fixed-basis sampling construction}
\label{II:U:sec:fixed-basis-sampling}
The following construction uses exact rational arithmetic for intermediate
linear algebra and a specified dyadic output. It is a finite-bit algorithm,
but it is not a theorem about unrestricted $b$-bit floating-point LU.

\begin{proposition}[Pathwise perturbation and normalized initial estimates]
\label{II:U:prop:pathwise-stability}
Let $p,\widehat q$ be proper row-normalized positive automata on the same
$R$ states, with $\tau\ge s\mathbf1>0$ for the first automaton. Then
\begin{equation}
 \UTV(p,\widehat q)\le\frac12\|\widehat\alpha-\alpha\|_1
 +\frac1{2s}\left(\|\widehat\tau-\tau\|_\infty+
                   \sum_a\|\widehat A_a-A_a\|_\infty\right).
 \label{II:U:eq:pathwise-general}
\end{equation}
If $\alpha_i,A_{a,ij}\ge\gamma>0$, $R\ge2$, and an unnormalized
estimate $a^0$ satisfies $\|a^0-\alpha\|_\infty\le\delta_\alpha
\le\gamma/[2(R+1)]$, its normalization $\bar\alpha=a^0/(a^0\mathbf1)$
is positive and
\begin{equation}
 \|\bar\alpha-\alpha\|_\infty
 \le\frac{(R-1)\delta_\alpha}{1-R\delta_\alpha}\le\gamma/2.
 \label{II:U:eq:initial-normalization}
\end{equation}
For transition estimates with per-letter maximum-entry errors $e_a$,
$\max_a e_a\le\gamma/2$ and $R\sum_a e_a\le s/2$ suffice for
$\bar A_{a,ij}\ge\gamma/2$ and
$\bar\tau=\mathbf1-\sum_a\bar A_a\mathbf1\ge(s/2)\mathbf1$.
\end{proposition}
\begin{proof}
For each $w=a_1\cdots a_m$, telescope
\begin{align*}
 \widehat\alpha\widehat A_w\widehat\tau-\alpha A_w\tau
 ={}&(\widehat\alpha-\alpha)\widehat A_w\widehat\tau
       +\alpha A_w(\widehat\tau-\tau)\\
 &+\sum_{k=1}^m\alpha A_{a_1\cdots a_{k-1}}
       (\widehat A_{a_k}-A_{a_k})
       \widehat A_{a_{k+1}\cdots a_m}\widehat\tau.
\end{align*}
The suffix sum $\widehat g=\sum_v\widehat A_v\widehat\tau$ is at most
$\mathbf1$ componentwise: its finite partial sums are
$\mathbf1-\widehat T^{n+1}\mathbf1$. Equality requires termination
from every state, and is not needed here. The original occupation row
$\nu=\alpha(I-T)^{-1}$ has $\nu\mathbf1\le1/s$.
Take absolute values, apply nonnegative matrix majorants, and sum over
all prefixes, the differing letter, and suffixes. Tonelli and these
two bounds give \eqref{II:U:eq:pathwise-general}. This also follows from
Lemma~\ref{II:U:lem:stopped-coupling}.

Write $a^0=\alpha+e$. In coordinate $i$ the normalized error is
$[e_i(1-\alpha_i)-\alpha_i\sum_{j\ne i}e_j]/(1+\sum_j e_j)$.
Its numerator is at most $(R-1)\delta_\alpha$ in absolute value;
$R\gamma\le1$ proves the last inequality in
\eqref{II:U:eq:initial-normalization}. Positivity of $a^0$ and of its sum
is immediate. Finally each row's total transition error is at most
$R\sum_a e_a$; recomputing STOP as the remainder proves the last claim.
\end{proof}


\begin{theorem}[Fixed rational basis: iid estimation, exact computation, dyadic output]
\label{II:U:thm:fixed-basis-sampling}
Fix the core in Theorem~\ref{II:U:thm:observable-certificate}, with
$|u_i|,|v_j|\le L_0$, $\sigma_R(H)\ge\sigma_0>0$, and
$q_i\ge q_*>0$. A \emph{supplied rational} invertible $V$ satisfies
$V\mathbf1=\mathbf1$ and certifies the true law with
$\alpha_i,A_{a,ij}\ge\gamma>0$, $\tau_i\ge s>0$, where the
supplied lower bounds $\gamma,s$ are rational.
The core and $V$ are fixed independently of the sample (or obtained
from an independent training sample with these true margins certified).
For $0<\delta<1$, observe $n$ iid words censored at $D=2L_0+1$, retaining whether and
when termination occurs by $D$, including at $D$ itself.

Estimate $H,M_a,h$ by exact completed-word counts and $q$ by prefix
counts, each divided by $n$. All subsequent intermediate operations
in this theorem are exact rational operations. Define
\begin{align*}
 t_n&=c_0/\sqrt n,&
 c_0&=\sqrt{\tfrac12\log(4R(R+2)/\delta)},\\
 b_n&=c_1/\sqrt n,&
 c_1&=\sqrt d+\sqrt{2\log(2R^2/\delta)},\\
 h_0&=\sqrt R/\sigma_0,&v_-&=\|V^{-1}\|_\infty,
 \qquad v_+=\|V\|_\infty,\\
 K_\alpha&=v_-[(2R+1)h_0+2Rh_0^2],\\
 K_0&=\frac{2Rh_0}{q_*^2}
          +\frac{4R^2h_0^2+4Rh_0}{q_*},\\
 E_\alpha&=K_\alpha c_0,&
 E_A&=v_+v_-\left(\frac{2h_0R^2}{q_*}c_1+K_0c_0\right).
\end{align*}
It suffices to take
\begin{equation}
 n\ge\max\left\{(2Rh_0c_0)^2,(2c_0/q_*)^2,
 (4E_\alpha/\gamma)^2,
 \left(2E_A/\min(\gamma,s)\right)^2\right\}.
 \label{II:U:eq:fixed-basis-n}
\end{equation}
Check that $\widehat H$ is invertible and all $\widehat q_i>0$;
otherwise return \texttt{UNRESOLVED}. Form
\begin{align*}
 a^0&=\widehat h\widehat H^{-1}D_{\widehat q}V^{-1},
 &\bar\alpha&=a^0/(a^0\mathbf1),\\
 \bar A_a&=V D_{\widehat q}^{-1}\widehat M_a
                  \widehat H^{-1}D_{\widehat q}V^{-1},
 &\bar\tau&=\mathbf1-\sum_a\bar A_a\mathbf1.
\end{align*}
Before dividing by $a^0\mathbf1$, check that this sum is positive.
Check normalization and the margins
$\bar\alpha_i,\bar A_{a,ij}\ge\gamma/2$ and
$\bar\tau_i\ge s/2$ exactly; on failure return
\texttt{UNRESOLVED}. The supplied checking thresholds can be taken
rational. Otherwise, with $h_b=2^{-b}\le\gamma/4$, round
every transition entry down to a multiple of $h_b$ and assign the
remaining row mass exactly to STOP. Round the first $R-1$ entries of
$\bar\alpha$ down and assign the remaining mass to its last entry.
The resulting dyadic law $q_b$ has exactly $R$ positive states and,
with probability at least $1-\delta$, is returned and satisfies
\begin{equation}
 \UTV(p,q_b)\le
 \frac{e_\alpha}{1-e_\alpha}+(R-1)h_b
          +\frac{e_A+dRh_b}{s},\qquad
 e_\alpha=E_\alpha/\sqrt n,\quad e_A=E_A/\sqrt n.
 \label{II:U:eq:fixed-basis-tv}
\end{equation}
In particular the rate is
$O_{R,\sigma_0,q_*,V,s,\gamma}
(\sqrt{(d+\log(R/\delta))/n}+d2^{-b})$.
The intermediate rational algorithm has polynomial bit complexity
in $n,D,d,R,b$ and the rational input bit length of $V$;
the output uses $O(dR^2b)$ bits for its probabilities.
\end{theorem}
\begin{proof}
For fixed $(i,j)$, the events $W=u_iav_j$ over $a$ are disjoint.
Cauchy--Schwarz bounds the expected sum of absolute frequency errors
by $\sqrt{d/n}$. Changing one observation changes that sum by at most
$2/n$. McDiarmid and a union bound over $R^2$ pairs, together with
Hoeffding for the $R^2+2R$ base entries, give a common event of
probability at least $1-\delta$ on which
\begin{equation}
 \|\widehat H-H\|_{\max},\ \|\widehat h-h\|_\infty,
 \|\widehat q-q\|_\infty\le t_n,
 \qquad
 \sum_a\|\widehat M_a-M_a\|_\infty\le R^2b_n.
 \label{II:U:eq:joint-count-event}
\end{equation}
The last inequality sums over both row and column indices. A
per-pair union bound cannot replace this step by exchanging a sum
over letters with a maximum over rows.

Since $\|H^{-1}\|_\infty\le h_0$, the first two sample thresholds give
\begin{align*}
 \|\widehat H^{-1}\|_\infty&\le2h_0,&
 \|\widehat H^{-1}-H^{-1}\|_\infty&\le2Rh_0^2t_n,\\
 \|D_{\widehat q}^{-1}\|_\infty&\le2/q_*,&
 \|D_{\widehat q}^{-1}-D_q^{-1}\|_\infty&\le2t_n/q_*^2.
\end{align*}
Distinct suffixes imply $\|h\|_1,\|\widehat h\|_1\le1$.
For each fixed core prefix, the words $u_iav_j$ over $(a,j)$ are
distinct, so
$\sum_a\|M_a\|_\infty,\sum_a\|\widehat M_a\|_\infty\le R$.
Also $\|D_q\|_\infty,\|D_{\widehat q}\|_\infty\le1$.
Telescoping the three factors in $hH^{-1}D_q$ gives
$\|a^0-\alpha\|_1\le K_\alpha t_n=e_\alpha$.
Telescoping $D_q^{-1}M_aH^{-1}D_q$ and summing over $a$ gives
\[
 \sum_a\|\bar A_a-A_a\|_\infty
 \le v_+v_-\left[
 \frac{2h_0}{q_*}\sum_a\|\widehat M_a-M_a\|_\infty+K_0t_n\right]
 \le e_A.
\]
The remaining thresholds imply $e_\alpha\le\gamma/4$ and
$e_A\le\min(\gamma,s)/2$. Thus $a^0$ is positive, and its sum is
within $e_\alpha$ of one. Its normalized entries are at least
$(\gamma-e_\alpha)/(1+e_\alpha)\ge\gamma/2$;
$\|\bar\alpha-\alpha\|_1\le2e_\alpha/(1-e_\alpha)$.
The transition entries are at least $\gamma/2$ and the recomputed
stopping probabilities at least $s/2$. All exact checks therefore pass.

Downward rounding changes at most $dR$ entries in each transition row,
with total change at most $dRh_b$; STOP absorbs this mass. Initial TV
changes by at most $(R-1)h_b$. The original source has occupation
sum at most $1/s$. Apply Lemma~\ref{II:U:lem:stopped-coupling}, or
\eqref{II:U:eq:pathwise-general} with the row-remainder identity, to obtain
\eqref{II:U:eq:fixed-basis-tv}. The rounding restriction leaves initial
and transition entries at least $\gamma/4$ and never reduces STOP.

For complexity, a trie of the $O(dR^2)$ queried words and prefixes
permits counting in $O(nD+dR^2D)$ symbol operations. Counts have
$O(\log n)$ bits. Fraction-free elimination computes the rational
inverses; determinants of the integer matrices have polynomial bit
length in $R,\log n$ and the input bit length of $V$. The remaining
$O(dR^3)$ dense rational arithmetic operations, normalization, and
integer floor divisions also have polynomial bit cost. This proves
the stated polynomial bound; it does not assert a small uniform
degree or a fixed-word-size implementation. Dyadic probabilities
are stored as their integer numerators, with remainders determined
exactly, giving the output bound.
\end{proof}


\paragraph{Scope of the numerical theorem.}
The source itself need not have rational coefficients: empirical
counts and the supplied basis are rational. Exact rational
intermediates avoid an unsupported universal backward-error constant
for partially pivoted LU. The result concerns the mathematical law
encoded by the dyadic probabilities, not bitwise equivalence to a
floating-point executable. A rational certified basis with strict
margins is a substantive input. Neither its existence for every
positive realization nor its recovery from data is asserted here.
Combining noisy SPA-selected rows with basis construction needs
control of the \emph{estimated basis}, not just selection of indices.
All these distinctions persist even though the fixed-basis sampling
and output-rounding result is now complete.

\subsection{Estimated bases and discovery in a bounded-depth promise class}
\label{II:U:sec:learned-basis}
The following perturbation argument closes the estimated-basis gap
for a stated finite-dimensional source class. All matrix norms in
this subsection are maximum absolute row sums.

\begin{lemma}[A perturbed normalized basis preserves a strict positive realization]
\label{II:U:lem:learned-basis}
Suppose $V\mathbf1=\mathbf1$ realizes a proper law with
$\alpha_j,A_{a,ij}\ge\gamma>0$ and $\tau_i\ge s>0$.
Let $\widehat V=V+E$, $E\mathbf1=0$, and
$\rho=\|E\|_\infty\|V^{-1}\|_\infty<1$.
In the new basis the exact coefficients of the \emph{same} law satisfy
\begin{align*}
 \|\alpha'-\alpha\|_1&\le\rho/(1-\rho),\\
 \sum_a\|A'_a-A_a\|_\infty&\le 2R\rho/(1-\rho),\\
 \|\tau'-\tau\|_\infty&\le\rho.
\end{align*}
Thus $\rho\le\min\{1/2,\gamma/(8R),s/2\}$ suffices for
$\alpha'_j,A'_{a,ij}\ge\gamma/2$ and $\tau'_i\ge s/2$.
Also $\|\widehat V^{-1}\|_\infty\le2\|V^{-1}\|_\infty$.
\end{lemma}
\begin{proof}
Set $K=EV^{-1}$, so $\widehat V=(I+K)V$, $K\mathbf1=0$.
Then $\alpha'=\alpha(I+K)^{-1}$,
$A'_a=(I+K)A_a(I+K)^{-1}$ and $\tau'=(I+K)\tau$.
The Neumann bound gives the first and third estimates, and
\[
 A'_a-A_a=(KA_a-A_aK)(I+K)^{-1}.
\]
Nonnegativity and normalization give
$\sum_a\|A_a\|_\infty\le R$, proving the second estimate.
The row identities remain exact; the margins imply positivity and
termination. The inverse bound is the same Neumann estimate.
\end{proof}


For an estimated observable row $z$ whose true row $f$ satisfies
$f\mathbf1=1$, use the affine normalization
\[
 \mathcal N(z)=z+(1-z\mathbf1)e_1^\top.
\]
Then $\mathcal N(z)\mathbf1=1$ exactly and
$\|\mathcal N(z)-f\|_1\le2\|z-f\|_1$. This works for signed
observable coordinates and avoids an unjustified probabilistic
interpretation of those coordinates.

\begin{theorem}[Frequent-prefix discovery, estimated basis, and dyadic law]
\label{II:U:thm:bounded-pool-learning}
Fix $R\ge2$, $d$, $L_0$, a supplied core of depth at most $L_0$,
and positive rational margins. Consider the following \emph{promise
class}: the source has a minimal proper $R$-state realization with
the Doeblin hypotheses of Theorem~\ref{II:U:thm:doeblin-sandwich},
stop margin $s_*$, $w_{\min}>0$, and $0<\mu<1/R$;
$\sigma_R(H)\ge\sigma_0$, core prefix masses are at least $q_*$;
for every relative vertex there is a witness prefix of depth at most
$L_0$, mass at least $p_*$, and relative purity at least $1-\eta_0$.
The realization and witness prefixes need not be supplied.

Choose $0\le\eta<1/2$ satisfying the strict inequality
\eqref{II:U:eq:sandwich-feasibility}, and a rational $m\in(m_-,m_+)$
with $C=1+(R-1)\eta$. Set
\begin{align*}
 a_\mu&=1-R\mu,& h_0&=\sqrt R/\sigma_0,& D_0&=1+2(R-1)m,\\
 s_G&=a_\mu/\sqrt R,& S_G&=\sqrt R h_0,
 & C_0&=1+80(S_G/s_G)^2,\\
 v_-&=\frac{1+(R-2)\mu}{a_\mu(1-2\eta)},&
 v_+&=D_0h_0,&
 \gamma_0&=\frac{\min(m-m_-,m_+-m)}{1+Rm}\min(1,w_{\min}).
\end{align*}
Observe $n$ iid words censored at $2L_0+1$, including termination
flags. Put $M_0=N_d(L_0)$,
\[
 K=M_0(R+1)+R^2+2R,\qquad
 t=\sqrt{\log(4K/\delta)/(2n)},\qquad
 b_n=\sqrt{d/n}+\sqrt{2\log(2R^2/\delta)/n}.
\]
Use the frequent observed-prefix pool
$\widehat{\mathcal P}=\{w:|w|\le L_0,\widehat q_w\ge3p_*/4\}$.
Compute its rows by Proposition~\ref{II:U:prop:pool-estimation}, then apply
$\mathcal N$. The following explicit sufficient tests are imposed:
\begin{align*}
 Rh_0t&\le1/2,& t&\le\min(p_*/4,q_*/2),\\
 e_0&=16Rh_0t/p_*+(2Rh_0^2+h_0)t,& e_F&=2e_0,\\
 e&=e_F+\sqrt2\eta_0S_G,&
 e&<s_G\theta_R/C_0,\qquad
 (e_F+C_0e)/s_G\le\eta,\\
 e_V&=D_0e_F,&
 e_Vv_-&\le\min\{1/2,\gamma_0/(8R),s_*/2\}.
\end{align*}
Here $\theta_R$ is as in Theorem~\ref{II:U:thm:relative-spa}.
Run SPA, let $\widehat U$ be its selected normalized rows, and set
$\widehat V=D_m\widehat U$. For the empirical realization define
\begin{align*}
 \bar v_-&=2v_-,&\bar v_+&=v_++e_V,\\
 e_\alpha&=\bar v_-[(2R+1)h_0+2Rh_0^2]t,\\
 K_0&=2Rh_0/q_*^2+(4R^2h_0^2+4Rh_0)/q_*,\\
 e_A&=\bar v_+\bar v_-
       [2h_0R^2b_n/q_*+K_0t].
\end{align*}
Require $e_\alpha\le\gamma_0/8$,
$e_A\le\min(\gamma_0,s_*)/4$, and $h_b=2^{-b}\le\gamma_0/8$.
Use exact rational arithmetic, normalize the initial vector, verify
initial/transition entries at least $\gamma_0/4$ and STOP at least
$s_*/4$, and round downward as in Theorem~\ref{II:U:thm:fixed-basis-sampling}.
On a failed check return \texttt{UNRESOLVED}.
With probability at least $1-\delta$ all checks pass and the returned
proper law on $R$ states satisfies
\begin{equation}
 \UTV(p,q_b)\le\frac{e_\alpha}{1-e_\alpha}+(R-1)h_b
                    +\frac{2(e_A+dRh_b)}{s_*}.
 \label{II:U:eq:learned-basis-tv}
\end{equation}
For fixed structural margins and sufficiently small $\eta_0$, the rate is
\[
 O\!\left(\sqrt{\frac{d+L_0\log d+\log(R/\delta)}n}+d2^{-b}\right).
\]
Counting uses a trie; the observed pool has at most $1+nL_0$ members.
With this supplied core the algorithm has polynomial bit cost in
$n,L_0,d,R,b$ and the rational parameter bit lengths.
\end{theorem}
\begin{proof}
Apply Hoeffding simultaneously to every prefix of depth at most
$L_0$ and its $R$ completed-word suffix queries, and to the core
queries. The union bound uses $K$ even though only observed frequent
prefixes are processed. Combine it with the joint-letter event in
\eqref{II:U:eq:joint-count-event}; the two failure budgets are $\delta/2$.
Every witness is retained, and every retained prefix has true mass
at least $p_*/2$. Proposition~\ref{II:U:prop:pool-estimation}, with this
lower mass, and affine normalization give row $\ell^1$ error $e_F$.

In native coordinates $\widetilde P=D_q^{-1}P$ is stochastic.
Thus $\sigma_R(J_\mu\widetilde P^{-1})\ge a_\mu/\sqrt R=s_G$.
The suffix matrix $C$ in $H=PC$ has nonnegative rows of mass at most
one, because suffixes are distinct. Hence
$\|V_*\|_\infty=\|CH^{-1}D_q\|_\infty\le h_0$ and
$\|J_\mu V_*\|_2\le S_G$.
The conservative SPA tests therefore imply that the selected
\emph{true} rows have relative purity $1-\eta$.

Their basis $V=D_mU$ has margins $\gamma_0,s_*$ by
Theorem~\ref{II:U:thm:doeblin-sandwich}. Every true observable residual
row has $\ell^1$ norm at most $h_0$, giving $\|V\|_\infty\le v_+$.
The identities
\[
 V^{-1}=\widetilde P J_\mu^{-1}\Pi^{-1}D_m^{-1},\quad
 \|\Pi^{-1}\|_\infty\le(1-2\eta)^{-1},\quad
 \|D_m^{-1}\|_\infty=1
\]
give the claimed $v_-$ bound. Moreover
$\|\widehat V-V\|_\infty\le e_V$ and both bases send
$\mathbf1$ to $\mathbf1$. Lemma~\ref{II:U:lem:learned-basis} implies
that the true law in the empirical basis has margins
$\gamma_0/2,s_*/2$ and the displayed inverse/norm bounds.

Now apply the deterministic product estimates from the proof of
Theorem~\ref{II:U:thm:fixed-basis-sampling} on the same uniform event.
They hold for every selected basis obeying these norm bounds; no
sample-independence claim for $\widehat V$ is needed. The raw initial
and transition errors are bounded by $e_\alpha,e_A$. The small-error
tests imply all stipulated margins after initial normalization.
Downward rounding preserves properness. Pathwise perturbation with
true STOP at least $s_*/2$ proves \eqref{II:U:eq:learned-basis-tv}.
A change of exact basis itself changes no word probability; it is
used to prove positivity, not counted as an additional distribution error.

All retained prefixes occur in the observed trie. Query lookup,
matrix products and SPA require polynomially many operations in its
size and the stated parameters. Squared-norm SPA comparisons can be
performed exactly on rationals; no square roots are needed to execute
the pivots. Fraction-free algebra and dyadic floor operations give
the polynomial bit bound as in Theorem~\ref{II:U:thm:fixed-basis-sampling}.
\end{proof}

This closes frequent-prefix discovery and the numerical basis bridge
\emph{within the stated promise class and with a supplied core}.
It does not certify the latent class assumptions from finite data,
discover the core in uniformly polynomial time, or use empirical SVD.
The finite exact positivity checks certify the output automaton;
the approximation guarantee additionally uses the source-class promise.

\begin{proposition}[Bounded-depth observations cannot certify an unrestricted source]
\label{II:U:prop:censored-obstruction}
For every censoring depth $D$, the laws $\delta_{a^{D+1}}$ and
$\delta_{a^{D+2}}$ induce identical censored observations, including
termination flags through $D$, but their TV distance is one.
For $\epsilon<1/2$ and $\delta<1/2$, no procedure using only these
observations can approximate both laws within $\epsilon$ with
probability at least $1-\delta$.
\end{proposition}
\begin{proof}
Every observation is the same nonterminated $D$-letter string.
The two TV balls of radius $\epsilon$ are disjoint, by the triangle
inequality. An output distribution that is identical under the two
sources cannot give both disjoint events probability greater than
$1/2$. Arbitrarily many such censored observations do not help.
\end{proof}


\begin{proposition}[Fixed-block empirical singular-space certificate]
\label{II:U:prop:fixed-block-svd}
Let $K$ be an $M\times N$ block of completed-word probabilities,
with its query words fixed before sampling, $\Urank K=r$, and
$\sigma_r(K)\ge s_0>0$. From $n$ iid observations let $\widehat K$
be the frequency matrix and put
\[
 e=\sqrt{MN}\sqrt{\log(2MN/\delta)/(2n)}.
\]
With probability at least $1-\delta$,
$\|\widehat K-K\|_2\le e$. If $e<s_0/2$, the leading $r$ left
and right singular projectors obey
\[
 \|\widehat P-P\|_2,\ \|\widehat Q-Q\|_2
          \le e/(s_0-e),\qquad
 \|\widehat K_r-K\|_2\le2e.
\]
If $3e<s_0$, thresholding empirical singular values at $2e$ selects
exactly $r$. The population rank and gap remain assumptions.
\end{proposition}
\begin{proof}
Hoeffding and a union bound give entry error at most
$e/\sqrt{MN}$ and thus the operator bound via the Frobenius norm.
The variational singular-value inequalities give
$\widehat\sigma_r\ge s_0-e$ and
$\widehat\sigma_{r+1}\le e$.
Writing $\widehat K\widehat V_r=\widehat U_r\widehat\Sigma_r$
and multiplying by $I-P$ gives
$(I-P)\widehat U_r=(I-P)(\widehat K-K)\widehat V_r
\widehat\Sigma_r^{-1}$. The norm is at most $e/(s_0-e)$.
For two projectors of equal rank this is their operator-norm distance;
the transposed argument gives the right projector. Finally
$\|\widehat K_r-K\|\le\|\widehat K_r-\widehat K\|+e\le2e$.
\end{proof}

This controls a fixed finite block. It neither estimates an infinite
Hankel tail nor supplies a positive cone or justifies adaptive block
selection outside a declared uniform family. Validated numerical
SVD errors must be added when the computed decomposition is inexact.

\subsection{A bounded-coordinate rank-two parametric construction}
\label{II:U:sec:parametric-rank-two}
\begin{theorem}[Estimated maps and endpoint clipping on a supplied valid interval]
\label{II:U:thm:parametric-rank-two}
Suppose $f_t$, $t\in I$, is a supplied affine family of probability
laws on a compact interval, with $p=f_{t_0}$,
$\|f_t-f_u\|_1\le K_f|t-u|$, true STOP at least $s>0$, and
letter masses $w_a(t)$ and projective updates $F_a(t)\in I$.
Let estimated affine-fractional maps $\widehat F_a:I\to I$, with
no pole on $I$, be contractions of modulus
at most $\bar\gamma<1$, and let nonnegative affine estimated masses
$\widehat w_a$ satisfy uniformly
\[
 \max_a|\widehat F_a(t)-F_a(t)|\le e_F,
 \qquad \sum_a|\widehat w_a(t)-w_a(t)|\le\zeta<s.
\]
Starting from $[t_0,t_0]$, perform $J$ all-letter endpoint-hull
iterations with $\widehat F_a$. At the two resulting endpoints,
use masses $\widehat w_a$, STOP as the remainder, and route to the
barycentric representation of $\widehat F_a(t)$ clipped to the
current interval. If the interval is degenerate, use one state.
Round probabilities down with step $h_b$ and initial normalization
by a remainder. The proper law on at most two states satisfies
\begin{equation}
 \UTV(p,q_{J,b})\le
 \frac{\zeta+\frac12K_f(e_F+\bar\gamma^J\operatorname{diam}I)}{s-\zeta}
       +h_b+\frac{2dh_b}{s-\zeta}.
 \label{II:U:eq:parametric-rank-two}
\end{equation}
Consequently fixed-depth coordinate estimates with
$e_F,\zeta=O(n^{-1/2})$ and $J=O(\log n)$ give a parametric
$O(n^{-1/2}+d2^{-b})$ TV rate. The interval and its validity are inputs.
\end{theorem}
\begin{proof}
The endpoint operator, taking minimum and maximum over $t_0$ and
all endpoint images, is $\bar\gamma$-Lipschitz in the maximum norm.
For an affine-fractional map without a pole on $I$, extrema on an
interval occur at its endpoints. Thus the interval expands monotonically and the
distance from any next image to the current interval is at most
$\bar\gamma^J\operatorname{diam}I$.
Initial mixing represents $f_{t_0}$ exactly. In a row, the STOP/letter
mass discrepancies contribute at most $2\zeta$ in $\ell^1$.
Since $\sum_a\widehat w_a\le1-s+\zeta\le1$, replacing true successors
by clipped estimated successors contributes at most
$K_f(e_F+\bar\gamma^J\operatorname{diam}I)$.
Every constructed STOP is at least $s-\zeta$. The occupation bound
and Theorem~\ref{II:U:thm:flows} give the first fraction. Initial dyadic TV
is at most $h_b$; per-row dyadic TV is at most $2dh_b$.
Apply the stopped-coupling bound for the remaining terms.
\end{proof}

For fixed-depth observable matrices, a certified nonzero denominator
on $I$ makes their fractional maps locally Lipschitz in the entries;
the inverse estimates already proved in Section~\ref{II:U:sec:fixed-basis-sampling}
then give the stated parametric errors. This does not assert that
the unknown valid interval, denominator margin or contraction bound
can always be inferred from bounded-depth data. Exact limiting rays
need not be evaluated: finite endpoint iteration and clipping suffice.

\section{Projective error, word laws and intermittent filtering}
\label{II:M:sec:projective-certificates}
Relative-error certificates can avoid an explicit
smallest-probability denominator, but their hypotheses must themselves
be established. They neither construct a positive cone from a Hankel
rank bound nor turn an average word-law guarantee into uniform factual
recall. Natural logarithms are used throughout this section.

\subsection{A sharp projective-to-KL bridge}
For positive vectors of the same dimension define
\[
 d_H(x,y)=\log\max_i\frac{x_i}{y_i}
          -\log\min_i\frac{x_i}{y_i}.
\]
This is a distance between rays: independent positive rescaling of
$x$ and $y$ does not change it. On distributions with common zeros,
the same definition may be used on their common positive support.

\begin{lemma}[Non-expansivity on the positive orthant]
\label{II:M:lem:projective-nonexpansive}
If $A$ is a nonnegative rectangular matrix with no zero row, then
$d_H(Ax,Ay)\le d_H(x,y)$ for all strictly positive $x,y$.
Normalizing either output by its sum leaves the conclusion unchanged.
\end{lemma}
\begin{proof}
Put $m=\min_i x_i/y_i$ and $M=\max_i x_i/y_i$. Multiplication by $A$
preserves $my\le x\le My$, giving $mAy\le Ax\le MAy$.
The outputs are strictly positive by the no-zero-row assumption.
Their coordinate ratios thus lie in $[m,M]$.
\end{proof}


\begin{lemma}[Sharp KL envelope at fixed projective diameter]
\label{II:M:lem:sharp-projective-kl}
Let $p,q$ be probability vectors positive on a common finite support.
If $d_H(p,q)\le e$, then
\begin{align}
 D_{\rm KL}(p\|q)&\le\Psi(e)\le e^2/8,\label{II:M:eq:sharp-projective-kl}\\
 \Psi(e)&=\frac{e}{\exp(e)-1}-1
          -\log\!\left(\frac{e}{\exp(e)-1}\right)\quad(e>0),
 \qquad \Psi(0)=0.\nonumber
\end{align}
The bound $\Psi(e)$ is attained by distributions on two points for each
$e>0$. Its small-$e$ expansion and large-$e$ behavior are
\[
 \Psi(e)=\frac{e^2}{8}-\frac{e^4}{576}
                 +\frac{e^6}{25920}+O(e^8),\qquad
 \Psi(e)=e-1-\log e+o(1)\quad(e\to\infty).
\]
\end{lemma}
\begin{proof}
Let $u_i=\log(p_i/q_i)$ and, with $I\sim p$, put
$W=\max_i u_i-u_I\in[0,e]$. Since $\mathbb E_p e^{-u_I}=1$,
\[
 D_{\rm KL}(p\|q)=\log\mathbb E e^W-\mathbb EW.
\]
For a fixed mean $\mu\in[0,e]$, convexity gives
$\mathbb E e^W\le1+(\mu/e)(\exp(e)-1)$, with equality for an
endpoint-supported law. Maximizing the resulting scalar expression
over $\mu$ gives $\mu_*=1-e/(\exp(e)-1)$ and the stated $\Psi(e)$.
To see attainability, give $W\in\{0,e\}$ mean $\mu_*$, call its
probabilities $p_i$, set $b=\log\mathbb E e^W$, and define
$q_i=p_i\exp(W_i-b)$. These $q_i$ sum to one, are positive, and
$\log(p_i/q_i)=b-W_i$ has oscillation $e$.

For the quadratic bound let $K(t)=\log\mathbb E e^{tW}$.
Its second derivative is the variance of $W$ under an exponential
tilt. Any variable in $[0,e]$ has variance at most $e^2/4$, since
$\operatorname{Var}(W)\le\mathbb E(W-e/2)^2\le e^2/4$.
Integrating $K''$ with weight $1-t$ on $[0,1]$ yields
$K(1)-K(0)-K'(0)\le e^2/8$.
Expanding the displayed elementary formula gives the asserted series;
using $e/(\exp(e)-1)\sim e\exp(-e)$ gives the large-$e$ expression.
\end{proof}


\subsection{What a relative matrix certificate actually asserts}
\begin{proposition}[Sufficient relative-error certificate]
\label{II:M:prop:relative-projective-certificate}
Let $A,B$ be nonnegative matrices with the same support and no zero
rows. If the oscillation of $\log(A_{ij}/B_{ij})$ on their support is
at most $\eta$, then $\sup_{x>0}d_H(Ax,Bx)\le\eta$.
In particular, if $0\le\delta<1$ and
$|B_{ij}-A_{ij}|\le\delta A_{ij}$ at every entry, then
\[
 \sup_{x>0}d_H(Ax,Bx)
 \le\log\frac{1+\delta}{1-\delta}
 =2\delta+O(\delta^3).
\]
The coefficient-oscillation condition is not necessary.
\end{proposition}
\begin{proof}
The first premise gives constants $m,M>0$ with $mB\le A\le MB$
and $\log(M/m)\le\eta$. Multiply by $x>0$ and compare output
ratios. In the relative-error case,
$(1-\delta)A\le B\le(1+\delta)A$, which proves the stated bound.
For failure of necessity take
\[
 A=\begin{pmatrix}1&1\\1&1\end{pmatrix},\qquad
 B=\begin{pmatrix}1&2\\1&2\end{pmatrix}.
\]
Both outputs are on the ray generated by $(1,1)^\top$, so their
projective distance is zero for every positive input. The coefficient
log-ratio oscillation is nevertheless $\log2$.
\end{proof}


With only the no-zero-row assumption, even equality of supports is
not necessary for a finite uniform projective defect: the matrices
\[
 \begin{pmatrix}1&0\\1&0\end{pmatrix},\qquad
 \begin{pmatrix}0&1\\0&1\end{pmatrix}
\]
again produce the same ray for every positive input. The nonzero-column
condition in Lemma~\ref{II:M:lem:projective-column-limit} below explains why
the normalized probabilistic setting is more restrictive.

\paragraph{Additive errors and zero entries.}
For known common support, minimum active entry $m>0$ and additive error
$\varepsilon<m$, the relative certificate gives
$\eta\le\log((1+\varepsilon/m)/(1-\varepsilon/m))$.
Dependence on $\varepsilon/m$ is unavoidable in worst cases:
$A=\operatorname{diag}(m,1)$ and
$B=\operatorname{diag}(m+\varepsilon,1)$ have defect
$\log(1+\varepsilon/m)$. This is not a universal equality in order:
$B=cA$ has zero projective defect despite a nonzero additive error.
Support-preserving quantization, multiplicative error bounds or other
separately proved restrictions are substantive premises. Smoothing is
one possible modification, not a free removal of these premises.

\begin{lemma}[Column limits of a uniform projective certificate]
\label{II:M:lem:projective-column-limit}
Suppose $A,B$ have no zero rows and every column of both has positive
sum. If $d_H(Ax,Bx)\le\eta$ for every $x>0$, then corresponding
normalized columns have identical supports and projective distance at
most $\eta$ on that support. Their KL divergence is at most $\Psi(\eta)$.
\end{lemma}
\begin{proof}
Apply the premise to $x=e_j+\epsilon\mathbf1$ and normalize both
outputs. The ratios of two normalized positive vectors with projective
distance at most $\eta$ lie in $[\exp(-\eta),\exp(\eta)]$.
Taking $\epsilon\downarrow0$ shows that each normalized column is
bounded above and below by these factors times the other. Thus their
supports coincide. Cross-ratio inequalities pass to the limit on that
support, so its projective distance is at most $\eta$.
Lemma~\ref{II:M:lem:sharp-projective-kl} proves the KL assertion.
\end{proof}


\subsection{Finite-horizon word laws without a mixing premise}
Use column-state conventions in this subsection. Two models have the
same finite latent state set $\{1,\ldots,N\}$, normalized positive
initial states $x_0,\widehat x_0$, and nonnegative matrices
$T_a,\widehat T_a$ with no zero rows. Define
\[
 C_{aj}=\mathbf1^\top T_a e_j>0,\qquad
 \widehat C_{aj}=\mathbf1^\top\widehat T_a e_j>0,
 \qquad \sum_a C_{aj}=\sum_a\widehat C_{aj}=1.
\]
From latent state $j$, the joint probability of emitting $a$ and moving
to $i$ is $(T_a)_{ij}$. Write $P_T,\widehat P_T$ for the induced
probability laws on $\Sigma^T$. These are fixed-horizon laws, not a
single normalized law on all finite words of all lengths.

\begin{theorem}[Word-law KL from aligned positive-model certificates]
\label{II:M:thm:projective-word-kl}
Assume for every positive $x$ and every letter $a$ that
\[
 d_H(T_ax,\widehat T_ax)\le\eta,
 \qquad d_H(Cx,\widehat Cx)\le\eta_{\rm obs}.
\]
Let $e_0=d_H(x_0,\widehat x_0)$. Then for every integer $T\ge1$,
\begin{align}
 D_{\rm KL}(P_T\|\widehat P_T)
 &\le\Psi(e_0)+T\bigl(\Psi(\eta_{\rm obs})+\Psi(\eta)\bigr),
 \label{II:M:eq:word-kl-relative}\\
 \frac1T D_{\rm KL}(P_T\|\widehat P_T)
 &\le\frac{e_0^2}{8T}+\frac{\eta_{\rm obs}^2+\eta^2}{8}.
 \label{II:M:eq:word-kl-quadratic}
\end{align}
Neither strict letter contraction nor intermittent mixing is required.
\end{theorem}
\begin{proof}
The emission distribution at latent state $j$ is $C_{\cdot j}$.
By Lemma~\ref{II:M:lem:projective-column-limit}, its KL divergence from
$\widehat C_{\cdot j}$ is at most $\Psi(\eta_{\rm obs})$.
Given $(j,a)$, the next-state distributions are
\[
 R_{a,j}=\frac{T_ae_j}{C_{aj}},\qquad
 \widehat R_{a,j}=\frac{\widehat T_ae_j}{\widehat C_{aj}}.
\]
The same lemma bounds their KL divergence by $\Psi(\eta)$ and
ensures matching support. The joint one-step kernels $K_j(a,i)$ and
$\widehat K_j(a,i)$ therefore satisfy, by the finite conditional
probability identity,
\[
 D_{\rm KL}(K_j\|\widehat K_j)
 =D_{\rm KL}(C_{\cdot j}\|\widehat C_{\cdot j})
  +\sum_a C_{aj}D_{\rm KL}(R_{a,j}\|\widehat R_{a,j})
 \le\Psi(\eta_{\rm obs})+\Psi(\eta).
\]
Expanding the log-likelihood ratio of the joint latent-state and letter
trajectory gives its initial KL plus the sum of the conditional kernel
KL divergences. The initial KL is at most $\Psi(e_0)$, so the displayed
bound follows for the joint trajectory. Marginalizing the latent states
cannot increase KL: grouping terms with the same observed word and
applying the log-sum inequality proves this directly. Finally apply
$\Psi(e)\le e^2/8$ to all three terms.
\end{proof}


\paragraph{Scope of the stronger word bound.}
The latent coordinates are aligned and both models have the same state
set. The uniform matrix certificates are stronger information than a
small empirical Hankel residual. The theorem does not construct such
coordinates, certify an arbitrary rank reduction, or infer the premises
from finitely many observed inputs. The total KL bound grows with $T$;
its rate bound does not imply a horizon-independent TV bound, uniform
error after conditioning on rare queries, or protected-fact persistence.
Termination and a law on $\Sigma^*$ require a separately specified
stopping model; they are not obtained by summing these normalized
fixed-length laws over $T$.

\subsection{Intermittent contraction controls the filter state}
For the models above set
$F_a(x)=T_ax/(\mathbf1^\top T_ax)$ and define $\widehat F_a$
similarly. Filters are evaluated on the same observed history drawn
under $P$. Non-expansivity holds even for identity or permutation
updates. It prevents amplification of an existing discrepancy but does
not, by itself, prevent accumulation of newly injected errors.

\begin{proposition}[Block filtering and the intermediate steps]
\label{II:M:prop:intermittent-projective-filter}
Assume the uniform update defect $\eta$ above. Suppose a set of
length-$L$ words has reference filter contraction coefficient at most
$\kappa_*<1$, and its conditional probability, under $P$ after every
supported history, is at least $p_*>0$. Put
\[
 q=\sqrt{1-p_*(1-\kappa_*^2)}<1,\qquad
 r_m=q^m e_0+L\eta\frac{1-q^m}{1-q}.
\]
With $e_t=d_H(x_t,\widehat x_t)$ one has, for $0\le j<L$,
\[
 \bigl(\mathbb E_P e_{mL}^2\bigr)^{1/2}\le r_m,
 \qquad
 \bigl(\mathbb E_P e_{mL+j}^2\bigr)^{1/2}\le r_m+j\eta.
\]
For $T=ML$, a valid filter-based word-law bound, including every
intermediate step, is
\[
 \frac1T D_{\rm KL}(P_T\|\widehat P_T)
 \le\frac1{8T}\sum_{m=0}^{M-1}\sum_{j=0}^{L-1}
                  (r_m+j\eta+\eta_{\rm obs})^2.
\]
The convention $q^0=1$ also covers $q=0$.
\end{proposition}
\begin{proof}
For a realized block $w$, telescoping its approximate updates against
reference updates and using non-expansivity after each substitution
gives $e_{(m+1)L}\le\kappa(w)e_{mL}+L\eta$.
Here $\kappa(w)\le\kappa_*$ on the mixing event and at most one
otherwise. Conditioning on the past yields
$\mathbb E[\kappa(w)^2\mid\mathcal F_{mL}]\le q^2$.
The $L^2(P)$ triangle inequality therefore gives
\[
 \|e_{(m+1)L}\|_{L^2(P)}
 \le q\|e_{mL}\|_{L^2(P)}+L\eta.
\]
Iteration proves the first bound; at an intermediate step use
$e_{mL+j}\le e_{mL}+j\eta$. The two next-letter distributions have
projective distance at most $e_t+\eta_{\rm obs}$. Apply
Lemma~\ref{II:M:lem:sharp-projective-kl}, the $L^2$ triangle inequality,
and the chain rule for the observed word law to obtain the full sum.
No independence of different blocks was used.
\end{proof}


Endpoint second-moment bounds do not on their own justify dropping the
$L-1$ intermediate steps from a word-law sum; the preceding proof includes
them. In the normalized,
aligned-model setting, Theorem~\ref{II:M:thm:projective-word-kl} gives the
simpler word-law estimate without any mixing premise. Existence and
efficient verification of $L,p_*,\kappa_*$ remain additional obligations
when a uniform filter-error bound is required.


\subsection{The columnwise hypothesis actually needed for word laws}
\begin{corollary}[Columnwise word-law certificate]\label{II:M:rev13:cor:column-word}
Consider two aligned $N$-state models with nonnegative matrices
$T_a,\widehat T_a$ and strictly positive normalized initial states
$x_0,\widehat x_0$. Let
\[
 C_{aj}=\mathbf1^\top T_ae_j>0,\quad
 \widehat C_{aj}=\mathbf1^\top\widehat T_ae_j>0,\quad
 \sum_a C_{aj}=\sum_a\widehat C_{aj}=1,
\]
and set $R_{a,j}=T_ae_j/C_{aj}$ and
$\widehat R_{a,j}=\widehat T_ae_j/\widehat C_{aj}$.
Assume corresponding $R_{a,j},\widehat R_{a,j}$ have common support and
\[
 d_H(R_{a,j},\widehat R_{a,j})\le\eta,\qquad
 d_H(Ce_j,\widehat Ce_j)\le\eta_{\rm obs}
 \quad\text{for every }a,j.
\]
With $e_0=d_H(x_0,\widehat x_0)$ and $T\ge1$,
\[
 D_{\rm KL}(P_T\Vert\widehat P_T)
 \le\Psi(e_0)+T\bigl(\Psi(\eta_{\rm obs})+\Psi(\eta)\bigr).
\]
These columnwise hypotheses suffice without a uniform projective
comparison of $T_ax$ and $\widehat T_ax$ for every positive $x$.
Zero rows of $T_a$ are allowed in this corollary. This weakening concerns
word-law KL; it does not replace the uniform update-error premise used
in the intermittent filter theorem.
\end{corollary}
\begin{proof}
For latent state $j$ the joint letter/next-state kernel is
$K_j(a,i)=C_{aj}R_{a,j}(i)$. Common support and the sharp envelope give
\[
 D_{\rm KL}(K_j\Vert\widehat K_j)
 =D_{\rm KL}(C_{\cdot j}\Vert\widehat C_{\cdot j})
   +\sum_a C_{aj}D_{\rm KL}(R_{a,j}\Vert\widehat R_{a,j})
 \le\Psi(\eta_{\rm obs})+\Psi(\eta).
\]
The initial KL is at most $\Psi(e_0)$. Apply the chain rule to the
length-$T$ latent trajectory and then marginalize the latent states.
\end{proof}


\section{Counterexamples for positive acquisition and conditional bounds}

\begin{proposition}[Uniform Doeblin positivity prevents arbitrary hidden-state purification]
\label{II:U:prop:doeblin-purity-obstruction}
If $A_{a,ij}\ge\mu(A_a\mathbf1)_i$ for every letter, each
nonempty-prefix posterior has all coordinates at least $\mu$.
The same is true of a normalized positive left Perron vector of a
nonempty word operator. Thus a collection of such anchors with
diagonal entries at least $1-\eta$ requires $\eta\ge(R-1)\mu$.
\end{proposition}
\begin{proof}
For a posterior $\pi$, $(\pi A_a)_j\ge\mu\pi A_a\mathbf1$.
Normalize and apply this observation to the last letter. A Perron
state returns to itself after the word and obeys the same inequality.
The other $R-1$ coordinates sum to at least $(R-1)\mu$.
\end{proof}

In particular $\eta<\mu s_*/(4R)$ is incompatible with these anchors
for $R\ge2$. An empty-prefix anchor supplies at most one exception.
The relative coordinates in Theorem~\ref{II:U:thm:doeblin-sandwich} avoid
this contradiction by using the inner simplex $J_\mu$.

\begin{proposition}[Noncommuting rank three with only one Perron ray]
\label{II:U:prop:common-perron}
Let $\alpha=(1,0,0)$ and
\[
 A_0=\frac1{30}\begin{pmatrix}4&1&1\\1&3&1\\1&2&4\end{pmatrix},
 \qquad
 A_1=\frac1{30}\begin{pmatrix}2&3&1\\3&1&2\\1&2&3\end{pmatrix}.
\]
With $\tau=(3/5,19/30,17/30)^\top$, this is a proper rank-three
law with noncommuting strictly positive transitions. Every nonempty
word operator has the same normalized left Perron vector
$(1/3,1/3,1/3)$.
\end{proposition}
\begin{proof}
The rows normalize and all stopping probabilities are positive.
For the core $\{\varepsilon,0,1\}$ on both sides, direct rational
calculation gives $\det H=-11/50625000\ne0$. Hence the rank is three.
The $(1,2)$ entry of $A_0A_1-A_1A_0$ is $1/450$.
Both matrices have column sums $1/5$, so their products have constant
column sum $5^{-|w|}$. Strict positivity makes the uniform left
Perron vector unique for each nonempty word.
\end{proof}


\paragraph{Dilation without a valid outer inclusion.}
Take $R=2$, $\eta=1/100$, $\mu=1/4$, $s_*=1/2$,
\[
 A_0=\begin{pmatrix}1/10000&1/10000\\1/8&1/8\end{pmatrix},\quad
 A_1=\tfrac14\mathbf1\mathbf1^\top-A_0,\quad
 \Pi=\begin{pmatrix}1&0\\1/100&99/100\end{pmatrix}.
\]
Although $\eta<\mu s_*/(4R)=1/64$, setting $m=1/98$ and
$W=D_m\Pi$ gives
\[
 W=\frac1{9800}\begin{pmatrix}9899&-99\\-1&9801\end{pmatrix},
 \qquad (WA_0W^{-1}\mathbf1)_1=-\frac{113851}{49000000}<0.
\]
This disproves the claimed algebraic positivity implication even
when a nearly diagonal stochastic anchor matrix is supplied.

\paragraph{The missing column-sum factor is substantive.}
With $R=3$, $\mu=1/5$, $\eta=1/10$, $m=49/100$, and
\[
 \Pi=\begin{pmatrix}1&0&0\\1/10&9/10&0\\0&0&1\end{pmatrix},
 \qquad
 W=D_m\Pi J_\mu
   =\frac1{2500}\begin{pmatrix}2431&59&10\\208&2282&10\\-39&59&2480\end{pmatrix},
\]
the old interval $1/8<m<1/2$ holds but $W_{31}<0$.
The corrected upper endpoint is $5/11$. This is realizable with
actual one-letter anchors, not only abstract matrices: put $Z=\Pi J_\mu$,
let the columns of
\[
 Q=\begin{pmatrix}1/4&1/1000&1/1000\\
                  1/10000&1/4&1/1000\\
                  1/10000&1/1000&1/4\end{pmatrix}
\]
be $w_a$, and define $A_a=w_aZ_{a,:}$, $a=0,1,2$, with uniform
initial distribution and STOP as the remainder. Every one-letter
posterior is the corresponding row of $Z$, all Doeblin inequalities
hold, and $\tau=(187/250,7489/10000,7489/10000)^\top$.
The one-letter prefix/suffix core has rank three. Nevertheless
$(WA_0W^{-1}\mathbf1)_3=-94961/25000000$.
Choosing $m=1/4$ in the corrected interval makes all transformed
coefficients strictly positive for this same source.

\subsection{Verified interior two-state example}
For the earlier source
$A_0=\frac1{20}\left(\begin{smallmatrix}3&2\\3&3\end{smallmatrix}\right)$,
$A_1=\frac1{20}\left(\begin{smallmatrix}2&2\\2&3\end{smallmatrix}\right)$,
$\alpha=(1/10,9/10)$, and $\tau=(11/20,9/20)^\top$, take
\[
 U=\begin{pmatrix}30/59&29/59\\20/49&29/49\end{pmatrix},\quad
 D(4,1)=\begin{pmatrix}2&-1\\-4&5\end{pmatrix},\quad
 V=\frac1{2891}\begin{pmatrix}1760&1131\\20&2871\end{pmatrix}.
\]
Then
\[
 \alpha_V=(897/5800,4903/5800),\qquad
 \tau_V=(29539/57820,26059/57820)^\top,
\]
\[
 A_{V,0}=\frac1{100606800}
       \begin{pmatrix}24761923&2357717\\24727123&5420117\end{pmatrix},
 \quad
 A_{V,1}=\frac1{100606800}
       \begin{pmatrix}16461862&5627438\\16427062&8689838\end{pmatrix}.
\]
These exact matrices certify $\gamma=1/50$ and $s=9/20$.
For the separate $R=2$ sandwich example with
$\mu=1/4$, $\Pi=\left(\begin{smallmatrix}1&0\\1/100&99/100\end{smallmatrix}\right)$,
and $m=1/98$, the correct values are
\[
 W=\frac1{19600}\begin{pmatrix}14799&4801\\4899&14701\end{pmatrix},
\]
\[
 WA_0W^{-1}=\frac1{970200000000}
       \begin{pmatrix}29484656149&30074228951\\
                       90086011249&91887363851\end{pmatrix}>0,
\]
where $A_0$ is the rare-letter matrix in the preceding obstruction.

\subsection{Why repeated whole-future KL cuts cannot be added}
\begin{proposition}[Counterexample to the proposed occupation-weighted KL bound]
\label{II:U:prop:repeated-kl-cuts}
The proposed lower bound $W_0\log(K/N)$ from $K>N$ disjoint future
signatures and cumulative prefix occupations is false without
nonoverlapping information increments.
\end{proposition}
\begin{proof}
On $\{0,a,b,c\}$ let
\[
 p=\tfrac13(\delta_{a0^{100}a}+\delta_{b0^{100}b}+\delta_{c0^{100}c}).
\]
For each $k\in\{a,b,c\}$ use prefixes $k0^j$, $0\le j\le100$,
and signature event $E_k=\{0^mk:m\ge0\}$. The 101 constant-length
cuts are prefix-free, the signatures are disjoint, and their
conditional probabilities equal one. Thus the proposed $N=2$, $K=3$
bound is $101\log(3/2)>40.95$. The one-state law with STOP $1/103$,
letter-$0$ probability $100/103$, and each other letter probability
$2/309$ is also representable by two identical states. Its divergence is
\[
 3\log103+100\log(103/100)+\log(3/4)<16.58.
\]
This contradicts the proposed bound. Each cut separately admits a
conditional KL lower bound, but its full future contains information
already present at the other cuts.
\end{proof}

If there are $M$ such cuts, summing their separate valid inequalities
and dividing by $M$ is legitimate. Adding them without this factor
requires signatures measurable in successive disjoint prediction
blocks and a corresponding chain rule. Neither version establishes
a general spectral KL discontinuity or directed conductance theorem.


\chapter{Executable Predictive Memory}\label{sec:memory}

\section{Which memory is being minimized?}
A rank-$r$ linear realization stores $r$ real coordinates. This statement
alone does not determine how many bits an executable approximation needs,
how much immutable code it uses, or how expensive its update and readout
are. The purpose of this section is to separate those quantities and give
exact results in interfaces where they can be compared rigorously.

The equality between finite Hankel rank and minimum linear realization
dimension is classical~\citep{carlyle1971,fliess1974}.
The underlying linear convention is
$f(w)=\alpha^\top A_w\beta$, with row state
$s(w)=\alpha^\top A_w$. Words are read chronologically; geometric
address strings are stored newest first and explicitly reversed for
replay. Logarithms used for bit counts are $\log_2$. Natural logarithms
appearing in pressure or analytic estimates are indicated by $\log$.
We distinguish recurrent dimension, persistent identifier bits,
immutable program size, workspace, numerical output precision and time.
An uncharged clock is not allowed in autonomous finite-state claims.

The operator norm in Section~\ref{sec:hankel} and the completed-word
divergences in Section~\ref{sec:probability} measure different errors
from uniform conditional prediction. We use the shared foundations in
Subsections~\ref{I:U:sec:preliminaries}, \ref{I:S:sec:exact},
and~\ref{I:S:sec:analytic} for the constructions below.

\section{Exact residual state behind a suffix window}
Fix a window length $N$ and write $q_N(w)=\operatorname{suff}_N(w)$,
retaining the whole word when $|w|<N$. Define
\[
 U_N=\operatorname{span}\{s(u)-s(v):q_N(u)=q_N(v)\}.
\]
The space $U_N$ consists of differences the external window fails to
identify. Its transition closure, rather than the ambient Hankel rank,
determines the residual state in the fixed affine interface
\[
 z(wa)=z(w)B_a+g(q_N(w),a),\qquad
 f(w)=z(w)c+h(q_N(w)).
\]
Here the matrices $B_a$ and readout $c$ are fixed, while the window and
forcing tables are supplied externally.

\begin{theorem}[Residual dimension for a fixed affine interface]
For a minimal realization,
\[
 U_N=\sum_{|q|=N}U_0A_q,\qquad
 U_{N+1}=\sum_a U_NA_a\subseteq U_N.
\]
The first plateau persists. The minimum recurrent dimension in the
specified interface is $\dim U_N$. The minimum Loewy length is that
of the actual letter-marked transition algebra acting on $U_N$.
\end{theorem}
Theorems~\ref{III:S:thm:filtration} and \ref{III:S:thm:window} give the exact
construction and lower bound. For the latter, differences within a
window fiber cancel both forcing and direct readout. Equality of all
future outputs then induces a surjective letter-intertwining map onto
the canonical difference space. Choosing its basis attains the bound.

If the homogeneous update and readout may depend on the window, the
correct interface has per-window dimensions $r_q$ and rectangular
$B_{q,a}\in\mathbb K^{r_q\times r_{q'}}$. Its charged physical
quantity is $\max_q r_q$, since only one window is active. Defining
$F_N(q)=\operatorname{span}\{s(u)-s(v):q_N(u)=q_N(v)=q\}$, the
minimum is $\max_q\dim F_N(q)$, with a marked path-image algebra
rather than the fixed-letter algebra. Definition~\ref{III:S:def:fiber-interface}
and Theorem~\ref{III:S:thm:fiber} make this contract explicit. The two
interfaces optimize different quantities.

An invertible neutral letter may preserve every $U_N$, even when
another letter has a nilpotent action. Thus single-letter forgetting does
not certify a memory saving under arbitrary ordered input sequences.
Conditional future laws introduce another distinction: normalization can
preserve differences that decay in the unnormalized linear state.

\section{Persistent bits and exact finite-state examples}
For a declared output error contract, an autonomous machine consists of
$(Q,q_0,(U_a),D)$. The charged persistent identifier length is
$B=\lceil\log_2|Q|\rceil$. Two histories with a future output difference
strictly larger than $2\epsilon$ cannot end in the same state of an
$\epsilon$-accurate machine. This classical packing principle becomes an
exact finite result only after a matching executable construction is given.

For the explicit rank-two family $f_a$ of
Section~\ref{III:M:sec:memory}, the manuscript proves
\[
 B_{f_a}^*(2^{-((a+1)m+2)})=m+2,
 \qquad a,m\in\mathbb Z,\quad a\ge1,\quad m\ge a+2.
\]
This is exact in integer bits. It does not identify the exact minimum
number of states; distinct state counts can have the same ceiling of
their base-two logarithm. The construction packs its clock into the state
identifier. The leading memory coefficient varies with $a$ although
Hankel rank stays two, so rank alone cannot determine that coefficient.

For the unequal-contraction system
\[
 \tau_0(x)=x/4,\qquad \tau_1(x)=x/2+1/2,\qquad h(x)=x,
\]
let $\theta=2^{-\ell}$ and let $\mathcal L_\ell$ be the stopped
newest-first address leaves. Put
$U_\ell=\bigcup_{v\in\mathcal L_\ell}\tau_v([0,1])$.
The all-history contract includes the empty history and uses no external clock.

\begin{theorem}[Fibonacci optimum on the leaf union]
For $\ell\ge1$, initial point $x_0\in U_\ell$, and error at most
$\theta/2$ at every finite history,
\[
 |Q|_{\min}=F_{\ell+2},\qquad
 B_{\min}=\lceil\log_2F_{\ell+2}\rceil.
\]
For every $x_0\in[0,1]$, the same state lower bound holds.
An exact count outside $U_\ell$ is not established here.
\end{theorem}
Theorem~\ref{III:M:rev13:thm:fib-general} proves this statement. Stopped
leaves supply the finite automaton and midpoint outputs. A family of
same-length histories approaching separated leaf endpoints forces the
matching lower bound. The initial-point restriction is necessary for the
stated leaf-machine initialization and is not removed by an asymptotic
dimension calculation.

\section{Rational addresses and irreversible modes}
For rational contractions, a stopped address can be ranked and unranked
without storing an exponential transition table. Exact integer counting
controls the persistent identifier. The proofs charge update arithmetic,
startup and workspace separately; a bit saving does not automatically
change byte-aligned resident memory.

Now let a finite mode obey $r_{t+1}=G(r_t,a_{t+1})$, allowing irreversible
resets, and let the contraction ratio depend on both letter and previous
mode. Storing only the retained letters can lose the information needed
to replay their dynamics. The construction stores a checkpoint
$b_t=r_{t-|v_t|}$ at the beginning of the retained window, together with
its newest-first address $v_t$. Forward replay recovers the applied mode
sequence. The algorithm then truncates using the actual product of
ratios and advances the checkpoint through discarded oldest letters.
It never requires inversion of $G$.

The pressure matrix is
\[
 M(s)_{m,m'}=\sum_{a:G(m,a)=m'}r_{a,m}^{s}.
\]
For the strongly connected mode graph, the positive root
$\rho(M(s^*))=1$ controls the symbolic address count. Theorem~\ref{III:M:rev13:thm:mode}
gives the construction and its conditional optimality statement with
all geometric and future-observability assumptions retained. Symbolic
pressure need not equal geometric dimension under arbitrary overlaps;
Hausdorff dimension by itself does not supply a finite executable with
constant overhead.

Numerical guards are sound only when the certified state enclosure lies
on one side of the decision threshold. If it intersects both sides, the
result is inconclusive. Refinement does not guarantee eventual resolution
at a genuinely ambiguous point. The feedback results certify cache
choices at the time of use and distinguish bounded returned errors from
unconditional discovery of persistent semantic facts.

\section{Execution, transfer approximations and actual coding cost}
An integer-execution certificate propagates roundoff through an explicit
invariant error bound. For transfer operators, the real branches must
preserve $[-1,1]$ and satisfy the stated complex two-radius bounds.
Repeated coefficient truncation then has the finite-horizon guarantee
\[
 \|\mathcal L_\sigma^n1-P_n\|_\infty
 \le n\delta_m K^{n-1},
 \qquad P_{n+1}=\Pi_m\mathcal L_\sigma P_n.
\]
This includes errors accumulated after each truncation; replacing
$P_n$ by a single final truncation is invalid in general.
Theorem~\ref{III:U:thm:oracle} and the rounding theorem state the additional
positivity, normalization and arithmetic hypotheses. Sound finite-horizon
bounds do not imply that a proposed compiler always terminates or that
its total bit complexity is optimal.

For recurrent matrix states, an exact main-state-plus-residual identity
is useful only after the serialized representation is charged. A smaller
residual norm, covariance or rank is not itself a smaller lossless code.
The final application develops reversible integer lifting with an actual
coding contract and separates complete encoded length from execution
workspace and update time. The application includes a controlled example
and states the conditions needed for a measured coding gain.

\section{Open scope and proof organization}
Open problems include exact Fibonacci counts outside the specified leaf
union, general nonconformal memory laws with effective updates, discovery
of protected facts from observations, and compiler completeness with a
sufficient precision rule and full cost bounds. A demonstrated advantage
in complete lossless code length and execution time on independent
recurrent trajectories is also still required for the codec application.

The remaining subsections cover exact residual interfaces,
conditional-future analysis, persistent-bit and address constructions,
mode-aware replay and pressure bounds, numerical execution certificates,
and the recurrent-state coding application. Every theorem is tied to its
own interface and error contract.
\label{III:III:main-end}

\section*{Detailed proofs and supporting results}

\section{Exact memory with an external suffix window}\label{III:S:sec:window}
This section is algebraic and does not assume bounded $H_f$. In a minimal row-state realization over $\SK$, let $q_N(w)=\Ssuff_N(w)$; for $|w|<N$, retain the entire word. The external window uses the original alphabet. Define
\[
 U_N=\SSpan\{s(u)-s(v):q_N(u)=q_N(v)\}.
\]
\begin{theorem}[Canonical difference filtration]\label{III:S:thm:filtration}
Each $U_N$ is transition-invariant, and
\[
 U_0=\SSpan\{s(w)-s(\Seps)\},\qquad
 U_N=\sum_{|q|=N}U_0A_q,\qquad
 U_{N+1}=\sum_aU_NA_a\subseteq U_N.
\]
The first plateau is permanent. If $\tau=\min\{N:U_N=U_{N+1}\}$, then $\tau\le\dim U_0-\dim U_\tau$. Also $U_N=0$ if and only if $f$ depends only on its $N$-suffix.
\end{theorem}
\begin{proof}
Equal windows remain equal after a common letter, proving invariance. Distinct words with a common full window have the form $xq,yq$, $|q|=N$, and their state difference is $(s(x)-s(y))A_q$. Shorter window fibers are singletons. This proves the closed form; group words by their last letter for the recursion. Nesting follows directly from equality of longer suffixes implying equality of shorter ones. A plateau persists under the same recursion, and each preceding proper inclusion loses at least one dimension. If $U_N=0$, outputs agree on fibers. Conversely if current outputs factor through the window, then outputs also agree after every continuation; column-span minimality makes the state difference zero.
\end{proof}

The initial space is the invariant closure of the finitely many rows $\alpha^\top(A_a-I)$. Indeed
\[
 s(a_1\cdots a_t)-\alpha^\top
 =\sum_{i=1}^t\alpha^\top(A_{a_i}-I)A_{a_{i+1}}\cdots A_{a_t}.
\]
Thus the entire filtration is computable by exact basis closure and row reduction, without enumerating all histories.

\begin{definition}[Charged affine interface]\label{III:S:def:affine-interface}
An order-$N$ implementation has a recurrent row $z(w)\in\SK^{1\times r}$ and
\[
 z(wa)=z(w)B_a+g(q_N(w),a),\qquad
 f(w)=z(w)c+h(q_N(w)).
\]
The matrices $B_a$ and recurrent readout $c$ are fixed; the forcing and direct-readout tables are unrestricted. The cost charges recurrent dimension and the actual unital algebra generated by $B_a$, while treating the window and tables as externally supplied. It does not count table bytes, arithmetic, or numerical precision.
\end{definition}

\begin{theorem}[Exact residual minimum and marked quotient]\label{III:S:thm:window}
Let $C_N$ be the image of the transition algebra on $U_N$, with $C_N=0$ when $U_N=0$. Within the stated interface the minimum recurrent dimension is $\dim U_N$, and the minimum homogeneous Loewy length is $\SLL(C_N)$. Every exact competitor has a letter-marked generated-algebra surjection onto $C_N$. One construction attains all these lower bounds and every graded radical dimension simultaneously.
\end{theorem}
\begin{proof}
Choose a row-basis matrix $E_N$ for $U_N$, and define $C_{N,a}$ by $E_NA_a=C_{N,a}E_N$. The coordinate relation
\[
 z(w)E_N=s(w)-s(q_N(w))
\]
is well-defined. For $q'=\Ssuff_N(qa)$, the forcing is defined uniquely by
$g(q,a)E_N=s(q)A_a-s(q')$, which belongs to $U_N$. With $c=E_N\beta$ and $h(q)=s(q)\beta$, this gives an exact attaining implementation.

For any competitor, let $V_N$ span its state differences within window fibers. Future windows match, so forcing and direct readout cancel:
\[
 f(ut)-f(vt)=(z(u)-z(v))B_tc.
\]
Send each competitor difference to $s(u)-s(v)$. A linear relation among competitor differences gives a target row annihilating every $A_t\beta$, hence zero by column-span minimality. The map $V_N\twoheadrightarrow U_N$ is therefore well-defined, onto, and intertwines all letters. This proves the dimension bound and a representation-kernel inclusion, hence the generated-algebra quotient. Apply Theorem~\ref{I:S:thm:marked} for the radical-layer bounds.
\end{proof}


\begin{theorem}[Saturation, one-symbol bound, and invertible padding]\label{III:S:thm:windowbounds}
Writing $\lambda_N=\SLL(C_N)$, one has
\[
 0\le\lambda_N-\lambda_{N+1}\le1,\qquad
 U_NJ(\ScA_f)^j\subseteq U_{N+j}.
\]
For nonzero $f$, $0\le\SLL(\ScA_f)-\lambda_0\le1$. If one transition is invertible, $U_N=U_0$ for every $N$. For a faithful finite-dimensional right $A$-module $M$, $L=\SLL(A)$, and $0\le r\le L$, the image algebra on $MJ(A)^r$ has Loewy length exactly $L-r$.
\end{theorem}
\begin{proof}
Restriction gives epimorphisms $C_N\twoheadrightarrow C_{N+1}$. The positive-word ideal $A_+=\SSpan\{A_w:|w|\ge1\}$ has quotient zero or $\SK$, so $J(A)\subseteq A_+$. Invariance and the recursion give $U_NA_+=U_{N+1}$, hence $U_NJ(A)\subseteq U_{N+1}$. Iteration proves the inclusion, and if $U_{N+1}J(A)^k=0$, then $U_NJ(A)^{k+1}=0$. This proves the one-symbol bound, including the zero endpoint. The quotient $M/U_0$ is at most one-dimensional and every letter acts as the identity; it is semisimple. Thus $MJ(A)\subseteq U_0$, giving the centering bound by faithfulness. If $A_a$ is invertible, its restriction to invariant $U_0$ is onto; hence $U_1\supseteq U_0A_a=U_0$ and the plateau persists. Finally, on $MJ^r$ the $k$th radical power of the image vanishes exactly when $MJ^{r+k}=0$, or equivalently $J^{r+k}=0$ by faithfulness. The least such $k$ is $L-r$, with the zero endpoint treated separately.
\end{proof}


\begin{example}[The neutral-letter counterexample]
For $f(w)=\mathbf1\{\#a(w)=1\}$ over $\{a,b\}$, take $A_a=\begin{psmallmatrix}0&1\\0&0\end{psmallmatrix}$, $A_b=I$, initial row $(1,0)$ and final column $(0,1)^\top$. The rank and Loewy length are two, and $U_0=\SR^2$. The invertible neutral letter forces $U_N=\SR^2$ at every order. Histories $b^N$ and $ab^N$ already have the same raw window and different current outputs. Finite algebraic nilpotence in one letter does not imply a saving from a raw token window when another letter can pad indefinitely.
\end{example}

\begin{example}[Sharp decrease and the limits of radical tails]
The unary length-$m$ impulse has $\dim U_N=\lambda_N=\max(m+1-N,0)$, since $U_0$ is the whole nilpotent shift space and $U_N=U_0A^N$. This attains the one-symbol bound. Adding $(-1)^{|w|}$ preserves a semisimple coordinate in every $U_N$: the dimension becomes $\max(m+1-N,0)+1$ and the Loewy cost is $\max(m+1-N,1)$. No positive radical tail contains that semisimple coordinate. Thus the exact radical-tail restriction theorem does not assert that such tails optimize every window problem.
\end{example}

\subsection{Window-conditioned homogeneous dynamics}
If the homogeneous matrices and recurrent readout may also depend on the window, define fibers
\[
 F_N(q)=\SSpan\{s(u)-s(v):q_N(u)=q_N(v)=q\}.
\]
They are zero for $|q|<N$, and equal $U_0A_q$ for $|q|=N$. Letters map $F_N(q)$ into $F_N(\Ssuff_N(qa))$. Let $\mathcal P_N$ be the image of the resulting path algebra on $\bigoplus_q F_N(q)$, including the vertex idempotents.

\begin{definition}[Window-indexed affine interface]
\label{III:S:def:fiber-interface}
Let $\mathcal Q_N=\{q_N(w):w\in\Sigma^*\}$, with successor
$q'=q_N(qa)$. Choose a dimension $r_q\ge0$ for each vertex $q$.
The recurrent row after history $w$ belongs to
$z(w)\in\mathbb K^{1\times r_{q_N(w)}}$, and the update/readout are
\[
 z(wa)=z(w)B_{q,a}+g_{q,a},\qquad
 f(w)=z(w)c_q+h_q,\qquad q=q_N(w),
\]
where $B_{q,a}\in\mathbb K^{r_q\times r_{q'}}$,
$g_{q,a}\in\mathbb K^{1\times r_{q'}}$ and
$c_q\in\mathbb K^{r_q\times1}$. The active window is external.
We charge peak recurrent dimension $\max_q r_q$ and, separately,
the Loewy length of the path-image algebra on the direct sum of
within-window reachable difference spaces, with vertex idempotents
and letter arrows marked. Window and forcing-table storage is not
included. This extends Definition~\ref{III:S:def:affine-interface} by
allowing the homogeneous update and readout to depend on the window.
\end{definition}

\begin{theorem}[Fibered minimum]\label{III:S:thm:fiber}
Within Definition~\ref{III:S:def:fiber-interface}, the minimum peak recurrent dimension is $\max_q\dim F_N(q)$. The minimum Loewy length of the actual marked path-image algebra is $\SLL(\mathcal P_N)$. These are different objectives from Theorem~\ref{III:S:thm:window}.
\end{theorem}
\begin{proof}
Within each fiber, the difference map used in Theorem~\ref{III:S:thm:window} is well-defined and onto, since future window paths and their forcing agree. The direct sum of these maps intertwines every arrow and every vertex idempotent, giving a path-algebra quotient and its radical lower bounds. Choosing bases for the canonical fibers attains the bounds, with translations supplied by the reference rows $s(q)$. Only one vertex is active in an execution, so the physical dimension charged here is the maximum, although the algebra acts on a direct sum. If desired, each fiber can be embedded into that common peak dimension; unused coordinates are set to zero.
\end{proof}

These different algebra costs need not decrease together. For the neutral-letter example at $N=1$, order the fiber coordinates as one coordinate at $a$ and two at $b$. The path radical is spanned by $E_{13},E_{21},E_{23}$; its square contains $E_{23}$ and its cube is zero. Its Loewy length is three, although the fixed-letter cost is two. This is not an extra saving or an extra physical register forced by the fixed-letter theorem: the marked algebra being charged has changed.

\section{Finite-energy forgetting on the difference module}\label{III:S:sec:metric}
For a real finite realization define the unnormalized future-energy diameter
\[
 \omega_N(f)=\sup_{q_N(u)=q_N(v)}
       \left(\sum_{t\in\Swords}|f(ut)-f(vt)|^2\right)^{1/2},
\]
allowing $+\infty$. It is nonincreasing in $N$ and vanishes exactly for suffix-definite behavior. If $H_f,H_g$ are bounded, then
\begin{equation}\label{III:S:eq:omegalipschitz}
 |\omega_N(f)-\omega_N(g)|\le\sqrt2\Snorm{H_f-H_g}.
\end{equation}
Indeed each row difference is tested by $e_u-e_v$, of norm $\sqrt2$ when $u\ne v$. The constant is sharp: a unary impulse of length $m\ge N+1$ has operator norm one and two such residual rows at distance $\sqrt2$.

Choose a basis matrix $E$ of $U_0$, and define $EA_a=C_aE$, $b=E\beta$, $g_aE=\alpha^\top(A_a-I)$. Center states by $s(w)=\alpha^\top+z(w)E$. Then
\[
 z(wa)=z(w)C_a+g_a,\qquad z(xq)=z(x)C_q+z(q).
\]
\begin{theorem}[Correct analytic domain for forgetting]\label{III:S:thm:diffstable}
For a minimal finite realization with $U_0\ne0$, these are equivalent:
\[
 \omega_0<\infty;\quad
 G_U=\sum_tC_tbb^\top C_t^\top\text{ finite};\quad
 \rho(\mathcal T)<1;\quad
 \exists R\succ0:R-\mathcal T(R)\succ0,
\]
where $\mathcal T(X)=\sum_aC_aXC_a^\top$. Then $G_U\succ0$ and
\[
 \omega_N^2=\max_{|q|=N}\sup_{x,y}
 (z(x)-z(y))C_qG_UC_q^\top(z(x)-z(y))^\top.
\]
No boundedness of the full $H_f$ is assumed. For $U_0=0$ all diameters vanish.
\end{theorem}
\begin{proof}
Choose a basis from the realized state differences spanning $U_0$. Finite $\omega_0$ makes each of its future sequences square summable, and finite-dimensional linearity gives a finite Gramian. Minimal future observability on $U_0$ makes it positive definite. Its tail identity $\mathcal T^j(G_U)\to0$ and the order-unit argument of Theorem~\ref{I:S:thm:stability} prove strict positive-map stability. Conversely strict stability yields $R=I+\mathcal T(R)$. Every $C_a$ is a strict contraction in the row norm $\Snorm x_R^2=xRx^\top$, with a common factor below one. The affine forcing belongs to a finite set, so all centered states are uniformly bounded by a geometric series. The finite Gramian then bounds their future diameter. Finally a common full window $q$ propagates a difference by $C_q$, proving the formula.
\end{proof}


\begin{theorem}[Finite computation with a certified remainder]\label{III:S:thm:finiteomega}
Choose $R\succ0$, $0\le\eta<1$, $c>0$ such that $C_aRC_a^\top\preceq\eta^2R$ and $G_U\preceq cR$. Put $M=\max_a\Snorm{g_a}_R/(1-\eta)$ and
\[
 D_{N,L}=\max_{\substack{|q|=N\\|x|,|y|\le L}}
             \Snorm{(z(x)-z(y))C_q}_{G_U}.
\]
Then
\[
 \omega_N\le2\sqrt c M\eta^N,\qquad
 D_{N,L}\le\omega_N\le D_{N,L}+2\sqrt c M\eta^{N+L}.
\]
With rational data this computes $\omega_N$ to arbitrary positive accuracy by finite exact arithmetic and certified bounds on square roots.
\end{theorem}
\begin{proof}
Induction gives $\Snorm{z(w)}_R\le M$. Replacing a long history by its last $L$ letters changes its centered state by at most $M\eta^L$. Replace both histories and propagate the common $N$-word, then dominate the future norm by $\sqrt c\Snorm\cdot_R$. This proves the remainder. Every finite pair is an admissible pair for the infinite supremum, proving the lower bound. Exact Stein systems give rational $R,G_U$; rational matrix inequalities certify upper constants, and $D_{N,L}^2$ is a finite maximum of rational numbers. Use $\eta^0=1$, including $\eta=0$.
\end{proof}


\begin{theorem}[Exact root rate and finite extinction]\label{III:S:thm:jsr}
Under Theorem~\ref{III:S:thm:diffstable}, let $r=\widehat\rho(\{C_a\})$ be the joint spectral radius. Then $0\le r<1$ and $\lim_N\omega_N^{1/N}=r$. If $0<r<1$ and $N_\epsilon=\min\{N:\omega_N\le\epsilon\}$, then
\[
 \lim_{\epsilon\downarrow0}\frac{N_\epsilon}{\log(1/\epsilon)}
  =\frac1{\log(1/r)}.
\]
For any finite-dimensional real or complex difference family, $r=0$ is equivalent to simultaneous vanishing of all sufficiently long products, to $U_N=0$ for some $N$, and to finite suffix-definiteness. Extinction occurs by $\dim U_0$ when nonzero.
\end{theorem}
\begin{proof}
Put $b_N=\max_{|q|=N}\Snorm{C_q}_{G_U}$. The defining state differences span, so choose a row basis $d_i$ among them. There is a fixed $\kappa$ such that $\Snorm T_{G_U}\le\kappa\max_i\Snorm{d_iT}_{G_U}$ for all $T$, by coordinate norm equivalence. Thus
$\kappa^{-1}b_N\le\omega_N\le\omega_0b_N$. Submultiplicativity gives $b_N^{1/N}\to r$, proving the rate. The common contraction gives $r<1$. Sandwich $\omega_N$ eventually between $r_-^N$ and $r_+^N$ for $r_-<r<r_+$ and use monotonicity to obtain the inverse logarithmic limit.

For extinction, $r=0$ implies every nonempty word product is nilpotent by repeating that word. In the generated algebra the positive-word ideal $A_+$ consists of nilpotent elements: traces of every power of $X\in A_+$ vanish because they are sums of traces of positive word products. Newton identities and Cayley--Hamilton imply $X^d=0$ in dimension $d$. For any $Y\in A$, $YX\in A_+$, so $I-YX$ is invertible. The radical criterion gives $A_+\subseteq J(A)$, hence $A_+$ is a nilpotent ideal. All sufficiently long products vanish. The descending chain $W_{j+1}=\sum_aW_jC_a$, $W_0=\SK^{1\times d}$, cannot have a nonzero plateau before zero because that would contradict this nilpotence. There are at most $d$ strict drops. The equivalence with $U_N=0$ and suffix-definiteness follows from Theorem~\ref{III:S:thm:filtration}. The reverse implication to zero JSR is immediate.
\end{proof}

This rate comparison is a specialization of finite-dimensional realization and switched-system stability, not a new general theorem about the JSR. For any $r<\eta<1$ there is a norm $\Snorm x_\eta=\sup_{j,w:|w|=j}\eta^{-j}\Snorm{xC_w}_{G_U}$ with contraction factor $\eta$. A finite version suffices once a block length $m$ satisfies $C_wG_UC_w^\top\preceq\eta^{2m}G_U$ for every $|w|=m$: take the maximum only over lengths $0,\ldots,m-1$. Multiplying once shifts each term to the next length, and the last length is controlled by the block inequality. This need not be a quadratic norm.

\begin{example}[Persistent drive; exact memory versus metric forgetting]
Let $z(wa)=\rho z(w)+1$, $z(wb)=\rho z(w)-1$, $z(\Seps)=0$, with $0<\rho<1/\sqrt2$ and $f=z$. The supremum distance between states is $2/(1-\rho)$, and $G_U=1/(1-2\rho^2)$. Hence
\[
 \omega_N=\frac{2\rho^N}{(1-\rho)\sqrt{1-2\rho^2}}.
\]
The full Hankel is unbounded since $f(a^n)\to1/(1-\rho)\ne0$. The homogeneous difference dimension remains one despite metric forgetting. In contrast, $C_a=C_b=3/4$ has JSR $3/4<1$ but all-word positive-map radius $9/8>1$. The JSR stability domain cannot replace the finite-energy hypothesis. A Jordan block at eigenvalue $1/2$ has powers with a factor $N2^{-N}$; a root rate of $1/2$ does not imply a bound $K2^{-N}$.
\end{example}

\section{Conditional futures: normalization changes the question}\label{III:S:sec:conditional}
\subsection{Completed-string probability semantics}
Assume $\alpha,\beta,A_a$ are nonnegative, $T=\sum_aA_a$ has $\rho(T)<1$, and
\[
 v=(I-T)^{-1}\beta\text{ is entrywise positive},\qquad \alpha^\top v=1.
\]
Then $f(w)=\alpha^\top A_w\beta$ is a distribution on \emph{completed finite strings}. For an admissible prefix $u$, define
\[
 F(u)=s(u)v=\sum_t f(ut)>0,\quad
 \bar s(u)=s(u)/F(u),\quad p(t\mid u)=\bar s(u)A_t\beta.
\]
The vectors $\bar s(u)$ lie in the compact simplex $\mathcal S_v=\{r\ge0:rv=1\}$. The future Gramian $Q=\sum_tA_t\beta\beta^\top A_t^\top$ is finite, since each coordinate future sequence has finite $\ell^1$ mass $v_i$. It is positive definite precisely when the future columns span. If $V=\Sdiag(v)$, the gauge $r\mapsto rV$, $A_a\mapsto V^{-1}A_aV$, $\beta\mapsto V^{-1}\beta$ gives $\widetilde\beta+\sum_a\widetilde A_a\mathbf1=\mathbf1$.

Define
\[
 \widehat\omega_N^2=
 \sup_{\substack{q_N(u)=q_N(w)\\F(u),F(w)>0}}
       (\bar s(u)-\bar s(w))Q(\bar s(u)-\bar s(w))^\top.
\]
It is a distance between probability laws, so $0\le\widehat\omega_N\le\sqrt2$, it is nonincreasing, and it vanishes exactly when conditional complete futures agree within each admissible window fiber. In an observable realization this is equality of normalized states; it is not $U_N=0$ for unnormalized states. Cylinder probabilities of nonterminating processes require a separately declared horizon or weighting; their unrestricted sum-of-squares metric is not substituted here.

\begin{theorem}[Exact projective difference identity]\label{III:S:thm:projective}
Let $\Psi_q(r)=rA_q/(rA_qv)$ on its positive-mass domain and $B(y)=I-vy$ for $yv=1$. Then $B(y)^2=B(y)$, and
\[
 \Psi_q(r)-\Psi_q(r')
 =\frac{(r-r')A_qB(\Psi_q(r'))}{rA_qv},\qquad
 D\Psi_q(r)[h]=\frac{hA_qB(\Psi_q(r))}{rA_qv}\quad(hv=0).
\]
Equivalently, the numerator can use $(s(x)-s(y))A_q$ and the denominator $F(xq)$ for histories $xq,yq$.
\end{theorem}
\begin{proof}
For $y=\Psi_q(r')$, $r'A_qB(y)=0$ and $rA_qB(y)=rA_q-(rA_qv)y$. Divide by $rA_qv$. The idempotence follows from $yv=1$. Differentiate the exact identity or the normalized quotient to obtain the derivative.
\end{proof}


\begin{theorem}[Finite row-simplex upper certificate]\label{III:S:thm:rowsimplex}
For $|q|=N$, let $b_i(q)=e_i^\top A_qv$ and $r_i(q)=e_i^\top A_q/b_i(q)$ for live rows with $b_i(q)>0$. Then
\[
 \widehat\omega_N^2\le
 \max_{|q|=N}\max_{i,j:\ b_i(q)b_j(q)>0}
           (r_i(q)-r_j(q))Q(r_i(q)-r_j(q))^\top.
\]
The right side is exactly the squared diameter when arbitrary inputs from $\mathcal S_v$ are allowed. It need not equal the reachable-history diameter. All its squared terms are rational for rational model data.
\end{theorem}
\begin{proof}
For any admissible $r$,
\[
 \Psi_q(r)=\sum_{i:b_i(q)>0}\frac{r_i b_i(q)}{rA_qv}\,r_i(q).
\]
This is a convex combination. A zero row mass implies a zero row of $A_q$ because $v>0$. The diameter of a convex hull in the $Q$ seminorm equals the largest vertex distance, by expanding the difference of two mixtures into a convex combination of pairwise differences. Each live vertex is reached from $e_i^\top/v_i\in\mathcal S_v$. If the requisite vertices lie in the closure of the normalized reachable set, continuity at live vertices also yields equality for reachable histories. The unique rational Stein system supplies rational $Q$.
\end{proof}


\begin{proposition}[A qualified mass-floor bound]
Let $D=\operatorname{diam}_2(\mathcal R)$ for the normalized reachable set, $\Pi=\sup_{y\in\mathcal S_v}\Snorm{I-vy}_2$, $a_N=\max_{|q|=N}\Snorm{A_q}_2$, and $m_N=\inf_{r\in\mathcal R,\ |q|=N,\ rA_qv>0}rA_qv$. If $m_N>0$, then
\[
 \widehat\omega_N\le\min\{\sqrt2,\sqrt{\lambda_{\max}(Q)}\Pi D a_N/m_N\}.
\]
\end{proposition}
\begin{proof}
Apply Theorem~\ref{III:S:thm:projective} and norm submultiplicativity. Compactness gives $\Pi<\infty$. A zero mass infimum makes this particular bound vacuous, not false or divergent.
\end{proof}


\subsection{The precise unary conditional exponent}
For a fixed letter $j$, let $\widehat\omega_{j,N}$ be the supremum over histories $xj^N,yj^N$, rather than over every length-$N$ window.
\begin{theorem}[Perron rate with visible modes and uniform positive overlap]\label{III:S:thm:perron}
Assume $A=A_j$ is primitive. Let $\lambda>0$ be its Perron root and choose positive $h,\ell$ with $Ah=\lambda h$, $\ell A=\lambda\ell$, $\ell h=1$. Put $P=h\ell$, $R=A/\lambda-P$, and $\theta=\rho(R)<1$. For $0<\theta<1$, let $b$ be the largest Jordan-block size on the spectral circle $|z|=\theta$. Then
\[
 \widehat\omega_{j,N}\le K(1+N)^{b-1}\theta^N.
\]
If the realization has both full row and full column spans, then
$\lim_N\widehat\omega_{j,N}^{1/N}=\theta$. If $\theta=0$, sufficiently long $j$-runs have identical normalized futures. Without the span assumptions the upper bound remains valid, but equality of the observable rate need not hold.
\end{theorem}
\begin{proof}
For $N\ge1$, $(A/\lambda)^N=P+R^N$. Put $t_*=\ell/(\ell v)$. Direct normalization yields
\[
 \Psi_{j^N}(r)-t_*=
 \frac{rR^N(I-vt_*)}{(rh)(\ell v)+rR^Nv}.
\]
On the compact simplex, $rh\ge\min_i h_i/v_i>0$. Thus denominators are uniformly positive for all large $N$, and the finitely many smaller $N$ cause no problem because $A^Nv>0$. The normalized-state diameter is bounded by a constant times $\Snorm{R^N}$. The Jordan power formula gives the asserted polynomial-exponential upper estimate.

For the lower estimate choose independent reachable normalized rows $r_1,\ldots,r_n$, possible under full row span. Set $w_i=r_i/(r_ih)$ and $d_i=w_i-w_1$ for $i\ge2$. These form a basis of $\{x:xh=0\}$; together with $\ell$ they span all rows. For $u_{i,N}=\Psi_{j^N}(r_i)$ one has
\[
 d_iR^N=\frac{u_{i,N}}{u_{i,N}h}-\frac{u_{1,N}}{u_{1,N}h},\qquad\ell R^N=0.
\]
The map $u\mapsto u/(uh)$ is uniformly Lipschitz on the simplex. Full future span gives $Q\succ0$ and norm equivalence, so $\Snorm{R^N}\le K'\widehat\omega_{j,N}$. Taking roots and using the spectral-radius formula proves equality. At $\theta=0$, $R$ is nilpotent and the displayed normalized error eventually vanishes.
\end{proof}

This is a specialization of Perron--Frobenius and projective-filtering theory~\cite{cohen,atar}. A simple dominant eigenvalue alone is insufficient, and a Jordan factor cannot in general be removed. Subsection~\ref{III:S:app:conditional} gives exact rational counterexamples to both omissions.

\begin{example}[Exact separation of mass rate and conditional rate]\label{III:S:ex:rates}
Take $\alpha^\top=(1/2,1/2)$,
\[
 A_j=\begin{pmatrix}6/25&4/25\\4/25&6/25\end{pmatrix},\quad
 g=\binom{1/5}{1/10},\quad A_a=ge_1^\top,\quad A_b=ge_2^\top,
 \quad\beta=\binom{1/5}{2/5}.
\]
Then $v=\mathbf1$, the realization is minimal, and the reset letters reach both simplex vertices. $A_j=(2/5)P+(2/25)(I-P)$ for $P=\mathbf1\mathbf1^\top/2$. Exact Stein calculation gives
\[
 (1,-1)Q(1,-1)^\top=700/16369,\qquad
 \widehat\omega_{j,N}^2=(700/16369)25^{-N}.
\]
Thus the conditional root rate is $1/5$. The unnormalized pair $aj^N,bbj^N$ has a nonzero Perron contrast because $F(a)=3/20$, $F(bb)=3/200$, so its future-energy root rate is $2/5$. These are exact rates on an explicitly specified fixture, not inferred rates from an unavailable numerical trace.
\end{example}

\begin{theorem}[Uniform positive-matrix certificate]
If every $A_a$ is strictly positive, define
\[
 \Delta_a=\log\max_{i,j,k,l}\frac{(A_a)_{ik}(A_a)_{jl}}{(A_a)_{il}(A_a)_{jk}},\quad
 \Delta_* =\max_a\Delta_a,\quad\tau_* =\max_a\tanh(\Delta_a/4)<1.
\]
For $N\ge1$,
\[
 \widehat\omega_N\le
 \min\{\sqrt2,2\tanh(\Delta_*\tau_*^{N-1}/4)\}
 \le(\Delta_*/2)\tau_*^{N-1}.
\]
When $\Delta_*=0$, interpret the bound as zero after the first letter.
\end{theorem}
\begin{proof}
The classical Birkhoff theorem gives projective image diameter $\Delta_a$ and contraction factor $\tanh(\Delta_a/4)$~\cite{cohen}. The first positive update sends boundary states into the cone interior, and each following update contracts the Hilbert distance. Thus any $N$-word has diameter at most $\Delta_*\tau_*^{N-1}$. For probability rows the sharp comparison is $\Snorm{\mu-\nu}_1\le2\tanh(H(\mu,\nu)/4)$, with $\ell^1$ rather than half-$\ell^1$ convention~\cite{cohen}. The gauge by $V$ preserves projective distance. Complete-future laws are mixtures of state future laws, so their $\ell^2$ distance is at most the mixture-weight $\ell^1$ distance. Combine these facts and $\tanh x\le x$.
\end{proof}


\begin{theorem}[Mass-qualified transfer; failure without a margin]\label{III:S:thm:conditionaltransfer}
Let $f,g$ be completed-string distributions and use one fixed history set $\mathcal H$ admissible in both, with $F_f(u),F_g(u)\ge\mu>0$. If $\Snorm{H_f-H_g}\le e$ and $\sup_{u\in\mathcal H}|F_f(u)-F_g(u)|\le e_F$, then
\[
 |\widehat\omega_N^{\mathcal H}(f)-\widehat\omega_N^{\mathcal H}(g)|
 \le2(e+e_F)/\mu.
\]
The hypothesis $\Snorm{f-g}_1\le\epsilon$ suffices for the alternative bound $4\epsilon/\mu$. No unqualified uniform continuity in Hankel norm holds for conditional diameters.
\end{theorem}
\begin{proof}
For a history row $r_u^f(t)=f(ut)$,
\[
 \Snorm{r_u^f/F_f-r_u^g/F_g}_2
 \le\frac{\Snorm{r_u^f-r_u^g}_2+|F_f-F_g|}{F_f},
\]
because $\Snorm{r_u^g}_2\le F_g$. A Hankel row error has norm at most $e$. Use the bound for both histories and the reverse triangle inequality, then take the same supremum. An $\ell^1$ error bounds both subtree mass and row errors by $\epsilon$.

For failure without a margin, assign $f_\epsilon(b)=1-\epsilon$, $f_\epsilon(ac)=f_\epsilon(bad)=\epsilon/2$, and $g_\epsilon(b)=1-\epsilon$, $g_\epsilon(ac)=g_\epsilon(bac)=\epsilon/2$, with other masses zero. In $f_\epsilon$ the histories $a,ba$ have deterministic distinct futures, so $\widehat\omega_1(f_\epsilon)=\sqrt2$. In $g_\epsilon$ all admissible equal-last-letter histories have identical futures, so its diameter is zero. Yet $\Snorm{f_\epsilon-g_\epsilon}_1=\epsilon$ and Proposition~\ref{I:S:prop:normbounds} gives $\Snorm{H_{f_\epsilon}-H_{g_\epsilon}}\le\epsilon\to0$.
\end{proof}

Hankel norm alone does not bound $e_F$: summing an infinite residual against the all-ones vector is not a bounded $\ell^2$ functional. Nor does an $\ell^2$ bound on completed futures imply a dimension-free total-variation or KL bound. In a finite set of $M$ outcomes, $\Snorm{p-q}_1\le\sqrt M\Snorm{p-q}_2$; using this inequality on an infinite completion space would be invalid.

\section{Behavioral equivalence and autonomous certification}
\label{III:M:sec:behavioral-contract}
This section supplies elementary bridges needed to read the later results
as a theory of agent memory. They are classical arguments, not asserted
novelty. Fix an event alphabet $\Sigma$, a query set $\mathcal Q$ and a
common answer space. Assume every history and continuation used below is
admissible; for a restricted language, specify its state and common
admissibility conditions as additional persistent information.

\paragraph{From agent records to a predictive contract.}
For an agent-memory application, an implementation must specify how
observations, feedback, and corrections become letters, which queries
are admissible, and what answer law or scalar observable is to be
preserved. Source records and provenance can support validation of this
encoding. The external-memory organization discussed by
Martin~\citep{martin2026memory} motivates retaining enough evidence to
revise a mistaken rule. The mathematical replay statements below assume
that the supplied events and mode information describe the intended
updates; they do not by themselves establish that a natural-language
extractor identified the right rule or correction.

FrontierMemory's task, user, and environment cases suggest complementary
tests of this interface~\citep{letta2026frontiermemoryv1}. A proposed
DRAGA evaluation should freeze the event and query interfaces, separate
memory construction from held-out tasks, and compare task performance
with and without the retained memory. It should also measure update
correctness, correction of stale records, complete stored bytes, and
online and offline computation. This accounting matters because the
benchmark's reported task costs exclude its offline memory-generation
cost. These are evaluation requirements for a future implementation;
this paper reports no FrontierMemory score. The pseudometric below
instead states a worst-case preservation target for a specified model,
including the effect of future continuations.

\begin{definition}[Future answer pseudometric]\label{III:M:def:agent-future}
For reference answer laws $P_h(\cdot\mid q)$ define
\[
 d_{\rm fut}^{\rm ans}(h,h')
 =\sup_{w\in\Sigma^*,\,q\in\mathcal Q}
 \MTV\!\left(P_{hw}(\cdot\mid q),P_{h'w}(\cdot\mid q)\right).
\]
For scalar observables replace TV by absolute difference. The scalar
metric in the later memory theorems is this latter convention; a scalar
readout theorem is not automatically a theorem for an answer distribution.
\end{definition}

\begin{proposition}[Exact quotient and behavioral packing]
\label{III:M:prop:agent-quotient}
Zero future distance defines a right congruence:
$h\sim h'$ implies $ha\sim h'a$ for every letter. If it has finitely
many classes, its quotient is an exact autonomous predictive machine.
Every exact deterministic machine has at least that many reachable states.
More generally, if $M$ histories have pairwise future distances greater
than $2\epsilon$, every deterministic machine with uniform answer error at
most $\epsilon$ has at least $M$ reachable states and needs
$\lceil\log_2M\rceil$ persistent bits.
\end{proposition}
\begin{proof}
The supremum defining $d_{\rm fut}^{\rm ans}(ha,h'a)$ ranges only over
continuations beginning with $a$, so it is at most
$d_{\rm fut}^{\rm ans}(h,h')$. This proves the congruence and makes
the quotient transition $[h]\mapsto[ha]$ well defined. Empty continuation
shows that the answer law is constant on each class. Conversely, equal
machine states remain equal after every common continuation, so an exact
machine cannot merge distinct classes. For approximate machines,
if two selected histories shared a state, their reference answers after
any common continuation would be at distance at most $2\epsilon$ by the
triangle inequality. This contradicts their separating future.
\end{proof}

The quotient result is an existence/minimality statement. An infinite
supremum, exact class equality or the quotient transitions need not be
computable from a finite data sample.

\begin{proposition}[A sufficient finite invariant-cover certificate]
\label{III:M:prop:agent-cover}
Let reference states evolve by $x^+=F_a(x)$. Suppose finitely many
nonempty sets $C_i$, a start index $i_0$, transitions $T(i,a)$ and
answer laws $\widehat P_i(\cdot\mid q)$ satisfy
\[
 x_0\in C_{i_0},\qquad F_a(C_i)\subseteq C_{T(i,a)},\qquad
 \sup_{x\in C_i,q\in\mathcal Q}\MTV(P_x(\cdot\mid q),
                     \widehat P_i(\cdot\mid q))\le\epsilon.
\]
Then storing only the index $i$ and applying $T$ gives uniform answer
error at most $\epsilon$ for all histories.
\end{proposition}
\begin{proof}
Induction gives $x_t\in C_{i_t}$. The final displayed bound applies
to every query at every time. Transitions and answers use only the index
and current input, so the implementation is autonomous.
\end{proof}

The finite cover alone does not supply its transition inclusion or an
efficient description. In particular, the relation
$d_{\rm fut}\le\epsilon$ is generally not an equivalence relation:
three real scalar states $0,3\epsilon/4,3\epsilon/2$ already violate
transitivity. Pairwise similarity clustering is therefore not a proof
of an update-compatible quotient.

\begin{proposition}[Bounded-loss decision transfer]\label{III:M:prop:decision-transfer}
Let $P,Q$ be answer laws with $\MTV(P,Q)\le\epsilon$. For actions
$a$ and loss $\ell(a,y)\in[0,1]$, suppose minimizers $a_P,a_Q$
exist. Then
\[
 \mathbb E_P\ell(a_Q,Y)-\mathbb E_P\ell(a_P,Y)\le2\epsilon.
\]
If the $P$-optimal action is unique with loss gap greater than
$2\epsilon$ to every other action, it remains the unique $Q$-optimal action.
\end{proposition}
\begin{proof}
Expectations of any $[0,1]$-valued function differ by at most TV.
Insert and subtract both $Q$ expectations. Their difference is
nonpositive by optimality of $a_Q$, leaving at most $2\epsilon$.
For the gap assertion, the difference between any two action risks
can change by at most $2\epsilon$.
\end{proof}

This result concerns the declared reference loss and query. It does not
convert a completed-word TV bound into a prefix-uniform guarantee, nor
does it certify that two agents generating different actions encounter
the same future observations.

\subsection{Where the multiletter aspect enters}
The order of events matters: $F_aF_b$ need not equal $F_bF_a$.
A direction invisible to the current readout may become observable
after another letter. A valid reduction must control all declared future
readouts and remain compatible with every allowed update.
Hankel theory encodes this past--future interaction linearly; the invariant
cover and stopped-address constructions handle finite executable states.
Neither linear Hankel rank nor spectral magnitude alone determines how
many persistent factual distinctions survive.

\paragraph{Convention for word order.}
The linear sections use $A_w=A_{a_1}\cdots A_{a_t}$ and a row state
updated on the right. Sequential geometric sections use
$x_t=\tau_{a_t}\circ\cdots\circ\tau_{a_1}(x_0)$.
The unequal-ratio address is written newest first and explicitly reversed
for chronological replay. These conventions describe the same input
ordering in their respective state orientations; they must not be mixed.

\section{An exact optimum in persistent bits}\label{III:M:sec:memory}
The following exact optimum does not depend on
the real dimension of the deployed representation.

\begin{definition}[Autonomous deterministic memory]\label{III:M:def:autonomous-memory}
A machine is a finite Moore machine $(Q,q_0,T_0,T_1,y)$ with deterministic,
time-independent transitions and output depending only on its current state.
All persistent information, including a clock, branch flag, cache, or
coroutine state, belongs to $Q$. Fixed code and immutable parameters may be
arbitrary, but cannot depend on the past input. Define
\[
 B_f^*(\epsilon)=\min\left\{\lceil\log_2|Q|\rceil:
       \sup_{w\in\{0,1\}^*}|f(w)-y(q(w))|\le\epsilon\right\}.
\]
An external clock or access to the deleted input history changes the model.
\end{definition}

Fix an integer $a\ge1$, put $\rho=2^{-a}$, and let
\[
 A_c=\begin{pmatrix}\rho&\rho c/2\\0&\rho/2\end{pmatrix},\quad
 \alpha=(1,0),\quad\beta=(0,1)^\top,\quad f_a(w)=\alpha A_w\beta.
\]
For a word $w=c_1\cdots c_t$, define $k(w)=\sum_{j=1}^tc_j2^{j-1}$.
The row-state recurrence and its closed form are
\[
 (z,x)\longmapsto(\rho z,\rho(x+cz)/2),\qquad
 z(w)=\rho^t,\quad f_a(w)=\frac{k(w)}{2^{(a+1)t}}.
\]
The letter matrices do not commute. The Hankel minor on
$\{\varepsilon,1\}$ has determinant $-\rho^2/4$, so rank is exactly
two. For $a=1$ its exact reachability and observability Gramians are
\[
 P=\begin{pmatrix}4/21&2/21\\2/21&8/7\end{pmatrix},\qquad
 Q=\begin{pmatrix}2&1/3\\1/3&4/21\end{pmatrix}.
\]
Thus $\|H_{f_1}\|_{\rm HS}^2=\operatorname{tr}(PQ)=292/441$.
Since $f_a(w)=(2\rho)^{|w|}f_1(w)$ and $2\rho\le1$, every source
has $\|H_{f_a}\|\le\sqrt{292/441}<5/6<1$.

\subsection{A finite witness for every future}
\begin{lemma}[Two suffixes determine the future metric]\label{III:M:lem:future}
For two states with differences $\Delta z,\Delta x$ and $0<\rho\le1/2$,
their supremum difference over all common suffixes is
\[
 d_{\rm fut}=\max\{|\Delta x|,\tfrac\rho2|\Delta x+\Delta z|\}.
\]
\end{lemma}
\begin{proof}
For a suffix of length $\ell$, the difference equals
$(\rho/2)^\ell\Delta x+\rho^\ell r_v\Delta z$ with
$0\le r_v\le1-2^{-\ell}$. A linear function on this interval is
bounded in absolute value by its endpoints. The first endpoint is at
most $|\Delta x|$. For $\ell\ge2$, the other can be written
$c_\ell(\rho/2)(\Delta x+\Delta z)-d_\ell\Delta x$ with
\[
 c_\ell=2\rho^{\ell-1}(1-2^{-\ell}),\qquad
 d_\ell=\rho^\ell(1-2^{1-\ell}),\qquad
 c_\ell+d_\ell\le5\,2^{-\ell}-6\,4^{-\ell}\le7/8.
\]
Lengths zero and one give the two asserted terms and attain them.
\end{proof}


\begin{lemma}[Behavioral packing]
If $M$ reachable prefixes have pairwise future distances greater than
$2\epsilon$, every admissible machine of error at most $\epsilon$
has at least $M$ reachable states.
\end{lemma}
\begin{proof}
Two merged prefixes give the same machine output after each common
suffix. Their exact outputs could then differ by at most $2\epsilon$,
contradicting the separating suffix.
\end{proof}

The packing principle is classical. The explicit construction and matching
integer bit count below are the specific result.

\subsection{The exact lower and upper bounds}
\begin{theorem}[Exact persistent-bit optimum]\label{III:M:thm:bits}
For $a\ge1$, $m\ge a+2$, and $\epsilon=2^{-((a+1)m+2)}$,
\[
 \boxed{B_{f_a}^*(\epsilon)=m+2.}
\]
The assertion is exact for $B=\lceil\log_2|Q|\rceil$, not for the minimum number of states. The lower bound allows arbitrary immutable tables. The upper construction
uses a closed integer update and includes its clock in the state identifier.
\end{theorem}
\begin{proof}[Lower bound]
Put $L=(a+1)m+2$ and take the prefix set
$\mathcal W_m=\{w:|w|\le m\}\cup\{0^{m+1},1^{m+1}\}$,
of size $2^{m+1}+1$. At one depth $t\le m$, distinct outputs differ
by at least $2^{-(a+1)t}\ge4\epsilon$.
The added pair has output gap
$2\rho(2^{m+1}-1)\epsilon\ge15\epsilon$.
At different depths, $|\Delta z|\ge\rho^m(1-\rho)$.
Lemma~\ref{III:M:lem:future} implies
$|\Delta z|\le(1+2/\rho)d_{\rm fut}$, hence
\[
 \frac{d_{\rm fut}}\epsilon\ge
 \rho 2^{m+2}\frac{1-\rho}{2+\rho}
 \ge16\frac{1-\rho}{2+\rho}\ge\frac{16}5.
\]
Every pair is separated by at least $16\epsilon/5>2\epsilon$.
Packing gives $|Q|\ge2^{m+1}+1$, and therefore at least $m+2$ bits.
\end{proof}

\begin{proof}[Upper bound]
Let $T=\lceil(L-1)/a\rceil$ and
$n_t=\min\{t,L-1-at\}$ for $0\le t<T$.
Create states $(t,b)$ for $0\le b<2^{n_t}$, plus an absorbing state
$\perp$. The integer $b$ records the most recent $n_t$ bits, newest
first. Encode the pair by consecutive blocks in a single identifier;
the clock is not stored separately. The initial state is $(0,0)$.
If $t+1<T$, on symbol $c$ update by
\[
 b'=c2^{n_{t+1}-1}
       +\left\lfloor b/2^{n_t-n_{t+1}+1}\right\rfloor,
 \qquad(t,b)\mapsto(t+1,b').
\]
All shifts have nonnegative integer exponents. The final layer enters
$\perp$, which loops forever. Output
\[
 y(t,b)=\begin{cases}b/2^{(a+1)t},&t\le m,\\
                      (2b+1)2^{-L},&m<t<T,
        \end{cases}\qquad y(\perp)=\epsilon.
\]
The update retains precisely the stated suffix. Before $m$ all bits are
retained. Later the discarded normalized binary tail is in
$[0,2^{-n_t}]$, so the midpoint error is at most
$\rho^t2^{-n_t-1}=\epsilon$. At and after $T$, the source is in
$[0,\rho^T)\subset[0,2\epsilon]$, so the absorbing output is valid.
This proves the all-horizon contract.

The constructed state count is $N=1+\sum_{t=0}^{T-1}2^{n_t}$.
The sum through $m$ equals $2^{m+1}-1$. For $t=m+j$, $j\ge1$,
one has $n_t=m+1-aj$, so the remaining finite geometric tail gives
\[
 N<2^{m+1}+\frac{2^{m+1}}{2^a-1}\le2^{m+2}.
\]
A single $(m+2)$-bit identifier therefore suffices. When $a=1$ the
exact constructed count is $2^{m+2}-2$.
\end{proof}

By monotonicity between successive tolerances, for each fixed $a$,
\[
 B_{f_a}^*(\epsilon)=\frac1{a+1}\log_2(1/\epsilon)+O(1).
\]
The coefficient varies while Hankel rank stays two. Thus rank alone does
not specify the optimal leading bit cost.

\begin{corollary}[Robustness and rank-four sources]\label{III:M:cor:robust-bits}
Under Theorem~\ref{III:M:thm:bits}, if $\|g-f_a\|_\infty\le\eta\le\epsilon/4$,
then $B_g^*(\epsilon+\eta)=m+2$. In particular
$B_g^*(5\epsilon/4)=m+2$ for every such $g$.
An explicit normalized rank-four example is
\[
 g(w)=f_a(w)+\frac\eta2(3^{-|w|}+5^{-|w|}),\qquad\eta=\epsilon/4.
\]
\end{corollary}
\begin{proof}
The same machine gives error $\epsilon+\eta$. Future distances decrease
by at most $2\eta$, leaving at least $16\epsilon/5-2\eta>
2(\epsilon+\eta)$. At tolerance $5\epsilon/4$, the lower margin is
at least $27\epsilon/10>5\epsilon/2$.

For the rank-four example append the scalar blocks $1/3,1/5$ to $A_c$.
Reachability rows and observability columns indexed by
$\varepsilon,0,00,1$ have full rank: the first three yield Vandermonde
minors in the distinct decay rates, and the last supplies the missing
coordinate. The extra series $h$ satisfies
$\|H_h\|_{\rm HS}\le9/7+25/23$ by the triangle inequality, since
$\sum_{t\ge0}(t+1)(2/c^2)^t=(1-2/c^2)^{-2}$ for $c=3,5$.
Consequently $(\eta/2)\|H_h\|_{\rm HS}\le(191/161)\eta<(9/7)\eta$.
The factor $1/2$ is required in this norm estimate;
the conclusion is unchanged. Since $\eta\le1/1024$,
$\|H_g\|<5/6+(9/7)\eta<1$.
\end{proof}


\subsection{What this solves, and what it leaves}
The optimum is exact in \emph{bits}, not in state count: the latter lies
between $2^{m+1}+1$ and the constructed $N$. Jointly encoding time is
essential. Decoding the layer by scanning block sizes takes $O(T)$
iterations and avoids an exponential table. A signed 64-bit implementation of these formulas
imposes $L\le62$; larger parameters need multiword arithmetic. Temporary
space and output precision may exceed the persistent state size.

Most importantly, $f_a$
decays uniformly with input length and the machine eventually absorbs.
It is a sharp executable-memory example, not evidence of sustained
language-memory retention. The rank-four perturbation depends on the
tolerance through $\eta$; it does not establish the same law for every
fixed rank-four source.
\subsection{A constant-overhead existence law on regular attractors}
\label{III:M:sec:ahlfors-memory}
For the following theorem, extend Definition~\ref{III:M:def:autonomous-memory}
to any fixed finite alphabet, with one deterministic transition for each
letter. Immutable transition and output tables are allowed and charged
separately from the persistent state identifier.

\begin{theorem}[Persistent-bit existence law under Ahlfors regularity]
\label{III:M:thm:ahlfors-memory}
Let $\mathcal X\subset\mathbb R^k$ be compact with positive diameter,
$x_0\in\mathcal X$, and maps $\tau_a:\mathcal X\to\mathcal X$ of
Lipschitz constant at most $0<\lambda<1$. Suppose the reachable states
are dense in $\mathcal X$. For $f(w)=h(\tau_w(x_0))$, assume
\[
 c_{\rm obs}\|x-y\|_2\le d_{\rm fut}(x,y)
 :=\sup_v|h(\tau_v(x))-h(\tau_v(y))|
 \le C_{\rm obs}\|x-y\|_2,
\]
with $0<c_{\rm obs}\le C_{\rm obs}<\infty$.
Let a finite Borel measure $\nu$ on $\mathcal X$ have total mass
$M\in(0,\infty)$ and satisfy, for open Euclidean balls and
$0<r<\operatorname{diam}\mathcal X$,
\[
 c_{\rm AD}r^s\le\nu(B(x,r)\cap\mathcal X)
                  \le C_{\rm AD}r^s,
 \qquad s,c_{\rm AD},C_{\rm AD}>0.
\]
Then, for $0<\varepsilon<c_{\rm obs}\operatorname{diam}\mathcal X/4$,
\begin{equation}
 s\log_2(1/\varepsilon)+C_-
 \le B_f^*(\varepsilon)
 \le s\log_2(1/\varepsilon)+C_+,
 \label{III:M:eq:ahlfors-bit-law}
\end{equation}
where
\[
 C_-=\log_2\!\left[\frac{M}{C_{\rm AD}}
                         \left(\frac{c_{\rm obs}}4\right)^s\right],
 \qquad
 C_+=1+\log_2\!\left[\frac{M}{c_{\rm AD}}
                   \left(\frac{2C_{\rm obs}}{1-\lambda}\right)^s\right].
\]
This is an existence theorem for persistent state bits. It does not
assert a computable net, a succinct immutable program, or an efficient
online state update.
\end{theorem}
\begin{proof}
Put $r_-=4\varepsilon/c_{\rm obs}$. A maximal finite
$r_-$-separated set has open $r_-$-balls covering $\mathcal X$;
maximality and compactness ensure both statements. Upper regularity
therefore gives at least
$M/(C_{\rm AD}r_-^s)$ centers. Their pairwise future distances are
at least $4\varepsilon$. By reachable density and the upper
Lipschitz bound, choose reachable states arbitrarily near the finitely
many centers so their future distances remain strictly greater than
$2\varepsilon$. Behavioral packing forces at least as many distinct
machine states and proves the lower bound.

For the upper bound, set
$\eta=(1-\lambda)\varepsilon/C_{\rm obs}$.
Choose a maximal $\eta$-separated set
$\{\widehat x_1,\ldots,\widehat x_N\}\subset\mathcal X$.
It is an $\eta$-net, and its open balls of radius $\eta/2$ are
pairwise disjoint. Lower regularity gives
\[
 N\le\frac{M}{c_{\rm AD}}
                   \left(\frac2\eta\right)^s.
\]
The net is chosen by separation; a minimum-cardinality net would not
automatically have the disjointness property used here.
Initialize at a nearest center to $x_0$, send center $i$ under letter
$a$ to a nearest center to $\tau_a(\widehat x_i)$, and output
$h(\widehat x_i)$. Tie-breaking is fixed. The tracking error satisfies
\[
 e_0\le\eta,\qquad e_{t+1}\le\lambda e_t+\eta,
 \qquad e_t\le\frac\eta{1-\lambda}
                          =\frac\varepsilon{C_{\rm obs}}.
\]
Since the empty suffix is allowed in $d_{\rm fut}$, output error is
at most $C_{\rm obs}e_t\le\varepsilon$. Finally
$\lceil\log_2N\rceil\le\log_2N+1$ yields $C_+$.
\end{proof}


The invariant-attractor condition
$\mathcal X=\bigcup_a\tau_a(\mathcal X)$ implies the required density
when $x_0\in\mathcal X$: every point lies in an iterated image of
diameter at most $\lambda^n\operatorname{diam}\mathcal X$, and the
same image of $x_0$ is reachable. The theorem's content is the
combination of packing and an invariant rounding tube; no novelty is
claimed for these standard ingredients. If finite-precision outputs are
required, reserve half the error for output quantization and halve the
tracking tolerance. This changes only the constant overhead, while all
output constants and transition tables still need to be charged.


\subsection{An optimal table-free executable on homogeneous separated attractors}
\label{III:M:sec:address-memory}
\begin{theorem}[Homogeneous addresses with all storage and update costs charged]
\label{III:M:thm:address-memory}
Suppose $\mathcal X\subset\mathbb R^k$ is the attractor of $d\ge2$
similarities $\tau_a(x)=\lambda O_ax+t_a$ of a common ratio
$0<\lambda<1$, with pairwise separated images. Let the dimension
$k$ and the rational affine maps be fixed. Assume $\mathcal X$ has
positive diameter, $x_0$ is the fixed point of $\tau_0$, and the
future-observability inequalities of Theorem~\ref{III:M:thm:ahlfors-memory}
hold for an output $h$ admitting certified $b$-bit evaluation in
polynomial time and $O_H(b)$ scratch bits. Then there is a table-free
finite-state executable with error at most $2^{-p}$ and
\[
 B(p)=\frac{\log d}{\log(1/\lambda)}p+O_H(1)
       =B_f^*(2^{-p})+O_H(1)
\]
persistent state bits. Its symbol update costs $O_H(p)$ bit operations.
For fixed-dimensional rational affine maps, readout takes
$O_H(p\mathsf M(p))$ bit operations plus the cost of evaluating $h$,
with $O_H(p)$ scratch and immutable parameter bits. No exponential
transition table is stored.
\end{theorem}
\begin{proof}
Put $D=\operatorname{diam}\mathcal X$ and choose
\[
 m=\max\left\{1,
 \left\lceil\frac{p+\log_2(2C_{\rm obs}D)}{\log_2(1/\lambda)}
                                      \right\rceil\right\}.
\]
The persistent state is a word $u$ of length $m$, initialized to $0^m$;
on a letter append it and discard the oldest symbol. Its mathematical
representative is $x(u)=\tau_{u_m}\circ\cdots\circ\tau_{u_1}(x_0)$.
Before $m$ inputs, the initial zeros fix $x_0$, so this representative
equals the true state. After $m$ inputs both states share the same
last $m$ maps, giving distance at most $D\lambda^m$ and output
error at most $2^{-p-1}$.

Pack the word as a base-$d$ integer $I\in\{0,\ldots,d^m-1\}$.
The exact update is $I\leftarrow(dI+a)\bmod d^m$.
It uses $\lceil m\log_2d\rceil$ persistent bits, without a startup
counter; $m,d^m$ and the program are immutable. Fixed-$d$ multiplication,
comparison and subtraction on this integer cost $O_H(p)$ bit
operations, not constant time on a machine with a fixed word width.
To read out, copy the index into scratch, recover its $m$ symbols,
and evaluate the $m$ affine maps. Per-step certified rounding error
$\epsilon_r$ contributes at most $\epsilon_r/(1-\lambda)$.
Choosing $O_H(p)$ fractional bits and a final output error at most
$2^{-p-2}$ leaves total numerical error at most $2^{-p-1}$.
The stated readout and storage costs follow for fixed-dimensional
affine maps and the stipulated output evaluator.

For completeness the separated homogeneous attractor is Ahlfors
regular of dimension $s=\log d/\log(1/\lambda)$: assign mass $d^{-m}$
to each depth-$m$ cylinder. A cylinder containing $x$ and of diameter
less than $r$ supplies the lower ball-mass bound. At a depth with
$\lambda^m$ comparable to $r$, distinct cylinders have separation
bounded below by a constant times $\lambda^m$; Euclidean packing
bounds the number intersecting a ball of radius $r$ by a model
constant, giving the upper bound. Reachable states are dense, since
the images of $x_0$ lie in every cylinder at every depth.
Theorem~\ref{III:M:thm:ahlfors-memory} supplies the matching lower bound.
\end{proof}

\begin{theorem}[An effective contractive coordinate chart]
\label{III:M:thm:chart-memory}
Suppose predictive states admit a supplied global chart
$\phi:[0,1]^k\to\mathcal X$ such that the conjugate updates
$F_a=\phi^{-1}\tau_a\phi$ map the cube to itself and are
$\lambda$-Lipschitz in maximum norm, $\lambda<1$.
Assume reachable states are dense and
\[
 c\|z-z'\|_\infty\le d_{\rm fut}(\phi z,\phi z')
                         \le C\|z-z'\|_\infty.
\]
Suppose the initial coordinate, the $F_a$ and $h\circ\phi$ have
certified $q$-bit evaluators using polynomial time, $O_H(q)$ scratch
bits and $O_H(q)$ immutable parameter bits, with fixed $k,d$.
Then an executable without a transition table attains error $2^{-p}$
with $kp+O_H(1)=B_f^*(2^{-p})+O_H(1)$ persistent bits and polynomial
compilation, update and readout bit cost.
\end{theorem}
\begin{proof}
Set $\epsilon=2^{-p}$ and choose a mesh $\delta=2^{-q}$ with
$\delta\le(1-\lambda)\epsilon/(4C)$, so $q=p+O_H(1)$.
Store the $k$ integer grid coordinates in $\{0,\ldots,2^q\}$.
Evaluate each update to maximum-norm error at most $\delta/2$,
clip to the cube and round to the nearest grid point. Clipping is
nonexpansive; the total one-step numerical error is at most $\delta$.
Initialize with error at most $\delta$. Thus
$e_{t+1}\le\lambda e_t+\delta$ and
$e_t\le\delta/(1-\lambda)$ for every history.
Output error from state tracking is at most $\epsilon/4$; evaluating
the readout to error at most $\epsilon/2$ meets the target.
Each coordinate needs $q+1$ bits, giving the asserted upper bound
and the stated computational costs by the supplied evaluators.
A grid packing of the cube at spacing greater than $2\epsilon/c$,
and reachable density, forces $\Omega_H(\epsilon^{-k})$ distinct
executable states by future observability. This gives the lower bound.
\end{proof}

The chart theorem assumes contraction in the executable coordinates;
bi-Lipschitz equivalence alone can enlarge a contraction constant
beyond one. Discovering such a chart is not part of the theorem.

Theorem~\ref{III:M:thm:packed-startup} below removes the fixed-point
initialization restriction for effective initial coordinates by encoding
startup lengths within the state index; it needs no address of $x_0$.

\paragraph{Why unequal ratios are a different target.}
For $\tau_0(x)=x/4$ and $\tau_1(x)=x/2+1/2$, the separated attractor
has dimension $s$ determined by $4^{-s}+2^{-s}=1$, hence
$s=\log_2((1+\sqrt5)/2)<1$.
A fixed-length register based on the worst contraction $1/2$ stores
$p+O(1)$ bits, rather than $sp+O(1)$. Thus the optimal
coefficient does not follow from that fixed-window construction for unequal ratios. Uniform bi-Lipschitz
bounds alone also do not prove Ahlfors regularity for arbitrary
nonconformal systems. Adaptive address coding or effective charts need separate theorems.
The rational-similarity case is supplied next in
Section~\ref{III:M:audit:sec:addresses}, with its own update and storage costs.

\subsection{Packed startup states remove the fixed-point initialization restriction}
\label{III:M:sec:packed-startup}
\begin{theorem}[Homogeneous addresses from an arbitrary effective initial state]
\label{III:M:thm:packed-startup}
Retain the separated homogeneous similarities and future-observability
assumptions of Theorem~\ref{III:M:thm:address-memory}, but replace its
fixed-point assumption by an arbitrary $x_0\in\mathcal X$ admitting
certified $q$-bit evaluation in polynomial time and $O_H(q)$ scratch
and immutable parameter bits. Assume also that $h$ has an effectively
evaluable $C_{\rm obs}$-Lipschitz extension to a fixed neighborhood
of $\mathcal X$, with the same evaluation resource bounds as before.
Then the same optimal persistent memory
\[
 B(p)=\frac{\log d}{\log(1/\lambda)}p+O_H(1)
       =B_f^*(2^{-p})+O_H(1)
\]
is achieved without a transition table or a supplied address of $x_0$.
Updates use $O_H(p)$ bit operations. Readout uses
$O_H(p\mathsf M(p))$ bit operations plus evaluation of $x_0$ and $h$,
and $O_H(p)$ scratch and immutable parameter bits.
\end{theorem}
\begin{proof}
Choose $m$ as in Theorem~\ref{III:M:thm:address-memory}, and retain the
entire input word until length $m$, and thereafter its last $m$
letters. The state set is $Q_m=\bigcup_{\ell=0}^m\Sigma^\ell$.
Its cardinality is $(d^{m+1}-1)/(d-1)$, so it uses
$m\log_2d+O_d(1)$ bits. Storing a separate length counter is
unnecessary: for a length-$\ell$ word of base-$d$ value $v$, put
\[
 S_\ell=\frac{d^\ell-1}{d-1},\qquad I=S_\ell+v.
\]
The length classes occupy the disjoint consecutive intervals
$[S_\ell,S_{\ell+1})$. Initialize $I=0$. On input
$a\in\{0,\ldots,d-1\}$, the exact update is
\[
 I\leftarrow
 \begin{cases}
 dI+a+1,&I<S_m,\\
 S_m+\big(d(I-S_m)+a\big)\bmod d^m,&I\ge S_m.
 \end{cases}
\]
This needs only the immutable $O_H(p)$-bit constants $S_m,d^m$
and fixed-$d$ integer arithmetic; hence the bit cost is $O_H(p)$.
No time-dependent external counter is used.

The representative of a stored word $u$ is $\tau_u(x_0)$, in the
sequential convention of Theorem~\ref{III:M:thm:address-memory}.
Before the window fills, this is the true state. Afterwards the
representative and the true state share their last $m$ contractions,
and so are within $\lambda^m\operatorname{diam}\mathcal X$.
The truncation output error is at most $2^{-p-1}$.

For readout, a scratch copy of $I$ decodes the letters backwards by
$a=(I-1)\bmod d$ and $I\leftarrow\lfloor(I-1)/d\rfloor$
until zero; a scratch word stores the reversed letters for replay.
There are at most $m=O_H(p)$ such operations, each on $O_H(p)$ bits.
Evaluate $x_0$ and replay the affine maps with initial error
$e_0\le 2^{-p}/(8C_{\rm obs})$ and per-step error
$e_r\le(1-\lambda)2^{-p}/(8C_{\rm obs})$.
The total state error is at most
$e_0+e_r/(1-\lambda)\le2^{-p}/(4C_{\rm obs})$.
A final readout error at most $2^{-p-2}$ completes the budget.
Fixed model constants ensure the computation stays in the stipulated
neighborhood; finitely many small precisions can use finer arithmetic.
The stated storage and readout bounds follow.
For every $x_0\in\mathcal X$, its finite images meet every cylinder,
so reachability is dense. The same packing lower bound as in
Theorem~\ref{III:M:thm:address-memory} applies.
\end{proof}

This removes an initialization restriction within the homogeneous
class. Section~\ref{III:M:audit:sec:addresses} now treats unequal rational
ratios; discovering effective coordinates for a general regular
attractor remains open.

\section{Unequal rational similarity addresses}\label{III:M:audit:sec:addresses}
This section gives a constructive answer to the
unequal-rational-ratio part of the effective-memory target. Its ingredients
include classical stopping trees, enumerative coding, and self-similar packing.
In particular, Reznik--Anisimov~\cite{reznik} already use composition groups and
combinatorial indexing to avoid complete variable-to-fixed coding trees.
Those coding ideas are not claimed as new here. The proposed addition is their
explicit use in an autonomous predictive-state update, with initialization,
finite precision, and bit costs stated together.

\begin{theorem}[Streaming rational-similarity addresses]
\label{III:M:audit:thm:addresses}
Fix a finite alphabet of size $d\ge2$, a dimension $k$, and rational affine
similarities
\[
 \tau_a(x)=r_aO_ax+t_a,\qquad O_a^\top O_a=I,\qquad
 r_a\in\mathbb Q\cap(0,1).
\]
Assume their compact attractor $X$ has positive diameter and strongly separated
first-level images. The affine coefficients and all model constants are fixed.
Let $x_0\in X$ be effectively evaluable. Suppose the scalar output $h$ has a
certified $C$-Lipschitz extension to a fixed neighborhood of $X$, and
\[
 c\Mnorm{x-y}\le d_{\mathrm{fut}}(x,y)
 :=\sup_{v\in\Sigma^*}|h(\tau_vx)-h(\tau_vy)|
 \le C\Mnorm{x-y},\qquad c,C>0.
\]
Suppose $x_0$ and $h$ have certified $q$-bit evaluators with polynomial time
and $O_H(q)$ scratch and immutable parameter bits. Here $H$ denotes the fixed
model, not an input whose size is being suppressed in a uniform claim.
Let $s>0$ be the unique solution of $\sum_a r_a^s=1$.
Then, for error $2^{-p}$, an autonomous finite-state executable attains
\[
 B(p)=sp+O_H(1)=B_f^*(2^{-p})+O_H(1)
\]
persistent bits, polynomial compilation/update/readout bit cost, and $O_H(p)$
scratch and immutable parameter bits. No exponential word or transition table
is stored. A conservative schoolbook implementation of the indexing operations
uses $O_H(p^{d+4})$ bit operations per update or readout, in addition to the
stipulated evaluator cost. This is not a uniform low-degree bound as $d$ varies.
\end{theorem}

\begin{proof}
\emph{Stopping address and update.}
Take an integer constant $c_0\ge1$ with $C\operatorname{diam}(X)2^{-c_0}
\le1/2$, and set $\theta=2^{-(p+c_0)}$. Write a history backwards, newest
letter first. Keep its shortest prefix $v$ with contraction product
$r_v=\prod_i r_{v_i}\le\theta$, or the entire reversed history if the
threshold has not been reached. The possible addresses are all internal nodes
and leaves of the finite ordered tree which stops a branch as soon as its
product is at most $\theta$.

On a new input $a$, prepend $a$ and apply this stopping rule again.
This uses only the old address and $a$. If the old address was already a leaf,
the new stopping prefix ends within that address; discarded older letters
are never needed. If it was internal, the full old history was retained.
This proves the update invariant and includes startup without a separate clock.
The representative replays the retained letters chronologically from $x_0$.
Before truncation it is exact. Thereafter the true state and representative
share all retained contractions, so
\[
 \Mnorm{x_t-\widehat x_t}\le r_v\operatorname{diam}(X)
 \le\theta\operatorname{diam}(X).
\]
The exact readout error is at most $2^{-p-1}$.

\emph{Number of addresses.}
Put $q_a=r_a^s$, so $\sum_a q_a=1$, and let $L_\theta$ be the number of
leaves. The stopped tree is finite and complete. Its leaf masses sum to one,
and, with $r_{\min}=\min_a r_a$,
\[
 r_{\min}\theta<r_v\le\theta
 \quad\Longrightarrow\quad
 (r_{\min}\theta)^s<q_v\le\theta^s.
\]
Hence $\theta^{-s}\le L_\theta<(r_{\min}\theta)^{-s}$. A full $d$-ary
tree with $L_\theta$ leaves has
\[
 N_\theta=\frac{dL_\theta-1}{d-1}
\]
total nodes, including all startup states. An ordinal in
$\{0,\ldots,N_\theta-1\}$ thus needs $sp+O_H(1)$ bits.

\emph{Counting without listing words.}
For a node of contraction product $R>\theta$, define
\[
 I(R)=\sum_{\substack{n\in\MN^d\\MR\prod_a r_a^{n_a}>\theta}}
       \binom{|n|}{n_1,\ldots,n_d},
 \qquad N(R)=1+dI(R).
 \tag{M.1}
\]
For $R\le\theta$, set $N(R)=1$. Every word whose product is above the
threshold is internal, since all shorter prefixes have a larger product.
The multinomial counts all such ordered words of a composition $n$ exactly.
Thus $I(R)$ is the number of internal nodes below this node; the full-tree
identity gives $N(R)$. Each coordinate of $n$ is $O_H(p)$ and there are
$O_H(p^d)$ compositions. This groups words for \emph{counting only}; distinct
letter orders retain distinct ordinals and distinct map compositions.

Use preorder indexing. At each internal node its empty continuation has
index zero, and its ordered child subtrees have sizes $N(Rr_a)$.
Subtracting or adding these sizes performs rank/unrank. There are $O_H(p)$
letters in every address. All threshold comparisons are exact rational
comparisons; neither $s$ nor a floating-point logarithm is needed by the codec.

For completeness, enumerate composition counters one at a time, compute
products and multinomials afresh, and retain only the sum. A multinomial can
be computed as a product of binomial coefficients using multiplicative exact
recurrences. All retained arithmetic values have $O_H(p)$ bits; full factorials
are unnecessary. Each composition costs at most $O_H(p^3)$ schoolbook bit
operations. Rank/unrank makes $O_H(p)$ counting calls, giving the stated
conservative time and linear scratch bounds. This deliberately trades time
for space; a faster implementation may use a separately charged cache.

\emph{Finite-precision readout.}
Replay $O_H(p)$ rational affine maps. With $r_{\max}<1$, initial and per-step
errors satisfying
\[
 e_0\le\frac{2^{-p}}{8C},\qquad
 e_r\le\frac{(1-r_{\max})2^{-p}}{8C}
\]
give state error at most $2^{-p}/(4C)$. A final readout error at most
$2^{-p-2}$ meets the total budget. The fixed neighborhood assumption makes
the rounded readout well defined. Finitely many small precisions may use
finer arithmetic with a fixed model-dependent overhead.

\emph{Lower bound.}
Strongly separated self-similar attractors are $s$-Ahlfors regular, also when
the ratios differ. One elementary proof assigns cylinder mass $r_v^s$ and
uses stopping cylinders with diameters comparable to a ball radius: separation
bounds how many meet the ball, and a cylinder through its center gives the
lower mass bound. The orbit of any $x_0\in X$ is dense, because it meets every
cylinder. Future observability and the packing argument of Section~\ref{III:M:sec:memory},
therefore force $\Omega_H(2^{sp})$ distinguishable executable states.
This gives $B_f^*(2^{-p})\ge sp-O_H(1)$ and completes the proof.
\end{proof}


\textbf{Scope.} Rational ratios and fixed alphabet size make the exact
polynomial counting argument effective. This proof supplies neither efficient
discovery of the model nor the general nonconformal or homographic state encoder.
Persistent bits are optimal to additive constant; the leading coefficient of
\emph{total} peak storage has not been optimized. The generic implementation
is a proof of polynomial cost, not a latency improvement.

\subsection{The unequal-ratio example already singled out in the preceding development}
For
\[
 \tau_0(x)=x/4,\qquad\tau_1(x)=x/2+1/2,\qquad h(x)=x,
\]
the attractor has diameter one, the first-level images are separated, and
$d_{\rm fut}(x,y)=|x-y|$. Its dimension solves
\[
 4^{-s}+2^{-s}=1,\qquad s=\log_2\frac{1+\sqrt5}{2}
 \simeq0.694241914.
\]
At threshold $\theta=2^{-\ell}$, the stopped tree has $F_{\ell+2}$ leaves
and $2F_{\ell+2}-1$ nodes, where $F_0=0,F_1=1$. The leading persistent-bit
coefficient is about 30.6\% below the worst-contraction fixed binary window.
This percentage concerns a mathematical state representation; no measured
memory or runtime improvement for a deployed system is asserted.

\begin{center}
\begin{tabular}{@{}rrr@{}}\toprule
Threshold exponent $\ell$ & Stopping-address bits & Fixed-window bits, with startup\\\midrule
8 & 7 & 9\\16 & 13 & 17\\32 & 24 & 33\\64 & 46 & 65\\128 & 91 & 129\\\bottomrule
\end{tabular}
\end{center}
These are exact finite address counts. A finite-precision output budget needs
the constant tolerance allocation used in the theorem.

\section{An exact finite optimum for the Fibonacci memory}
\label{III:M:sec:leaf-optimum}
The 91-bit entry in the preceding table counts every node of the stopped
tree, including startup nodes. It is a correct cost for that construction.
This section changes the startup contract: the prescribed approximation
error is allowed from time zero. For the explicit initial state $x_0=0$,
leaf addresses alone then give an exact finite minimax result. No claim
about an arbitrary initial point, exact startup, or total peak bytes is
implicit in this improvement.

Fix an integer $\ell\ge1$, put $\theta=2^{-\ell}$, and consider
\[
 x_{t+1}=\tau_{a_{t+1}}(x_t),\qquad
 \tau_0(x)=x/4,\quad \tau_1(x)=x/2+1/2,\quad x_0=0.
\]
The readout is $h(x)=x$. A deterministic memory consists of a finite set
$Q$, a fixed initial state, updates $U_a:Q\to Q$ and a readout
$D:Q\to\MR$, all independent of an uncharged external clock or transcript.
Its required error is $|D(q_t)-x_t|\le\theta/2$ for every finite input
history, including the empty history. Persistent model constants and the
fixed precision parameter are treated as in the earlier fixed-model
contract; temporary arithmetic is accounted for separately below.

Write addresses from newest to oldest. For a word $v=v_1\cdots v_n$, let
\[
 \tau_v=\tau_{v_1}\circ\cdots\circ\tau_{v_n},\qquad
 c(v)=2\#_0(v)+\#_1(v),\qquad I_v=\tau_v([0,1]).
\]
Let $\mathcal L_\ell$ be the words whose cost first reaches or exceeds
$\ell$: $c(v)\ge\ell$ and $c(v^-)<\ell$, where $v^-$ deletes the last
letter. These are leaves, not all prefixes of leaves.

\begin{lemma}[Exact leaf count and separated endpoints]
\label{III:M:lem:fibonacci-leaf-packing}
With $F_0=0,F_1=1$, one has $|\mathcal L_\ell|=F_{\ell+2}$.
Every leaf interval has length $\theta$ or $\theta/2$.
The right endpoints $b_v=\tau_v(1)$, ordered from left to right, have
successive separation at least $3\theta/2$.
Consequently the minimum number of closed intervals of length $\theta$
covering the attractor $K$ is exactly $F_{\ell+2}$.
\end{lemma}
\begin{proof}
The leaf count $L(n)$ satisfies $L(n)=L(n-2)+L(n-1)$ for $n>0$ and
$L(n)=1$ for $n\le0$. Thus $L(\ell)=F_{\ell+2}$. The last increment
is either one or two, so a leaf has cost $\ell$ or $\ell+1$.

Consider adjacent leaves and their longest common prefix $u$, of
contraction ratio $r_u$. The left leaf is the rightmost descendant of
$u0$, so its right endpoint is $\tau_u(1/4)$, because $\tau_1(1)=1$.
The right leaf is the leftmost descendant of $u1$, with some number of
zeros afterwards. Its left endpoint is $\tau_u(1/2)$ and its length is
some $r_v\in\{\theta,\theta/2\}$. The right-endpoint separation is
$r_u/4+r_v$. Since $u$ is internal, its dyadic ratio satisfies
$r_u\ge2\theta$. If $r_u=2\theta$, the right child is already a leaf
and $r_v=\theta$. Otherwise $r_u\ge4\theta$ and $r_v\ge\theta/2$.
Both cases give the asserted $3\theta/2$ bound.
Each endpoint belongs to $K$. A length-$\theta$ interval contains at
most one of these endpoints, while the $F_{\ell+2}$ leaf intervals
themselves cover $K$ and have length at most $\theta$.
\end{proof}


\begin{theorem}[Exact autonomous optimum with approximate startup]
\label{III:M:thm:fibonacci-exact-optimum}
Under the stated contract and $x_0=0$, the minimum number of executable
memory states and the minimum fixed-length persistent bit budget are
\[
 |Q|_{\min}=F_{\ell+2},\qquad
 B_{\min}(2^{-\ell-1})=\left\lceil\log_2 F_{\ell+2}\right\rceil.
\]
In particular, $B_{\min}(2^{-129})=90$ bits at $\ell=128$.
\end{theorem}
\begin{proof}
For the upper bound take $Q=\mathcal L_\ell$, initialize at
$v_0=0^{\lceil\ell/2\rceil}$, update by
\[
 U_a(v)=\operatorname{trim}_\ell(av),\qquad
 D(v)=\tau_v(1/2),
\]
where trimming retains the shortest prefix of cost at least $\ell$.
Such a prefix always exists in $av$. Initially $x_0\in I_{v_0}$.
If $x_t\in I_v$ and $av=v'u$, invariance of $[0,1]$ implies
$x_{t+1}\in\tau_{v'}(\tau_u([0,1]))\subseteq I_{v'}$.
Induction proves this inclusion for all histories. The midpoint readout
therefore has error at most $|I_{v'}|/2\le\theta/2$, including startup.
Every leaf is reachable by supplying its letters in chronological order,
since $\operatorname{trim}_\ell(vv_0)=v$.

For the lower bound one must use reachable histories, rather than assume
that every point of $K$ occurs at finite time. Set $T=\ell+3$.
For each leaf $v$ first supply $T-|v|$ copies of letter 1 and then the
letters of $v$ in reverse order. The final true state is
\[
 y_v=\tau_v\bigl(\tau_1^{T-|v|}(0)\bigr),\qquad
 |b_v-y_v|=r_v2^{-(T-|v|)}\le\theta/8,
\]
because $|v|\le\ell$. Hence two different $y_v$ are separated by at
least $3\theta/2-2\theta/8=5\theta/4>\theta$.
All these histories have the same length. If two ended in the same
stored state, their common readout would be within $\theta/2$ of both,
contradicting this separation. Thus at least $F_{\ell+2}$ states are
necessary. Finally
$F_{130}=659034621587630041982498215$ lies strictly between $2^{89}$
and $2^{90}$, proving the numerical bit optimum.
\end{proof}


\begin{proposition}[Leaf identifiers without an exponential table]
\label{III:M:prop:fibonacci-leaf-cost}
The preceding memory can store only a leaf ordinal. With elementary
integer arithmetic its update and exact midpoint readout use $O(\ell)$
temporary bits and at most $O(\ell^3)$ bit operations per update.
No list of the $F_{\ell+2}$ leaves or transition table is required.
\end{proposition}
\begin{proof}
At remaining cost $n>0$, the left subtree contains $L(n-2)$ leaves.
For ranking, add this number when the next letter is 1, then decrease
$n$ by the letter cost. For unranking, compare the current ordinal with
$L(n-2)$ and choose the corresponding child, subtracting that count on
the right. There are at most $\ell$ letters. Each needed Fibonacci
number can be recomputed by a rolling two-integer recurrence, with
$O(\ell)$ additions of $O(\ell)$-bit integers. Recomputing rather than
caching all counts gives $O(\ell^3)$ bit operations and $O(\ell)$ live
bits for the ordinal, counters and retained word. Prepend, trim and rank
give the autonomous update. The midpoint is a dyadic rational with at
most $\ell+2$ fractional bits, obtainable by exact affine replay within
the same space and a smaller time bound.
\end{proof}


\paragraph{What changes, and what does not.}
The earlier $2F_{130}-1$-node representation still uses 91 bits and
retains exact startup histories. The new 90-bit representation uses a
midpoint readout and permits the declared error immediately. The
91-bit count was never a lattice quantization of occupation parameters.
The asymptotic coefficient remains
$\log_2\varphi\simeq0.694241914$, with
$\varphi=(1+\sqrt5)/2$; improving startup changes an additive constant.
Hausdorff dimension alone would only give a rate with unspecified
finite constants. The exact packing above supplies the missing finite
argument. Independently byte-aligned 90-bit and 91-bit payloads both
occupy 12 bytes; saving one bit is not automatically a resident-byte gain.
Python object sizes, working storage, output buffers and time are not
included in the persistent-bit number.


\subsection{Extension to the leaf union and non-dyadic tolerances}\label{III:M:rev13:sec:fib}

Notation as in Section~\ref{III:M:sec:leaf-optimum}: $\tau_0(x)=x/4$, $\tau_1(x)=x/2+1/2$, $h(x)=x$, $\theta=2^{-\ell}$, leaves
$\mathcal L_\ell$ (newest-first words whose cost first reaches $\ell$),
$I_v=\tau_v([0,1])$, $b_v=\tau_v(1)$, and $K$ the attractor. Put
\[
U_\ell=\bigcup_{v\in\mathcal L_\ell}I_v\supseteq K .
\]
The preceding theorem proves $|Q|_{\min}=F_{\ell+2}$ for $x_0=0$ and states that no claim for an
arbitrary initial point is implicit. The proof in fact needs only $x_0\in U_\ell$.

\begin{theorem}[Exact optimum for every initial point of the leaf union]
\label{III:M:rev13:thm:fib-general}
Fix $\ell\ge1$ and the contract $|D(q_t)-x_t|\le\theta/2$ for every finite history
including the empty one, for an autonomous deterministic memory $(Q,q_0,(U_a),D)$ without
an uncharged clock. For every $x_0\in U_\ell$ (in particular every $x_0\in K$),
\[
|Q|_{\min}=F_{\ell+2},\qquad B_{\min}(2^{-\ell-1})=\lceil\log_2F_{\ell+2}\rceil .
\]
For every $x_0\in[0,1]$ the lower bound $|Q|_{\min}\ge F_{\ell+2}$ holds. An exact state count for $x_0\notin U_\ell$ is not asserted here.
\end{theorem}
\begin{proof}
\emph{Upper bound.} Choose (or supply to the compiler) $v_0\in\mathcal L_\ell$ with $x_0\in I_{v_0}$; this is the only
place where $x_0$ enters. Take $Q=\mathcal L_\ell$, $U_a(v)=\Mtrim_\ell(av)$,
$D(v)=\tau_v(1/2)$. If $x_t\in I_v$ and $av=v'u$ with $v'=\Mtrim_\ell(av)$, then
$x_{t+1}=\tau_a(x_t)\in\tau_{av}([0,1])=\tau_{v'}(\tau_u([0,1]))\subseteq I_{v'}$, so by
induction $x_t\in I_{v_t}$ for all histories and the midpoint readout has error at most
$|I_{v_t}|/2\le\theta/2$.

\emph{Lower bound.} Let $T=\ell+3$. For each leaf $v$ feed $T-|v|$ letters $1$ followed by
the letters of $v$ in chronological order (possible because $|v|\le\ell$). The final state
is $y_v=\tau_v\bigl(\tau_1^{T-|v|}(x_0)\bigr)$ and, since $\tau_1^j(x_0)=1-2^{-j}(1-x_0)$,
\[
|b_v-y_v|=r_v2^{-(T-|v|)}|1-x_0|\le\theta\,2^{-3}=\theta/8
\]
for every $x_0\in[0,1]$. The right endpoints $b_v$ are pairwise at least $3\theta/2$ apart
(Lemma~\ref{III:M:lem:fibonacci-leaf-packing}), hence distinct $y_v$ are
at least $3\theta/2-\theta/4=5\theta/4>\theta$ apart. All $F_{\ell+2}$ histories have the
same length $T$; two of them ending in the same memory state would force a common
readout within $\theta/2$ of both final states, contradicting the separation.
\end{proof}


\begin{remark}
The leaf union $U_\ell$ is a strict superset of $K$ (it contains the gaps of $K$ at
depth $\ge\ell$), so the theorem covers more initial points than the attractor. For
$x_0\in[0,1]\setminus U_\ell$ (for example $x_0=3/10$ at $\ell\ge2$) the leaf machine
cannot be initialised, the exact value lies in $[F_{\ell+2},\,2F_{\ell+2}-1]$ (the
all-node construction gives the upper end). The dependence of the precise
count on $x_0$ in this residual case remains open.
\end{remark}

\begin{proposition}[Non-dyadic tolerances]\label{III:M:rev13:prop:fib-eps}
Let $x_0\in K$, let $0<\epsilon\le1/4$ and let $\ell\ge1$ be the integer with
$2^{-\ell-1}\le\epsilon<2^{-\ell}$, $\theta=2^{-\ell}$. For the contract
$|D(q_t)-x_t|\le\epsilon$:
\[
\text{if }\ \tfrac\theta2\le\epsilon<\tfrac{3\theta}4:\quad |Q|_{\min}(\epsilon)=F_{\ell+2};
\qquad
\text{if }\ \tfrac{3\theta}4\le\epsilon<\theta:\quad F_{\ell+1}\le|Q|_{\min}(\epsilon)\le F_{\ell+2}.
\]
\end{proposition}
\begin{proof}
The leaf machine at threshold $\theta\le2\epsilon$ has error $\le\theta/2\le\epsilon$, giving
the upper bound. For the lower bound repeat the construction of
Theorem~\ref{III:M:rev13:thm:fib-general} with $T=\ell'+j$ at threshold $\theta'=2^{-\ell'}$: the
final states are within $\theta'2^{-j}$ of the endpoints, hence pairwise more than
$3\theta'/2-\theta'2^{1-j}$ apart, and this exceeds $2\epsilon$ for large $j$ whenever
$\epsilon<3\theta'/4$. Take $\ell'=\ell$ in the first case and $\ell'=\ell-1$ (so
$3\theta'/4=3\theta/2>\epsilon$) in the second.
\end{proof}


\section{Persistent facts, checkpoint replay and certified guards}\label{III:M:audit:sec:persistent}
The new address construction still forgets old distinctions under a sufficiently
long common continuation. It does not solve permanent factual recall.

\begin{proposition}[Uniform contraction cannot preserve a non-fading fact]
Suppose two histories encode a binary fact $b\in\{0,1\}$ that remains unchanged
after $t$ repetitions of an irrelevant event. The two stored states initially
have distance at most $D$. Suppose that event is $\lambda$-contractive,
$0<\lambda<1$, and the answer distribution is $L$-Lipschitz in state distance
under total variation. If both fact answers have TV error at most $\epsilon$
against the correct point-mass answers, then
\[
 1\le 2\epsilon+LD\lambda^t.
\]
Consequently no fixed $\epsilon<1/2$ can hold for all $t$ under these assumptions.
\end{proposition}
\begin{proof}
The correct answers have TV distance one. The triangle inequality bounds this
by the two answer errors plus the distance between the model's two answer
distributions. Lipschitz readout and contraction bound the latter by
$LD\lambda^t$.
\end{proof}

This is an elementary obstruction, not a new novelty claim. Discontinuous
exact-real readout is excluded by the Lipschitz hypothesis; finite robust
storage cannot silently use infinitely precise coordinates.

If all $N$ independently supplied bits must be individually recalled with
worst-case error below $1/2$ by a deterministic-state encoder, at least $N$
persistent bits are necessary. Otherwise two of the $2^N$ fact assignments
share a state and disagree on a query. Thus growing factual content cannot be
hidden inside a fixed total byte budget with unrestricted recall guarantees.

\subsection{The persistent-plus-predictive construction target}
\textbf{Target: a certified streaming memory with exact persistent distinctions
and a compressed predictive component.} The deployed state has the form
\[
 m_t=(r_t,z_t),\qquad
 r_{t+1}=G(r_t,a_{t+1}),\qquad z_{t+1}=F_{a_{t+1},r_t}(z_t).
\]
Here $r_t$ retains designated non-fading facts or exceptions, while $z_t$
captures compressible dynamics. Exact storage refers to the chosen typed
facts, not to infallible extraction from natural language. Every byte in
$r_t$, provenance and indexes is counted. A guarantee must also account for
coupling from $z_t$ into any update of $r_t$.

\subsection{Constructive extension: replay from the mode at the window base}
\begin{proposition}[Rational addresses with a mode checkpoint]
\label{III:M:audit:prop:checkpoint}
Let $\mathcal R$ be a nonempty finite set, $G:\mathcal R\times\Sigma\to
\mathcal R$ a deterministic computable transition, and
\[
 r_{t+1}=G(r_t,a_{t+1}),\qquad
 x_{t+1}=\tau_{a_{t+1},r_t}(x_t),\qquad
 \tau_{a,r}(x)=r_aO_{a,r}x+t_{a,r}.
\]
The ratios $r_a\in\mathbb Q\cap(0,1)$ depend only on the letter; the
orthogonal maps and translations may depend on the mode. Assume a common
nonempty compact invariant set $X$ contains $x_0$, and let $D$ be a known
upper bound on its diameter. Fix $0<\theta<1$ and use the stopped address
tree of Theorem~\ref{III:M:audit:thm:addresses}, including its internal startup nodes.
An autonomous representation is
\[
 m_t=(b_t,I_t),\qquad b_t=r_{\mathrm{base},t}
       =r_{t-|v_t|},\qquad I_t=\operatorname{rank}(v_t),
\]
where $v_t$ is the retained word, newest letter first. It uses at most
\[
 \lceil\log_2|\mathcal R|\rceil+\lceil\log_2N_\theta\rceil
\]
persistent bits. Start with $(b_0,I_0)=(r_0,0)$. On receiving $a$, form
\[
 w=(a)\operatorname{unrank}(I_t)=v_{t+1}u_{\rm drop},\qquad
 v_{t+1}=\operatorname{trim}(w),
\]
and update, with $G^*$ iterating $G$ in the displayed word order,
\[
 b_{t+1}=G^*(b_t,\operatorname{rev}(u_{\rm drop})),\qquad
 I_{t+1}=\operatorname{rank}(v_{t+1}).                 \tag{M.2}
\]
To read, initialize $(\rho,y)=(b_t,x_0)$ and for each letter $a$ of
$\operatorname{rev}(v_t)$ apply the simultaneous assignment
\[
 (\rho,y)\longleftarrow\bigl(G(\rho,a),\tau_{a,\rho}(y)\bigr),
                                                        \tag{M.3}
\]
whose right-hand side uses the old mode. The final mode is exactly $r_t$.
If all history is retained, $y=x_t$; otherwise
\[
 \Mnorm{x_t-y}\le D r_{v_t}\le D\theta.               \tag{M.4}
\]
Neither invertibility nor strong connectivity of $G$ is required.

For executable numerical bounds, assume additionally rational affine
coefficients and the certified evaluation hypotheses of
Theorem~\ref{III:M:audit:thm:addresses}, uniformly over the fixed finite mode catalogue.
Then the same polynomial indexing cost and $O_H(p)$ scratch bounds apply
at $\theta=2^{-(p+c_0)}$. For an output uniformly $C$-Lipschitz in the
continuous coordinate within each exact mode, the same error allocation
gives output error $2^{-p}$ and
\[
 B(p)\le sp+\lceil\log_2|\mathcal R|\rceil+O_H(1),
 \qquad \sum_a r_a^s=1.                              \tag{M.5}
\]
All mode/transition descriptions and numerical evaluators are charged.
This is a constructive upper bound, not a joint optimality theorem.
\end{proposition}
\begin{proof}
Let $\ell=|v_t|$. Its oldest endpoint is time $t-\ell$, whose mode is $b_t$.
After appending one letter, discarding $j=|u_{\rm drop}|$ oldest letters
advances that endpoint to $t-\ell+j=t+1-|v_{t+1}|$. Those discarded letters
are precisely $\operatorname{rev}(u_{\rm drop})$ in chronological order.
Equation (M.2) therefore advances the checkpoint to the correct new mode,
proving its invariant by induction. This only evaluates $G$ forwards.

Starting at this checkpoint, (M.3) reconstructs every subsequent mode exactly,
because $G$ depends on that mode and the input letter alone. The true state
at the window base and $x_0$ both lie in $X$. The two replays apply identical
mode-dependent similarities, whose composition has contraction ratio
$r_{v_t}$. This proves (M.4). If the whole history is retained, its base state
is $x_0$ and the error is zero. An internal tree node necessarily retains the
whole history. A leaf can also occur at the first threshold crossing before
any deletion; (M.4) is still a safe upper bound there.

The ratio depends only on the letter, so the address tree and composition
counts are exactly those of Theorem~\ref{III:M:audit:thm:addresses}. The one extra mode
uses $\lceil\log_2|\mathcal R|\rceil$ bits. Replaying or dropping letters
adds only $O_H(p)$ mode transitions. With certified rounded affine replay,
the numerical error satisfies the same geometric recurrence as before;
the discrete mode remains exact. This proves the computational claims.
\end{proof}


Strong separation and future observability are not needed for this upper
bound. They cannot be silently transferred from the unmixed family to the
mode-dependent one to claim a matching lower bound. Some pairs $(b,v)$ can
be unreachable, and different modes can have indistinguishable futures.
The model is fixed in (M.5); a growing mode catalogue must be charged explicitly.

\textbf{Feedback boundary.} If instead
$r_{t+1}=G(r_t,a_{t+1},\sigma_{t+1})$, where $\sigma_{t+1}$ is determined
by the continuous trajectory, the letters and checkpoint need not recover
the branch choices. One must prove that each choice is reconstructible,
or retain the branch information and charge its storage and update costs.
A log of one independent binary choice per retained step costs
$O_H(p)$ additional bits; no additive-constant overhead follows in general.
Proposition~\ref{III:M:audit:prop:checkpoint} does not resolve this feedback case.

\subsection{Numerical refinement, geometric envelopes and actual ambiguity}
Fix a stored state $m=(b,v)$, a query/event, and its scalar guard $g$.
Let $\Phi_{b,\operatorname{rev}(v)}$ denote chronological replay with modes
initialized at $b$. For a leaf address put
\[
 c_m=\Phi_{b,\operatorname{rev}(v)}(x_0),\qquad
 S_m=\Phi_{b,\operatorname{rev}(v)}(X),\qquad R_m=Dr_v.
\]
The set of states attainable by admissible histories with this \emph{full}
stored state is
\[
 \mathcal C_m=\{x(h):\operatorname{Enc}(h)=m\}
       \ \subseteq\ S_m\ \subseteq\ B(c_m,R_m).       \tag{M.6}
\]
In general the first inclusion is strict: arbitrary points of $X$ need not
be reachable with the required base mode and encoding. For an internal
startup node use $S_m=\{c_m\}$ and $R_m=0$ instead. All these histories
have the same current mode, which is recovered by (M.3).

\begin{lemma}[Termination of a centered Lipschitz sign certificate]
\label{III:M:audit:lem:guard}
Suppose $g$ is $L_g$-Lipschitz on a region containing the relevant states and
evaluation points. At precision $q$ let $\widetilde g_q$ be a certified
evaluation of the exact replay center with
\[
 |\widetilde g_q-g(c_m)|\le\delta_q,
 \qquad \delta_q=L_gE_{\rm num}(q)+\eta_g(q)\longrightarrow0.
\]
Put $A=L_gR_m$ and $\Delta_m=|g(c_m)|-A$. The test
\[
 |\widetilde g_q|>A+\delta_q                          \tag{M.7}
\]
certifies a constant strict sign on $\mathcal C_m$. Eventual acceptance
by this particular test is equivalent to $\Delta_m>0$; a sufficient
precision is $2\delta_q<\Delta_m$.
The alternative test
\[
 |\widetilde g_q|+\delta_q\le A                      \tag{M.8}
\]
certifies $\Delta_m\le0$, so further numerical refinement cannot make
(M.7) succeed. If $\Delta_m<0$, (M.8) eventually succeeds. If $\Delta_m=0$,
neither finite-time exit is guaranteed.
\end{lemma}
\begin{proof}
For any $x\in\mathcal C_m$, $|g(x)-\widetilde g_q|\le A+\delta_q$,
which proves the sign claim. Acceptance implies
$|g(c_m)|\ge|\widetilde g_q|-\delta_q>A$. Conversely,
$|\widetilde g_q|\ge|g(c_m)|-\delta_q>A+\delta_q$ whenever
$2\delta_q<\Delta_m$. The upper triangle inequality proves (M.8);
if $\Delta_m<0$, its left side is at most $|g(c_m)|+2\delta_q<A$
eventually. For $g(c_m)=A$ and
$\widetilde g_q=A+\delta_q/2$ with $\delta_q>0$, both (M.7) and (M.8)
fail for every $q$, even though the errors tend to zero.
\end{proof}


\textbf{Counterexample: numerical refinement is not a truncation certificate.}
Take $L_g=1$, $R_m=4/256$, $\delta_q=2/256$,
$\widetilde g_q=5/256$, and $g(c_m)=3/256$.
Then $|\widetilde g_q|>L_gR_m$, but $\Delta_m=-1/256$.
The weaker inequality omitting numerical uncertainty is invalid.
For instance, $c_m=0$, $g(x)=3/256+x$ on
$[-4/256,4/256]$ realizes a sign change despite that weaker test.

\textbf{Interval certificate.} Any certified enclosure
$[\underline g_m,\overline g_m]\supseteq g(\mathcal C_m)$ yields a
positive decision when $\underline g_m>0$, and a negative decision when
$\overline g_m<0$. An enclosure of $g(S_m)$ is sufficient.
Zero in the enclosure is inconclusive. Likewise $\Delta_m\le0$ only
excludes the centered test (M.7); it does not establish a sign change.
For $g(x)=R^2+x^2$ on $[-R,R]$, $R>0$, one has $L_g=2R$ and
$g(0)=R^2<2R^2=L_gR$, whereas the exact range
$[R^2,2R^2]$ is strictly positive.

\textbf{An actual information-loss certificate.} Exhibit two admissible
histories $h^+,h^-$ with identical \emph{full} stored state, the same
query/event, and $g(x(h^+))>0>g(x(h^-))$. No deterministic decision from
that state alone can be correct for both. This is stronger than finding
zero in an outer interval. A concrete example uses a single mode,
$\tau_0(x)=x/2$, $\tau_1(x)=x/2+1/2$, $x_0=0$, and threshold $1/2$.
Chronological histories $(0,0)$ and $(1,0)$ retain the same last-letter
address $(0)$, but their states are $0$ and $1/4$; the guard $x-1/8$
has opposite signs. Equality with zero requires a separately declared
tie convention and is not used in this example.

\textbf{Branching is a set-valued guarantee.} With feedback, unresolved
guards can be handled by keeping joint hypotheses about the discrete mode,
continuous compatible set and required replay/branch payload. Intersect each
compatible set with its guard region, then propagate it by that branch's
transition. This preserves containment of the reference trajectory when the
enclosures are sound. It does not provide a unique exact mode. A subset of
the original mode--address product alone need not encode this information,
because different branch histories can share its elements. Explicit binary
branching can double the number of hypotheses per unresolved step; merging
or symbolic compression requires a separate correctness and cost argument.

\subsection{Persistence is an observability property, not a spectral shortcut}
For $F(x)=x/(1+x)$ on $[0,1]$, direct induction gives
\[
 F^t(1)=\frac1{t+1}\longrightarrow0,\qquad
 \lim_{t\to\infty}(F^t(1))^{1/t}=1,\qquad F'(0)=1.
\]
With identity readout and this single neutral transition,
$d_{\rm fut}(F^t(1),0)=1/(t+1)$. Thus unit asymptotic exponential rate
does not imply a positive surviving separation.

For $P_0=\operatorname{diag}(1,0)$ and
$P_1=\operatorname{diag}(0,1)$, the joint spectral radius is one:
all products have norm at most one and repeated $P_0$ has norm one.
Nevertheless $P_1P_0=P_0P_1=0$, so a mixed continuation can erase every
distinction. These projectors \emph{commute}; noncommutativity is not needed
for the counterexample. Joint spectral radius controls the largest product
growth, not a lower bound for every continuation.

A word-Hankel operator connects pasts to future observables. Its nonzero
singular values count linear predictive rank, and their magnitudes measure
spectral approximation strength. They do not by themselves provide a
common center--stable decomposition for all letter transitions.

One \emph{sufficient specification} for finitely many protected facts is a
finite history congruence $\pi:\Sigma^*\to\mathcal R$ satisfying
$\pi(ha)=G(\pi(h),a)$, together with a designated fact label $f(r)$ and
a declared neutral alphabet. Require fact invariance under neutral events
and the uniform observational margin
\[
 \inf_{w\in\Sigma_{\rm neutral}^*}
 d_{\rm fut}(hw,h'w)\ge2\delta>0
 \quad\text{whenever } f(\pi(h))\ne f(\pi(h')).       \tag{M.9}
\]
The admissible histories and future queries in (M.9) must be specified.
Different auxiliary control modes carrying the same fact need not be
separated. This is a conditional design requirement, not a proved universal
decomposition or a learning theorem. A fixed finite register cannot retain
an unbounded number of independently queryable facts.

\subsection{What remains open, and what existing proofs need}
The difficult contribution is learning or synthesizing this decomposition
with an efficiently checkable future-error invariant and competitive cost.
A supplied decomposition plus a geometric error series is useful machinery,
but does not establish a general discovery theorem.

Theorem~\ref{III:M:audit:thm:addresses} and Proposition~\ref{III:M:audit:prop:checkpoint} do not
depend on solving the feedback, learned-congruence or general Hankel
decomposition targets. Their hypotheses and proofs are explicit above.
Lemma~\ref{III:M:audit:lem:guard} also does not assume that an inconclusive test can
always be resolved. Resolving those open targets is necessary only for the
corresponding stronger claims: general closed-loop agent memory, automatic
discovery of protected facts, or global optimality beyond the stated class.

\textbf{Proof obligations.} Specify the event/query language and reference
semantics; establish all-letter update compatibility; transport numerical and
model error to the declared answers; and prove the storage/update bound.
The ordinary counting-Hankel estimate does not automatically certify rare
conditional queries. The conditional bound divides by prefix mass.
A controlled agent whose actions change subsequent observations needs a
controlled-system or policy-level argument in addition to a fixed-input bound.

\textbf{Decisive experiment.} Use a frozen answering model and typed events
with a known oracle: insert facts, apply corrections and retractions, interleave
long irrelevant sequences, and ask delayed or counterfactual queries. Compare
an exact key--value baseline, retrieval/summary baselines, a homogeneous fading
state, the unequal-ratio state, and the persistent-plus-predictive construction.
Match total memory rather than just injected tokens. Report factual accuracy,
order sensitivity, certificate coverage/violations, update and readout latency,
and peak allocation. Afterwards test paraphrased natural-language interfaces,
separating extraction failures from memory failures. No such agent benchmark
is reported in this paper.

\section{Certified feedback, joint memory and consolidation}
\label{III:M:sec:agent-extensions}
This section gives conditional constructions for feedback, joint memory
and consolidation. All input alphabets, mode catalogues and structural constants
are fixed unless a bound explicitly varies them. The elementary packing,
coupling, quotient and concentration arguments below are not claimed as
new general principles. Their role is to state the exact executable contract.

\subsection{Cache decisions at the time of their certification}
Use the following time convention throughout this section:
\begin{equation}\label{III:M:eq:feedback-time}
 \sigma_j=\mathbf1\{g_{a_j}(x_{j-1})\ge0\},\qquad
 x_j=\tau_{a_j,r_{j-1}}(x_{j-1}),\qquad
 r_j=G(r_{j-1},a_j,\sigma_j).
\end{equation}
Both right-hand sides use the old mode; the guard uses the state before
the continuous update. A model whose guard uses the updated state requires
a different, explicitly stated convention. No guard is recomputed at an
intermediate replay state.

\begin{lemma}[Maximum stopped-window length]\label{III:M:lem:feedback-length}
For $0<r_a<1$, put $r_{\max}=\max_a r_a$ and
\[
 L_\theta=\left\lceil\frac{\log(1/\theta)}{\log(1/r_{\max})}\right\rceil.
\]
Every word stopped at its first product at most $\theta$, and every
internal startup word, has length at most $L_\theta$. For
$\theta=2^{-(p+c_0)}$, this is $O_H(p)$ for a fixed model.
\end{lemma}
\begin{proof}
The product of any $L_\theta$ ratios is at most
$r_{\max}^{L_\theta}\le\theta$; its first crossing occurs no later.
An internal word has product greater than $\theta$, so is shorter.
\end{proof}

The slowest contraction is decisive. For ratios $1/4,3/4$ and
$\theta=1/16$, the all-$3/4$ word first crosses at length $10$:
$(3/4)^9>1/16\ge(3/4)^{10}$. Replacing $r_{\max}$ by $r_{\min}$
would incorrectly give $2$.

\begin{theorem}[Certified branch-cache replay]\label{III:M:thm:branch-cache}
Assume the rational similarities, common compact $X$, initial point $x_0$,
and effective evaluators of Proposition~\ref{III:M:audit:prop:checkpoint}, but
replace its mode update by \eqref{III:M:eq:feedback-time}. Let $D\ge\operatorname{diam}(X)$.
Suppose at each incoming event a finite certified computation produces
an interval containing $g_{a_j}(x_{j-1})$ that excludes zero, using only
the deployed state, the new event and charged model data. Store its exact
sign bit $\sigma_j$. The hypothesis is an explicit restriction on the
admissible trajectories and guard algorithm, not an automatic consequence
of contraction.

Let $v_t=(a_t,\ldots,a_{t-\ell+1})$ be the stopped word and store
\[
 m_t=(b_t,I_t,\beta_t),\qquad
 b_t=r_{t-\ell},\quad I_t=\Mrank(v_t),\quad
 \beta_t=(\sigma_t,\ldots,\sigma_{t-\ell+1}).
\]
Replay the reversed letter--bit pairs from $(b_t,x_0)$ using the
simultaneous update $(\rho,y)\leftarrow
(G(\rho,a,\sigma),\tau_{a,\rho}(y))$.
Then the final mode is exactly $r_t$, and exact-arithmetic replay satisfies
\[
 \Mnorm{x_t-y_t}\le D r_{v_t}\le D\theta
\]
whenever deletion has occurred; if all history is retained, $y_t=x_t$.
Certified numerical replay adds its explicit $E_{\rm num}$.

On arrival of $a$, first certify the new bit on the current state enclosure,
then prepend $(a,\sigma)$, and trim letters and bits together. Advance $b_t$
through the dropped pairs in chronological order. This is an autonomous
update, including noninvertible mode resets. A fixed-width representation
uses at most
\begin{equation}\label{III:M:eq:branch-bits}
 \lceil\log_2|\mathcal R|\rceil+
 \lceil\log_2N_\theta\rceil+L_\theta
\end{equation}
persistent bits. The decoded address supplies the active bit-string length;
unused bits in the fixed field are zeroed. Addressing and mode replay keep
the previous polynomial bounds and $O_H(p)$ scratch. Total event time also
includes the guard's actual certification cost, for which no universal
bound is asserted here.
\end{theorem}
\begin{proof}
Initially the address and bit string are empty and $b_0=r_0$.
Inductively every cached bit is the reference event-time bit by the
certified enclosure. If $j$ oldest pairs are dropped after appending,
their forward mode transitions advance the base from $t-\ell$ to
$t-\ell+j=t+1-|v_{t+1}|$. No inverse transition is used.
Replaying cached pairs therefore follows exactly the reference mode path.
The true state at the retained base and $x_0$ both lie in $X$ and undergo
the same similarities. Their distance contracts by $r_{v_t}$. Before
deletion they start from the same $x_0$. Numerical errors follow the
geometric recurrence of the preceding checkpoint proposition.
Lemma~\ref{III:M:lem:feedback-length} bounds the extra bit field and the number
of added mode operations. Guard work is separate from these operations.
\end{proof}


The guard can use the image of $X$ under the cached mode path, enlarged
for numerical error. This outer enclosure need not be tight. At zero
margin the theorem does not promise a successful event: halt or abstention
must be declared, or a separate justified mechanism must supply the bit.
An external decision oracle is not an uncharged implementation.

\paragraph{Optional augmented-alphabet encoding.}
If only letters in a declared $\Sigma_{\rm gate}$ require a bit, encode
the alphabet consisting of one copy of each other letter and two copies
of each gated letter. Both copies have ratio $r_a$. The address counting
argument gives an upper rate $\widetilde s p+O_H(1)$, in addition to the
mode checkpoint, where
\[
 \sum_{a\notin\Sigma_{\rm gate}}r_a^{\widetilde s}
 +2\sum_{a\in\Sigma_{\rm gate}}r_a^{\widetilde s}=1.
\]
This is an upper bound on all encoded pair words. It is not an optimality
claim: many bit words can be unreachable and the copied geometric branches
need not be separated. Neither this construction nor \eqref{III:M:eq:branch-bits}
preserves the earlier coefficient $s$ without further information.

\subsection{Sound coalescence with its actual cardinality bound}
Here hypotheses are an outer representation of possible states, not a
probability distribution. Soundness means that the actual state remains
covered. It does not authorize selecting an arbitrary hypothesized mode
as an exact answer.

\begin{proposition}[A shared continuous address and a mode subset]
\label{III:M:prop:oneway-coalescence}
Suppose $\tau_{a,r}=\tau_a$ and the ratios satisfy the address assumptions.
Maintain the single address $I_t$ and a subset $S_t\subseteq\mathcal R$
containing the actual mode. For the new letter let $J_t\subseteq\{0,1\}$
contain every guard value on the sound continuous enclosure; if unresolved,
use $J_t=\{0,1\}$. Set
\[
 S_{t+1}=\{G(r,a,\sigma):r\in S_t,\ \sigma\in J_t\}.
\]
This representation remains sound and has at most $|\mathcal R|$ mode
labels at every time. A bit-mask implementation uses
$|\mathcal R|+\lceil\log_2N_\theta\rceil$ persistent bits.
\end{proposition}
\begin{proof}
The continuous path depends only on the input word, so every branch shares
its continuous address and sound error enclosure. The actual pair
$(r_t,\sigma_{t+1})$ occurs in the displayed Cartesian product; its successor
mode belongs to $S_{t+1}$. Taking a set merges duplicate labels and leaves
at most $|\mathcal R|$ elements. A common enclosure may discard correlations
and admit extra possibilities, but cannot discard the actual trajectory.
\end{proof}

The proposition controls representation size, not enclosure sharpness or
factual ambiguity. A singleton answer is certified only when all retained
possibilities justify it under the declared output contract.

\begin{proposition}[Separated gates: settled and transient counts]
\label{III:M:prop:dwell-coalescence}
Allow mode-dependent similarities with letter-only ratios. Let gated
events have binary successor choices; ordinary events keep the mode fixed.
Between consecutive gated events require at least $L_\theta$ ordinary
events. Start from the known initial state, or a settled representation
with at most one canonical stopped-window hypothesis per current mode.
Use common letter addressing and sound enclosures, and discard mode-path
details only after they leave the stopped window.

Then a sound implementation can maintain at most $|\mathcal R|$
hypotheses immediately before each gated event, and at most
$2|\mathcal R|$ at arbitrary intermediate times. Once the previous gate
has left the active window, histories with the same current mode can use
the common replay from $(r,x_0)$ and its $D\theta$ error enclosure.
This is canonical re-encoding with a certified error, not equality of
the original continuous states.
\end{proposition}
\begin{proof}
After $L_\theta$ ordinary events, the stopped word is contained wholly
in those events. Its mode is constant and equal to the current mode;
thus the mode at its base is the current mode. For each such mode, replay
from $x_0$ covers all true paths in that mode within $D\theta$, because
their base states lie in $X$. Common ratios make the retained letter word
identical across hypotheses. Coalescing to one canonical enclosure per
mode is therefore sound at this settled time.
The next gate creates at most two successors per settled hypothesis.
Ordinary events do not split them. The next gate occurs only after a new
settled time, which proves both counts by induction. One common address,
the remaining mode paths (at most one active gate each), and a finite
hypothesis catalogue suffice; their storage and guard work are charged.
\end{proof}


\paragraph{Why the all-time $|\mathcal R|$ bound is false.}
Take $\mathcal R=\{0,1\}$, $X=[0,1]$, ratio $1/4$ and
$\tau_r(x)=x/4+3r/4$. Before a gate, the enclosures
$[0,1/4]$ and $[3/4,1]$ are compatible with the two modes.
For $g(x)=(x-1/8)(x-7/8)$ both enclosures contain both signs.
With $G(r,a,\sigma)=\sigma$, splitting at the gate produces distinct
old-mode/new-mode pairs $(0,0),(0,1),(1,0),(1,1)$.
The new continuous enclosures from the two old modes lie respectively
in $[0,1/16]$ and $[15/16,1]$; sharing a new mode does not identify them.
There can therefore be four candidates for two modes immediately after
the gate, even though a sufficiently long subsequent dwell permits
canonical coalescence. This is a counterexample to the claimed inference
for the maintained sound hypothesis representation; it does not assert
that all its outer candidates must be reachable from one exact history.

\subsection{A matching lower bound for separated observable fibres}
\begin{theorem}[Conditional joint persistent-bit optimum]
\label{III:M:thm:joint-mode-optimum}
Use the letter-driven model and effective hypotheses of
Proposition~\ref{III:M:audit:prop:checkpoint}; this theorem does not include
unlogged continuous feedback. Let $\mathcal R_{\rm reach}$ be its finite
reachable mode set. Assume the closures $X_r$ of the reachable continuous
fibres are nonempty $s$-Ahlfors regular, with common positive regularity
constants and scales, where $\sum_a r_a^s=1$.
For actual reachable states assume
\[
 d_{\rm fut}((r,x),(r,y))\ge c\Mnorm{x-y},\qquad
 d_{\rm fut}((r,x),(r',y))\ge2\delta_R\quad(r\ne r'),
\]
with $c,\delta_R>0$. Assume also the uniform within-mode readout
$C$-Lipschitz bound and effective numerical error allocation of the
checkpoint proposition for every declared query. Futures are the
declared common input continuations. For sufficiently large $p$ with
$2^{-p}<\delta_R$,
\[
 \log_2|\mathcal R_{\rm reach}|+sp-O_H(1)
 \ \le B^*(2^{-p})\ \le
 \log_2|\mathcal R_{\rm reach}|+sp+O_H(1).
\]
Here $B^*$ minimizes deterministic autonomous persistent memory under
uniform future-answer TV error. Constants may depend on the fixed model,
including regularity, observability and readout constants.
\end{theorem}
\begin{proof}
Regularity gives, in each compact fibre closure, a packing at separation
$4\varepsilon/c$ of cardinality at least $a_H\varepsilon^{-s}$ at small
scales. Approximate its finitely many points by actual reachable states
within $\varepsilon/(2c)$; their separation remains greater than
$2\varepsilon/c$. Within a mode, observability makes the future distance
of any two selected histories greater than $2\varepsilon$. Across modes
the same conclusion follows from $\varepsilon<\delta_R$.
If two such histories shared a stored state, autonomous updates would
give identical stored states under every common continuation. The two
uniform approximation inequalities and TV's triangle inequality would
then give future distance at most $2\varepsilon$, a contradiction.
Consequently at least $|\mathcal R_{\rm reach}|a_H\varepsilon^{-s}$
codes are required. Taking logarithms proves the lower bound.

Every retained base mode is reachable. Indexing only those modes in the
checkpoint construction and choosing $\theta$ and numerical precision
for output error $\varepsilon$ gives the upper bound. Under a common
future input, equal modes remain equal and identical similarities do not
increase continuous discrepancies. The uniform query Lipschitz bound
therefore supplies the declared future contract as well as the current
readout contract. The original online addressing estimates apply.
\end{proof}

The upper Lipschitz hypothesis cannot be replaced by lower observability.
For example, the family of threshold queries $h_q(x)=\delta_{\mathbf1\{x\ge q\}}$
separates every distinct pair by TV $1$, however close their coordinates.
On an infinite reachable set no finite code achieves uniform error below
$1/2$ for that family. The theorem is a supplied-model result; it neither
discovers the fibres nor proves their regularity for a learned agent.
For a growing catalogue, dependence hidden in $O_H(1)$ must be made uniform.

\subsection{A conditional-query contract without a query prior}
\begin{proposition}[Uniform query kernels]\label{III:M:prop:query-kernels}
Let $h_q(r,x)$ be specified probability kernels for every declared query,
including queries outside the support of a sampling distribution. Suppose
\[
 \sup_q\MTV(h_q(r,x),h_q(r,y))\le C\Mnorm{x-y}
\]
for all admissible within-mode pairs. If the deployed mode is exact and
its state error is at most $e$, then its plug-in kernel has uniform query
error at most $Ce$, independently of any query prior.
\end{proposition}
\begin{proof}
Apply the assumed inequality to the exact and decoded states and take the
supremum over $q$. No conditioning of a joint word law is performed.
\end{proof}

For a null-probability query, the kernel is part of the specified reference
semantics, not an inferred conditional probability on a null event.
This proposition does not remove the rare-prefix obstruction for an
unconditional Hankel or completed-word TV approximation. It uses a
stronger, separately verified family of conditional readouts.

\subsection{External trajectory guarantees in a controlled environment}
Compare two agents with the same initial external-state distribution and
the same environment/observation kernels. The reference and deployed agents
receive the same history until their actions differ. Assume standard
measurable kernels (finite actions suffice below). ``Trajectory'' means
external states, observations and actions; it does not include the two
different internal memory representations. The memory error assumptions
below must hold on every shared history in the specified support, not just
on one realized test rollout.

\begin{theorem}[Guarded actions and probability of guard failure]
\label{III:M:thm:guarded-policy}
Let the reference policy choose the largest of finitely many scores
$g_a(r,x)$ with a fixed tie rule. At a shared history, suppose the mode is
exact and certified computed scores obey
$|\widetilde g_a-g_a(r,x)|\le\beta_t$ for every action, where for example
$\beta_t=L_g e_t+\eta_t$. Choose a computed maximizer $a^*$ only when
\[
 \widetilde g_{a^*}-\max_{b\ne a^*}\widetilde g_b>2\beta_t.
\]
Then the selected action equals the unique reference maximizer.
Let $\tau$ be the first failed action or memory certificate in a coupling
that keeps the two external trajectories identical before failure. If the
deployed policy is defined after failure as well, then
\[
 \MTV(P_{\rm ref}^{1:T},P_{\rm dep}^{1:T})
 \le\Pr(\tau\le T).
\]
In particular the laws coincide if certificates succeed almost surely
throughout the supported horizon. If failure at each next step, conditional
on no earlier failure, has probability at most $\alpha_t$, the right-hand
side is at most $\min(1,\sum_{t=1}^T\alpha_t)$.
\end{theorem}
\begin{proof}
For every $b\ne a^*$, the true score difference exceeds the computed
difference minus $2\beta_t$, hence is positive. Identical actions allow
the environment to use the same transition and observation samples; the
next external histories remain equal. Induction constructs the coupling
until $\tau$. The probability that its full external paths differ is at
most $\Pr(\tau\le T)$, which bounds TV. Summing the probabilities of
first failure gives the last assertion. Observed guard success on one
sample path does not bound this failure probability.
\end{proof}


\begin{theorem}[Stochastic-policy coupling]\label{III:M:thm:stochastic-policy}
Suppose on every shared history the two action kernels have TV distance
at most $d_t\le1$. Then
\[
 \MTV(P_{\rm ref}^{1:T},P_{\rm dep}^{1:T})
 \le1-\prod_{t=1}^T(1-d_t)
 \le\min\left(1,\sum_{t=1}^Td_t\right).
\]
In particular take $d_t=\min(1,L_\pi e_t)$ if exact modes and a uniform
$L_\pi$-Lipschitz stochastic policy are independently established on
shared histories. Uncertified mode errors are not covered by this choice.
\end{theorem}
\begin{proof}
While histories agree, maximally couple the next finite-action kernels;
their conditional probability of disagreement is at most $d_t$.
On agreement couple the environment identically. The probability of no
disagreement through time $T$ is at least $\prod_t(1-d_t)$. The coupling
inequality proves the first bound, and the union/product bound the second.
The environment need not be contractive: any amplification after a first
disagreement is already included in its probability.
\end{proof}

These theorems transfer a supplied uniform memory/policy certificate to
closed-loop behavior. They do not establish its premises for arbitrary
trained policies or infer total-variation guarantees from benchmark scores.

\subsection{Exact active-map quotients and certified approximate merging}
\begin{proposition}[Active-map quotient with a complete finite code]
\label{III:M:prop:active-map}
Let $K,V$ be fixed finite sets of sizes $k,v\ge1$. Events are deterministic
set/update, delete and no-op operations on $K\to V\cup\{\bot\}$.
Assume every future discrete output and transition depends on the past
only through the current map, and every continuous update depends only
on that map, its continuous state and the next event. No historical,
provenance, scheduled activation or expiration query is included unless
its required metadata is added to the state.
Then replacing the operation history by its active map preserves all
future behavior exactly. All maps can be encoded in
$\lceil\log_2 (v+1)^k\rceil$ bits. If every map is reachable and every
key including absence can be queried, this number is the exact minimum
fixed-width bit budget for the discrete map alone.

A simpler enumerative code for a map with $n$ active keys uses at most
\[
 \lceil\log_2(k+1)\rceil+
 \lceil\log_2\tbinom{k}{n}\rceil+n\lceil\log_2 v\rceil
\]
bits, including the variable cardinality. If $n$ is fixed and known as
part of the contract, the first term can be omitted.
\end{proposition}
\begin{proof}
Equal active maps and continuous states produce equal next maps and states
under every common event, and equal outputs. Induction proves exact
future equivalence. Each key has $v+1$ possibilities, yielding $(v+1)^k$
maps and a finite radix code. Any two different reachable maps differ at
a key; its exact query distinguishes them, proving the fixed-width lower
bound. Alternatively encode $n$, the ranked $n$-element key subset, and
its values in sorted-key order, which gives the displayed upper length.
\end{proof}

Canonical semantics does not imply that this serialization is shorter than
every previous representation. For $k=v=1$ there are two maps, empty and
present: omitting the cardinality from the conditional-size formula would
assign zero bits to both. An exact old/new codec selector costs one extra
bit and permits length $1+\min(L_{\rm old},L_{\rm map})$ if both decoders
and their data are available and charged. Conversion workspace, retained
logs and peak allocation must also be counted. Future-dated activations
cannot be erased merely because they are inactive now.

\begin{proposition}[Merging in a certified future pseudometric]
\label{III:M:prop:future-merge}
Let $d_{\rm fut}$ be the declared common-continuation pseudometric.
Suppose a possible true state $z$ belongs to at least one certified ball
$d_{\rm fut}(z,z_i)\le e_i$. Partition the candidate centres into groups
and choose admissible representatives $\bar z_j$ with certified distances
$d_{\rm fut}(z_i,\bar z_j)\le\delta_i$ for each member of group $j$.
Replacing the groups by their representatives with radii
$\max_{i\in j}(e_i+\delta_i)$ preserves sound coverage.

Under common deterministic inputs, if the transition is nonexpansive in
this future metric, successive certified replacements of radii at most
$\delta_1,\ldots,\delta_n$ increase a single-representative error by at
most $\sum_j\delta_j$. This sum, or a sharper proved transport bound,
must be charged against the total allowed error.
\end{proposition}
\begin{proof}
For the group containing a centre that covers $z$, the triangle inequality
gives $d_{\rm fut}(z,\bar z_j)\le e_i+\delta_i$. The representative thus
covers the true state with the stated radius. For repeated replacements,
apply this inequality after each common update; nonexpansiveness prevents
the intervening updates from enlarging the existing bound. In the usual
supremum-over-continuations metric, nonexpansiveness follows because
continuations beginning with a fixed letter form a subset of all
continuations. This does not compare agents taking different actions.
\end{proof}

The result certifies coverage, not a posterior weight or permission to
select an arbitrary remaining hypothesis as true. Computationally useful
certificates of the stated future distances remain a separate obligation.

\paragraph{When a midpoint estimate is justified.}
For an admissible convex within-mode domain and a proved bound
$d_{\rm fut}((r,x),(r,y))\le C\Mnorm{x-y}$ on all its points,
the midpoint of $x_i,x_j$ has future distance at most
$C\Mnorm{x_i-x_j}/2$ from either. This controls the centres only;
their existing uncertainty radii must be added. Convexity and the
\emph{future} bound are essential. A current-output Lipschitz constant
is insufficient when a later guard can switch a persistent mode.

\paragraph{A rational counterexample to Euclidean branch merging.}
Take $X=[-1,1]$, current mode $0$, and states $x=\pm\eta$, $0<\eta<1$.
At the next gated input set $r^+=\mathbf1\{x\ge0\}$ and $x^+=x/2$;
ordinary subsequent inputs preserve the mode. Querying that mode gives
opposite Dirac laws, hence future TV $1$. The continuous distance is
$2\eta\to0$. Merging into $x=0$ selects the positive branch and loses
the negative one. Even a perfectly contractive continuous map does not
make this merge safe. A nonconvex invariant set introduces the additional
problem that the midpoint need not be an admissible state.

\subsection{Verification from supplied probes}
Neither numerical closeness of trajectories nor an eigenvalue gap alone
establishes a statistical identification theorem. With exact noiseless
observations of $x_1=\lambda x_0$ and known $x_0=1$, a single observation
distinguishes $\lambda=1$ from every specified $\lambda<1$. Conversely,
hidden models with identical observable laws cannot be distinguished by
any sample size, even if their latent eigenvalues have a gap.

\begin{proposition}[Finite supplied probes under explicit sampling promises]
\label{III:M:prop:probe-test}
Fix $M$ supplied probes, $0<\alpha<1$, $0<\delta\le1$, $D\ge\delta$,
and $0<\gamma<1$. For each probe $i$ and chosen lag $L$, suppose one can
obtain $n$ independent identically distributed observations in $[0,1]$
with mean $\mu_i(L)$. Each probe is promised to be either persistent,
with $\mu_i(L)\ge\delta$, or fading, with
$0\le\mu_i(L)\le D(1-\gamma)^L$. Choose
\[
 L\ge\left\lceil\gamma^{-1}\log(4D/\delta)\right\rceil,
 \qquad
 n\ge\left\lceil8\delta^{-2}\log(2M/\alpha)\right\rceil.
\]
Classify a probe as persistent exactly when its sample mean exceeds
$\delta/2$. All $M$ classifications are correct with probability at
least $1-\alpha$. The count is $Mn$ probe observations at lag $L$,
not $n$ unrestricted training examples for an unknown memory architecture.
\end{proposition}
\begin{proof}
The lag makes every fading mean at most $\delta/4$, since
$(1-\gamma)^L\le e^{-\gamma L}$.
For a bounded observation $Y$ of mean $\mu$, convexity gives
$\mathbb E e^{tY}\le1-\mu+\mu e^t$. The logarithm of the centred
right-hand side has value and derivative zero at $0$ and second derivative
at most $1/4$ (the variance of a tilted Bernoulli variable). Hence
$\mathbb E e^{t(Y-\mu)}\le e^{t^2/8}$ for either sign of $t$.
Independence and exponential Markov, optimized at $t=4a$, give
$\Pr(|\bar Y-\mu|\ge a)\le2e^{-2na^2}$.
Taking $a=\delta/4$ and a union bound over $M$ probes gives failure
probability at most $2M e^{-n\delta^2/8}\le\alpha$.
On its complement the fading sample means are below $\delta/2$ and
persistent means above $3\delta/4$, so the decision rule is correct.
\end{proof}

The proposition assumes observable probes with separated means and
independent sampling. It does not identify arbitrary latent modes,
construct a protected fact congruence, validate a semantic extractor, or
prove persistence against every neutral word. Correlated trajectories
require their own concentration and effective-sample-size analysis.

\subsection{Nonconformality is not itself an impossibility theorem}
Take $A_0=A_1=\Mdiag(1/3,1/5)$ with translations $(0,0),(2/3,0)$.
The attractor is the middle-third Cantor set times $\{0\}$.
For $0\le s\le1$, every length-$n$ product has singular-value function
$3^{-ns}$. Both norm and singular-value pressures are therefore
$\log2-s\log3$, with the same zero $\log2/\log3<1$:
nonconformality does not imply a strict exponent gap.
With zero translations and distinct linear parts
$\Mdiag(1/3,1/5),\Mdiag(1/3,1/6)$, that pressure calculation remains valid
but the attractor is a singleton. Affinity dimension need not equal
covering dimension in unrestricted overlapping systems.
Memory rates require future-metric packing/covering bounds and effective
updates; dimension alone does not give $O_H(1)$ overhead. General
nonconformal constructions remain open here, not proved impossible.

\section{Mode-dependent ratios: pressure bounds and conditional optimality}\label{III:M:rev13:sec:mode}

Proposition ``Rational addresses with a mode checkpoint'' requires the ratio $r_a$ to depend
on the letter only. This section removes that restriction, gives a constructive pressure coefficient and a conditional matching
lower bound, and keeps the table-free polynomial update.

\subsection{Setting}
Let $\mathcal R$ be a finite mode set with deterministic transition
$G:\mathcal R\times\Sigma\to\mathcal R$ (resets allowed), initial mode $r_0$, and
\[
r_{t+1}=G(r_t,a_{t+1}),\qquad x_{t+1}=\tau_{a_{t+1},r_t}(x_t),\qquad
\tau_{a,m}(x)=r_{a,m}O_{a,m}x+t_{a,m},
\]
with rational ratios $r_{a,m}\in\mathbb Q\cap(0,1)$ depending on \emph{both} the letter
and the mode, orthogonal $O_{a,m}$, a common compact invariant set $X\ni x_0$ of known
diameter bound $D$, and all model data fixed. Write $r_{\min},r_{\max}$ for the extreme
ratios. Define the pressure matrix
\begin{equation}\label{III:M:rev13:eq:pressure}
M(s)_{m,m'}=\sum_{a\in\Sigma:\,G(m,a)=m'}r_{a,m}^{\,s},\qquad s\ge0 .
\end{equation}
Call the mode graph (edges $m\to G(m,a)$) \emph{strongly connected} if $M(0)$ is
irreducible. Then $s\mapsto\rho(M(s))$ is continuous and strictly decreasing, with
$\rho(M(0))=d\ge2$ and $\rho(M(s))\to0$, so there is a unique $s^\ast>0$ with
\begin{equation}\label{III:M:rev13:eq:sstar}
\rho(M(s^\ast))=1,\qquad M(s^\ast)u=u\ \text{ for some }u>0 .
\end{equation}
This is a symbolic pressure exponent. Its identification with geometric
Hausdorff dimension additionally requires a separation hypothesis, as in
\citet{mauldin_williams}. Overlapping or identical branches need not have
that dimension. In this section $d=|\Sigma|\ge2$ and $p\ge1$.

\subsection{The mode-aware stopped address}
For a history $a_1\cdots a_t$ with modes $r_0,\dots,r_t$, let
$\pi_j=\prod_{i=t-j+1}^{t}r_{a_i,r_{i-1}}$ be the product of the ratios \emph{actually
applied} by the $j$ newest letters ($\pi_0=1$). Let
$j^\ast=\min\{j:\pi_j\le\theta\}$ if this set is nonempty, else $j^\ast=t$. The stored state is
\[
m_t=(b_t,v_t),\qquad v_t=(a_t,\dots,a_{t-j^\ast+1})\ \text{(newest first)},\qquad
b_t=r_{t-j^\ast}.
\]

\begin{lemma}[Autonomous update and exact replay]\label{III:M:rev13:lem:mode-update}
From $(b_t,v_t)$ and a new letter $a$, replay $\operatorname{rev}(v_t)$ from $(b_t,x_0)$ to
obtain the modes $r_{t-j^\ast},\dots,r_t$ exactly, form $w=a\,v_t$, compute its
newest-first prefix products with these modes, cut at the first product $\le\theta$ (or
keep all of $w$ if none), and advance $b_t$ chronologically through the dropped oldest
letters. The result is $(b_{t+1},v_{t+1})$. Replaying $\operatorname{rev}(v_t)$ from
$(b_t,x_0)$ returns the exact current mode $r_t$ and a point $y_t$ with
\[
\Mnorm{x_t-y_t}\le D\,\pi_{j^\ast}\le D\theta\ \text{ if a deletion has occurred},\qquad
y_t=x_t\ \text{ otherwise}.
\]
Every retained word has length at most
$L_\theta=\lceil\log(1/\theta)/\log(1/r_{\max})\rceil$.
\end{lemma}
\begin{proof}
If $v_t$ is a leaf ($\pi_{j^\ast}\le\theta$) then $w$ has full product
$r_{a,r_t}\pi_{j^\ast}\le\theta$, so the first product $\le\theta$ among the newest-first
prefixes of the full new history occurs within $w$; these prefix products coincide with
those computed from the replayed modes, which are exact because $G$ is deterministic and
$b_t$ is the true mode at the window base. If $v_t$ is internal it equals the whole history,
so $w$ is the whole new history. Dropping the $j$ oldest letters moves the base from
$t-j^\ast$ to $t-j^\ast+j=t+1-|v_{t+1}|$ and $G$ is only evaluated forwards. The true state
at the window base and $x_0$ both lie in $X$ and undergo the same mode-dependent
similarities, whose composition has ratio $\pi_{j^\ast}$. The length bound is the
Lemma~\ref{III:M:lem:feedback-length}.
\end{proof}


\subsection{Counting the addresses}
For a mode $m$ and $x>0$ let
\[
W_m(x)=\#\Bigl\{c_1\cdots c_n\ (n\ge0)\text{ chronological from }m:\ 
\textstyle\prod_{i=1}^n r_{c_i,m_{i-1}}>x\Bigr\},
\]
where $m_0=m$ and $m_i=G(m_{i-1},c_i)$.
\begin{lemma}[Address set and growth]\label{III:M:rev13:lem:mode-count}
Identify a newest-first stored word with its chronological reversal for
counting and ranking. Every reachable stored state lies in
\[
\begin{aligned}
\mathcal A_\theta={}&\{(r_0,w):\ w\text{ internal from }r_0\}\\
&\cup\bigcup_{b\in\mathcal R,\ c\in\Sigma}\{(b,\,c\,w'):\ w'\text{ chronological from }
G(b,c)\text{ with product in }(\theta,\theta/r_{c,b}]\},
\end{aligned}
\]
so that $|\mathcal A_\theta|=W_{r_0}(\theta)+\sum_{b,c}\bigl[W_{G(b,c)}(\theta)-
W_{G(b,c)}(\theta/r_{c,b})\bigr]$. If the mode graph is strongly connected then, with
$C_u=\max_mu_m/\min_mu_m$ for the Perron vector in \eqref{III:M:rev13:eq:sstar},
\[
W_m(x)\le\frac{C_u\,r_{\min}^{-s^\ast}}{d-1}\,x^{-s^\ast}\quad(m\in\mathcal R,\ 0<x<1),
\qquad
|\mathcal A_\theta|\le\frac{(1+d|\mathcal R|)\,C_u\,r_{\min}^{-s^\ast}}{d-1}\,\theta^{-s^\ast}.
\]
\end{lemma}
\begin{proof}
A leaf $v_t=(c_{j^\ast},\dots,c_1)$ with base $b$ and oldest letter $c_1$ satisfies
$\pi_{j^\ast}\le\theta<\pi_{j^\ast-1}$, and $\pi_{j^\ast}=r_{c_1,b}\,\pi'$ where $\pi'$
is the product of $c_2\cdots c_{j^\ast}$ read from $G(b,c_1)$; thus
$\pi'\in(\theta,\theta/r_{c_1,b}]$. Shorter newest suffixes have larger products, so no
further condition is needed. Internal states retain the whole history, so their base is
$r_0$. This proves the inclusion and the count.

For the growth, consider the forward stopped tree from $m$ at threshold $x$ (a branch
stops at its first node of product $\le x$). It is a full $d$-ary tree whose internal
nodes are counted by $W_m(x)$, so it has $(d-1)W_m(x)+1$ leaves, each of product in
$(r_{\min}x,x]$. Induction from the leaves upward, using $M(s^\ast)u=u$ at every internal
node of product $R$ and mode $m'$ (its children contribute
$\sum_aR^{s^\ast}r_{a,m'}^{s^\ast}u_{G(m',a)}=R^{s^\ast}u_{m'}$), gives
$\sum_{\text{leaves }\lambda}\rho_\lambda^{s^\ast}u_{\mathrm{end}(\lambda)}=u_m$.
Hence $(\#\text{leaves})(r_{\min}x)^{s^\ast}\min_mu_m\le\max_mu_m$, which gives the bound
on $W_m(x)$; the bound on $|\mathcal A_\theta|$ follows from the count with
$W_{G(b,c)}(\theta/r_{c,b})\ge0$.
\end{proof}


\begin{lemma}[Polynomial counting without listing words]\label{III:M:rev13:lem:mode-dp}
For a fixed model, every $W_m(x)$ with $x=\theta/R$, $R$ a product of ratios, and every
rank/unrank operation on $\mathcal A_\theta$ can be computed with $O_H(p^{d|\mathcal R|+O(1)})$
bit operations on $O_H(p)$-bit integers, using exact rational comparisons only.
\end{lemma}
\begin{proof}
The product of a chronological word depends only on its type vector
$n\in\MN^{\Sigma\times\mathcal R}$, $n_{a,m'}$ being the number of positions at which
letter $a$ is read in mode $m'$: the product is $\prod r_{a,m'}^{n_{a,m'}}$. Let
$N(n,m,m'')$ be the number of words from $m$ of type $n$ ending in mode $m''$. Then
$N(0,m,m'')=[m=m'']$ and
\[
N(n,m,m'')=\sum_{(m',a):\,G(m',a)=m'',\ n_{a,m'}\ge1}N(n-e_{a,m'},m,m').
\]
Words have length at most $L_\theta=O_H(p)$, so there are
$\binom{L_\theta+d|\mathcal R|}{d|\mathcal R|}=O_H(p^{d|\mathcal R|})$ type vectors, each
with a count of $O_H(p)$ bits, and
$W_m(x)=\sum_{n:\ \prod r^{n}>x}\sum_{m''}N(n,m,m'')$. Rank and unrank use preorder in the
forward tree from $G(b,c)$: the subtree below a node of product $R$ and mode $m'$ has
$W_{m'}(\theta/R)$ internal nodes, of which $W_{m'}(\theta/R)-W_{m'}(\theta/(r_{c,b}R))$
have product in the required window; summing the sizes of earlier sibling subtrees along
the $O_H(p)$ letters of the address gives the ordinal, and the inverse operation chooses
children by comparison. Offsets for $(b,c)$ blocks are the corresponding block counts.
\end{proof}


\begin{proposition}[Streaming type counts with linear scratch]\label{III:M:rev13:prop:euler-count}
Let $n=(n_{a,m})$ have total length $L$, and let $N(n,s,t)$ count
chronological words from $s$ to $t$ with those edge multiplicities.
Let $d_m=\sum_a n_{a,m}$ and $d_m^{\rm in}=\sum_{(a,b):G(b,a)=m}n_{a,b}$.
For $L>0$ first require
$d_m-d_m^{\rm in}=[m=s]-[m=t]$. Restrict to the active vertices
(incident to a positive edge, together with $s,t$), and let $T_t(n)$ be
the weighted number of directed spanning trees oriented toward $t$,
where an edge $(m,a)$ has weight $n_{a,m}$ and loops cannot enter a tree.
If the balance condition fails, or an active $m\ne t$ has $d_m=0$,
then $N=0$. Otherwise
\[
 N(n,s,t)=T_t(n)\,
 \frac{\displaystyle\prod_m\binom{d_m}{(n_{a,m})_{a\in\Sigma}}}
      {\displaystyle\prod_{m\ne t}d_m}.
\]
The empty-word value is $N(0,s,t)=[s=t]$.
For fixed model size this formula can be evaluated in polynomial bit time
using $O_H(p)$ scratch when $L=O_H(p)$. Consequently streaming the type
vectors gives polynomial-time rank/unrank and updates with linear scratch,
without memoising a polynomial-size DP table.
\end{proposition}
\begin{proof}
Label temporarily all copies of every edge. In an Euler trail ending at $t$,
the last exit from each vertex other than $t$ forms an arborescence toward
$t$. Conversely, choose such a tree, place its edge last in the outgoing
list at each non-root vertex, and order all remaining outgoing edges.
Following these lists from $s$ uses every edge: if it stopped too early,
the balance equations force a stop at $t$. The remaining edge degrees
are then balanced. Any unused outgoing list has its tree edge unused;
following these last-exit edges, the remaining degree balances would force
an unused exit at $t$, contradicting the stop. Thus the labelled count is
$T_t(n)\prod_m(d_m-1+[m=t])!$.
Each unlabelled word has exactly $\prod_{a,m}n_{a,m}!$ labelled versions.
Dividing and rewriting factorial ratios as multinomials proves the formula,
including the case $s=t$. A zero tree count gives zero walks.

For fixed $|\mathcal R|$, $T_t(n)$ can be evaluated by explicitly summing
all finitely many choices of one outgoing edge type at every non-root
vertex and testing whether they form a tree. Each summand is a product of
at most $|\mathcal R|-1$ integers bounded by $L$. This avoids any reliance
on a numerical determinant. Multinomials are formed as products of
binomial coefficients using exact incremental divisions. Their product is
at most $d^L$; the tree factor and denominator have only $O_H(\log L)$
bits. Hence every retained integer has $O_H(L)$ bits. In contrast, storing
uncancelled $L!$ would require $\Theta(L\log L)$ bits.
Enumerating the $O_H(p^{d|\mathcal R|})$ type vectors needs only the
current vector, count and sum. The sums $W_m(x)$ and the subtree differences
in Lemma~\ref{III:M:rev13:lem:mode-dp} therefore use $O_H(p)$ scratch. Recomputing them
along an $O_H(p)$-letter path preserves polynomial time. With elementary
integer arithmetic an ample bound is $O_H(p^{d|\mathcal R|+6})$ per
rank/unrank or update. This exponent depends on the fixed model and is not
a uniform low-degree bound.
\end{proof}



\subsection{The theorem}
\begin{theorem}[Pressure upper bound and conditional persistent-bit optimum]\label{III:M:rev13:thm:mode}
Assume the setting above, a strongly connected mode graph, and the effective evaluation
hypotheses of Proposition ``Rational addresses with a mode checkpoint'' (certified
$q$-bit evaluators for $x_0$ and the mode-dependent affine data, a readout uniformly
$C$-Lipschitz in the continuous coordinate within each mode). Let $s^\ast$ be defined by
\eqref{III:M:rev13:eq:sstar} and $\theta=2^{-(p+c_0)}$ with $CD2^{-c_0}\le1/2$.
\begin{enumerate}[label=(\roman*)]
\item The memory $(b_t,\operatorname{rank}(v_t))$ of Lemma~\ref{III:M:rev13:lem:mode-update} is an
autonomous finite-state executable with exact mode recovery, output error $2^{-p}$, and
\[
B(p)\le s^\ast p+\lceil\log_2|\mathcal R|\rceil+O_H(1)
\]
persistent bits; its update and readout are polynomial in $p$ (Lemma~\ref{III:M:rev13:lem:mode-dp}
plus the stipulated evaluators) with $O_H(p)$ scratch by Proposition~\ref{III:M:rev13:prop:euler-count} and streaming replay. No exponential
word or transition table is stored.
\item Let $\mathcal R_{\rm reach}$ be the reachable modes and assume, as in Theorem
``Conditional joint persistent-bit optimum'', that the closures $X_r$ of the reachable
continuous fibres are $s^\ast$-Ahlfors regular with common constants, that
$d_{\rm fut}((r,x),(r,y))\ge c\Mnorm{x-y}$ within modes and
$d_{\rm fut}((r,x),(r',y))\ge2\delta_{\mathcal R}$ across modes. Then for
$2^{-p}<\delta_{\mathcal R}$,
\[
\log_2|\mathcal R_{\rm reach}|+s^\ast p-O_H(1)\ \le\ B^\ast(2^{-p})\ \le\
\log_2|\mathcal R_{\rm reach}|+s^\ast p+O_H(1).
\]
\item $s^\ast\le\bar s$, where $\bar s$ solves $\sum_a\bar r_a^{\,\bar s}=1$ with
$\bar r_a=\max_mr_{a,m}$ a pointwise dominating letter-only ratio; the inequality is strict whenever some ratio depends on the mode along a cycle of
the mode graph.
\end{enumerate}
\end{theorem}
\begin{proof}
(i) Lemma~\ref{III:M:rev13:lem:mode-update} gives exact modes and continuous error $\le D\theta$,
hence output error $\le CD\theta\le2^{-p-1}$ by the choice of $c_0$; the certified
replay error allocation of Proposition~\ref{III:M:audit:prop:checkpoint} contributes the other
$2^{-p-1}$ because only the mode-dependent constants change. The address ordinal ranges over
$\mathcal A_\theta$, so Lemma~\ref{III:M:rev13:lem:mode-count} gives
$\lceil\log_2|\mathcal A_\theta|\rceil\le s^\ast(p+c_0)+\log_2\frac{(1+d|\mathcal R|)C_u}{d-1}
+s^\ast\log_2\frac1{r_{\min}}+1$; the base mode is stored separately or folded into the
ordinal. Lemma~\ref{III:M:rev13:lem:mode-dp} gives the costs.
(ii) Indexing only reachable base modes in (i) gives the upper bound. For the lower bound
the proof of Theorem ``Conditional joint persistent-bit optimum'' applies verbatim: in
each reachable fibre closure, Ahlfors regularity of exponent $s^\ast$ yields a
$(4\varepsilon/c)$-separated set of cardinality $\ge a_H\varepsilon^{-s^\ast}$, approximated
by reachable states of the same mode within $\varepsilon/(2c)$; within-mode observability
and the inter-mode margin make all these histories pairwise future-separated by more than
$2\varepsilon$, and behavioural packing forces $|\mathcal R_{\rm reach}|a_H\varepsilon^{-s^\ast}$
distinct codes.
(iii) $M(s)\le\overline M(s)$ entrywise, where
$\overline M(s)_{m,m'}=\sum_{a:G(m,a)=m'}\bar r_a^{\,s}$ has all row sums
$\sum_a\bar r_a^{\,s}$, equal to $1$ at $s=\bar s$; hence
$\rho(M(\bar s))\le\rho(\overline M(\bar s))=1$ and $s^\ast\le\bar s$. If some
$r_{a,m}<\bar r_a$ on an edge of a cycle, irreducibility makes the inequality strict.
\end{proof}


\begin{remark}[Separation and startup are substantive hypotheses]
The invariant graph-directed fibres satisfy
\[
 K_r=\bigcup_{(m,a):G(m,a)=r}\tau_{a,m}(K_m).
\]
Their construction matrix is $M(s)^\top$, which has the same spectral
radius. A mass construction for these incoming cylinders uses a positive
Perron vector of $M(s^\ast)^\top$, rather than the outgoing-tree vector
$u$ used in Lemma~\ref{III:M:rev13:lem:mode-count}. With incoming strong separation
and strong connectivity the usual cylinder argument gives
$s^\ast$-regular fibres; compare \citet{mauldin_williams}.
If $x_0\in K_{r_0}$, these fibres are the closures of reachable states:
to approximate a chosen length-$n$ incoming cylinder, first reach its base
mode from $r_0$ and then follow that cylinder; contraction makes the initial
point immaterial as $n\to\infty$.

For $x_0\notin K_{r_0}$ the extra startup orbit need not be finite. For
example the one-mode Fibonacci system at $x_0=3/10$ has infinitely many
reachable points $(3/10)4^{-n}$ outside $K$. They belong to scaled copies
of the gap $(1/4,1/2)$. One must therefore not infer Ahlfors regularity of
the entire reachable closure from ``finitely many startup points''.
For the lower bound, it also suffices to assume regular observable compact
subsets of the reachable fibre closures, rather than regularity of the
whole closures.

Without separation/observability there is no matching lower bound from
pressure alone. For one mode with two identical maps $\tau_a(x)=x/2$ and
$x_0=0$, $s^\ast=1$ but one state reproduces the continuous state exactly.
\end{remark}


\begin{remark}[Reducible mode graphs]
Without strong connectivity let $s^\ast=\sup\{s:\rho(M(s))\ge1\}$ over the modes reachable
from $r_0$. For every $s>s^\ast$ the vector $u=(I-M(s))^{-1}\mathbf 1>0$ satisfies
$M(s)u\le u$, and the induction of Lemma~\ref{III:M:rev13:lem:mode-count} gives
$W_m(x)\le C_s x^{-s}$, hence $B(p)\le sp+O_{H,s}(1)$ for every $s>s^\ast$; two
components with the same exponent in sequence can add a $\log p$ term, so the clean
$O_H(1)$ statement is the irreducible one.
\end{remark}



\begin{remark}[What this does not settle]
Feedback-driven modes (guards on the continuous state) still need the branch cache or a
certified reconstruction; the mode-dependent counting above assumes letter-driven $G$,
exactly as the checkpoint proposition. Determining the optimal branch-cache coefficient in the general feedback
setting remains open.
\end{remark}

\section{From exact algebra to a finite executable}\label{III:U:sec:execution}
The general certificate returns a realization in exact algebraic arithmetic.
Separate the three error contributions as follows:
\[
 \sup_w|f(w)-\widehat g(w)|
 \le\epsilon_{\rm reduction}+\epsilon_{\rm coefficients}
                                      +\epsilon_{\rm execution}.
\]
The two linear errors can be certified using the joint Stein construction.
The rounded executable has a separate nonlinear state recurrence.

\subsection{A sufficient integer-execution certificate}
Use a row state initialized at $c^*$, updated by right multiplication by
$T_a$, and read against $b$; this respects the product order in
$g(w)=c^*T_wb$. The following standard positive-majorant bound is a useful
reference certificate.
\begin{proposition}[Componentwise execution tube]\label{III:U:prop:tube}
Let $M_{ij}=\max_a|(T_a)_{ij}|$ and suppose $\rho(M)<1$.
Use column vectors for componentwise absolute values. If each machine
step incurs local state error at most $\eta\ge0$, and the initial error
is at most $e_0$, then
\[
 e_t\le (M^\top)^te_0+\sum_{j=0}^{t-1}(M^\top)^j\eta.
\]
If $e_0\le E=(I-M^\top)^{-1}\eta$, then $e_t\le E$ for all $t$,
and the state contribution to output error is at most $|b|^\top E$.
An invariant state-radius vector is
$R=(I-M^\top)^{-1}(|x_0|+\eta)$.
\end{proposition}
\begin{proof}
Taking componentwise absolute values after one update gives
$e_{t+1}\le M^\top e_t+\eta$. Iterate and use the nonnegative Neumann
series for the inverse. Also $R\ge|x_0|$ and
$M^\top R+\eta=R-|x_0|\le R$, proving the state invariant.
The output bound is the triangle inequality.
\end{proof}

Unbounded verification of quantized filters has established precedents
\citep{cox2014}. The local bound must cover the whole invariant region, including products,
accumulation, coefficient representation, and output evaluation. For a fixed
integer kernel, the accumulator width is checked against all products and
sums on that region. Tie handling and operation order are part of the
machine specification. Saturation is a different machine and requires its
own proof. Initial state and readout quantization are charged explicitly.

This sufficient condition is stronger than mean-square stability and may
fail for a stable source. It is not a general replacement for the Hankel
certificate. Nor is its geometric error accumulation a new principle.
Searching for a useful invariant, choosing a base, and allocating bits are
separate optimization tasks.

\section{Polynomial transfer computation: the remaining completeness gap}\label{III:U:sec:oracle}
Positive contracting two-state predictive models represented by homographic
maps admit a different computational question: when can word enumeration in
a causal state encoder be replaced by a polynomial transfer calculation?
The following theorem gives a sufficient analytic certificate. Its existence
does not follow from the exact-bit family.

\subsection{A two-radius analytic certificate}
On normalized coordinates $x\in[-1,1]$, write
\[
 G_b(x)=\frac{\alpha_b+\beta_bx}{1+\xi_bx},\qquad
 g_b(x)=\kappa_b(1+\xi_bx)^{-2},\qquad 0<\kappa_b<1.
\]
Suppose rational $1<w<R$ satisfy, for each branch,
\[
 |\xi_b|R<1,\qquad \max(|G_b(-R)|,|G_b(R)|)\le w.
\]
A real homography maps the boundary circle to a circle centered on the
real axis; its real endpoints give the maximum modulus. Thus
$G_b(\overline D_R)\subset\overline D_w$. Strict invariance on the
real interval and poles outside it guarantee that such radii exist near
one, though their quantitative margins remain model-dependent.

For $\sigma\in[s_-,s_+]\subset(0,\infty)$ define
\[
 (\mathcal L_\sigma P)(x)=\sum_{b=0}^1g_b(x)^\sigma P(G_b(x)).
\]
Use the analytic logarithm equal to zero at $x=0$ for $1+\xi_bx$.
With $A_* =\lceil2s_+\rceil$ one may take
\[
 W=\left\lceil\sum_b(1-|\xi_b|R)^{-A_*}\right\rceil,
 \qquad\Lambda=\frac W{1-w/R},\qquad
 \delta_m=\frac{WR^{-m}}{R-1}.
\]
For $\|P\|_w=\sum_j|c_j|w^j$ and coefficient truncation $\Pi_m$,
Cauchy's bound gives
\[
 \|\Pi_m\mathcal L_\sigma P\|_w\le\Lambda\|P\|_w,
 \qquad\|(I-\Pi_m)\mathcal L_\sigma P\|_{\infty,[-1,1]}
                                      \le\delta_m\|P\|_w.
\]
The section matrix can be formed without enumerating words. Its entries are
\[
 M_{kj}=\sum_b\kappa_b^\sigma
 \sum_{h=0}^{\min(j,k)}\binom jh\alpha_b^{j-h}\beta_b^h
 (-\xi_b)^{k-h}\frac{(2\sigma+j)_{k-h}}{(k-h)!}.
\]
For one branch, compute column zero recursively and subsequent columns by
\[
 a_{k+1}^{(0)}=-\xi_b\frac{2\sigma+k}{k+1}a_k^{(0)},\qquad
 a_k^{(j+1)}=\alpha_ba_k^{(j)}+\beta_ba_{k-1}^{(j)}
                                      -\xi_ba_{k-1}^{(j+1)}.
\]
This costs $O(m^2)$ scalar operations. Real powers and rounding still
require certified enclosures in a finite implementation.

\begin{theorem}[Correct propagation of repeated truncations]\label{III:U:thm:oracle}
Assume the real branch maps satisfy $G_b([-1,1])\subseteq[-1,1]$ for $b=0,1$, and the two-radius analytic bounds above hold for fixed $1<w<R$.
Let $V_n=\mathcal L_\sigma^n1$, $P_0=1$, and
$P_{n+1}=\Pi_m\mathcal L_\sigma P_n$. Assume also $g_b(x)\le1$
on $[-1,1]$. With $K=\max(2,\Lambda)$,
\[
 \|V_n-P_n\|_\infty\le n\delta_mK^{n-1}.
\]
If $g_b(x)\ge g_->0$, $g_-<1$, put $a_*=2g_-^{s_+}$.
Then $V_n\ge a_*^n$, and relative error at most $\eta$ follows when
\[
 R^m\ge\frac{nWK^{n-1}}{(R-1)\eta a_*^n}.
\]
Consequently $m=O_H(n+\log(1/\eta)+\log(n+1))$ suffices for fixed
model $H$ and its certified margins.
\end{theorem}
\begin{proof}
The polynomial norm grows at most as $\Lambda^n$.
Positivity and the two branches give
$\|\mathcal L_\sigma f\|_\infty\le2\|f\|_\infty$.
Hence $E_{n+1}\le2E_n+\delta_m\Lambda^n$, with $E_0=0$.
Summing
$\sum_{j=0}^{n-1}2^{n-1-j}\Lambda^j\le n\max(2,\Lambda)^{n-1}$
yields the displayed bound. Each real branch weight is at least
$g_-^{s_+}$, so induction gives $V_n\ge a_*^n$.
\end{proof}


The distinction $M^n1\ne\Pi_m\mathcal L_\sigma^n1$ is essential:
truncation and iteration generally do not commute. The propagation bound
above is therefore necessary.

\subsection{What is established and what remains conditional}
The exact-arithmetic theorem removes exponential word enumeration for this
finite-horizon transfer calculation. In the intended regime
$n=O_H(\log p)$, it gives degree $O_H(\log p)$ under the stated relative
accuracy requirements. It does not by itself prove that a numerical
candidate generator always returns a polynomial accepted by the downstream
rational encoder certificate.

The desired $O_H(p^2\log p)$ bit cost depends on fast arithmetic and
specific subroutine cost bounds. Ordinary arithmetic gives a different
bound. The compiler, stored polynomial, precision, and online update
must all be charged. The unresolved obligation is a terminating, uniformly certified construction over the
declared model class with quantitative constants, including a proof that
the proposed acceptance conditions can always be met. Successful instances
do not establish that obligation.

\paragraph{A second logical boundary: dimension versus bounded overhead.}
A box dimension $h$ generally gives
$\log N(\epsilon)=h\log(1/\epsilon)+o(\log(1/\epsilon))$, not an
$O(1)$ remainder. For $K=\{0\}\cup\{2^{-j}:j\ge1\}$, the dimension
is zero but the covering number at scale $2^{-p}$ grows on the order of
$p$, giving $\log_2N\asymp\log_2p$. Stronger two-sided covering or
separation hypotheses are required for a constant-overhead memory law.
\subsection{Certified arithmetic error without a completeness claim}
\label{III:U:sec:transfer-rounding}
Strictly contracting branches do not make an unnormalized transfer
operator a contraction toward a positive fixed function.
For example, take
\[
 G_0(x)=x/4-1/2,\quad G_1(x)=x/4+1/2,\quad
 g_0=g_1=1/4,\quad \sigma=1.
\]
These branches are separated on $[-1,1]$ and satisfy the two-radius
assumptions with $R=2$, $w=3/2$, and $\xi_b=0$. Nevertheless,
$\mathcal L_\sigma1=1/2$ and $V_n=2^{-n}$, so there is no positive
limit of $V_n$. A positive
eigenfunction is not, without normalization, a fixed function.
Even after normalizing the leading eigenvalue, a quantitative
convergence rate and its normalization must be proved separately.

\begin{theorem}[Finite-horizon rounding and conditional eigenfunction error]
\label{III:U:thm:transfer-rounding}
Under Theorem~\ref{III:U:thm:oracle}, let
$A_m=\Pi_m\mathcal L_\sigma$ and compute degree-at-most-$m$
polynomials by
\[
 \widehat P_0=1,\qquad
 \widehat P_{j+1}=A_m\widehat P_j+e_j,\qquad
 \|e_j\|_w\le r_j.
\]
The certified errors $e_j$ must include matrix-coefficient approximation,
arithmetic rounding, and serialization. Then
\begin{equation}
 \|V_n-\widehat P_n\|_{\infty,[-1,1]}
 \le n\delta_m K^{n-1}
                   +\sum_{j=0}^{n-1}\Lambda^{n-1-j}r_j.
 \label{III:U:eq:transfer-rounded-error}
\end{equation}
If, in addition, a specified positive function $h\ge h_{\min}>0$
and constants $C<\infty$, $0<\vartheta<1$ satisfy the independently
proved normalization and convergence bound
$\sup_x|V_n(x)-h(x)|/h(x)\le C\vartheta^n$, then
\begin{equation}
 \sup_x\frac{|\widehat P_n(x)-h(x)|}{h(x)}
 \le C\vartheta^n+
 \frac{n\delta_mK^{n-1}
                 +\sum_{j=0}^{n-1}\Lambda^{n-1-j}r_j}{h_{\min}}.
 \label{III:U:eq:transfer-conditional-eigenfunction}
\end{equation}
No such convergence assumption is inferred from branch contraction alone.
\end{theorem}
\begin{proof}
For the exact truncated iteration $P_j=A_m^j1$, repeated substitution
gives $\widehat P_n-P_n=\sum_{j=0}^{n-1}A_m^{n-1-j}e_j$.
Since $\|A_m\|_{w\to w}\le\Lambda$ and
$\|P\|_{\infty,[-1,1]}\le\|P\|_w$, its norm is bounded by the
sum in \eqref{III:U:eq:transfer-rounded-error}. Add
Theorem~\ref{III:U:thm:oracle}'s finite-horizon error. Dividing by
$h_{\min}$ and adding the separately assumed horizon-convergence error
proves \eqref{III:U:eq:transfer-conditional-eigenfunction}.
\end{proof}


An $O(m^2)$ construction of the section matrix does not account for
$n$ repeated dense matrix--vector products, which cost $O(nm^2)$
scalar operations without further structure. Bit precision, evaluation
of real powers, interval dependency, and verified inversion also need
their own bounds. A small polynomial approximation error is not by
itself proof of acceptance by the complete state compiler.
Accordingly, universal compiler termination, a sufficient precision rule, and uniform offline/online bit bounds remain unestablished.


\subsection{Normalized transfer convergence on a strict log-Lipschitz cone}
\label{III:U:sec:normalized-transfer}
The endpoint value $L=c/(1-\gamma)$ gives
only cone invariance. Strict inclusion requires a strictly larger
$L$. Moreover the relevant order is that of the regularity cone,
not only pointwise positivity in $C([-1,1])$.

\begin{theorem}[Quantitative normalized positive eigenfunction]
\label{III:U:thm:normalized-transfer}
On $I=[-1,1]$, let
$(\mathcal L f)(x)=\sum_{a=1}^d w_a(x)f(G_a(x))$, where
$G_a:I\to I$ have Lipschitz constants at most $\gamma<1$,
$0<w_-\le w_a(x)\le w_+$, and
$|\log w_a(x)-\log w_a(y)|\le c|x-y|$.
Choose $L>c/(1-\gamma)$ and put $L'=c+\gamma L<L$.
Define
\[
 t_* =\frac{L-L'}{L+L'}e^{-2(L+2L')},\qquad
 \Delta=2\log(1/t_*),\qquad \vartheta=\tanh(\Delta/4)<1.
\]
There is a unique eigenfunction $h>0$ in the log-Lipschitz cone of
constant $L$, normalized by $h(0)=1$. It satisfies
\[
 \mathcal L h=\lambda h,\quad
 e^{-L'}\le h(x)\le e^{L'},\quad
 d w_-e^{-L'}\le\lambda\le d w_+e^{L'}.
\]
Writing $V_n=\mathcal L^n1$, for $n\ge1$,
\begin{equation}
 \sup_{x\in I}\left|\frac{V_n(x)/V_n(0)}{h(x)}-1\right|
       \le e^{\Delta\vartheta^{n-1}}-1.
 \label{III:U:eq:normalized-transfer}
\end{equation}
\end{theorem}
\begin{proof}
Let $K_L$ be the closed convex cone of nonnegative Lipschitz functions
with $f(x)\le e^{L|x-y|}f(y)$ for all $x,y$; its only function
vanishing at a point is zero. Each positive summand of $\mathcal L f$
has log-Lipschitz constant at most $c+\gamma L$, and their sum has
the same bound. Thus $\mathcal L(K_L\setminus\{0\})\subset K_{L'}$.

For $f,g\in K_{L'}$ normalized at zero, $g(y)/f(y)\ge e^{-2L'}$,
and hence also $g(y)/f(y)\ge e^{-4L'}$. For $r=|x-y|\in(0,2]$,
\[
 \frac{e^{Lr}-e^{L'r}}{e^{Lr}-e^{-L'r}}
 \ge\frac{L-L'}{L+L'}e^{-2L}.
\]
This follows by expressing numerator and denominator as integrals
of $r e^{ur}$. Therefore
$e^{Lr}g(y)-g(x)\ge
t_*[e^{Lr}f(y)-f(x)]$. Pointwise positivity also holds, so
$g-t_*f\in K_L$, and symmetrically $f-t_*g\in K_L$.
The image diameter in the Hilbert metric of $K_L$ is at most
$2\log(1/t_*)=\Delta$.

For projective contraction, use the Lipschitz Banach space: the strict
constant $L'<L$ places these positive image functions in the interior
of $K_L$. The normalized slice $f(0)=1$ of $K_L$ is convex and compact in the
uniform norm by uniform bounds and equicontinuity. The continuous
map $f\mapsto\mathcal L f/(\mathcal L f)(0)$ has a fixed point
by Schauder. Birkhoff's projective contraction theorem
\citep{gaubert2015} gives contraction factor at most $\vartheta$ on
comparable rays. Every eigenfunction in $K_L$ lies in the image,
whose diameter is finite, so this also gives uniqueness.
Both $\mathcal L1$ and $\mathcal Lh$ are image points; iteration
gives projective distance at most $\Delta\vartheta^{n-1}$.
Cone order implies pointwise order. Since both representatives have
value one at zero, their pointwise ratios are between $e^{-D}$ and
$e^D$ when their projective distance is at most $D$. This proves
\eqref{III:U:eq:normalized-transfer}. The fixed point lies in $K_{L'}$;
its normalization gives the bounds on $h$, and evaluation at zero
gives those on $\lambda$.
\end{proof}


\begin{corollary}[A certified polynomial eigenfunction approximator]
\label{III:U:cor:normalized-polynomial}
Assume the two-radius analytic hypotheses of Theorem~\ref{III:U:thm:oracle}
and the hypotheses of Theorem~\ref{III:U:thm:normalized-transfer}, with
fixed effective margins. Suppose the section coefficients can be
enclosed to $B$ bits in time polynomial in $m,B$ and their input
description. Then, for every $p\ge1$, a terminating algorithm returns
a dyadic polynomial $\widehat h_p$, positive on $I$, together with
certified error $\|\widehat h_p-h\|_\infty\le2^{-p}$.
Its horizon, degree and coefficient precision are $O_H(p)$.
Compilation has polynomial bit cost. Horner evaluation takes
$O_H(p\mathsf M(p))$ bit operations and $O_H(p)$ scratch bits,
with $O_H(p^2)$ bits of stored polynomial coefficients; here
$\mathsf M(p)$ is the cost of multiplying $O(p)$-bit integers.
These costs concern the eigenfunction approximator, not the full state compiler.
\end{corollary}
\begin{proof}
Choose $n=O_H(p)$ so that the right side of
\eqref{III:U:eq:normalized-transfer}, multiplied by $e^{L'}$, is at most
$2^{-p-2}$. Let $a_*=2g_-^{s_+}$ be the lower growth bound of
Theorem~\ref{III:U:thm:oracle}. For a rounded polynomial iterate with
$E_n=\|\widehat P_n-V_n\|_\infty\le a_*^n/2$, normalization gives
\[
 \left\|\frac{\widehat P_n}{\widehat P_n(0)}-
                  \frac{V_n}{V_n(0)}\right\|_\infty
 \le 2(1+e^{L'})E_n/a_*^n.
\]
Indeed $V_n/V_n(0)\in K_{L'}$ normalized at zero and
$\widehat P_n(0)\ge a_*^n-E_n$.
Take $E_n\le a_*^n2^{-p-2}/[2(1+e^{L'})]$.
Equation~\eqref{III:U:eq:transfer-rounded-error}, with
$\delta_m=O_H(R^{-m})$, then permits $m=O_H(p)$ and per-step
certified arithmetic error $2^{-O_H(p)}$. The polynomial coefficient
norms grow at most exponentially in $n$ and normalization divides
by a quantity at least $a_*^n/2$, also exponential in $n$.
Thus $O_H(p)$ coefficient bits suffice after allowing for the
$O(p)$ operations per row and final dyadic rounding of the normalized
coefficients, whose error is included in the budget. This can be
implemented by interval arithmetic with directed dyadic rounding
and the stipulated coefficient enclosures; double precision until
the explicit error tests pass. The preceding bounds prove termination.

If necessary request the smaller final error
$\min(2^{-p},e^{-L'}/2)$; this changes precision only by a model
constant and proves positivity from $h\ge e^{-L'}$.
There are $O(nm^2)=O_H(p^3)$ dense scalar operations after section
construction, on $O_H(p)$-bit numbers; this and the coefficient oracle
give polynomial compilation cost. There are $O(p)$ coefficients of
$O(p)$ bits. Horner evaluation at a rational argument known to the
required $O_H(p)$ bits uses $O(p)$ multiplications; the coefficient
norm bound controls error accumulation. Its working accumulator uses
$O_H(p)$ bits. Immutable coefficients are charged separately.
\end{proof}

For the homographic weights of Section~\ref{III:U:sec:oracle}, one can use
$c=\max_a 2s_+|\xi_a|/(1-|\xi_a|)$ and a verified branch derivative
bound $\gamma<1$. Rational fixed exponents and rational positive
parameters give effective power enclosures by standard interval
root or logarithm/exponential algorithms. Alternatively their
certified coefficient oracle is an explicit input to the corollary.

\paragraph{Implications for state compilation.}
The normalized target, a positive lower bound, geometric convergence,
and its polynomial approximation are now supplied under explicit
effective hypotheses. A state compiler has additional acceptance
conditions, inverse-coordinate operations and inductive update
relations. The above theorem does not verify them. In particular its
fast Horner cost cannot be relabeled the cost of every compiled state update.

\section{Controlled recurrent states and exact residual coding}\label{III:U:sec:gdn}
A word-Hankel theorem for exogenously supplied letters is not automatically
a theorem for a trained decoder whose controls depend on its current
activations. This section isolates the valid matrix identity from that
unproved transfer. Gated DeltaNet combines a decay gate with a delta-rule
state update~\citep{yang2025gated}. The controlled example below uses
the unit-decay delta rule; it does not assume that the learned gates of
a trained GDN layer realize that special case.

\subsection{An executable core and a reduced residual}
Write a controlled recurrence in key-by-value orientation as
$S^+=AS+U$. Let $c$ be a retained core, $D(c)$ a computable matrix
predictor, $B^\top B=I_r$, and $Z=(S-D(c))B$.
\begin{proposition}[Closed projected residual update]\label{III:U:prop:core}
Under the same controls and the same supplied core trajectory,
\[
 Z^+=AZ+[AD(c)+U-D(c^+)]B,
 \qquad \widehat y=q^\top D(c)+q^\top ZB^\top.
\]
The readout error is
$q^\top(S-D(c))(I-BB^\top)$.
\end{proposition}
\begin{proof}
Substitute $S^+=AS+U$ into $(S^+-D(c^+))B$ and add and subtract
$AD(c)B$. The readout follows by projecting the residual on the value
side. No invariant-subspace assumption on $\operatorname{range}D$ is
needed for this identity.
\end{proof}

The core update must itself be computed from available deployed quantities.
It cannot silently query a deleted reference state. Further, the identity
is over exact arithmetic; matching FP32 bits requires an explicit
operation-order analysis or a reversible word-level residual.

\begin{proposition}[Weighted core-conditioned prediction]
Let $O$ be a specified frozen future-readout operator and $W=O^\top O$.
Assume the needed weighted second moments exist, and take a finite-valued
core. Define
\[
 D_*(c)=\mathbb E[W\mid c]^\dagger\mathbb E[WS\mid c],\qquad
 M_c=\mathbb E[(S-D_*(c))^\top W(S-D_*(c))].
\]
Then
\[
 \inf_{D,\,B^\top B=I_r}
 \mathbb E\|O(S-D(c))(I-BB^\top)\|_F^2
 =\sum_{j>r}\lambda_j(M_c).
\]
\end{proposition}
\begin{proof}
Conditional weighted normal equations eliminate the cross terms. For
any $D$, its residual Gram equals $M_c$ plus the positive semidefinite
matrix
$\mathbb E[(D-D_*)^\top W(D-D_*)]$. Null directions of
$\mathbb E[W\mid c]$ are also null for $\mathbb E[WS\mid c]$, so the
pseudoinverse solves the normal equations. Choosing $D_*$ and the top
$r$ eigenvectors of $M_c$ minimizes the omitted trace.
\end{proof}

This is conditional weighted regression plus principal components, not a
new general compression theorem. The operator $O$ must cover the future
queries and intervention cuts for which a guarantee is claimed. A richer
core can lower the optimal residual Gram, but its bits, predictor, basis,
and execution cost must be paid for.

\subsection{A constructive controlled example}
For key and value dimensions both $n\ge2$, take
$S=cI_n+a e_1^\top$, $D(c)=cI_n$, $B=e_1$, with a finite-valued
stored scalar $c$. This state is reachable from zero with $n$
rank-one delta writes: $k=e_i$, $\beta=1$, and $v=ce_i+a_i e_1$.
Subsequently allow a unit key $k$, $\beta\in[0,1]$, and
$v=ck+\xi e_1$. The unit-decay delta update gives
\[
 a^+=a+\beta k(\xi-k^\top a),\qquad
 y=cq^\top+(q^\top a)e_1^\top.
\]
Thus $c$ and $n$ residual coordinates suffice for exact continuation of
this controlled family without storing a dense $cI_n$ or dense basis.
If $c,a$ are independent centered variables, $\operatorname{Var}c=\sigma_c^2$,
and $\mathbb E\|a\|^2=V$, then for $W=I$ the raw value Gram is
$\sigma_c^2I+Ve_1e_1^\top$. A raw rank-$r$ projection loses
$(n-r)\sigma_c^2$ for $r\ge1$, whereas the core-conditioned residual
has rank one. This is an executable synthetic separation, not a finding
that Qwen exhibits this family.

\subsection{The missing mathematical bridge to encoded length}
An exact residual supplies injectivity. The predictor changes residual
word statistics, not exactness. A future-readout norm can ignore state
components that a lossless full-state codec must retain. Small floating
values can still contain many independent mantissa bits. A positive coding
theorem therefore needs assumptions on the encoded integer residual,
small-magnitude behavior, predictor description length, and adaptation
cost. A small Hankel singular tail alone does not provide them.

The next useful experiment is a frozen comparison of a deployable dynamic
core, an energy predictor, cross approximation, and existing exact residual codecs, with the
same residual format, complete bytes, and complete update/read cost.
Selecting on total encoded length is a new candidate that needs new held-out
data. For a trained-model claim, controls, gates, normalization, all caches,
logits, and free-generation token divergence must be covered, and the
deleted native state must actually be absent from resident memory.
\subsection{A covariance surrogate and what exact coding actually permits}
\label{III:U:sec:codec-covariance-audit}
Covariance does not determine exact expected code length. The valid
algebraic inequality is the following surrogate bound.

\begin{proposition}[Classical diagonal-covariance surrogate inequality]
\label{III:U:prop:covariance-surrogate}
For $\Gamma\succeq0$, $a,b>0$, and a real orthogonal matrix $U$,
\begin{equation}
 \sum_j\log_2\bigl(a+b(U^\top\Gamma U)_{jj}\bigr)
 \ge\log_2\det(aI+b\Gamma)
 =\sum_j\log_2(a+b\lambda_j(\Gamma)).
 \label{III:U:eq:covariance-surrogate}
\end{equation}
Equality holds exactly when $U^\top\Gamma U$ is diagonal, including
arbitrary rotations within repeated-eigenvalue eigenspaces.
This is an inequality for a covariance surrogate, not an exact
lossless-entropy formula.
\end{proposition}
\begin{proof}
Apply Hadamard's determinant inequality to the positive definite matrix
$aI+bU^\top\Gamma U$. Its determinant is invariant under the orthogonal
change of basis. Equality in Hadamard's inequality holds exactly when
the Gram vectors are orthogonal, equivalently when this matrix is
diagonal. Taking base-two logarithms proves the result.
\end{proof}


\begin{proposition}[A valid variance bound for discrete lossless symbols]
\label{III:U:prop:discrete-entropy-bound}
Let $Z=(Z_1,\ldots,Z_D)$ be an integer-valued residual vector with
finite coordinate variances $v_j$. Then
\begin{equation}
 H(Z)\le\sum_jH(Z_j)
 \le\frac12\sum_j\log_2\!\bigl(2\pi e(v_j+1/12)\bigr).
 \label{III:U:eq:discrete-entropy-bound}
\end{equation}
An ideal prefix code on the whole vector has expected length between
$H(Z)$ and $H(Z)+1$; coding coordinates separately instead admits
length less than $\sum_jH(Z_j)+D$. All model, predictor, basis,
framing, and coding-distribution descriptions are additional costs.
The same statements apply conditional on available side information,
by applying the bounds to each conditional law and averaging.
\end{proposition}
\begin{proof}
The first inequality is entropy subadditivity. For an integer variable
$Z_j$, add an independent uniform variable on $[-1/2,1/2]$.
The resulting density has differential entropy $H(Z_j)$ in bits and
variance $v_j+1/12$. The Gaussian maximum-entropy bound gives the
second inequality; this is the classical discrete variance bound
\citep{massey1988}. The coding claims are the prefix-code entropy
lower bound and the Shannon lengths $\lceil-\log_2p\rceil$.
\end{proof}


\paragraph{Counterexamples and missing codec assumptions.}
Covariance alone cannot determine discrete entropy. The symmetric
variable $Z=\pm1$ has variance one and entropy one bit. In contrast,
\[
 P(Y_m=0)=1-m^{-2},\qquad
 P(Y_m=m)=P(Y_m=-m)=\frac1{2m^2}
\]
also has variance one but entropy
$h_2(m^{-2})/\log2+m^{-2}\to0$.
More directly, the two-point variable $Z=\pm1$ remains exactly
representable on grids $\Delta_q=1/m$ with entropy one bit, whereas
$\tfrac12\log_2(1+2\pi e/\Delta_q^2)$ diverges. An
$O(\Delta_q^2)$ remainder cannot repair this claimed equality for
such a discrete source.
Thus an equality inferred solely from covariance, or an exact saving
threshold based on it, is false without further distributional and coding assumptions.

Nor does a nonidentity readout weight force a different basis. For
example, $\Gamma=\Udiag(1,2)$ and a diagonal weight
$W=\Udiag(2,3)\ne I$ share the coordinate eigenbasis with their
corresponding weighted covariance. No strict surrogate gap follows
merely from $W\ne I$. A measured saving would require empirical evidence
or additional structural assumptions.

An arbitrary orthogonal change of basis followed by coordinate
quantization is not an exact integer-lattice codec. For example, rotate
$(1,0)$ and $(2,0)$ by $45$ degrees; their coordinates are respectively
$(1/\sqrt2,1/\sqrt2)$ and $(\sqrt2,\sqrt2)$, both rounding to
$(1,1)$ on the unit grid. Exact reconstruction requires an injective
representation, such as an appropriate reversible integer transform
or an additional exact residual. A finite-precision basis must itself
be serialized and its actual decoder used.

For a residual matrix $R$ with $n_k$ rows,
$\mathbb E[R^\top R]$ aggregates row second moments. It is not the
per-row covariance used by a scalar marginal-entropy formula; identical
centered row laws would require division by $n_k$. Likewise, the
dimension $n_vr-r(r+1)/2$ of a Stiefel manifold is not the bit length
of an exact basis header. Actual encodings, including a completion
basis when needed, determine that length.

For a fixed exact codec the genuine saving condition is simply
\[
 L_{\rm metadata}+L_{\rm core}+\mathbb E L_{\rm residual}
                         <\mathbb E L_{\rm baseline},
\]
with the same scope and all shared costs amortized explicitly.
Equation~\eqref{III:U:eq:covariance-surrogate} can guide a heuristic or bound
a specified surrogate; it cannot replace these code lengths. For an
injective discrete representation $S\mapsto(C,Z)$,
$H(S)=H(C,Z)=H(C)+H(Z\mid C)$: an invertible representation does
not lower the fundamental joint lossless entropy, although it may
improve a restricted practical coder or exploit already available side
information.


\subsection{Reversible integer lifting with an actual coding contract}
\label{III:U:sec:integer-lifting}
Integer lifting is classical \citep{calderbank1998}. Its relevance here
is that it repairs injectivity of the proposed residual transform;
it does not by itself establish a Hankel-specific compression gain.

\begin{lemma}[Constructing lifting factors from a supplied orthogonal basis]
\label{III:U:lem:rotation-lifts}
A supplied computable $U\in SO(D)$ can be approximated arbitrarily
closely by a product of at most $3D(D-1)$ rational shears. A planar
rotation of angle $\theta$ is
\[
 \begin{pmatrix}1&-\tan(\theta/2)\\0&1\end{pmatrix}
 \begin{pmatrix}1&0\\\sin\theta&1\end{pmatrix}
 \begin{pmatrix}1&-\tan(\theta/2)\\0&1\end{pmatrix}.
\]
Split angles if necessary so that $|\theta|\le\pi/2$; each shear
coefficient then has magnitude at most one. If all $\ell$ coefficients
are rounded within $\delta$ while staying in $[-1,1]$, the product
error in maximum-row-sum norm is at most $\ell2^{\ell-1}\delta$.
\end{lemma}
\begin{proof}
Givens elimination expresses a special orthogonal matrix as at most
$D(D-1)/2$ planar rotations, allowing angles up to $\pi$.
Splitting each into at most two proves the shear count by the displayed
identity. Telescope the product; all factors before and after rounding
have norm at most two, and each factor error has norm at most $\delta$.
\end{proof}


\begin{theorem}[Bijective lifting and a covariance bound including correlated errors]
\label{III:U:thm:integer-lifting}
Let $S_r=I+s_re_{i_r}e_{j_r}^\top$, $i_r\ne j_r$, be rational
shears, $1\le r\le\ell$, and use a fixed deterministic nearest-integer
rule. Apply $z_{i_r}\leftarrow z_{i_r}+\operatorname{round}(s_rz_{j_r})$
in the stated order. The resulting map $T:\mathbb Z^D\to\mathbb Z^D$
is a bijection, whose inverse applies the same rounded values with
subtraction in reverse order. Put $U_*=S_\ell\cdots S_1$ and
\[
 b_j=\frac12\sum_{r=1}^\ell
       |(S_\ell\cdots S_{r+1})_{j,i_r}|.
\]
For a random integer vector $Z$ of finite second moment and covariance
$M$, the output obeys
\begin{equation}
 \operatorname{Var}(T(Z)_j)
 \le\left(\sqrt{(U_*MU_*^\top)_{jj}}+b_j\right)^2.
 \label{III:U:eq:lifting-variance}
\end{equation}
If a desired orthogonal $U$ diagonalizes $M$ with diagonal entries
$\lambda_j$, the right side is at most
\[
 \left(\sqrt{\lambda_j}+
    \|(U_*-U)_{j,:}\|_2\sqrt{\Utr M}+b_j\right)^2.
\]
These assertions apply separately to conditional laws given a core.
\end{theorem}
\begin{proof}
The donor coordinate $j_r$ is unchanged by a shear, so subtraction
recovers its integer target exactly. Reverse composition proves
bijectivity, for any deterministic rounding convention.
Writing the rounding error at step $r$ as $\epsilon_r e_{i_r}$,
$|\epsilon_r|\le1/2$, gives
\[
 T(z)=U_*z+\sum_{r=1}^\ell
            S_\ell\cdots S_{r+1}e_{i_r}\epsilon_r(z).
\]
Thus the $j$th error lies in $[-b_j,b_j]$ and its standard deviation
is at most $b_j$. The triangle inequality in centered $L^2$ proves
\eqref{III:U:eq:lifting-variance}, without independence. For the last
bound, split $U_*=U+(U_*-U)$ and apply the same inequality, using
$\operatorname{Var}(vZ)\le\|v\|_2^2\Utr M$.
\end{proof}

The claim that each lift adds at most $1/4$ to the \emph{output}
variance is false. Let $Z_1=Z_2$ be Bernoulli$(1/2)$ and lift with
$s=51/100$. The exact linear target has variance $(151/100)^2/4$,
whereas the integer target is $2Z_1$, with variance one. The increase
is $17199/40000>1/4$, because the rounding error is correlated with
the signal. The propagated-error terms above are indispensable.

\begin{theorem}[Complete expected length against a specified baseline]
\label{III:U:thm:actual-lifting-code}
Let $S$ be a finite-bit integer vector and $C=c(S)$ a core. Given $C$,
an integer predictor and a specified lifting transform produce
$Z'=T_C(S-D(C))$. The decoder receives the core and every model,
predictor, transform and framing description needed for inversion;
their total nonpayload length is at most $B_{\rm meta}$ bits.
Fix the coding models before coding the evaluated source; any dependence
on the transmitted core is explicitly conditional. Let prefix codes
for $C$ and for each $Z'_j$ conditional on $C$ have
Shannon lengths for specified full-support models $Q_C,Q_{j|C}$.
Use base-two KL and put
\[
 D_C=D_2(P_C\|Q_C),\qquad
 D_j(c)=D_2(P_{Z'_j|c}\|Q_{j|c}).
\]
If $v_j(c)$ is any valid bound on the conditional variance, the
expected complete length is at most
\begin{align}
 \mathbb E L\le B_{\rm meta}+H(C)+D_C+1+D
 +\mathbb E_C\sum_{j=1}^D\left[
 \tfrac12\log_2(2\pi e(v_j(C)+1/12))+D_j(C)\right].
 \label{III:U:eq:actual-lifting-code}
\end{align}
If this explicit upper bound is smaller than a proved lower bound
on the expected length of a \emph{specified} baseline coder, the new
coder improves that baseline in expectation at exact reconstruction.
\end{theorem}
\begin{proof}
The inverse lifting steps and the integer predictor reconstruct $S$
exactly from $C,Z'$. Expected Shannon lengths are bounded by cross
entropy plus one for each code: cross entropy is entropy plus the
displayed model mismatch. Proposition~\ref{III:U:prop:discrete-entropy-bound}
bounds each marginal conditional entropy by its variance term.
Sum these inequalities, average over the core, and add every
nonpayload description. The final comparison is then immediate.
\end{proof}

This is a coding theorem for the declared models, not an assertion
that their probabilities or description lengths are free to obtain.
An implementation must also specify finite integer widths. For
bounded source support, the lifting recurrence bounds every intermediate
magnitude by $\|S_r\|_\infty$ times the previous bound plus $1/2$;
these explicit bounds determine exact numerator, product and accumulator
widths. In-place lifting stores $D$ integers plus multiplication scratch;
holding additional input/output buffers must be charged. No reduction
of update latency or peak memory is implied by \eqref{III:U:eq:actual-lifting-code}.

\begin{proposition}[The universal saving threshold is impossible]
\label{III:U:prop:no-universal-code-gain}
For the deterministic exact representation above,
$H(C,Z')=H(S)$. Consequently a valid complete-code upper bound cannot
lie below a valid lower bound on $H(S)$. In particular a
criterion for beating \emph{every} lossless coder by placing its
complete length below an entropy lower bound has no realizable instance.
\end{proposition}
\begin{proof}
The encoder is deterministic and injective, and has a deterministic
inverse on its image; entropy is invariant. Any uniquely decodable
code has expected length at least $H(S)$. If $U$ bounds its complete
length above and $L$ bounds $H(S)$ below, then
$U\ge\mathbb E L_{\rm code}\ge H(S)\ge L$, so $U<L$ is impossible.
\end{proof}

For comparison, let $S=(Z,Z)$ with $Z$ uniform on $b$-bit integers.
The lift $(z_1,z_2)\mapsto(z_1,z_2-z_1)$ yields $(Z,0)$.
With the distribution and lift fixed, $b$ payload bits suffice instead
of a raw $2b$-bit layout; including a header of length $H_0$ still
saves bits if $H_0<b$. The joint entropy was already $b$, so this is
a gain over the specified raw layout, not over an optimal lossless code.

For a practical fixed frame, a one-bit selector between a new exact
codec and an exact baseline gives length at most
$1+\min(L_{\rm new},L_{\rm base})$, with all headers included.
This provides a deterministic fallback guarantee; it does not supply
held-out evidence for learned Hankel predictors.

\section{The homogeneous norm identity}
Here $d=|\Sigma|$ and $C_{i,j}(f)=(f(uv))_{|u|=i,|v|=j}$.
\begin{lemma}[Homogeneous norm and spectrum]\label{III:S:lem:homogeneous}
If $f\ne0$ is supported at length $m$, then $\Snorm{H_f}=a:=\Snorm f_2$. If $m\ge2$, then $\sigma_1=\sigma_2=a$ and $\sigma_3\ge a/\sqrt d$.
\end{lemma}
\begin{proof}
The word-length decomposition makes each nonzero block map degree $q$ to degree $m-q$, with block $C_{m-q,q}(f)$. Their input and output spaces are mutually orthogonal, so their singular-value lists combine. Each Frobenius norm is $a$. The two endpoint blocks have singular value $a$, and every other block has norm at most $a$. The block $C_{1,m-1}$ has at most $d$ rows and Frobenius norm $a$, hence has a singular value at least $a/\sqrt d$.
\end{proof}


\section{Conditional counterexamples and exact sanity checks}\label{III:S:app:conditional}
All examples below use completed-string semantics with $v=\mathbf1$ and initial row uniform. The matrices are nonnegative and satisfy $\beta+\sum_aA_a\mathbf1=\mathbf1$, with strictly positive termination probabilities. They therefore define proper finite-string distributions.

\paragraph{A simple Perron root without positive overlap is insufficient.}
Take
\[
 A_j=\Sdiag(1/2,1/4),\quad
 A_a=\tfrac18\mathbf1e_1^\top,\quad
 A_b=\tfrac18\mathbf1e_2^\top,\quad\beta=(1/4,1/2)^\top.
\]
Both pure normalized states are reached by reset letters. Repeated $j$ leaves them pure, with current termination probabilities differing by $1/4$ for every run length. Thus $\widehat\omega_{j,N}\ge1/4$, although the eigenvalue ratio of $A_j$ is $1/2$. The dominant right eigenvector has a zero coordinate; the second state has zero Perron overlap. This is exactly the hypothesis missing from a simple-eigenvalue-only claim.

\paragraph{A Jordan factor is necessary even with a strictly positive matrix.}
In dimension three put $P=\mathbf1\mathbf1^\top/3$, $u=(1,-1,0)^\top$, $w=(1,1,-2)^\top$, $L=uw^\top$, and
\[
 A_j=\tfrac14I+\tfrac14P+\tfrac1{48}L,\quad
 \beta=(1/8,1/4,3/8)^\top,\quad
 g=(3/16,1/8,1/16)^\top,\quad A_a=ge_1^\top,\ A_b=ge_2^\top.
\]
Here $A_j>0$, $A_j\mathbf1=\mathbf1/2$, $L^2=PL=LP=0$, and both reset states are reachable. Expanding the Jordan power gives
\[
 (e_1^\top-e_2^\top)(2A_j)^N\beta
       =-\frac{N+2}{16}\,2^{-N}.
\]
This is a coordinate of the conditional future difference. Hence no bound $K2^{-N}$ with a fixed $K$ holds, although the exact root rate is $1/2$.

\paragraph{Absolute contraction may leave conditional memory unchanged.}
For
\[
 A_j=I/2,\quad A_a=E_{11}/4,\quad A_b=E_{22}/4,\quad
 \beta=\mathbf1/4,
\]
one has $Q=\begin{psmallmatrix}1/11&1/12\\1/12&1/11\end{psmallmatrix}$ and
$\widehat\omega_N^2=1/66$ for all $N$. The window $j^N$ retains the two pure normalized states, and the simplex has no larger diameter. The raw transition JSR is $1/2$. This example independently checks that the metric used in Section~\ref{III:S:sec:metric} cannot simply be relabeled as a conditional one.



\backmatter
\chapter*{Companion paper}\addcontentsline{toc}{chapter}{Companion paper}
\emph{When Nearly Rank-One Predictions Require Large Sequential Models}
studies a single proper binary trace law with two distinct approximation barriers.
Its proof combines the geometric waiting process of Chapter 1 with the rare
block-enumerating word of Chapter 2. The mixture keeps the operator obstruction
at a fixed rank while enforcing a growing positive-state budget for vanishing
log-loss. The accompanying notebook regenerates the figures and the source archive.
The component proofs, additional hypotheses, and broader results on executable
memory remain available in the three chapters of this volume.
\begin{thebibliography}{99}
\bibitem{littman2001}
Michael L. Littman, Richard S. Sutton, and Satinder Singh.
Predictive Representations of State.
\emph{Advances in Neural Information Processing Systems}, 14:1555--1561, 2001.
\href{https://proceedings.neurips.cc/paper/2001/file/1e4d36177d71bbb3558e43af9577d70e-Paper.pdf}{Proceedings paper}.

\bibitem{carlyle1971}
J.~W. Carlyle and Azaria Paz.
\newblock Realizations by stochastic finite automata.
\newblock \emph{Journal of Computer and System Sciences}, 5\penalty0
  (1):\penalty0 26--40, 1971.
\newblock \doi{10.1016/S0022-0000(71)80005-3}.

\bibitem{fliess1974}
M. Fliess.
Matrices de Hankel. \emph{Journal de Math\'ematiques Pures et Appliqu\'ees}, 53:197--222, 1974; corrigendum, 54:481, 1975.

\bibitem{balle2017}
Borja Balle and Odalric-Ambrym Maillard.
Spectral Learning from a Single Trajectory under Finite-State Policies.
\emph{Proceedings of Machine Learning Research}, 70:361--370, 2017.
\url{https://proceedings.mlr.press/v70/balle17a.html}.

\bibitem{martin2026memory}
Lance Martin.
The components of agent memory.
\emph{Anthropic Engineering}, 6 November 2026.
\url{https://www.anthropic.com/research/components-of-agent-memory}.
Accessed 9 November 2026.

\bibitem{letta2026frontiermemoryv1}
Letta Research.
FrontierMemory 1.0: Benchmarking continual learning through memory.
\emph{FrontierMemory Blog}, October 2026.
\url{https://frontiermemory.com/blog/v1/}.
Accessed 9 November 2026.

\bibitem{yang2025gated}
Songlin Yang, Jan Kautz, and Ali Hatamizadeh.
Gated Delta Networks: Improving Mamba2 with Delta Rule.
\emph{International Conference on Learning Representations}, 2025.
\url{https://arxiv.org/abs/2412.06464v3}.

\bibitem{lacroce2024}
Clara Lacroce, Borja Balle, Prakash Panangaden, and Guillaume Rabusseau.
\newblock Optimal approximate minimization of one-letter weighted finite
  automata.
\newblock \emph{Mathematical Structures in Computer Science}, 34\penalty0
  (8):\penalty0 807--833, 2024.
\newblock \doi{10.1017/S0960129524000276}.
\newblock URL \url{https://arxiv.org/abs/2306.00135}.

\bibitem{sva2019}
Borja Balle, Prakash Panangaden, and Doina Precup.
\newblock Singular value automata and approximate minimization.
\newblock \emph{Mathematical Structures in Computer Science}, 29\penalty0
  (9):\penalty0 1444--1478, 2019.
\newblock \doi{10.1017/S0960129519000094}.
\newblock URL \url{https://arxiv.org/abs/1711.05994}.

\bibitem{lacroce2022}
Clara Lacroce, Prakash Panangaden, and Guillaume Rabusseau.
\newblock Towards an {AAK} theory approach to approximate minimization in the
  multi-letter case, 2022.
\newblock URL \url{https://arxiv.org/abs/2206.00172}.
\newblock Presented at LearnAut 2022.

\bibitem[Brevig and Saksman(2020)]{brevig2020}
O. F. Brevig and E. Saksman.
Coefficient estimates for $H^p$ spaces with $0<p<1$.
\emph{Proceedings of the American Mathematical Society}, 148(9):3911--3924, 2020.
\doi{10.1090/proc/14995}.
\url{https://arxiv.org/abs/1905.09547}.

\bibitem[Brevig and S\o vik(2026)]{outer2026}
O. F. Brevig and \O. S\o vik. \emph{Extremals for linear functionals on $H^p$ are outer when $0<p<1$}. arXiv:2610.21788v1 (15 October 2026). Theorem 1 and Corollary 3. \url{https://arxiv.org/abs/2610.21788}.

\bibitem{rational2026}
L. Baratchart, A. Borichev, S. Chevillard, C. Coiffard Marre, and R. Zarouf.
Lower bounds in $H^2$-rational approximation to Blaschke products.
\emph{Research in the Mathematical Sciences}, 13, article 80, 2026.
Revised author version: arXiv:2603.16295v2, 29 September 2026.
\url{https://arxiv.org/abs/2603.16295}.

\bibitem{popescu2003}
Gelu Popescu.
\newblock Multivariable {Nehari} problem and interpolation.
\newblock \emph{Journal of Functional Analysis}, 200\penalty0 (2):\penalty0
  536--581, 2003.
\newblock \doi{10.1016/S0022-1236(03)00078-8}.

\bibitem{ball2005}
Joseph~A. Ball, Gilbert Groenewald, and Tanit Malakorn.
\newblock Structured noncommutative multidimensional linear systems.
\newblock \emph{SIAM Journal on Control and Optimization}, 44\penalty0
  (4):\penalty0 1474--1528, 2005.
\newblock \doi{10.1137/S0363012904443750}.
\newblock URL \url{https://vtechworks.lib.vt.edu/handle/10919/48136}.

\bibitem[Schaefer et al.(2024)]{schaefer2024}
Marcus Schaefer, Jean Cardinal, and Tillmann Miltzow.
\newblock The existential theory of the reals as a complexity class: A
  compendium, 2024.
\newblock URL \url{https://arxiv.org/abs/2407.18006}.

\bibitem{basu2006}
Saugata Basu, Richard Pollack, and Marie-Fran\c{c}oise Roy.
\newblock \emph{Algorithms in Real Algebraic Geometry}, volume~10 of
  \emph{Algorithms and Computation in Mathematics}.
\newblock Springer, Berlin, Heidelberg, 2nd edition, 2006.
\newblock \doi{10.1007/3-540-33099-2}.

\bibitem{renegar1992}
James Renegar.
\newblock On the computational complexity and geometry of the first-order
  theory of the reals. {Part I}: Introduction. preliminaries. the geometry of
  semi-algebraic sets. the decision problem for the existential theory of the
  reals.
\newblock \emph{Journal of Symbolic Computation}, 13\penalty0 (3):\penalty0
  255--299, 1992.
\newblock \doi{10.1016/S0747-7171(10)80003-3}.

\bibitem[Bultheel and Lasarow(2007)]{schur_nevanlinna}
A. Bultheel and A. Lasarow. Schur--Nevanlinna sequences of rational functions.
\emph{Proceedings of the Edinburgh Mathematical Society}, 50(3):571--596, 2007.
\url{https://doi.org/10.1017/S0013091504001063}.

\bibitem{parrott}
S. Parrott.
On a quotient norm and the Sz.-Nagy--Foia\c{s} lifting theorem.
\emph{Journal of Functional Analysis}, 30(3):311--328, 1978.
\href{https://doi.org/10.1016/0022-1236(78)90060-5}{doi:10.1016/0022-1236(78)90060-5}.

\bibitem{parrott2025}
A. Renard. \emph{A criterion of contractivity for four by four upper-triangular matrices}. arXiv:2506.10485v2 (2025), Theorem 1.3 (Parrott completion). \url{https://arxiv.org/abs/2506.10485}.

\bibitem{bourin2019}
J.-C. Bourin and E.-Y. Lee.
\newblock On the Russo--Dye theorem for positive linear maps.
\newblock \emph{Linear Algebra and its Applications}, 571:92--102, 2019.
\newblock \href{https://doi.org/10.1016/j.laa.2019.02.017}{doi:10.1016/j.laa.2019.02.017};
\href{https://arxiv.org/abs/1911.10573}{arXiv:1911.10573}.

\bibitem{cohen}
S. N. Cohen and E. Fausti.
Hyperbolic contractivity and the Hilbert metric on probability measures.
2023, revised 2024. \href{https://arxiv.org/abs/2309.02413}{arXiv:2309.02413}.

\bibitem{atar}
R. Atar and O. Zeitouni.
Exponential stability for nonlinear filtering.
\emph{Annales de l'Institut Henri Poincar\'e, Probabilit\'es et Statistiques}, 33(6):697--725, 1997.
\href{https://www.numdam.org/item/AIHPB_1997__33_6_697_0/}{Original article}.

\bibitem{denis2008}
F. Denis and Y. Esposito.
On Rational Stochastic Languages.
\emph{Fundamenta Informaticae}, 86(1--2):41--77, 2008.
\doi{10.3233/FUN-2008-861-203}.
\url{https://arxiv.org/abs/cs/0602093}.

\bibitem{gu1996}
Ming Gu and Stanley C. Eisenstat.
\newblock Efficient algorithms for computing a strong rank-revealing QR
factorization.
\newblock \emph{SIAM Journal on Scientific Computing}, 17(4):848--869, 1996.
\newblock \doi{10.1137/0917055}.
\newblock URL \url{https://math.berkeley.edu/~mgu/MA273/Strong_RRQR.pdf}.

\bibitem{gaubert2015}
St\'ephane Gaubert and Zheng Qu.
\newblock Dobrushin's ergodicity coefficient for Markov operators on cones.
\newblock \emph{Integral Equations and Operator Theory}, 81(1):127--150, 2015.
\newblock \doi{10.1007/s00020-014-2193-2}.
\newblock URL \url{https://arxiv.org/abs/1307.4649}.

\bibitem[Gillis and Vavasis(2014)]{gillis2014}
Nicolas Gillis and Stephen A. Vavasis.
\newblock Fast and robust recursive algorithms for separable nonnegative
matrix factorization.
\newblock \emph{IEEE Transactions on Pattern Analysis and Machine
Intelligence}, 36(4):698--714, 2014.
\newblock \doi{10.1109/TPAMI.2013.226}.
\newblock URL \url{https://arxiv.org/abs/1208.1237}.

\bibitem{reznik}
Y. A. Reznik and A. V. Anisimov. \emph{Enumerative
encoding/decoding of variable-to-fixed-length codes for memoryless sources}.
\emph{Proceedings of the 9th International Symposium on Communication Theory
and Applications (ISCTA)}, Ambleside, UK, 16--20 July 2007.
\url{https://reznik.org/papers/ISCTA07_VF-codes.pdf}.

\bibitem[Mauldin and Williams(1988)]{mauldin_williams}
R. D. Mauldin and S. C. Williams.
Hausdorff dimension in graph directed constructions.
\emph{Transactions of the American Mathematical Society}, 309(2):811--829, 1988.
\url{https://doi.org/10.1090/S0002-9947-1988-0961615-4}.

\bibitem{cox2014}
A. Cox, S. Sankaranarayanan, and B.-Y. E. Chang.
A bit too precise? Verification of quantized digital filters.
International Journal on Software Tools for Technology Transfer
16(2):175--190, 2014.
\doi{10.1007/s10009-013-0279-9}.
\url{https://plv.colorado.edu/bec/papers/filters-sttt13.html}.

\bibitem{massey1988}
James L. Massey.
\newblock On the entropy of integer-valued random variables.
\newblock \emph{Proceedings of the Beijing International Workshop on
Information Theory}, pages C1.1--C1.4, July 1988.
\newblock URL \url{https://www.isiweb.ee.ethz.ch/archive/massey_pub/pdf/BI527.pdf}.

\bibitem{calderbank1998}
A. R. Calderbank, Ingrid Daubechies, Wim Sweldens, and Boon-Lock Yeo.
\newblock Wavelet transforms that map integers to integers.
\newblock \emph{Applied and Computational Harmonic Analysis},
5(3):332--369, 1998.
\newblock \doi{10.1006/acha.1997.0238}.
\newblock URL \url{https://cm-bell-labs.github.io/who/wim/papers/integer.pdf}.
\end{thebibliography}

\end{document}
